\documentclass[11pt,a4paper]{article}
\usepackage[T1]{fontenc}
\usepackage{lmodern,microtype,amsmath,amssymb,amsthm,mathtools}
\usepackage[margin=27mm]{geometry}
\usepackage[authoryear,round]{natbib}
\usepackage[colorlinks=true,linkcolor=blue,citecolor=blue,urlcolor=blue]{hyperref}
\usepackage{enumitem}
\setlist{itemsep=3pt,topsep=5pt}
\newtheorem{theorem}{Theorem}[section]
\newtheorem{proposition}[theorem]{Proposition}
\newtheorem{lemma}[theorem]{Lemma}
\newtheorem{corollary}[theorem]{Corollary}
\theoremstyle{definition}\newtheorem{definition}[theorem]{Definition}
\theoremstyle{remark}\newtheorem{remark}[theorem]{Remark}
\newtheorem{example}[theorem]{Example}
\newcommand{\Z}{\mathbb Z}\newcommand{\Q}{\mathbb Q}
\newcommand{\F}{\mathbb F}\newcommand{\Lat}{\operatorname{Lat}}
\newcommand{\tr}{\operatorname{tr}}\newcommand{\Pic}{\operatorname{Pic}}
\newcommand{\workbenchurl}{https://github.com/KokunoYumeto/algebraic-groups-and-representations-workbench}
\newcommand{\recordurl}{https://github.com/KokunoYumeto/mathematics-workbenches}
\title{Grothendieck's question on symmetric-group\newline lattices over local rings}
\author{}
\date{23 November 2026}
\begin{document}
\maketitle
\begin{abstract}
In his manuscript on $\lambda$-rings, Grothendieck asks which coefficient
rings make the induction ring of symmetric-group lattices polynomial
on the trivial representations.  We prove that coherent regular local
rings of characteristic zero have this property, while every
positive-characteristic ring satisfies an explicit nonzero polynomial
relation.  In characteristic zero the proposed polynomial ring always
embeds as a direct summand; the question is generation.
For a monic polynomial over any characteristic-zero DVR in which $p$
is a uniformizer, whose reduction is a power of its variable,
generation holds exactly when the coefficient algebra is regular.
This criterion also holds for every complete $p$-torsion-free local
ring whose special fibre is Artinian with a principal maximal ideal,
and for finite free local algebras over any mixed-characteristic DVR
with that special fibre.  Every singular member has a countable split
sum of finite cyclic classes in its generic kernel, with a retraction
on the whole lattice Grothendieck group that kills the Young span.
We determine their exact order from the first coefficient of valuation
one; when the completed generic algebra is a field, that order is the
polynomial degree.  We also classify arbitrary square-zero extensions
of $\Z_{(p)}$ and quadratic $\Z_p$-orders.  Arbitrary local
$\Q$-algebras, filtered local colimits and rational nil thickenings
give further positive cases; singular normal domains give negative ones.
Every nontrivial truncated polynomial extension of every
mixed-characteristic DVR fails generation, without a restriction on
ramification, residue field, or completeness.  Prime-to-characteristic
exponents have exact-order obstruction classes; for divisible exponents
a different ideal lattice gives a class of exact order $p$, which
splits off when $p^2$ does not divide the exponent.
Two pure-power examples show that a standard witness family can vanish
in the lattice Grothendieck group.  Completion is injective on the full lattice Grothendieck group and
identifies the generic kernels over any characteristic-zero DVR.
We also prove the integral tensor-power formula and an explicit failure
of the special $\lambda$-identities for the split Witt--Green ring.
The classification of all singular mixed-characteristic local rings
remains open.  The mathematics and text were produced by GPT-6 Astra
(ChatGPT Pro mode), with checks by separate instances of the same model;
no human has checked them.
\end{abstract}

\section{The question and the answers}

The coefficient ring is part of the mathematical content of a
representation ring.  A characteristic-zero fraction field forgets
extensions and the geometry of integral lattices; whether its familiar
generators still suffice is therefore a question about integral
$K$-theory, not merely ordinary characters.

Grothendieck asks on cote 133, scan 29 (manuscript page 28), for which
rings $k$ the ring $\bigoplus_n R_k(S_n)$ is polynomial over $\Z$ in
the trivial classes $\rho_i$, and in particular whether locality, or
the condition $K_0(k)\simeq\Z$, suffices.  Scan 28 defines the product
through induction and uses the same generators in a tensor-power
formula.  We use the exact convention: $\Lat_A(G)$ consists of
$A[G]$-modules finite free over the commutative local ring $A$, and
its exact sequences are all short exact sequences with three such
terms.  They split over $A$, but need not split equivariantly.
When a comparison uses a nonlocal coefficient algebra, the same
notation means modules isomorphic to $A^r$ for a single integer $r$;
the exact sequences are coefficient-split.  Thus freeness there
requires a common rank across all components.
All rings in the paper have an identity and are nonzero.

Put
\[
 R(A)=\bigoplus_{n\geq0}K_0(\Lat_A(S_n)),\qquad
 \Phi_A:\Z[h_1,h_2,\ldots]\longrightarrow R(A),\quad h_n\mapsto\rho_n.
\]
The degree of $h_n$ is $n$ and $\rho_0=1$.  The multiplication is
\[
 [M][N]=\left[\operatorname{Ind}_{S_a\times S_b}^{S_{a+b}}
                  (M\otimes_A N)\right].
\]
Tensor product and induction preserve the stated exact sequences,
so this is a well-defined graded ring.  Its degree-$n$ proposed
basis consists of the Young permutation lattices, one for each
partition of $n$.

\begin{theorem}\label{thm:answers}
The following statements hold.
\begin{enumerate}[label=\textup{(\roman*)}]
\item If $A$ is a coherent local ring of characteristic zero such that
every finitely presented $A$-module has finite projective dimension,
or if $A$ is any local $\Q$-algebra, then $\Phi_A$ is an isomorphism.
This includes all noetherian regular local rings and all
characteristic-zero valuation domains.  The property is preserved
by filtered colimits of local rings along local homomorphisms,
and is unchanged by quotienting by a uniquely divisible nil ideal.
\item If $m\cdot1_A=0$ for an integer $m\geq2$, the $m$th Newton
power-sum polynomial satisfies $P_m(\rho_1,\ldots,\rho_m)=0$.
In particular characteristic zero is necessary for $\Phi_A$ to be
an isomorphism.
\item For every characteristic-zero local ring $A$, $\Phi_A$ is
split injective as a graded ring map.
\item Locality is not sufficient for surjectivity, even with
characteristic zero, reducedness and flatness over $\Z$.  The rings
in Propositions~\ref{prop:nilpotent-trace} and
\ref{prop:reduced-fiber-product} give explicit failures.  For odd
residue characteristic they have no square roots of one other than
$\pm1$.
\item For every prime $p$, generation fails over
$\Z_{(p)}[\epsilon]/(\epsilon^2)$, already in degree $p$; the same
holds with $\Z_p$ in place of $\Z_{(p)}$.
For every $p$ the reduction kernel in degree $p$ contains a direct
summand isomorphic to $\bigoplus_{\mathbb N}\Z/2$.
For $p=2$, the involution
$\left(\begin{smallmatrix}1&\epsilon\\\epsilon&-1\end{smallmatrix}\right)$
minus the regular representation is a nonzero class of exact order two.
The same direct-summand conclusion holds over
$O[\epsilon]/(\epsilon^2)$ for every characteristic-zero DVR $O$
of residue characteristic two, with arbitrary ramification and residue
field.  For odd $p$, it holds in degree $p$ whenever $O$ is a
DVR of residue characteristic $p$ and
$e=v_O(p)$ satisfies $v_2(e)\le v_2(p-1)$; $v_2$ denotes the
exponent of two.  This includes $O=\Z_3[\zeta_3]$ and
$O=\Z_3[\sqrt3]$.
For every odd-residue DVR, without any ramification or residue-field
restriction, there is a nonzero exact-order-two class in both the
degree-$p$ reduction and generic kernels, detected by an additive
homomorphism annihilating the Young span
(Corollary~\ref{thm:all-dvr-odd-dual}).  No split-summand conclusion
is asserted by that unrestricted theorem.
For every $n\ge2$ and $O=\Z_{(p)}$ or $\Z_p$, the degree-$p$
reduction kernel for $O[e]/(e^n)$ contains a split countable direct
sum of copies of $\Z/n$.  The explicit cyclic classes and their
inductions have exact order $n$ and vanish after inverting $p$.  More generally, for any mixed-characteristic DVR $O$ with
$h=v_O(p)$, the same split conclusion holds when $(n,h)=1$,
and generation fails whenever $n\nmid h$.
\item If $M$ is any $\Z_{(p)}$-module and
$A=\Z_{(p)}\oplus M$ with $M^2=0$, then $\Phi_A$ is an isomorphism
if and only if $M$ is a $\Q$-vector space.  This also characterizes
every square-zero extension with quotient $\Z_{(p)}$; the extension
has a canonical splitting.  There is no finite-generation hypothesis.
\item For $A=\Z_{(p)}[\epsilon]/(p\epsilon,\epsilon^2)$ the degree-$p$
reduction kernel contains a direct summand isomorphic to
$\bigoplus_{\mathbb N}\Z$.  For odd $p$, an explicit infinite-order
class has trace zero at every element of $S_p$.
\item For $A=\Z_{(2)}[t]/(t^2+4)$, and also with $\Z_2$ in place of
$\Z_{(2)}$, the ordinary-character kernel in degree two contains a
direct summand isomorphic to $\bigoplus_{\mathbb N}\Z/2$.
More generally, for every prime $p$, a $\Z_p$-order in a quadratic field extension of $\Q_p$
satisfies polynomiality
exactly when it is maximal; every nonmaximal order has the same infinite
two-torsion conclusion in degree $p$.  These are one-dimensional noetherian local
domains with field generic fibre.
\item Let $p$ be any prime and let $O$ be any characteristic-zero
DVR whose uniformizer is $p$, with arbitrary residue field and no
completeness hypothesis.  Let
\[
 F=t^n+\sum_{i=0}^{n-1}a_it^i,\qquad n\ge2,\quad a_i\in pO.
\]
For $A=O[t]/F$, the following are equivalent:
\[
 \Phi_A\text{ is an isomorphism},\qquad A\text{ is regular},
 \qquad a_0\notin p^2O.
\]
If $a_0\in p^2O$, let $r$ be the least index $1\le i<n$
with $a_i\notin p^2O$, or $r=\infty$ if none exists, and put
$N=\min(n,pr)$, with $N=n$ when $r=\infty$.
The degree-$p$ generic kernel contains a countable direct sum of
exact order-$N$ classes that splits from the whole lattice
Grothendieck group.  Its coordinate retraction annihilates the
Young span.  The analogous assertion holds for $C_p$.
This includes every prime, $N<n$, and nonreduced generic algebras.
\item Under the singular hypotheses of the preceding assertion,
if the completed generic algebra
$\operatorname{Frac}(\widehat O)[t]/F$ is a field, then
$N=n$.  Thus every such field order has a split countable sum of
exact order-$n$ generic-kernel classes.  In particular this covers
every distinguished monogenic field order of generic residue degree
greater than one, as well as the totally ramified ones whose
constant term lies in $p^2O$.
\item For every mixed-characteristic DVR $O$ of residue characteristic
$p$ and every $n\ge2$, generation fails over $O[t]/t^n$ in degree
$p$.  There is no restriction on ramification, residue field or
completeness.  If $p\nmid n$, there is an exact-order-$n$ class in
both the reduction and generic kernels, detected by a whole-group
homomorphism annihilating the Young span.  If $p\mid n$, such a
homomorphism to $\Z/n$ takes the value $n/p$ on a class of
exact order $p$ in both kernels.  When $p^2\nmid n$, its cyclic
subgroup splits from the whole Grothendieck group, with a retraction
annihilating the Young span.  If $h=v_O(p)$ is prime to $p$,
the standard mesh instead gives a split cyclic subgroup of exact
order $p^{v_p(n)}$ containing the vertical class.  When $p^2\mid n$,
the order-$p$ subgroup generated by that vertical class is not a
direct summand of the whole group.
\item Let $A$ be a complete noetherian mixed-characteristic local
ring such that $p$ is a nonzerodivisor and $A/pA$ is Artinian with
embedding dimension at most one.  Then $\Phi_A$ is an isomorphism
if and only if $A$ is regular, equivalently a DVR.  In the singular
case a Cohen coefficient presentation as in
Theorem~\ref{thm:curvilinear-special-fibre} gives a whole-group split
sum $\bigoplus_{\mathbb N}\Z/N$ in the degree-$p$ generic kernel,
with a Young-killing retraction.  If $A[1/p]$ is a field, then
$N=\operatorname{length}_{A/pA}(A/pA)$.
The criterion and the split obstruction also apply to every local
finite free algebra over any mixed-characteristic DVR with this
special fibre.  In particular they apply to every finite free local
monogenic algebra over an unramified DVR, without any completeness,
residue-cardinality or residue-separability restriction.
\end{enumerate}
\end{theorem}

For every prime $p$, Theorem~\ref{thm:normal-domain-counterexample}
gives a two-dimensional normal local domain of residue characteristic
$p$ over which generation fails.  It is flat over $\Z_{(p)}$ and
can be chosen complete.  Dimension two is the smallest possible for
a noetherian normal local domain counterexample: a one-dimensional
noetherian normal local domain is a discrete valuation ring.
Thus the remaining classification concerns generation over singular
mixed-characteristic rings; no failure of polynomial independence
can occur in characteristic zero.  The positive result does not
require an unramified residue characteristic.  Conversely, a
determinant test is useful but cannot characterize generation.  The nonzero-primary Tate length-parity detectors used for odd
ramification all vanish when $v_2(e)>v_2(p-1)$,
by Theorem~\ref{thm:complementary-detector-limit}.
Corollary~\ref{thm:all-dvr-odd-dual} proves failure of generation
throughout that range using a divisor class.  The vanishing theorem
therefore describes a limitation of length parity, rather than a
positive class of coefficient rings.

The integral polynomial bialgebra is classical:
\citet[Proposition~1.1, p.~543, and Theorem~1.2, p.~544]{Burroughs1974}
proves it over $\Z$.  Swan's comparison theorem for semilocal
Dedekind domains supplies the traditional discrete-valuation case:
\citet[Addendum, p.~242]{Gersten1973} reports
\citet[Theorem~3, p.~87]{Swan1963} in exactly this form.  Our proof for regular local
rings uses the resolution and localization theorems and James's
unitriangular modular decomposition.  We state exactly where those
published results enter.  We make no assertion that an arbitrary
singular coefficient ring has the same Grothendieck group as its
fraction field.

Recent work of \citet[Theorem A, p.~2912]{BarthelEtAl2025}
stratifies derived representation categories over commutative
noetherian rings by group cohomology.  This gives a structural
setting for lattice categories over singular bases; their exact
Grothendieck groups and the induction generators considered here
are a separate invariant.

\citet[Section~2, p.~105; Theorem~3, p.~108]{Roggenkamp1971}
computes exact Grothendieck groups of local Bass-order lattices,
including two-torsion cases.  Those lattices are projective over an
underlying discrete valuation ring.  Here freeness is required over
the specified, possibly singular, coefficient ring, which defines
a different exact category.


A recent note of \citet[Theorem~1.2(c), p.~4, and
Section~3.2, p.~12]{GPT6Grinberg2026} proves integrality of the
anti-arithmetic product of symmetric functions using exterior powers
and Boorman's generation theorem for the rational representation
ring of a fixed symmetric group.  Its virtual-character setting
describes the rational specialization of the lattice groups here.
The coefficient-ring question asks which additional exact lattice
classes survive before that specialization.

The comparison of congruent coefficient algebras in
Theorem~\ref{thm:congruent-coefficients} uses the classical
averaging and precision estimates of
\citet[Theorems~2--3, pp.~347--351]{Maranda1953}.
We prove the change of coefficient algebra and the preservation
of every coefficient-split exact relation, which transfer
Grothendieck-group obstructions from nilpotent thickenings to
nonzero deformations.  The congruence bound depends on the group
order; a fixed modulus does not compare all symmetric-group
degrees at once.

Theorem~\ref{thm:completion-generic-kernel} also compares a finite
free algebra over any characteristic-zero DVR with its completion.  It proves injectivity
on the full lattice Grothendieck group and an isomorphism on the
kernels of generic specialization.  The algebra need be neither
local nor reduced; the underlying modules must be free of a common
coefficient rank.  The published $K_0$ input is
\citet[Chapter~II, Lemma~7.2, p.~62 of the author chapter version]{Weibel2013}:
an extension-closed cofinal exact subcategory induces an injection
on $K_0$.  The arithmetic argument here constructs such subcategories
by descending regular generic representations and their exact relations.
This comparison gives the monic-deformation assertion over the
localized integers as well as the complete integers.

The monic theorem at depth $p^2$ also uses the triangulated
stable category of the Frobenius algebra $\mathbb F_p[t]/t^n$.
The category and its exact triangles are those of Happel's theorem,
as stated by \citet[Section~13.4, Theorem~13.7, author p.~50]{Buhler2010}.
The passage to a selected summand of an exact triangle uses
\citet[Lemma~1.6, pp.~821--822]{BalmerSchlichting2001}.
The coefficient operators, their compatibility with suspension,
and the resulting length invariant are proved below.

The manuscript also asks whether the group completion of orthogonal
sums of nondegenerate quadratic modules is a special $\lambda$-ring,
even when two is invertible (cote 133, scan 6, manuscript page 5).
The answer to that split question is negative.  The modern theorem
of \citet[main theorem, p.~1056; Definition~1.2, p.~1060]{Zibrowius2015}
concerns the Grothendieck--Witt group with its metabolic relations;
our calculation is consistent with that theorem.

\section{Proof overview}

The positive theorem separates the ordinary and modular fibres.
Over a regular local base, every finite equivariant module has a
finite resolution by modules free over the base.  On inverting the
residue characteristic the group algebra becomes a product of
matrix algebras.  The only possible obstruction is therefore a
class supported on the closed characteristic.  Its modular simple
factors can be expressed as integral combinations of Young
permutation modules.  Resolving their multiplicity modules over the
regular base gives Euler rank zero and kills the obstruction.
This is where regularity is used; the special fibre itself need
not be regular.

A different positive mechanism applies to nil ideals whose additive
groups admit unique division by every integer.  Matrix defects in a
lifted finite-group action then form cocycles that can be averaged
inside the ideal, even though the group order need not be a unit of
the coefficient ring.  The same calculation lifts exact sequences.
Finite sets of defects lie in a nilpotent rational subideal, so the
argument works for nil ideals without a uniform nilpotence bound.
It produces positive mixed-characteristic rings that are not coherent.

For the negative results, the induction generators have much more
rigid invariants than arbitrary lattices.  They have integral traces,
and after inverting the residue characteristic their invariant
modules have trivial determinant line.  Inducing a nonconstant
root of unity violates the trace condition.  A rank-two cyclic
representation over a singular normal local domain, induced to a
symmetric group, violates the invariant-line condition.  In positive characteristic a
different mechanism applies: the diagonal vector lies in the
augmentation representation, and its exterior powers give a
Newton relation.  The split quadratic example is a finite Jordan
block calculation.

The flat square-zero examples require an invariant beyond traces and
generic eigenbundles.  The nilpotent extension gives an odd boundary
on two-periodic Tate cohomology.  Its square has primary components
whose composition lengths modulo two are additive on admissible
exact sequences.  Induction from a cyclic subgroup gives the obstruction
for every prime.  An explicit exact sequence annihilates each flat
class by two, and companion matrices produce infinitely many
independent classes in every residue characteristic.  For odd-prime
ramification, a finite coefficient algebra makes the cyclotomic
distance odd in the stated range.  A companion representation supplies
two Tate groups of that odd length.  A spectral parameter selects one
residue factor of the coefficient algebra, so completeness is unnecessary.  A separation map
from any torsion-free module into the $p$-adic integers then gives the
exact square-zero idealization criterion.

For arbitrary odd-prime ramification, a second invariant replaces
parity by a divisor class on the normal surface
$y^2=(1+x)^p-1$.  Its ideal $(x,y)$ is nonprincipal and has
order two.  A lifted square-zero differential has square equal to
the surface equation times an operator; on the primary-one part,
a square root normalizes this to a module on the surface.  An
integral averaged intertwiner proves independence of the lift,
and upper-triangular lifts prove additivity on every exact triple.
Stable curvature first extracts a cyclic-group-free summand from
an arbitrary lattice.  The companion witness maps to $(x,y)$,
and normalizer twists separate its value before induction.

For every truncation length, lift the group matrix to a formal power
series ring and divide its failure to have order $p$ by the defining
power of the parameter.  The resulting curvature is canonical modulo
maps through regular lattices.  Its nonzero primary components have
an exact integral Tate complex.  Their invariant lattices therefore
have additive covolumes, and a change of generic basis alters the
valuation by a multiple of the truncation length.  The asymmetric
cyclic lattice has invariant module $(p,e)$, of covolume one.
This gives exact torsion of every truncation order.  Completion is
used for primary decomposition; composition with completion supplies
the invariant over the localized integers as well.

When the truncation exponent is divisible by the residue
characteristic, the ramification factor in the first companion
family is removed by changing the integral lattice.  Write $n=pr$
and translate the group coordinate by $t^r$.  The ideal generated
by this translated coordinate and a uniformizer is free over the
coefficient polynomial ring.  Its cyclic curvature is a unit, and
its invariant lattice is $(\pi,t^r)$ inside the regular invariant
line.  Its covolume is $r$, irrespective of the ramification index.
The corresponding primary coordinate survives induction with one
normalizer term.  This settles every divisible exponent, while the
surface argument settles the exponents prime to $p$.

For torsion ideals, a map to the divisible module $\Q/\Z_{(p)}$
separates an element of order $p$.  Composing the ideal-extension
boundary with the rational coefficient boundary gives an even
operator.  Its primary Euler characteristics take integer values,
producing infinite-order classes.  For the singular quadratic order,
the two integral coefficients of an involution instead give operators
on residual homology satisfying a matrix-algebra relation.  The same
relation survives connecting maps, and half-length parity detects
infinitely many classes of order two over a domain.
At odd primes the canonical homotopy for multiplication by $p$ gives
a contracting square-zero operator on each primary component.  This
forces even lengths, including on connecting images; half-length parity
then covers every nonmaximal quadratic field order.
The distinguished monic classification uses the generic hypersurface
of the two equations defining the coefficient algebra and the cyclic
group algebra.  Bounded periodic group cohomology restricts every
lattice to finitely many singular points of a projective parameter
line.  At a chosen point, an integral root gives a regular immersion
into a truncated polynomial ring.  The truncation length is precisely
$N=\min(n,pr)$: the first coefficient of valuation one competes with
the $p$th power in the cyclic-group equation.  A finite filtration
defines length modulo $N$ on the supported singularity category.
This produces an additive invariant for every admissible sequence.
The witness lattices form coefficient-split meshes with regular
endpoints, proving $N$-annihilation.  Their ideals $(p,t^i)$ have
detected value $i$, proving the exact orders.  Unramified parameters
separate both Galois conjugates and normalizer terms, so induction
to $S_p$ preserves value one.  For a completed generic field,
uniqueness of the extended valuation forces $N=n$.
When the residue field is infinite, rational parameter orbits give
the same separation; finite fields instead use finite unramified
extensions of distinct degrees.  A Cohen coefficient ring then
reduces a complete ring with an Artinian principal special fibre
to the distinguished monic presentation.  Lifting a residue power
basis and cancelling powers of $p$ proves the required finite
freeness directly.  Its cotangent space makes regularity equivalent
to the constant coefficient having valuation one.

For a divisible truncation exponent, the ideals $(\pi,Y^j)$,
where $Y=X+t^{n/p}$, form an admissible finite mesh with regular
endpoints.  This proves $p$-annihilation of the vertical-ideal
class; its primary covolume then proves exact order $p$.
Reducing the covolume modulo $p$ supplies a retraction when
$v_p(n)=1$.

The norm calculation for field orders gives an additional description
of a quotient of this obstruction.


Finally, transport action matrices through a sufficiently precise
isomorphism of coefficient reductions.  Their multiplicative defects
are small.  Averaging corrects the defects with a controlled loss
of divisibility, and iteration converges when the initial precision
exceeds twice the valuation of the group order.  Averaged
intertwiners make the resulting isomorphism class unique.  Performing
the correction inside upper triangular matrix algebras preserves
all admissible exact relations.  Thus the full truncated-polynomial
obstruction survives deep deformations of its equation.
To pass from a complete DVR to the original DVR, add the coefficient-free
kernel of a regular representation mapping onto a lattice.  The resulting
generic representation is free over the group algebra.  Approximation
of a coefficient basis then descends the lattice, and finite group
cohomology descends each coefficient-split extension.  Cofinality
preserves the exact relations on $K_0$.  Thus completion identifies
all generic-kernel classes, including their split torsion subgroups.
For binomials, finite companion formulas exhibit these generators
directly over either base.

An explicit binomial construction gives a second description through
stable modules over the residual truncated polynomial ring.  The
coefficient Bockstein and cyclic homotopy select summands whose
lengths are additive modulo the truncation degree.  Its companion
matrices and the shallow-coefficient calculation provide small
arithmetic examples of the general theorem.

\section{Polynomial independence and rational coefficients}

This section isolates the part of the question which has an answer
over every characteristic-zero coefficient ring.

\begin{proposition}\label{prop:split-injection}
For any commutative ring $A$ of characteristic zero, the homomorphism
$\Phi_A$ is split injective.  If $A$ is local and contains $\Q$,
it is an isomorphism.
\end{proposition}
\begin{proof}
Since no positive integer annihilates $1_A$, the ring
$A\otimes_\Z\Q$ is nonzero.  Choose a prime ideal in it and let
$F$ be the fraction field of the corresponding domain.  Then
$F$ is a characteristic-zero field with a unital map from $A$.
Base change along this map is exact on lattice exact sequences,
because their underlying $A$-module sequences split.  The composite
\[
 \Z[h_1,h_2,\ldots]\xrightarrow{\Phi_A}R(A)
           \longrightarrow R(F)
\]
is the ordinary characteristic-zero isomorphism, with $h_n$ sent
to the trivial $S_n$-module.  Its inverse composed with the second
arrow is the required retraction.

If $A$ is a local $\Q$-algebra, tensor the split Wedderburn
decomposition of $\Q[S_n]$ with $A$.  A module finite free over
$A$ decomposes into the standard matrix modules tensored with
finite projective multiplicity modules over $A$.  Those
multiplicity modules are free because $A$ is local.  Hence the
ordinary irreducible representations form a basis for its exact
Grothendieck group.  The Young permutation basis and induction
multiplication give the assertion, in each degree and as a ring.
\end{proof}

% Self-contained proof module. Citation keys: James1978, Quillen1973.
\section{Regular local coefficient rings}

The positive result over a discrete valuation ring extends to every
commutative noetherian regular local ring of characteristic zero.  The
argument uses ordinary and modular representation theory only to identify
two triangular systems of generators; the passage between the fibres is
an argument with finite resolutions.

\begin{theorem}\label{thm:regular-local}
Let $A$ be a commutative noetherian regular local ring of characteristic
zero.  For every $n\geq 0$, extension of scalars to a characteristic-zero
field containing $A$ induces an isomorphism
\[
 K_0(\operatorname{Lat}_A(S_n))\simeq R_{\mathbb Q}(S_n).
\]
Consequently the homomorphism
\[
 \mathbb Z[\rho_1,\rho_2,\ldots]\longrightarrow
 \bigoplus_{n\geq 0}K_0(\operatorname{Lat}_A(S_n)),
\]
with multiplication given by induction and $\rho_n$ the trivial
$A[S_n]$-lattice, is an isomorphism.
\end{theorem}

\begin{proof}
Fix $n$ and write $G=S_n$.  We first compare lattices with arbitrary
finitely generated modules.  Put $d=\dim A$.  Any finitely generated
$A[G]$-module $N$ has a resolution
\[
 0\longrightarrow L\longrightarrow A[G]^{a_{d-1}}
 \longrightarrow\cdots\longrightarrow A[G]^{a_0}
 \longrightarrow N\longrightarrow 0
\]
in which $L$ is finite free over $A$ (if $d=0$, take $L=N$): choose successive finite free
$A[G]$-covers, and use $\operatorname{gl.dim}(A)=d$ for the last kernel.
The category of lattices is closed under extensions and under kernels
of surjections between its objects, since such surjections split over
$A$.  The resolution theorem therefore gives
\begin{equation}\label{eq:lattice-G0}
 K_0(\operatorname{Lat}_A(G))\simeq G_0(A[G]).
\end{equation}
Here $G_0$ denotes the Grothendieck group of finitely generated modules.
At the level used here, the inverse sends a module to the alternating
sum of any finite lattice resolution.  Comparison of resolutions and
the horseshoe construction prove independence and additivity; an
acyclic bounded complex of $A$-free modules has $A$-free cycles, by
successively splitting from its right end.  This is also the
$K_0$ case of the resolution theorem in
\cite[\S4, Theorem~3 and Corollary~1]{Quillen1973}.

If the residue field of $A$ has characteristic zero, every positive
integer is invertible in $A$.  The split Wedderburn decomposition
\[
 \mathbb Q[S_n]\simeq
 \prod_{\lambda\vdash n}\operatorname{Mat}_{d_\lambda}(\mathbb Q)
\]
then gives $A[S_n]\simeq\prod_{\lambda\vdash n}
\operatorname{Mat}_{d_\lambda}(A)$.  Morita equivalence and
$G_0(A)=\mathbb Z[A]$ prove the assertion in this case.

Suppose instead that the residue characteristic is $p>0$, and put
$B=A/pA$ and $C=A[1/p]$.  Every prime different from $p$ is already a
unit in $A$, so $C$ is a $\mathbb Q$-algebra.  Localization of
finitely generated modules gives an exact sequence
\begin{equation}\label{eq:regular-localization}
 G_0(\mathcal T)\longrightarrow G_0(A[G])
 \longrightarrow G_0(C[G])\longrightarrow 0,
\end{equation}
where $\mathcal T$ consists of modules killed by a power of $p$.
Filtering a module by powers of $p$ shows that the image of the first
map is generated by the classes of finitely generated $B[G]$-modules.
We prove that every such class is zero in $G_0(A[G])$.

Let $J$ be the radical of the finite-dimensional algebra
$\mathbb F_p[G]$.  By the splitting-field theorem for symmetric groups,
\[
 \mathbb F_p[G]/J\simeq
 \prod_{\mu\ p\text{-regular}}\operatorname{Mat}_{e_\mu}(\mathbb F_p).
\]
The ideal $B\otimes_{\mathbb F_p}J$ is nilpotent in $B[G]$.
Its powers filter every finitely generated $B[G]$-module.  Morita
equivalence for the displayed product shows that each successive
quotient is a direct sum of modules
\[
 M\otimes_{\mathbb F_p}D^\mu,
\]
where $M$ is a finitely generated $B$-module and $D^\mu$ is a simple
$\mathbb F_p[G]$-module.  The group acts trivially on $M$.

For a partition $\nu$ of $n$, let $Y^\nu_{\mathbb Z}$ denote the
permutation lattice on the cosets of its Young subgroup.  James's
splitting and triangularity theorems
\cite[Theorem~11.5, p.~40, and Theorem~12.1, p.~42]{James1978}
imply that the classes of $Y^\nu_{\mathbb F_p}$, with $\nu$
$p$-regular, form a $\mathbb Z$-basis of
$G_0(\mathbb F_p[G])$.  Indeed, their composition multiplicities
against the $D^\mu$, ordered by dominance, form a unitriangular
matrix.  Write
\[
 [D^\mu]=\sum_\nu a_{\mu\nu}[Y^\nu_{\mathbb F_p}].
\]
Tensoring an exact sequence of $\mathbb F_p[G]$-modules with the
$\mathbb F_p$-vector space $M$ is exact.  Therefore
\begin{equation}\label{eq:kill-modular-class}
 [M\otimes_{\mathbb F_p}D^\mu]
 =\sum_\nu a_{\mu\nu}
 [M\otimes_A Y^\nu_A]
 \quad\hbox{in }G_0(A[G]).
\end{equation}
The finitely generated $A$-module $M$ has a finite free resolution
$F_\bullet$.  Tensoring this resolution with the $A$-free lattice
$Y^\nu_A$ gives
\[
 [M\otimes_A Y^\nu_A]
 =\left(\sum_i(-1)^i\operatorname{rank}_A F_i\right)
 [Y^\nu_A]=0.
\]
The last equality holds because $pM=0$ and $p\ne0$ in the domain
$A$, so $M$ has rank zero.  Equation~\eqref{eq:kill-modular-class}
proves the required vanishing.  Thus
\begin{equation}\label{eq:regular-invert-p}
 G_0(A[G])\simeq G_0(C[G]).
\end{equation}

Finally, localization gives a surjection $G_0(A)\to G_0(C)$.
Since $G_0(A)=\mathbb Z[A]$, and rank over the fraction field sends
$[C]$ to $1$, this map is an isomorphism and
$G_0(C)=\mathbb Z[C]$.  Tensoring the rational split Wedderburn
decomposition with $C$ now yields
\[
 G_0(C[G])\simeq\bigoplus_{\lambda\vdash n}G_0(C)
 \simeq\mathbb Z^{\{\lambda\vdash n\}}.
\]
These identifications agree with extension to the fraction field of
$A$ and hence with the ordinary representation group of $S_n$.
The ordinary Young permutation classes form a unitriangular integral
basis, and the monomial $\rho_{\lambda_1}\cdots\rho_{\lambda_r}$
is precisely the class induced from the corresponding Young subgroup.
The degreewise isomorphisms are compatible with induction, proving
the final assertion.
\end{proof}

\begin{remark}
No regularity of $A/pA$ was used.  In particular, the theorem includes
ramified regular local rings.  Regularity of $A$ enters in the
comparison with $G_0$ and in the finite free resolution of the
multiplicity module $M$.
\end{remark}

\section{Further positive coefficient rings}

Two elementary extensions make clear that the positive cases are
not confined to noetherian regular rings.

\begin{proposition}\label{prop:filtered-colimits}
Let $(A_i)$ be a filtered system of commutative local rings and
unital local homomorphisms.  Suppose that, for every $i$, the
trivial symmetric-group classes freely generate the induction
ring over $A_i$.  Then they do so over $A=\varinjlim_iA_i$.
\end{proposition}

\begin{proof}
The colimit of a filtered system of local rings with local maps
is local.  For a fixed finite group $G$, every finite free
$A$-module with $G$-action is specified by finitely many matrices
and finitely many equations: matrix inverses and the group
multiplication equations.  It therefore descends to a finite
stage.  An equivariant isomorphism descends after increasing
the stage, since it too is specified by finitely many entries
and equations.

An exact sequence of lattices is split as a sequence of modules
over its coefficient ring.  Choose bases compatible with such
a splitting.  Its representation matrices have block triangular
form, and the diagonal blocks specify its submodule and
quotient.  This finite matrix description descends as an exact
sequence of lattices.  Conversely all relations in $K_0$ use
only finitely many isomorphisms and exact sequences.  Hence
\[
 \varinjlim_i K_0(\operatorname{Lat}_{A_i}(G))
 \simeq K_0(\operatorname{Lat}_A(G)).
\]
These isomorphisms commute with induction.  In each degree the
maps between the groups on the left are the identity after
their identifications with the free group on Young induction
monomials.  Taking the colimit proves the assertion.
\end{proof}

\begin{proposition}\label{prop:regular-coherent}
Let $A$ be a commutative coherent local ring of characteristic
zero such that every finitely presented $A$-module has finite
projective dimension.  Then the trivial symmetric-group
classes freely generate the induction ring over $A$.
In particular, this holds for every valuation domain of
characteristic zero.
\end{proposition}

\begin{proof}
Use $G_0^{\mathrm{fp}}$ for the Grothendieck group of finitely
presented modules.  If $G$ is finite, a finitely presented
$A[G]$-module is finitely presented over $A$.  Conversely, an
$A[G]$-module finitely presented over $A$ is finitely presented
over $A[G]$: take a finite free $A[G]$-cover, whose kernel is
finitely presented over $A$ by coherence, and hence finitely
generated over $A[G]$.  It follows as well that $A[G]$ is
coherent.  Thus these finitely presented module categories
are abelian.

Successive finite free $A[G]$-covers provide a finite resolution
by lattices of any finitely presented $A[G]$-module.  The
process stops because the underlying $A$-module has finite
projective dimension; the last kernel is finite projective
over $A$, hence free.  Therefore the resolution argument gives
\[
 K_0(\operatorname{Lat}_A(G))\simeq
 G_0^{\mathrm{fp}}(A[G]),\qquad
 G_0^{\mathrm{fp}}(A)=\mathbb Z[A].
\]
If the residue characteristic is zero, Proposition~\ref{prop:split-injection} applies.  Otherwise let it be
$p$.  Localization of the coherent module category in $p$
has kernel the finitely presented modules killed by a power
of $p$, and has quotient the finitely presented module
category over $A[1/p][G]$.  The same finite presentations and
clearing-denominators argument as for noetherian modules
proves this assertion.

The filtration and modular Young triangularity argument in
the proof of Theorem~\ref{thm:regular-local} applies verbatim
to these finitely presented modules.  All its kernels,
images, and multiplicity modules are finitely presented,
by coherence.  The only required vanishing is that, for a
finitely presented module $M$ killed by $p$, a finite free
resolution has alternating rank zero.  Choose a prime of
$A[1/p]$, which exists because characteristic zero implies
$A[1/p]\ne0$, and tensor the resolution with its residue
field.  The resolution stays exact after inverting $p$ and
splits there, since it is an acyclic bounded complex of
projective modules.  Its Euler rank is consequently zero.
Thus localization induces
\[
 G_0^{\mathrm{fp}}(A[G])\simeq
 G_0^{\mathrm{fp}}(A[1/p][G]).
\]
Also $G_0^{\mathrm{fp}}(A[1/p])=\mathbb Z[A[1/p]]$:
localization from $\mathbb Z[A]$ is surjective, and every finitely presented $A[1/p]$-module has a finite
projective resolution, obtained by descending a presentation to $A$
and localizing.  Thus $G_0^{\mathrm{fp}}(A[1/p])=K_0(A[1/p])$;
rank at the same residue field proves that the image of $[A]$
has infinite order.  The rational split matrix decomposition
finishes the proof as before.

A valuation domain is coherent.  More explicitly, elementary
row and column operations diagonalize every finite matrix
over it: among its finitely many nonzero entries choose one
dividing all the others, clear its row and column, and
continue.  Thus every finitely presented module is a finite
direct sum of free modules and cyclic modules $A/(a)$.
For $a\ne0$, the latter has the length-one free resolution
$0\to A\xrightarrow{a}A\to A/(a)\to0$.
This verifies the hypotheses for valuation domains.
\end{proof}

\section{Nil ideals with uniquely divisible additive groups}

There are positive mixed-characteristic examples outside the coherent
case.  The reason is elementary rigidity: finite-group cohomology
vanishes on a module on which division by the group order is possible.
The coefficient ring itself need not contain that inverse.

\begin{theorem}\label{thm:rational-nil}
Let $A$ be a commutative local ring and $I$ a nil ideal whose
additive group is uniquely divisible.  For every finite group $G$,
reduction induces an isomorphism
\[
 K_0(\Lat_A(G))\xrightarrow{\ \sim\ }
 K_0(\Lat_{A/I}(G)).
\]
These isomorphisms commute with tensor products, restriction and
induction.  Consequently $\Phi_A$ is an isomorphism if and only
if $\Phi_{A/I}$ is an isomorphism.
\end{theorem}

Here a nil ideal means that every element is nilpotent; a common
nilpotence exponent for the ideal is not required.  Unique divisibility
means precisely that its additive group is a $\Q$-vector space.

\begin{proof}
We first suppose $I^2=0$ and put $B=A/I$ and $N=|G|$.
Every matrix in $\mathrm{GL}_r(B)$ lifts to one in
$\mathrm{GL}_r(A)$, since an element invertible modulo a nil ideal
is invertible.  Given a representation $\bar\rho$ on $B^r$,
choose lifts $s_g\in\mathrm{GL}_r(A)$, with $s_1=1$.
The kernel of the reduction map on general linear groups is
$1+M_r(I)$; its multiplication is addition of the matrices in
$M_r(I)$.  Conjugation by $s_g$ gives an action of $G$ on
$M=M_r(I)$, independent of the chosen lifts.

Define $c(g,h)\in M$ by
\[
 s_gs_h=(1+c(g,h))s_{gh}.
\]
Associativity gives
\[
 g\,c(h,k)-c(gh,k)+c(g,hk)-c(g,h)=0.
\]
Division by $N$ makes sense in $M$.  Put
$b(g)=N^{-1}\sum_{k\in G}c(g,k)$.  Summing the last identity
over $k$ gives
\[
 c(g,h)=b(g)+g\,b(h)-b(gh).
\]
Replacing $s_g$ by $(1-b(g))s_g$ therefore gives a representation
lifting $\bar\rho$.

Two such lifts are isomorphic by an isomorphism reducing to the
identity.  Indeed write $\rho'(g)=(1+d(g))\rho(g)$.
The homomorphism identities give
$d(gh)=d(g)+g\,d(h)$.  If
$b=N^{-1}\sum_h d(h)$, summation gives $d(g)=b-gb$.
Conjugation by $1+b$ carries $\rho$ to $\rho'$.
An arbitrary isomorphism between reductions can first be lifted
as an invertible matrix and then corrected in this manner.
Thus reduction is a bijection on isomorphism classes of lattices.

The same argument lifts every exact-sequence relation.  Choose
bases over $B$ adapted to a split underlying sequence
$0\to B^a\to B^{a+b}\to B^b\to0$.
Its group action lies in the group of invertible block upper
triangular matrices.  Reduction of that group is surjective;
its kernel is the additive group of block upper triangular
matrices with entries in $I$.  This is again uniquely divisible,
and the preceding two-cocycle calculation applies unchanged.
The resulting block representation over $A$ gives an admissible
exact sequence lifting the original one.  Conversely reduction
preserves every admissible exact sequence, because it splits over
$A$.  A bijection on isomorphism classes together with lifting
of these defining relations proves the asserted $K_0$ isomorphism.
No assertion of full faithfulness on morphisms is needed.

If $I$ is nilpotent, apply the square-zero case successively to
$A/I^{j+1}\to A/I^j$.  Each quotient $I^j/I^{j+1}$ is
uniquely divisible.  To see this, division by an integer in $I$
is $A$-linear, and
$(x_1\cdots x_j)/m=(x_1/m)x_2\cdots x_j$.
Thus every $I^j$ is a $\Q$-subspace of $I$.

Finally suppose merely that $I$ is nil.  Any finite set of its
elements generates a nilpotent ideal $J$, since the ring is
commutative.  Its rational saturation
\[
 J_{\Q}=\{x\in I:mx\in J\text{ for some integer }m>0\}
\]
is an $A$-ideal, is uniquely divisible, and is nilpotent with
the same bound: a product of sufficiently many elements becomes
zero after multiplication by a positive integer, and $I$ is
torsion-free.
For a proposed lift of a finite-group representation, only
finitely many entries occur in its multiplication defects.
Put them in such a $J_{\Q}$.  The matrices then define a
representation over $A/J_{\Q}$, and the nilpotent case lifts it
to $A$.  For uniqueness, put the finitely many entries of the
difference of two lifts, after a lifted change of basis, into
such an ideal.  The same argument works in the block upper
triangular groups and lifts every exact-sequence relation.
This proves the general nil case.

Reduction commutes with all three stated operations.  Applying
the result to every symmetric group and its induction product
therefore gives the final assertion.
\end{proof}

\begin{example}\label{ex:noncoherent-positive}
For a prime $p$ let
\[
 A=\Z_{(p)}\oplus\Q\varepsilon,\qquad\varepsilon^2=0,
\]
with the usual scalar action of $\Z_{(p)}$ on $\Q$.
This is a local ring with maximal ideal
$p\Z_{(p)}\oplus\Q\varepsilon$.  Its square-zero ideal
$\Q\varepsilon$ is uniquely divisible, and its quotient is
$\Z_{(p)}$.  Hence its induction ring is polynomial on the
trivial classes.  It is not coherent: the annihilator of
$\varepsilon$ is $\Q\varepsilon$, which is not finitely
generated as an $A$-module, whereas annihilators of individual
elements are finitely generated in a coherent ring.
This example also shows why positive coefficient rings cannot
be characterized simply by coherence and regularity.
\end{example}

% Self-contained proof module; needs no modular representation theory.
\section{An obstruction in every positive characteristic}

An exterior-power identity makes characteristic zero a necessary
condition, including for coefficient rings of characteristic $p^a$.

\begin{proposition}\label{prop:positive-characteristic}
Let $A$ be a nonzero commutative ring, and suppose that $m\cdot1_A=0$
for some integer $m\geq2$.  Let $P_m$ be the $m$th power-sum
symmetric function, regarded as an integral polynomial in the
complete symmetric functions $h_1,\ldots,h_m$.  Then
\[
 P_m(\rho_1,\ldots,\rho_m)=0
 \quad\text{in }K_0(\operatorname{Lat}_A(S_m)).
\]
In particular, the trivial classes cannot freely generate over any
nonzero local ring of positive characteristic.
\end{proposition}

\begin{proof}
For each $r\geq0$, write $\epsilon_r$ for the sign lattice of $S_r$,
with $\epsilon_0=\rho_0=1$.  For the permutation lattice
$A^r$, contraction by the sum of coordinates makes its exterior
algebra an exact complex
\[
 0\longrightarrow\bigwedge^r A^r\longrightarrow\cdots
 \longrightarrow\bigwedge^1 A^r\longrightarrow A
 \longrightarrow0.
\]
This complex is $S_r$-equivariant and split exact over $A$: contraction
has a contracting homotopy given by exterior multiplication by a
basis vector.  Its cycles are finite free, as can also be seen by
splitting $A^r$ into the kernel of the sum map and that basis line.
Since
\[
 \bigwedge^j A^r\simeq
 \operatorname{Ind}_{S_j\times S_{r-j}}^{S_r}
 (\epsilon_j\boxtimes\rho_{r-j}),
\]
its Euler relation is
\[
 \sum_{j=0}^r(-1)^j\rho_{r-j}\epsilon_j=0
 \qquad(r>0).
\]
Thus, if $H(t)=\sum_{r\geq0}\rho_rt^r$ and
$E(t)=\sum_{r\geq0}\epsilon_rt^r$, then $H(t)E(-t)=1$.
In particular $\epsilon_j$ is the usual elementary symmetric
polynomial in the $\rho_i$.

Now set $U=A^m$ with its permutation action, and let
$V=\ker(U\xrightarrow{\sum} A)$.  This is free of rank $m-1$,
with basis $e_i-e_m$ for $1\leq i<m$.  The invariant vector
$v=e_1+\cdots+e_m$ belongs to $V$ because $m=0$ in $A$.
Moreover,
\[
 v=\sum_{i=1}^{m-1}(e_i-e_m),
\]
so it is primitive and $W=V/Av$ is finite free of rank $m-2$.
The exterior-power filtration of
$0\to Av\to V\to W\to0$ gives
\[
 [\textstyle\bigwedge^i V]
 =[\textstyle\bigwedge^i W]+[\textstyle\bigwedge^{i-1}W].
\]
With exterior powers outside the possible range set to zero,
alternation gives
\begin{equation}\label{eq:augmentation-exterior-zero}
 \sum_{i=0}^{m-1}(-1)^i[\textstyle\bigwedge^iV]=0.
\end{equation}
On the other hand, $0\to V\to U\to A\to0$ yields
\[
 [\textstyle\bigwedge^iV]
 =\sum_{j=0}^i(-1)^{i-j}[\textstyle\bigwedge^jU].
\]
Substitution into~\eqref{eq:augmentation-exterior-zero} gives
\[
 0=\sum_{j=0}^{m-1}(m-j)(-1)^j
       \rho_{m-j}\epsilon_j.
\]
The right side is the coefficient of $t^m$ in
$tH'(t)E(-t)=tH'(t)/H(t)$, hence is
$P_m(\rho_1,\ldots,\rho_m)$ by the Newton identity.  The
polynomial $P_m$ is nonzero over $\mathbb Z$: its coefficient of
$h_m$ is $m$.  This proves the obstruction.
\end{proof}

\begin{example}
In characteristic $2$, the identity is
$2\rho_2-\rho_1^2=0$.
In characteristic $4$, the argument gives the nonzero polynomial
relation
\[
 4\rho_4-4\rho_1\rho_3-2\rho_2^2
       +4\rho_1^2\rho_2-\rho_1^4=0.
\]
The second example does not require a computation of the lattice
Grothendieck group of $S_2$ over $\mathbb Z/4\mathbb Z$.
\end{example}


\section{Determinants and traces}

The following invariants distinguish generation from independence.
For an element $g\in G$, trace is additive on lattice exact
sequences, and determinant is multiplicative on them: choose an
$A$-basis adapted to the underlying split sequence, and use its
block triangular action matrix.  Consequently
\[
 \det:K_0(\Lat_A(G))\longrightarrow\operatorname{Hom}(G,A^\times)
\]
is a homomorphism, with the target written multiplicatively.

\begin{proposition}\label{prop:determinant}
If $\Phi_A$ is surjective, every unit $u\in A$ with $u^2=1$
equals $1$ or $-1$.  The ring
$\Z_{(2)}[\varepsilon]/(2\varepsilon,\varepsilon^2)$ violates this
condition and does not have the polynomial-generation property.
\end{proposition}
\begin{proof}
The rank-one $S_2$-lattice with transposition acting by $u$ has
determinant $u$.  The two degree-two proposed generators are the
trivial lattice $\rho_2$ and the regular permutation lattice
$\rho_1^2$; their determinants are $1$ and $-1$.  Any integral
combination has determinant $\pm1$.
In the indicated local ring, $1+\varepsilon$ squares to one and
is neither $1$ nor $-1$.  Its nonzero nilpotent part proves the
first inequality; reduction modulo $\varepsilon$ proves the
second.  The ring has characteristic zero, maximal ideal
$(2,\varepsilon)$, and $K_0(A)=\Z$ because finite projectives over
a local ring are free.  It therefore also answers the
parenthetical question in the manuscript negatively.
\end{proof}

\begin{proposition}\label{prop:nilpotent-trace}
For every prime $p$, the characteristic-zero local ring
\[
 A_p=\Z_{(p)}[\varepsilon]/(p\varepsilon,\varepsilon^p)
\]
does not have the polynomial-generation property.  For odd $p$,
its only square roots of one are $\pm1$.
\end{proposition}
\begin{proof}
The additive group is
$\Z_{(p)}\oplus\bigoplus_{j=1}^{p-1}\F_p\varepsilon^j$,
and the maximal ideal is $(p,\varepsilon)$.  The unit
$u=1+\varepsilon$ satisfies $u^p=1$, by the binomial theorem.
Let $H=\langle g\rangle$ be generated by a $p$-cycle in $S_p$,
and induce the rank-one $A_p[H]$-lattice with action $u$ to $S_p$.
In a coset basis only cosets in $N_{S_p}(H)/H$ contribute to the
trace of $g$.  They give each nonzero exponent modulo $p$ once.
Thus
\[
 \tr(g\mid\operatorname{Ind}_H^{S_p}L_u)
  =\sum_{a=1}^{p-1}(1+\varepsilon)^a
  =p-1+\varepsilon^{p-1}.
\]
For $1\leq j\leq p-1$, the coefficient of $\varepsilon^j$
in the sum is $\binom p{j+1}$.  It vanishes in $A_p$ if
$j<p-1$ and equals one if $j=p-1$.  The trace therefore lies
outside the image of $\Z$.  Young permutation traces are integers,
so additivity of trace proves non-generation.
For odd $p$, two is invertible.  In a local ring the two factors
of $(v-1)(v+1)=0$ cannot both be nonunits if their difference is
a unit; hence $v=\pm1$.
\end{proof}

\section{Reduced counterexamples in mixed characteristic}

Gluing two characteristic-zero branches along a positive-characteristic
point can obstruct generation, even when the coefficient ring is
reduced and torsion-free over the integers.

\begin{proposition}\label{prop:reduced-fiber-product}
Let $p$ be a prime, let $\zeta_p$ be a primitive $p$th root of unity,
and put $O=\mathbb Z_{(p)}[\zeta_p]$.  Use the residue maps
$O\to\mathbb F_p$ and $\mathbb Z_{(p)}\to\mathbb F_p$ to form
\[
 A=O\times_{\mathbb F_p}\mathbb Z_{(p)}.
\]
Then $A$ is a noetherian reduced local ring of characteristic zero,
flat over $\mathbb Z_{(p)}$, and the trivial symmetric-group classes
do not generate its induction ring.  If $p$ is odd, every square
root of $1$ in $A$ is $1$ or $-1$.
\end{proposition}

\begin{proof}
The ring $O$ is local with maximal ideal $(\zeta_p-1)$ and residue
field $\mathbb F_p$: the cyclotomic polynomial becomes
$(X-1)^{p-1}$ modulo $p$, and $O$ is finite over $\mathbb Z_{(p)}$.
The fibre product is finite over $\mathbb Z_{(p)}$, since it is the
kernel of the difference of residue maps from
$O\oplus\mathbb Z_{(p)}$ to $\mathbb F_p$.  An element of $A$ is a
unit precisely when its common residue is nonzero.  This proves
noetherianity and locality.  As a subring of a product of
characteristic-zero domains, $A$ is reduced and torsion-free over
$\mathbb Z_{(p)}$; torsion-free modules over a discrete valuation
ring are flat.

The unit $u=(\zeta_p,1)$ belongs to $A$ and satisfies $u^p=1$.
Let $H=\langle g\rangle\subseteq S_p$ be generated by a $p$-cycle,
and let $L_u$ be the rank-one $A[H]$-lattice on which $g$ acts by
$u$.  Its induced lattice
$N=\operatorname{Ind}_H^{S_p}L_u$ is finite free over $A$.
All the nonidentity elements of $H$ are conjugate in $S_p$, and
$|C_{S_p}(g)|=p=|H|$.  The induced trace formula, which follows
directly from the coset basis of $N$, therefore gives
\[
 \operatorname{tr}(g\mid N)=\sum_{a=1}^{p-1}u^a=(-1,p-1).
\]
This element is not the image of an integer in $A$.  On every
Young permutation lattice the trace of $g$ is the integer counting
the fixed cosets.  Trace is additive on exact sequences of lattices,
as those sequences split over $A$ and the matrices become block
triangular.  It follows that $[N]$ is outside the subgroup generated
by all Young permutation classes, which is the degree-$p$ part of
the subring generated by the $\rho_i$.

If $p$ is odd, $2$ is a unit in the local ring $A$.  For any
$a\in A$ with $a^2=1$, the factors $a-1$ and $a+1$ generate the
unit ideal.  At least one is a unit, and their product is zero;
therefore the other is zero.  Hence $a=1$ or $a=-1$.
\end{proof}

\section{Normal local domains in every residue characteristic}

Generation can fail over a two-dimensional normal local domain that
is flat over the integers.  The obstruction is an invertible module
of order two on the characteristic-zero fibre, detected as the
invariant part of an induced representation.

\begin{theorem}\label{thm:normal-domain-counterexample}
\label{thm:normal-surface-counterexample}
For every prime $p$ there is a two-dimensional noetherian normal local
domain $A$ of characteristic zero and residue characteristic $p$ such
that the trivial symmetric-group classes do not generate its induction
ring.  The ring $A$ is flat over $\Z_{(p)}$, and its only square roots
of $1$ are $1$ and $-1$.
\end{theorem}

\begin{proof}
Choose a primitive $p$th root of unity $\zeta$, and put
\[
 O=\Z_{(p)}[\zeta],\qquad \pi=\zeta-1.
\]
The ring $O$ is a discrete valuation ring with uniformizer $\pi$ and
residue field $k=\F_p$: the polynomial $\Phi_p(1+X)$ is Eisenstein at
$p$, and $p$ is a unit times $\pi^{p-1}$.
Choose $u\in O^\times$ and $v\in O$ such that
\[
 z^2+\overline v z+\overline u\in k[z]
\]
is irreducible.  For odd $p$, take $v=0$ and $-\overline u$ a nonsquare;
for $p=2$, take $u=v=1$.  Set
\begin{equation}\label{eq:normal-surface-ring}
 A=\left(O[x,b]/(x(\pi-x)-b(vx+ub))\right)_{(\pi,x,b)}.
\end{equation}
We simultaneously verify the complete construction
\begin{equation}\label{eq:complete-normal-surface-ring}
 \widehat A=\widehat O[[x,b]]/(x(\pi-x)-b(vx+ub)).
\end{equation}
In the argument below, $A$ denotes either ring and $O$ is replaced by
$\widehat O$ in the complete case.

Write $f=x(\pi-x)-b(vx+ub)$.  In either regular local ambient ring,
the associated graded quotient is
\[
 k[\Pi,X,B]/(X\Pi-X^2-\overline v XB-\overline u B^2).
\]
The displayed polynomial is primitive and linear in $\Pi$, because
$X$ is relatively prime to $X^2+\overline v XB+\overline u B^2$.
It is therefore irreducible.  The associated graded ring is a domain,
so separatedness of the maximal-ideal filtration proves that $A$ is
a domain.  It is a Cohen--Macaulay hypersurface of dimension two and
embedding dimension three.
Its quotient by $\pi$ is also a domain: the binary quadratic
$X^2+\overline v XB+\overline u B^2$ is irreducible by the choice of
$u,v$, and the associated graded argument also treats the formal
power-series quotient.  Thus $(\pi)$ is a single height-one prime,
whose localization is regular because its maximal ideal is principal.

Put $C=A[1/\pi]$.  This ring is regular.  Indeed,
\[
 f_x=\pi-2x-vb,\qquad f_b=-vx-2ub,
 \qquad xf_x+bf_b=-\pi x\pmod f.
\]
If a prime not containing $\pi$ contained both partial derivatives,
it would contain $x$, then $b$ by $f=0$ and $u$ a unit, and then
$\pi$ by $f_x$, a contradiction.  Hence one partial derivative is a
unit at every such localization.  This proves regularity of the
hypersurface localization: a derivation sends the square of a prime
into that prime, so $f$ does not belong to the square of the localized
maximal ideal in the regular ambient ring.  The same argument uses
formal partial derivatives for~\eqref{eq:complete-normal-surface-ring}.
All height-one localizations of $A$ are therefore regular, and Serre's
criterion proves normality \citep[Lemma 10.157.4, Tag 031S]{Stacks2026}.  Evaluation at $x=b=0$ gives a retraction
onto $O$, so $A$ has characteristic zero.  Its torsion-freeness over
$\Z_{(p)}$ gives flatness.  In a characteristic-zero domain,
$(w-1)(w+1)=0$ forces $w=1$ or $w=-1$.

Consider $I=(x,b)$.  This is a height-one prime, since $A/I=O$.
It is nonprincipal: $x,b$ are independent in the cotangent space
$\mathfrak m/\mathfrak m^2$, and hence in $I/\mathfrak m I$, so two
generators are necessary.  We claim that $I_C$ is still nonprincipal.
If it had a generator, clearing powers of $\pi$ would give a generator
$r\in A$.  Strip off the largest power of $\pi$ dividing $r$, which
exists by the Krull intersection theorem
\citep[Lemma 10.51.4, Tag 00IP]{Stacks2026}, and write the result as
$r_0\notin\pi A$.  The ideal $(r_0)$ is $\pi$-saturated: if
$\pi s=r_0t$, primality of $(\pi)$ gives $t=\pi t_1$, and cancellation
gives $s=r_0t_1$.  Iteration treats every power of $\pi$.
The prime $I$ is also $\pi$-saturated, because $\pi\notin I$.
Contracting $I_C=r_0C$ to $A$ would therefore give $I=(r_0)$,
a contradiction.

The defining equation gives
\[
 I_C^2=x\bigl(x,b,u^{-1}(\pi-x-vb)\bigr)C=xC.
\]
The last equality uses that $u,\pi$ are units in $C$.  Thus $I_C$
is invertible and its Picard class has exact order two.
Put $J=(b,\pi-x)C$.  Similarly,
\[
 I_CJ=b(b,x,\pi-x,vx+ub)C=bC.
\]
Hence $J$ is a nontrivial invertible module, also of order two.

On $L=A^2$ define
\[
 M=\begin{pmatrix}x&b\\vx+ub&\pi-x\end{pmatrix},
 \qquad T=I+M.
\]
The trace and determinant give $M^2=\pi M$.  Therefore
\[
 T^p
 =I+\left(\sum_{j=1}^{p}\binom pj\pi^{j-1}\right)M=I,
\]
because the coefficient times $\pi$ is $(1+\pi)^p-1=0$ and $O$ is
a domain.  Let a generator of $H=C_p$ act by $T$, embed $H$ in $S_p$
as the subgroup generated by a $p$-cycle, and set
\[
 N=\operatorname{Ind}_{H}^{S_p}L.
\]
This is a finite free $A$-lattice.

Over $C$, the $H$-invariant part of $L$ is the image of the idempotent
\[
 I-\pi^{-1}M
 =\pi^{-1}\begin{pmatrix}\pi-x&-b\\-vx-ub&x\end{pmatrix}.
\]
Projection to its first coordinate identifies this image with $J$.
The image ideal is $(\pi-x,b)C$; injectivity follows because an
invariant vector $(0,z)$ satisfies $bz=0$, with $C$ a domain and
$b\ne0$.  Induction for finite groups is also coinduction, so
\[
 N_C^{S_p}\simeq L_C^H\simeq J.
\]
Every positive integer is invertible in $C$.  The averaging projector
therefore makes invariants an exact functor to finite projective
$C$-modules.  Extension of scalars, invariants, and determinant give
an additive homomorphism
\[
 K_0(\Lat_A(S_p))\longrightarrow K_0(C)
                       \xrightarrow{\det}\operatorname{Pic}(C).
\]
For a Young subgroup $S_\lambda$, the invariants of
$C[S_p/S_\lambda]$ form a free line, since invariant coefficients are
constant on the transitive coset set.  The homomorphism vanishes on
every degree-$p$ monomial in the trivial classes, but its value on
$[N]$ is $[J]\ne0$.  This proves failure of generation for both rings.
\end{proof}

\begin{example}
For $p=2$ one may take $u=v=1$, $\zeta=-1$, and $\pi=-2$.  The example
is
\[
 A=\left(\Z_{(2)}[x,b]/(x(x+2)+xb+b^2)\right)_{(2,x,b)},
 \qquad
 T=\begin{pmatrix}1+x&b\\x+b&-1-x\end{pmatrix}.
\]
Thus a rank-two involution already detects the obstruction for $S_2$.
\end{example}

\begin{corollary}\label{cor:complete-normal-domain-counterexample}
The counterexamples in Theorem~\ref{thm:normal-domain-counterexample}
can be chosen complete and two-dimensional.
\end{corollary}

\begin{proof}
Use~\eqref{eq:complete-normal-surface-ring}.  The preceding proof
verified its domain property, normality, surviving nonprincipal
invertible module, and induced representation together with those of
the polynomial-local example.
\end{proof}

\begin{corollary}\label{cor:minimal-normal-dimension}
Dimension two is the least possible dimension of a noetherian normal
local characteristic-zero counterexample to surjectivity of $\Phi$.
\end{corollary}

\begin{proof}
A zero-dimensional normal local domain is a field, and a
one-dimensional noetherian normal local domain is a discrete valuation
ring \citep[Lemma 10.119.7, Tag 00PD]{Stacks2026}.  In characteristic zero both satisfy the regular-local positive
theorem.  Theorem~\ref{thm:normal-domain-counterexample} supplies
counterexamples in dimension two.
\end{proof}

\section{Flat square-zero obstructions at every prime}
\label{sec:dual-numbers-all-primes}

A flat square-zero thickening can destroy generation in every residue
characteristic.  The obstruction is a boundary on periodic cohomology;
its square has primary components with additive lengths modulo two.

\begin{theorem}
\label{thm:dual-numbers-all-primes}
For every prime $p$, put
\[
 O_p=\mathbb Z_{(p)}\text{ or }\mathbb Z_p,\qquad
 A_p^{\mathrm{dual}}=O_p[\epsilon]/(\epsilon^2).
\]
There is an explicit $A_p^{\mathrm{dual}}[S_p]$-lattice $M_p$ of rank $p!$
whose class is not an integer combination of Young permutation classes.
Moreover
\[
 [M_p]-[A_p^{\mathrm{dual}}[S_p]]
\]
reduces to zero over $O_p$ and is detected by an additive
$\mathbb Z/2$-valued invariant.  In particular the polynomial map generated
by the trivial symmetric-group classes is not surjective over
$A_p^{\mathrm{dual}}$.
\end{theorem}

We give all details, including the effect of changing a generator of the
cyclic subgroup.  Proposition~\ref{prop:truncation-torsion} below
proves that the displayed class has exact order two for every prime;
Corollary~\ref{cor:all-prime-two-torsion} gives an infinite split family.

\subsection{Primary-length invariants for the cyclic group of order \texorpdfstring{$p$}{p}}

Fix $H=C_p=\langle s\rangle$, write $O=O_p$ and $A=A_p^{\mathrm{dual}}$,
and let $N$ be an $O[H]$-lattice with action matrix $S$.  Set
\[
 \mathcal N(S)=I+S+\cdots+S^{p-1}.
\]
Use the two-periodic complex with even differential $S-I$ and odd
differential $\mathcal N(S)$.  Its cohomology is
\[
 H^0(N)=\frac{\ker(S-I)}{\mathcal N(S)N},\qquad
 H^1(N)=\frac{\ker\mathcal N(S)}{(S-I)N}.
\]
Both groups are finite dimensional over $\mathbb F_p$.  For the even
group, this follows from $\mathcal N(S)v=pv$ on invariants.  For the odd
group use
\begin{equation}
\label{eq:dual-prime-D-identity}
 (S-I)D_p(S)=\mathcal N(S)-pI,
 \qquad
 D_p(t)=\sum_{j=0}^{p-2}(p-1-j)t^j.
\end{equation}
If $\mathcal N(S)v=0$, this identity expresses $pv$ as a boundary.

For an $A[H]$-lattice $L$, the canonical extension
\[
 0\longrightarrow L_0\xrightarrow{\epsilon}L\longrightarrow L_0
 \longrightarrow0,\qquad L_0=L/\epsilon L,
\]
gives a degree-one boundary $\delta_L$ on this periodic cohomology and an
even endomorphism $u_L=\delta_L^2$.  For every monic irreducible
$f\in\mathbb F_p[t]$ with $f\ne t$, define
\[
 \nu_f(L)=
 \ell_{\mathbb F_p[t]}
 \bigl(H^0(L_0)_{f\text{-primary}},\ t\text{ acting as }u_L\bigr)
 \pmod2.
\]
We prove additivity directly, including the sign issue when $p$ is odd.

For an admissible exact sequence $0\to L'\to L\to L''\to0$, choose an
$A$-linear splitting.  In each cochain degree write its differential as
\[
 d_L^j=\begin{pmatrix}d_{L'}^j&B^j\\0&d_{L''}^j\end{pmatrix},
\]
and expand every block into its constant and $\epsilon$ coefficients.
The coefficient of $\epsilon$ in the upper-right block of
$d_L^{j+1}d_L^j=0$ is
\begin{align*}
 d_{L',1}^{j+1}B_0^j+B_0^{j+1}d_{L'',1}^j
 ={}&-d_{L',0}^{j+1}B_1^j-B_1^{j+1}d_{L'',0}^j.
\end{align*}
On a degree-$j$ cocycle for $L''_0$, the last term is zero and the first
term on the right is a boundary.  The left side represents
\[
 \delta_{L'}^{j+1}\partial_j+
 \partial_{j+1}\delta_{L''}^j.
\]
Thus the extension boundaries anticommute with the long-exact-sequence
boundaries, and all maps commute with $u=\delta^2$.

Taking an $f$-primary component is exact, as is seen by a common finite
Bezout projector.  Since $f\ne t$, $u$ is invertible there, and $\delta$
identifies the two parities by $\mathbb F_p[t]$-linear isomorphisms.  The
two boundary images in the resulting exact six-term cycle have equal
composition length.  The minus sign in their intertwining identity is a
unit and does not affect that equality.  If $a,b,c$ are the three
degree-zero lengths and $r_0,r_1$ are the two boundary-image lengths, then
\[
 a-b+c=r_0+r_1=2r_0.
\]
This proves that $\nu_f$ is a homomorphism on
$K_0(\operatorname{Lat}_A(H))$.  It vanishes on every lattice obtained by
scalar extension from $O$, because its canonical $\epsilon$ extension then
splits as an $O[H]$-extension and $\delta=0$.

\begin{proposition}\label{prop:dual-bockstein-invariants}
For every monic irreducible $f\in\mathbb F_p[t]$ other than $t$, the
primary-length formula above defines a homomorphism
\[
 \nu_f:K_0(\operatorname{Lat}_{O[\epsilon]/(\epsilon^2)}(C_p))
       \longrightarrow\mathbb Z/2.
\]
It vanishes on scalar extensions from $O$.  For any fixed class only
finitely many $\nu_f$ are nonzero.
\end{proposition}
The final assertion follows from finite dimension of the cohomology and
then from finite support of a finite integer combination of lattice classes.


\subsection{An explicit rank-\texorpdfstring{$p$}{p} cyclic lattice}

Put $r=p-1$.  Let $C$ be the integral companion matrix of
$\Phi_p(t)=1+t+\cdots+t^{p-1}$, with basis $e_1,\ldots,e_r$ and action
\[
 Ce_i=e_{i+1}\ (i<r),\qquad
 Ce_r=-\sum_{i=1}^r e_i.
\]
Thus $C^p=I$ and $\mathcal N(C)=0$.  Set
\[
 b=e_1,\qquad a=e_r^{\mathsf T},\qquad
 T_p=\begin{pmatrix}1&\epsilon a\\\epsilon b&C\end{pmatrix}.
\]
The diagonal blocks of $T_p^p$ are $1$ and $C^p$, since every off-diagonal
round trip contains $\epsilon^2$.  Its off-diagonal blocks are respectively
\[
 \epsilon a\mathcal N(C),\qquad
 \epsilon\mathcal N(C)b,
\]
and therefore vanish.  Hence $T_p^p=I$, giving an $A[H]$-lattice
$L_p=A^p$.

The reduced representation is $O\oplus I$, where the companion module $I$
is isomorphic to the augmentation ideal of $O[H]$: identify $e_i$ with
$s^{i-1}(s-1)$ for $1\le i\le p-1$.  The map $C-I$ is injective over $O$
and has cokernel $\mathbb F_p$, with the common class of all $e_i$ as a
generator.  Therefore
\[
 H^0((L_p)_0)=\mathbb F_p\quad\text{from the trivial summand},
 \qquad
 H^1((L_p)_0)=\mathbb F_p\quad\text{from }I.
\]

The even extension boundary sends the trivial generator to the class of
$b=e_1$, hence is the identity in these coordinates.  For the odd boundary,
the upper-right $\epsilon$ coefficient of $\mathcal N(T_p)$ is
$aD_p(C)$.  Since $C^jb=e_{j+1}$ for $0\le j\le p-2$,
\[
 aD_p(C)b=1.
\]
Thus the odd boundary is also the identity.  We conclude that
\begin{equation}
\label{eq:dual-prime-u-one}
 u_{L_p}=1\quad\text{on both cohomology spaces},\qquad
 \nu_{t-1}(L_p)=1.
\end{equation}

\subsection{Changing the generator: exact chain maps and the scaling factor}

For $1\le q\le p-1$, let $L^{(q)}$ denote the twist in which $s$ acts as
$T^q$ rather than $T$.  Put
\[
 S_q(T)=I+T+\cdots+T^{q-1}.
\]
There is a cochain map from the periodic complex for $T$ to the periodic
complex for $T^q$, in the following direction and degrees:
\begin{equation}
\label{eq:dual-prime-comparison}
 F_q^{2j}=q^j I,\qquad
 F_q^{2j+1}=q^j S_q(T),\qquad j\in\mathbb Z.
\end{equation}
Negative powers of $q$ are allowed because $q\in O^\times$.  The two chain
identities are
\[
 T^q-I=S_q(T)(T-I),\qquad
 \mathcal N(T^q)S_q(T)=\mathcal N(T)S_q(T)=q\mathcal N(T).
\]
Here $\mathcal N(T^q)=\mathcal N(T)$ because multiplication by $q$ permutes
the residues modulo $p$.

The same maps apply before and after reduction modulo $\epsilon$ and
commute with the canonical extension.  On degree-zero cohomology,
$F_q^0=I$ identifies the two groups: $T$ and $T^q$ have the same invariants
and the same norm image.  On degree two, however, the map is $F_q^2=qI$.
Naturality of the extension boundaries gives
\[
 \delta_{L^{(q)}}^0F_q^0=F_q^1\delta_L^0,
 \qquad
 \delta_{L^{(q)}}^1F_q^1=F_q^2\delta_L^1.
\]
Composing and using the canonical period-two identification yields
\begin{equation}
\label{eq:dual-prime-scaling}
 u_{L^{(q)}}=q\,u_L
 \quad\text{on the commonly identified }H^0.
\end{equation}
This is why it would be incorrect to treat the periodicity identification
as invariant under the generator change without its scalar factor.

For $L=L_p$, equations \eqref{eq:dual-prime-u-one} and
\eqref{eq:dual-prime-scaling} give
\begin{equation}
\label{eq:dual-prime-twist-detector}
 \nu_{t-1}(L_p^{(q)})=
 \begin{cases}1,&q=1\text{ in }\mathbb F_p,\\0,&q\ne1.
 \end{cases}
\end{equation}

\subsection{Induction to the symmetric group and the Mackey multiplicities}

Embed $H$ in $S_p$ as the subgroup generated by a $p$-cycle and define
\[
 M_p=\operatorname{Ind}_H^{S_p}L_p.
\]
Its $A$-rank is $[S_p:H]p=p!$.  To distinguish its class, apply
$\nu_{t-1}$ after restriction back to $H$.

For a double coset in $H\backslash S_p/H$, the intersection
$H\cap xHx^{-1}$ is either $H$ or the trivial group, since $H$ has prime
order.  The first case occurs exactly when $x\in N_{S_p}(H)$.  For every
other double coset, the Mackey summand is induced from the trivial
subgroup, hence is free over $A[H]$.  Its reduced periodic cohomology is
zero, so it contributes zero to every $\nu_f$.

Inside the normalizer, the double cosets are exactly its cosets modulo
$H$.  The centralizer of a $p$-cycle in $S_p$ is $H$, and every automorphism
$s\mapsto s^q$, $q\in\mathbb F_p^\times$, is realized by a permutation of
the $p$ cyclically labeled letters.  Thus
\[
 N_{S_p}(H)/H\cong\mathbb F_p^\times.
\]
The remaining Mackey summands are precisely one copy of each
$L_p^{(q)}$, with $q$ running through the nonzero residues.  Changing a
Mackey convention replaces $q$ by $q^{-1}$ and does not change either the
multiplicity statement or the unique occurrence of $q=1$.  Consequently
\[
 \nu_{t-1}\!\left(\operatorname{Res}^{S_p}_H M_p\right)
 =\sum_{q\in\mathbb F_p^\times}\nu_{t-1}(L_p^{(q)})=1.
\]

Every Young permutation lattice over $A$ is obtained by scalar extension
from its integral $O$-model.  Its restriction to $H$ is also such a scalar
extension, so its invariant is zero.  Therefore $[M_p]$ is not in their
integer span.

Finally, $[(L_p)_0]=[O]+[I]=[O[H]]$ by the augmentation exact sequence.
Induction gives $[(M_p)_0]=[O[S_p]]$.  After inverting $p$, the matrix
$C-I$ is invertible and a change of basis $I+\epsilon X$, with
$X_{12}=a(C-I)^{-1}$ and $X_{21}=(I-C)^{-1}b$, identifies $L_p$ with its
constant reduction.  That reduction over $\operatorname{Frac}(O)$ is the regular
cyclic representation.  Thus the same equality holds after inverting
$p$, even with $\epsilon$ retained.  The class
$[M_p]-[A[S_p]]$ is thus in the asserted reduction kernel and is nonzero by
the displayed invariant.  This completes the proof of
Theorem~\ref{thm:dual-numbers-all-primes}.

\section{An exact relation for cyclic deformations}
\label{sec:uniform-torsion}

An explicit exact sequence annihilates a cyclic deformation class by
the truncation exponent.  The next section constructs an invariant
that proves this bound is exact for every exponent.

\begin{proposition}\label{prop:truncation-torsion}
Let $p$ be a prime, $O=\Z_{(p)}$ or $\Z_p$, $n\geq2$, and
$A=O[e]/(e^n)$.
Let $C$ be the companion of $\Phi_p(z)$ in the basis
$1,z,\ldots,z^{p-2}$, let $a=e_{p-1}^{\mathsf T}$ and $b=e_1$, and put
\[
 T_j=\begin{pmatrix}1&e^j a\\e^{n-j}b&C\end{pmatrix},
 \qquad L_j=(A^p,T_j),\qquad 0\leq j\leq n.
\]
If $x_j=[L_j]-[A[C_p]]$, then $x_j=jx_1$ and $nx_1=0$.
In particular, the flat dual-number obstruction of
Theorem~\ref{thm:dual-numbers-all-primes} has exact order two for every
prime $p$, both before and after induction to $S_p$.
\end{proposition}

\begin{proof}
The zero cross products and $\Phi_p(C)=0$ give $T_j^p=I$.
The endpoint lattices are regular.  Indeed, multiplication by a
generator $g$ on $A[C_p]$ has matrix $T_0$ in the basis
\[
 1+g+\cdots+g^{p-1},\quad 1,g,\ldots,g^{p-2},
\]
and matrix $T_n$ in the basis
\[
 1,\quad g-1,(g-1)g,\ldots,(g-1)g^{p-2}.
\]
Both lists are bases by triangular comparison with the powers of $g$.

For $0\leq j<n$, define an equivariant injection
\[
 L_j\longrightarrow L_{j+1}\oplus L_0,\qquad
 (v_0,v_c)\longmapsto
 \bigl((ev_0,v_c),(v_0,e^jv_c)\bigr).
\]
Its two blocks are $\operatorname{diag}(e,I_{p-1})$ and
$\operatorname{diag}(1,e^jI_{p-1})$; direct multiplication verifies
equivariance.  It is split over $A$.  Quotient coordinates are
\[
 q_0=x_0-ey_0,\qquad q_c=y_c-e^jx_c,
\]
where $x$ and $y$ denote the two target summands.  In these coordinates
the induced action is
\[
 \begin{pmatrix}1&-e a\\-e^{n-1}b&C\end{pmatrix},
\]
which is conjugate to $T_1$ by $\operatorname{diag}(1,-I_{p-1})$.
Thus there is an admissible exact sequence
\[
 0\longrightarrow L_j\longrightarrow L_{j+1}\oplus L_0
   \longrightarrow L_1\longrightarrow0.
\]
It gives $x_{j+1}=x_j+x_1$.  Since $x_0=x_n=0$, the two assertions
follow.  For $n=2$, the previously constructed detector is nonzero
on $x_1$ and on its induction to $S_p$, whereas the displayed relation
kills both classes by two.  They therefore have exact order two.
\end{proof}

Theorem~\ref{thm:all-truncations} proves that this annihilation bound
is exact, both for the cyclic classes and after induction.

\section{Every truncated polynomial thickening}
\label{sec:all-truncations}

The truncation exponent itself occurs as an exact torsion order.
On the periodic Tate complex, the curvature below is the
lifted-differential operator of
\citet[Section~1, p.~39; Proposition~1.3 and Corollary~1.4,
pp.~40--41]{Eisenbud1980}.  We require its realization as a stable
lattice endomorphism and the additive covolume that follows.
The invariant takes primary components of this curvature
endomorphism.  Their cyclic-group cohomology vanishes, so their
invariant lattices have an additive covolume, well-defined modulo
the truncation exponent.

\begin{theorem}\label{thm:all-truncations}
Let $p$ be a prime, $n\ge2$, $O=\Z_p$ or $\Z_{(p)}$, and
$A=O[e]/(e^n)$.  The degree-$p$ reduction kernel
\[
 \ker\bigl(K_0(\Lat_A(S_p))\longrightarrow
                 K_0(\Lat_O(S_p))\bigr)
\]
contains a split countably infinite direct sum of copies of $\Z/n$.
The displayed generators also vanish after inverting $p$.
In particular $\Phi_A$ is not surjective.

Let $C$ be the companion of $\Phi_p(z)$ in the monomial basis,
$a=e_{p-1}^{\mathsf T}$, and $b=e_1$.
Choose monic irreducibles $f_i\in\F_p[z]$, different from $z$,
of pairwise distinct degrees $r_i$ prime to $n(p-1)$, and integral
unit companion lifts $D_i$.  Explicit generators are
\[
 \begin{split}
 L_i&=\left(A^{pr_i},
 \begin{pmatrix}
 I_{r_i}&e(a\otimes I_{r_i})\\
 e^{n-1}(b\otimes D_i)&C\otimes I_{r_i}
 \end{pmatrix}\right),\\
 x_i&=[\operatorname{Ind}_{C_p}^{S_p}L_i]-r_i[A[S_p]].
 \end{split}
\]
The cyclic precursors $[L_i]-r_i[A[C_p]]$ likewise have exact
order $n$ and generate a split direct sum.
\end{theorem}

We first work over the complete DVR $O=\Z_p$, with fraction field
$K=\Q_p$.  A map will be called stably zero if it factors through
a finite free $A[C_p]$-lattice.  Sums and composites of such maps
have the same property.

\subsection{Curvature and integral primary components}

Let $E$ be an $A$-free $C_p$-lattice, with generator matrix $T$.
Choose any lift $\widetilde T$ over $O[[e]]$ and put
\[
 \widetilde U=\frac{\widetilde T^p-I}{e^n},\qquad
 U_E=\widetilde U\pmod{e^n}.
\]
This is an $A[C_p]$-endomorphism: the numerator is divisible
by $e^n$, and $\widetilde U$ commutes with $\widetilde T$.

\begin{lemma}\label{lem:truncation-stable-curvature}
The stable class of $U_E$ is independent of the lift and basis,
and is natural in $E$.  The groups of maps modulo stably zero maps
have characteristic $p$.
\end{lemma}

\begin{proof}
For an equivariant map $F:E\to E'$, choose lifts and write
$\widetilde T'\widetilde F-\widetilde F\widetilde T=e^n\widetilde R$.
Telescoping the difference of $p$th powers gives
\[
 U_{E'}F-FU_E
  =\sum_{i=0}^{p-1}T'^{\,p-1-i}RT^i
  =\sum_{j=0}^{p-1}T'^j H T^{-j},\qquad H=RT^{-1}.
\]
Every such group trace factors through the regular lattice:
\[
 E\longrightarrow A[C_p]\otimes_A E_{\mathrm{underlying}}
       \longrightarrow E',
\]
where the first map sends $v$ to
$\sum_jg^j\otimes T^{-j}v$, and the second sends
$g^j\otimes w$ to $T'^jH(w)$.
Changing a lift by $e^n\widetilde R$ gives the same trace formula;
a basis change is the case of an invertible $F$.
Finally the trace of an equivariant map $F$ is $pF$, proving the
last assertion.
\end{proof}

An endomorphism $U$ of a finite $O$-free module has an integral
primary decomposition
\[
 E=\bigoplus_f E_f,
\]
indexed by the irreducible factors $f\in\F_p[z]$ of its residue
characteristic polynomial.  Each $E_f/pE_f$ is $f$-primary for $U$.
This decomposition commutes with intertwined maps and is exact on
intertwined short exact sequences.  To obtain it, lift the residue
factors to $O[z]$.  A product of suitable powers of these factors
evaluated at $U$ is divisible by $p$, so a power annihilates
$E/p^kE$ for every $k$.  The factors are pairwise comaximal modulo
$p^k$.  Their Chinese remainder projectors are compatible as $k$
varies, and completeness gives the required idempotents on $E$.
Using a common annihilating product for a sequence makes the
projectors commute with its maps; projecting a chosen preimage
proves surjectivity on each component.

For $U_E$, these projectors are $A[C_p]$-linear.  Their summands are
$A$-free, since they are finite projective over the local ring $A$.
If $f\ne z$, the restriction of $U_E$ to $E_f$ is an automorphism.

\begin{lemma}\label{lem:truncation-primary-acyclic}
Every component $E_f$ with $f\ne z$ has zero Tate cohomology as an
$O[C_p]$-lattice.
\end{lemma}

\begin{proof}
We first verify that the restricted curvature can be realized by
a lift of the summand.  Choose an $A$-basis adapted to $E_f$ and
lift the change of basis to $O[[e]]$.  After conjugation, the
off-diagonal blocks of $\widetilde T$ are divisible by $e^n$.
In a diagonal block of $\widetilde T^p$, every term leaving the
summand and returning contains two such blocks, hence is divisible
by $e^{2n}$.  Thus the diagonal block $\widetilde T_f$ satisfies
\[
 \frac{\widetilde T_f^p-I}{e^n}\pmod{e^n}=U_E|_{E_f}.
\]
This curvature lift is a unit over $O[[e]]$, because its reduction
modulo the maximal ideal is a unit.

Write $\widetilde U_f$ for it.  If $(T_f-I)x=0$ modulo $e^n$,
lift $x$ and write
$(\widetilde T_f-I)x=e^ny$.  Multiplication by
$N(\widetilde T_f)$ and cancellation of $e^n$ give
\[
 \widetilde U_f x=N(\widetilde T_f)y.
\]
Therefore $x$ is a norm boundary modulo $e^n$.
If instead $N(\widetilde T_f)x=e^ny$, multiplication by
$\widetilde T_f-I$ gives
$x=(\widetilde T_f-I)\widetilde U_f^{-1}y$.
These are the two exactness assertions for the Tate complex.
\end{proof}

\subsection{The invariant lattice and its covolume}

Suppose $M$ is an $A$-free $C_p$-lattice with zero Tate cohomology.
Its invariant module $M^{C_p}$ is finite free over $O$.
After tensoring with $K$, the projector $N/p$ makes it a projective
module over the local ring $A_K=K[e]/(e^n)$, hence free of some
rank $r$.
Choose an $A_K$-basis and identify it with $A_K^r$.  The determinant
of an $O$-basis of $M^{C_p}$, relative to
$1,e,\ldots,e^{n-1}$ in these coordinates, has an integral valuation.
Define
\[
 \iota(M)=v_p(\det M^{C_p})\pmod n.
\]
This is independent of the choices.  An $O$-basis change has unit
determinant.  If $B\in\operatorname{GL}_r(A_K)$ changes the generic
basis, it induces its constant matrix $B_0$ on each of the $n$
successive $e$-adic quotients.  Consequently
\[
 \det_K(B:A_K^r\to A_K^r)=\det(B_0)^n,
\]
so the valuation changes by a multiple of $n$.

The index is additive on admissible sequences of cohomologically
trivial lattices.  Their invariants form a short exact sequence,
because the first cohomology of the left term vanishes.  Compatible
generic bases and $O$-bases give a block triangular determinant.
Also $\iota$ vanishes on finite free $A[C_p]$-lattices, whose
invariants have norm elements as an $A$-basis.

\begin{lemma}\label{lem:truncation-stable-index}
Two cohomologically trivial lattices that are stably isomorphic
have the same index $\iota$.
\end{lemma}

\begin{proof}
Let $f:M\to N$ and $g:N\to M$ be stable inverse maps, and write
$gf-I=ba$ through a projective $P$.
Then $(f,a):M\to N\oplus P$ is a split inclusion with retraction
$(g,-b)$.  Its complementary projection is
\[
 \begin{pmatrix}I-fg&fb\\-ag&I+ab\end{pmatrix}.
\]
Every entry factors through projectives: this holds for $I-fg$
by hypothesis and for the other entries because their source or
target is $P$.  The complement is therefore projective, since its
identity factors through a projective.
The ring $A[C_p]$ is finite over the local ring $A$, and is local:
its residue algebra is
$\F_p[C_p]=\F_p[g-1]/((g-1)^p)$.
Thus the complement and $P$ are free $A[C_p]$-modules.
We have an actual isomorphism after adding such modules, and their
indices vanish.
\end{proof}

\subsection{An additive primary index}

For each irreducible $f\in\F_p[z]$, $f\ne z$, set
\[
 \eta_f(E)=\iota(E_f)\in\Z/n.
\]
Lemma~\ref{lem:truncation-primary-acyclic} ensures that this is
defined.  We verify both independence and additivity.

In the stable quotient, the curvature is natural by
Lemma~\ref{lem:truncation-stable-curvature}.  On an $f$-component,
some power of $f(U)$ is divisible by $p$ and hence is stably zero.
The division by $p$ is again an equivariant $A$-linear map, since
the lattices have no $p$-torsion.  For distinct $f,g$, Bezout's
identity over $\F_p[z]$ therefore makes every stable map between
the two components zero.  An isomorphism comparing two choices of
lifts or bases, together with its inverse, becomes block diagonal
in the stable quotient.  Its diagonal blocks are stable inverses.
Lemma~\ref{lem:truncation-stable-index} proves independence of the
resulting indices.  A component occurring only on one side is
stably projective and has index zero.

For an admissible sequence, choose an $A$-linear splitting and
lift its generator matrix in upper triangular blocks.  The actual
curvature is then upper triangular with the endpoint curvatures
on its diagonal.  It is an endomorphism of the short exact sequence,
so its primary projectors give a short exact sequence
\[
 0\longrightarrow E'_f\longrightarrow E_f
       \longrightarrow E''_f\longrightarrow0.
\]
All three terms have zero Tate cohomology.  Their indices add,
and choice independence allows these compatible lifts to be used
for every original sequence.  We have proved homomorphisms
\begin{equation}\label{eq:all-truncation-index}
 \eta_f:K_0(\Lat_A(C_p))\longrightarrow\Z/n.
\end{equation}
They vanish on constant $O$-lattices, which have constant lifts
and curvature zero.  For each virtual class only finitely many
coordinates are nonzero.

\subsection{Evaluation and exact torsion orders}

\begin{proof}[Proof of Theorem~\ref{thm:all-truncations}]
For a unit companion matrix $D$, regard the displayed generator
first as a $p$ by $p$ matrix over the commutative ring
$O[D][[e]]$.  The companion cofactor identity
$a\operatorname{adj}(zI-C)b=1$ gives characteristic polynomial
\[
 (z-1)\Phi_p(z)-e^nD=z^p-1-e^nD.
\]
Cayley--Hamilton gives the exact identity
\begin{equation}\label{eq:all-truncation-companion-curvature}
 \widetilde T_D^p=I+e^n(I_p\otimes D).
\end{equation}
In particular it defines a $C_p$-lattice after reduction, and its
actual curvature is $I_p\otimes D$.
If $\overline D$ has irreducible characteristic polynomial $f$,
the entire lattice is $f$-primary.

Over $A_K$, its invariant vectors are determined by the upper
coordinate $v_0$, since
\[
 (C-I)v_c=-e^{n-1}(b\otimes D)v_0.
\]
The cokernel of $C-I$ is $O/p$, generated by the class of $b$.
Since $e^{n-1}$ retains only the constant coefficient of $v_0$,
the solution is integral precisely when that constant coefficient
is divisible by $p$.  The unit matrix $D$ does not change this
condition.  The upper row equation follows from $e^n=0$.
Projection onto $v_0$ therefore identifies the invariant lattice
with
\[
 (p,e)A^r\subset A_K^r.
\]
Its quotient in $A^r$ has $O$-length $r$.  Thus
\[
 \eta_f(L_D)=r\pmod n,
 \qquad \eta_g(L_D)=0\quad(g\ne f).
\]

Changing the generator to its $j$th power scales the curvature by
$j$.  Indeed, choosing $\widetilde T^j$ as the lift gives
\[
 \frac{\widetilde T^{jp}-I}{e^n}
 =\widetilde U\sum_{a=0}^{j-1}(\widetilde T^p)^a
 \equiv jU\pmod{e^n}.
\]
The invariant module itself is unchanged, since the two elements
generate the same cyclic group.  If $j\overline D_i$ has the same
characteristic polynomial as $\overline D_i$, comparison of
constant terms gives $j^{r_i}=1$.  The condition
$(r_i,p-1)=1$ forces $j=1$, and distinct degrees give disjoint
scalar orbits.

Mackey restriction from $S_p$ to $C_p$ gives one copy of each
normalizer twist.  Every other double coset has trivial intersection
and gives a regular cyclic lattice.  Hence
\[
 (\eta_{f_j}\circ\operatorname{Res}_{C_p}^{S_p})(x_i)
  =\begin{cases}r_i\pmod n,&i=j,\\0,&i\ne j.\end{cases}
\]
The exact sequence of Proposition~\ref{prop:truncation-torsion}
applies with the unit block $D_i$, giving $nx_i=0$ and the same
annihilator for the cyclic precursor.  The unit coordinate values
prove exact order $n$.  Multiplication of each coordinate by
$r_i^{-1}\in\Z/n$, together with finite primary support, gives
the asserted retractions onto direct sums of copies of $\Z/n$.

Reduction to $O$ gives trivial plus cyclotomic classes, with regular
class by the augmentation sequence.  After inverting $p$, the reduction is isomorphic to the regular
representation.  Lift such an isomorphism to the free modules over
$K[e]/(e^n)$, average the map over $C_p$ using $1/p$, and apply
Nakayama's lemma.  The averaged lift is an equivariant isomorphism.
This proves both kernel assertions.
Every Young lattice is constant over $O$, whereas the coordinate
invariants on the $x_i$ are nonzero.

Finally, for $O=\Z_{(p)}$, extend coefficients to
$\Z_p[e]/(e^n)$ and compose with the invariants just constructed.
This functor is exact on admissible sequences.  The explicit matrices
and their exact annihilation relations already have coefficients in
$\Z_{(p)}$, and their complete-base coordinate values are unchanged.
Thus the exact orders and the split injections descend to the stated
localized-integer base as well.
\end{proof}

\begin{corollary}\label{cor:ramified-all-truncations}
Let $O$ be an arbitrary mixed-characteristic DVR with residue
characteristic $p$, normalized valuation $v$, and absolute
ramification index $h=v(p)$.  Let $n\ge2$ and put
$A=O[e]/(e^n)$.
If $n\nmid h$, then $\Phi_A$ is not surjective in degree $p$.
If $(n,h)=1$, the degree-$p$ reduction kernel contains a split
countably infinite direct sum of copies of $\Z/n$, with generators
that also vanish after inverting $p$.
The same split assertion holds for the cyclic precursors.
No completeness, finiteness, or perfection assumption on the
residue field is required.
\end{corollary}

\begin{proof}
First assume that $O$ is complete, with uniformizer $\pi$ and
residue field $k$.  We adapt the proof of
Theorem~\ref{thm:all-truncations}, retaining only primary
polynomials $f\in\F_p[z]$, $f\ne z$.
These are \emph{coarse} primary components: if $f$ factors over $k$,
we take the sum of all the resulting primary components.  Equivalently,
reduce $f(U)$ modulo $\pi$ and take its nilpotent Fitting summand;
its complement is the invertible Fitting summand.  The characteristic
polynomial factorization over $k$, followed by the same complete-DVR
projector construction, lifts this decomposition integrally.
It is exact for intertwined sequences.  Distinct irreducibles over
$\F_p$ remain coprime over $k$, so these coarse components are disjoint,
and only finitely many occur on a fixed lattice.

The curvature trace formula is unchanged.  Stable maps are killed
by $p$, which is a unit times $\pi^h$.  On a selected component,
some power of $f(U)$ is divisible by $\pi$; its $h$th power is
therefore divisible by $p$ and is stably zero.  Division by $p$
preserves equivariance because the lattices have no $p$-torsion.
On the complement $f(U)$ is invertible.  This proves the same
stable independence of the coarse primary summand.  Since $f(0)\ne0$,
$U$ is a unit on it, and the lifted matrix-factorization argument
of Lemma~\ref{lem:truncation-primary-acyclic} proves vanishing of
its integral Tate cohomology.

Measure the invariant-lattice covolume using $v$, modulo $n$.
A change of generic $A_K$-basis still has $K$-determinant
$\det(B_0)^n$, so the ambiguity is a multiple of $n$.
Exact additivity and projective stabilization are unchanged.
We thus obtain homomorphisms
\[
 \eta_f:K_0(\Lat_A(C_p))\longrightarrow\Z/n
\]
which vanish on constant lattices and have finite primary support.
For the companion lattice $L_D$ of
Theorem~\ref{thm:all-truncations}, the curvature is again
$I_p\otimes D$, and its invariant lattice is $(p,e)A^r$.
Consequently its coarse primary coordinate is
\[
 \eta_f(L_D)=hr\pmod n.
\]

For $n\nmid h$, take $r=1$ and $D=(1)$.
Only the identity normalizer twist has coarse primary polynomial
$z-1$; the other scalar twists have distinct nonzero eigenvalues
in $\F_p$.  Mackey restriction therefore gives coordinate $h\ne0$
modulo $n$ on the induced difference from the regular lattice.
Every Young lattice is constant, so this proves non-surjectivity.

For $(n,h)=1$, choose the irreducibles $f_i$ and integral companions
$D_i$ exactly as in Theorem~\ref{thm:all-truncations}: their degrees
$r_i$ are pairwise distinct and prime to $n(p-1)$.
They exist over the fixed finite field $\F_p$, irrespective of $k$.
The scalar-orbit argument takes place over $\F_p$; its coprimality
identities remain valid after extension to $k$.
Thus the induced coordinate matrix is diagonal with entries
$hr_i\in(\Z/n)^\times$.
The same exact sequence still annihilates each displayed class
by $n$.  Normalizing the coordinates by $(hr_i)^{-1}$ gives exact
order $n$ and finite-support retractions, both cyclically and after
induction.  The reduction and generic-vanishing arguments are
unchanged.

Finally, for a noncomplete DVR, extend scalars to its completion.
This preserves the residue field and normalized valuation of $p$,
and is exact on the admissible sequences of free lattices.
Compose the complete-base invariants with this functor.  All the
explicit companion matrices and their annihilation sequences are
already defined over $O$, and their coordinate values remain $hr_i$.
The preceding conclusions therefore hold over $O$ itself.
When $(n,h)\ne1$, the value $hr$ alone does not determine the exact
order of the displayed class; no stronger torsion-order assertion
is made here.
\end{proof}

\begin{corollary}\label{cor:truncation-adams}
Let $O=\Z_{(2)}$ or $\Z_2$, $A=O[e]/(e^n)$, and
\[
 x=\left[\left(A^2,
       \begin{pmatrix}1&e\\e^{n-1}&-1\end{pmatrix}\right)\right]
       -[A[C_2]].
\]
Then $x$ has exact order $n$, $x^2=0$, and
\[
 \psi^k(x)=\begin{cases}kx,&k\text{ odd},\\0,&k\text{ even}.
                         \end{cases}
\]
For $n=3$, this gives $\psi^3(x)=0\ne x$, although $3$ is a unit
in $A$ and the cube map on $C_2$ is the identity.
\end{corollary}

\begin{proof}
Let $\sigma$ be the sign line and $P=1+\sigma$ the regular class.
For any involution lattice $(M,T)$, the maps
\[
 0\longrightarrow M\longrightarrow
 A[C_2]\otimes_A M_{\mathrm{underlying}}
 \longrightarrow\sigma M\longrightarrow0,
 \qquad
 v\longmapsto(v,Tv),\quad (u,v)\longmapsto v-Tu,
\]
form an admissible exact sequence, where the middle generator
interchanges its two coordinates.  Applied to the rank-two lattice
in the statement, it gives $[M]+[\sigma M]=2P$, and therefore
$\sigma x=-x$.
As the displayed lattice has determinant $\sigma$,
\[
 \lambda_t(x)=
 \frac{1+(P+x)t+\sigma t^2}{(1+t)(1+\sigma t)}
 =1+\frac{xt}{1-t^2}.
\]
Comparing $\lambda_t(\sigma x)=\lambda_{\sigma t}(x)$ with
$\lambda_t(-x)=\lambda_t(x)^{-1}$ gives
\[
 \left(1-\frac{xt}{1-t^2}\right)
 \left(1+\frac{xt}{1-t^2}\right)=1;
\]
its coefficient of $t^2$ is $-x^2$, so $x^2=0$.
Consequently the integral logarithmic-derivative identity becomes
\[
 \sum_{k\ge1}(-1)^{k-1}\psi^k(x)t^k
 =t\lambda_t(x)^{-1}\frac{d\lambda_t(x)}{dt}
 =\sum_{j\ge0}(2j+1)x\,t^{2j+1}.
\]
This proves the displayed values without division in the
Grothendieck ring.  Exact order follows from the theorem with
$f=z+1$.
\end{proof}

\begin{corollary}\label{cor:all-prime-two-torsion}
For every prime $p$, the degree-$p$ reduction kernel for
$O[\epsilon]/(\epsilon^2)$, with $O=\Z_{(p)}$ or $\Z_p$,
contains a split countably infinite direct sum of copies of $\Z/2$.
The displayed classes also vanish after inverting $p$.
\end{corollary}
\begin{proof}
Apply Theorem~\ref{thm:all-truncations} with $n=2$.
\end{proof}

\begin{corollary}\label{cor:adams-not-power}
Let $p$ be odd and $A=\Z_{(p)}[\epsilon]/(\epsilon^2)$.
On $K_0(\Lat_A(C_p))$ the second Adams operation defined by exterior
powers is not pullback along the automorphism $g\mapsto g^2$.
This failure occurs although two is invertible in $A$.
\end{corollary}

\begin{proof}
Exterior powers of admissible exact sequences have the usual finite
filtration.  They therefore define an additive-to-multiplicative map
$\lambda_t$ on the exact Grothendieck group, and
$\lambda_t(-x)=\lambda_t(x)^{-1}$.
Let $x=[L_1]-[A[C_p]]$ be the nonzero order-two class of
Proposition~\ref{prop:truncation-torsion}, with $n=2$.
As $-x=x$, comparison of coefficients of $t^2$ gives
\[
 \lambda^2(x)=x^2-\lambda^2(x),\qquad
 \psi^2(x)=x^2-2\lambda^2(x)=0.
\]
Pullback along $g\mapsto g^2$ is an automorphism of the exact category,
so its induced group automorphism cannot kill $x$.
\end{proof}

Thus a formula on ordinary characters cannot determine the Adams
operations on this integral Grothendieck group.  The covolume argument supplies the independent information needed
to determine the truncation classes.
The annihilation identity used in the proof is the degree-two instance
of the general torsion--Adams relation of
\citet[Proposition~2, p.~2]{Clauwens2010}; the proof above derives the
needed identity directly from the exterior-power addition formula.

In the projective group-ring setting,
\citet[Propositions~4.2.3 and~4.3.2, pp.~1490--1492]{Pappas2015}
compares Adams operations with Cassou-Nogu\`es--Taylor operations.
The examples here concern lattices free over singular coefficient
rings, whose exact classes need not be determined by ordinary characters.

\section{The smallest flat example}
\label{sec:dual-number-obstruction}

The general construction has a rank-two example that can be checked
without a companion matrix.  Put $O=\Z_{(2)}$ and
$A=O[\epsilon]/(\epsilon^2)$.

For $p=2$, the cyclotomic matrix is $(-1)$.  The smallest lattice is
\[
 L=A^2,\qquad
 T=\begin{pmatrix}1&\epsilon\\\epsilon&-1\end{pmatrix}.
\]
Its reduced even Tate class is the trivial first coordinate and its odd
Tate class is the sign second coordinate.  The two extension boundaries
send these basis classes to one another, so $u=I$ and
$\nu_{t+1}(L)=1$.  Every scalar-extended integral lattice has invariant zero.
In particular
\[
 [L]\notin\mathbb Z[A]+\mathbb Z[A[C_2]],\qquad
 0\ne[L]-[A[C_2]],\qquad
 2([L]-[A[C_2]])=0.
\]
Since $S_2=C_2$, the degree-two Young permutation classes are precisely
$\rho_2=[A]$ and $\rho_1^2=[A[S_2]]$.  Consequently
\[
 \mathbb Z[\rho_1,\rho_2,\ldots]\longrightarrow R(A)
\]
is not surjective.

The ring $A$ is noetherian, local, of characteristic zero, free of rank two
over $\mathbb Z_{(2)}$, and has regular reduction $A/(\epsilon)=O$.  Its
only square roots of one are $\pm1$: writing a unit as $a+b\epsilon$ gives
$a^2=1$ and $2ab=0$, forcing $a=\pm1$ and $b=0$.  Moreover
$A[1/2]=\mathbb Q[\epsilon]/(\epsilon^2)$ is local, so its Picard group is
zero.  The rank-two example becomes trivial plus sign after inverting two;
explicit eigenvectors are $(1,\epsilon/2)$ and $(-\epsilon/2,1)$.
Thus neither extra roots of unity nor a nontrivial generic line bundle is
responsible for this obstruction.  It is detected by the square of the
boundary attached to the nilpotent extension.

\subsection{A relative exact sequence for involutions}

The following elementary sequence will also be useful over ramified bases.
Let $P=A[C_2]$.  For any $A[C_2]$-lattice $M$ of
$A$-rank $d$, with involution $T$, there is an admissible sequence
\begin{equation}
\label{eq:dual-relative-resolution}
 0\longrightarrow M\otimes\mathrm{sgn}
 \longrightarrow P\otimes_A M_{\mathrm{underlying}}
 \longrightarrow M\longrightarrow0.
\end{equation}
Here the group acts only on the first factor of the middle term, which is
$P^d$.  In the two copies of $M$ corresponding to $1,s$, the maps are
\[
 v\longmapsto(-Tv,v),\qquad
 (u,v)\longmapsto u+Tv.
\]
The latter has the $A$-linear section $m\mapsto(m,0)$; the group action on
the kernel parameter $v$ is $-T$, as required.  

\section{Flat dual numbers over ramified discrete valuation rings}
\label{sec:ramified-dual}

Ramification does not restore generation in residue characteristic two.
The cohomology of a suitable integral involution has length one, even
when $O/2O$ has even length.

\begin{theorem}\label{thm:ramified-dual}
Let $O$ be a characteristic-zero discrete valuation ring with residue
characteristic two, and set $A=O[\epsilon]/(\epsilon^2)$.
The kernel of reduction
\[
 K_0(\Lat_A(C_2))\longrightarrow K_0(\Lat_O(C_2))
\]
contains a direct summand isomorphic to a countable direct sum of copies
of $\Z/2$.  In particular $\Phi_A$ is not surjective.  This applies to
$O=\Z_2[\sqrt2]$ and to every ramification index and residue field.
\end{theorem}

\begin{proof}
Choose a uniformizer $\pi$ and put $h=2/\pi\in O$.
For an $A$-free involution write $T=S+\epsilon U$ with matrices over
$O$.  Then $S^2=I$ and $SU+US=0$.  On the two Tate groups
\[
 H^0=\ker(S-I)/\operatorname{im}(S+I),\qquad
 H^1=\ker(S+I)/\operatorname{im}(S-I),
\]
$U$ induces an odd operator $\delta$, interchanging the parities.
These groups have finite $O$-length and are killed by two: the usual
norm and difference maps give the latter assertion, and after
inverting two their complex is exact.

The operator is natural for maps that depend on $\epsilon$.
Indeed, an equivariant map $F=F_0+\epsilon F_1$ satisfies
\[
 S'F_0=F_0S,\qquad U'F_0-F_0U=F_1S-S'F_1.
\]
On a cycle $Sv=\pm v$, the second right-hand side is an appropriate
Tate boundary.  Therefore it induces zero in the next parity.
In an $A$-split exact sequence choose bases with
\[
 S=\begin{pmatrix}S'&a\\0&S''\end{pmatrix},\qquad
 U=\begin{pmatrix}U'&b\\0&U''\end{pmatrix}.
\]
The connecting map in the Tate exact sequence is $[v]\mapsto[av]$.
The off-diagonal relation
\[
 U'a+aU''=-S'b-bS''
\]
shows on each cycle that $\delta$ anticommutes with the connecting map
modulo boundaries.  All these cohomology groups are killed by two,
so it also commutes there.  Hence the even operator $u=\delta^2$ is
natural on the whole six-term periodic exact sequence.

Fix a polynomial $q\in\F_2[z]$ with $q(0)=1$.
On each finite-length Tate group take the nilpotent Fitting summand
for $q(u)$, namely $\ker(q(u)^N)$ for sufficiently large $N$.
This operation is exact and natural for maps commuting with $u$:
the Fitting decomposition into eventual kernel and image is a direct
sum, and both summands are preserved by every such map.  On the
chosen summand $u$ is invertible, by the coprimality of $z$ and $q$;
thus $\delta$ identifies its two parities.  The six-term exact
sequence becomes an exact three-cycle after this identification.
For any exact three-cycle of finite-length modules, the sum of the
three term lengths is twice the sum of the three image lengths.
Consequently
\[
 \nu_q(T)=\operatorname{length}_O
       \ker\bigl(q(\delta^2)^N:H^0\longrightarrow H^0\bigr)\pmod2
\]
is additive on the stated exact category.  It vanishes on every
scalar extension of an $O$-lattice, since $U=0$ there.

Choose distinct monic irreducibles $q_i\in\F_2[z]$ of odd degrees
$d_i$, with $q_i\ne z$; one may take one of each odd prime degree.
Let $D_i$ be the $0,1$ matrix lifting its companion matrix.  Define
\[
 T_i=\begin{pmatrix}
 I+\epsilon D_i&\pi I\\
 -h\epsilon D_i&-I-\epsilon D_i
 \end{pmatrix}.
\]
Direct multiplication gives $T_i^2=I$.  At $\epsilon=0$, its Tate
invariants are $H^0=k^{d_i}$ generated by the first block and
$H^1=k^{d_i}$ generated by the columns of $(I,-hI)$, where
$k=O/(\pi)$.  Indeed the norm image in the first block is
$(2,\pi)O^{d_i}=\pi O^{d_i}$, and the difference image in the other
eigenlattice is $\pi(I,-hI)O^{d_i}$.
The operator $\delta$ is $\overline D_i$ in both directions.  Hence
$u=\overline D_i^{\,2}$ and
\[
 \nu_{q_j}(T_i)=\begin{cases}1,&j=i,\\0,&j\ne i.\end{cases}
\]
Here squaring preserves the irreducible polynomial $q_i$ over
$\F_2$, and different $q_i$ stay coprime over $k$; no assumption that
$k$ is finite or perfect is needed.  The detected length is $d_i$,
which is odd.

Let $L_i$ be this lattice and put
$\gamma_i=[L_i]-d_i[A[C_2]]$.
Its reduction has a trivial free sublattice and sign free quotient,
so $\gamma_i$ lies in the stated reduction kernel.  The matrix
\[
 J_i=\begin{pmatrix}I&0\\-h(I+\epsilon D_i)&-I\end{pmatrix}
\]
is invertible and satisfies $J_iT_i=-T_iJ_i$.
The standard exact sequence
\[
 0\longrightarrow L_i\otimes\mathrm{sgn}
 \longrightarrow A[C_2]\otimes_A L_i
 \longrightarrow L_i\longrightarrow0
\]
therefore gives $2\gamma_i=0$.
The invariants $\nu_{q_j}$ prove independence and nonvanishing.
Only finitely many $q_j(\delta^2)$ can have nonzero nilpotent Fitting
summands on any given finite-length module, since the $q_j$ are
pairwise coprime.  Thus $(\nu_{q_j})_j$ takes values in the direct sum
and retracts the inclusion of the $\gamma_i$.
Finally every degree-two monomial in the trivial generators is a
scalar extension from $O$, whereas these classes are detected by
$\nu_{q_j}$.  This proves non-surjectivity of $\Phi_A$.
\end{proof}

\section{Odd primes and the ramification index}
\label{sec:odd-ramification}

A finite coefficient algebra makes the cyclotomic distance odd in a
precise range of ramification indices.  Keeping the entire algebra,
and selecting one of its residue factors by an endomorphism, makes
the argument work without a completeness assumption.

\begin{theorem}\label{thm:odd-ramification}
Let $p$ be an odd prime and let $O$ be any characteristic-zero discrete
valuation ring with residue field $k$ of characteristic $p$.
Write $e=v_O(p)$ for its absolute ramification index, and let $v_2$
denote the exponent of two in a positive integer.  Suppose
\[
 v_2(e)\le v_2(p-1).
\]
For $A=O[\epsilon]/(\epsilon^2)$, the kernel of
\[
 K_0(\Lat_A(S_p))\longrightarrow K_0(\Lat_O(S_p))
\]
contains a split countable direct sum of copies of $\Z/2$.
The displayed classes also vanish after inverting $p$.
In particular $\Phi_A$ is not surjective.
There is no completeness, finiteness, or perfection hypothesis on $k$.
\end{theorem}

\begin{corollary}\label{thm:even-ramified-cubic}
Let $\Z_3$ denote the complete ring of $3$-adic integers.
For $O=\Z_3[\zeta_3]$ or $\Z_3[\sqrt3]$, the flat dual numbers
$O[\epsilon]/(\epsilon^2)$ have a split countable degree-three
reduction-kernel subgroup isomorphic to $\bigoplus_{\mathbb N}\Z/2$.
\end{corollary}

\begin{proof}[Proof of the corollary]
Both rings have absolute ramification index two, so the theorem applies.
\end{proof}

\begin{proof}[Proof of Theorem~\ref{thm:odd-ramification}]
\emph{Primary composition lengths.}
For an $A$-lattice $E$ restricted to $C_p$, the two Tate groups of
$E_0=E/\epsilon E$ have finite $O$-length and are killed by $p$.
The canonical coefficient boundary $\delta$ is natural and
anticommutes with the connecting maps of admissible exact sequences,
by the coefficient calculation in
Proposition~\ref{prop:dual-bockstein-invariants}.
Thus $u=\delta^2$ commutes with the entire Tate exact sequence.
For a monic irreducible $f\in k[U]$, $f\ne U$, take the primary part
for the maximal ideal $(\pi,f)$ of $O[U]$, where $U$ acts by $u$.
Primary extraction is exact on finite-length modules: apply the
Chinese remainder decomposition for powers of their finitely many
supporting maximal ideals.
On this part $u$ is invertible, so $\delta$ identifies both parities
and the two connecting-image modules.  The six-term length formula
therefore proves that
\[
 \nu_f(E)=\operatorname{length}_{O[U]}H^0(E_0)_f\pmod2
\]
is additive.  The length is composition length over $O[U]$, not
dimension over $k$ or length over $O$.
Every constant coefficient extension has $\delta=0$.
The comparison~\eqref{eq:dual-prime-scaling} shows that changing the
generator by $g\mapsto g^a$ scales the residual primary parameter
by $a\in\F_p^\times$.

\emph{A finite coefficient algebra.}
Put $d=\gcd(e,p-1)$, $t=(p-1)/d$, and $m=e/d$.
The hypothesis makes $m$ odd.
Choose a uniformizer $\pi$ and form
\[
 O_1=O[s]/(s^t-\pi).
\]
The polynomial is Eisenstein, so this is a finite free domain over
$O$ of rank $t$.  Its unique maximal ideal is $(s)$; hence it is a
discrete valuation ring, with $v_s(p)=(p-1)m$.
Write $b=p/s^{(p-1)m}$, a unit, and define
\[
 G(Y)=b+\sum_{j=1}^{p-2}\frac{\binom p{j+1}}p\,
              b s^{mj}Y^j+Y^{p-1},\qquad C=O_1[Y]/G(Y).
\]
This is a finite free coefficient algebra over $O$.
Its residue polynomial $\bar b+Y^{p-1}$ is separable with nonzero
roots, so $C/sC$ is a product of finite separable fields over $k$.
The element $Y$ is a unit, and $s$ is a nonzerodivisor.
The defining polynomial gives literal identities
\[
 \zeta=1+s^mY,\qquad \Phi_p(\zeta)=0,\qquad
 \zeta^p=1,\qquad \zeta-1=s^mY.
\]
The algebra $C$ may be semilocal.  No selection of a local coefficient
factor, and no Hensel lifting, is required.

\emph{A companion lattice and its two boundaries.}
The following construction applies to this $C$ and to the analogous
algebras after further finite unramified coefficient extensions.
Write
\[
 h(Z)=\frac{Z^p-1}{Z-\zeta}=(Z-1)q(Z),\qquad
 N(Z)=1+Z+\cdots+Z^{p-1}=(Z-\zeta)q(Z).
\]
These are polynomial identities.  For example $h(1)=0$ follows from
$(1-\zeta)h(1)=0$ because $1-\zeta$ is a nonzerodivisor.
Let $V=C[Z]/h(Z)$, let $B$ be multiplication by $Z$ in its monomial
basis, let $v=1\in V$, and let $w$ extract the coefficient of
$Z^{p-2}$.  The polynomial $q$ is monic of degree $p-2$, so $w(q)=1$.
For any $c\in C$, define
\[
 T_c=\begin{pmatrix}B&\epsilon v\\\epsilon c w&\zeta\end{pmatrix}
\]
on a free $C[\epsilon]/(\epsilon^2)$-module of rank $p$.
Its first-order off-diagonal terms in $T_c^p$ are multiples of
$h(B)=0$; hence $T_c^p=I$.
There is no unit condition on $c$.

For the constant reduction $M=V\oplus C_\zeta$, the invariant
submodule of $V$ is $Cq$, and its norm image is $(1-\zeta)Cq$.
Indeed $N(B)f=(1-\zeta)f(1)q$.
Its norm kernel is $(B-1)V$, because $1-\zeta$ is a nonzerodivisor.
The scalar $\zeta$ summand has zero even Tate group and odd group
$C/(\zeta-1)$.  Consequently
\[
 H^0(M)=H^1(M)=C/(s^m).
\]
The even coefficient boundary sends $q$ to $cw(q)=c$.
The odd coefficient boundary sends $1$ to $q$, by the polynomial
divided-difference identity
\[
 \frac{N(B)-N(\zeta)}{B-\zeta}\,v=q(B)v=q.
\]
No inverse of $B-\zeta$ is taken in this expression.
Thus $u=\delta^2$ is exactly multiplication by $c$.

The reduced class is regular: there is a $C$-split exact sequence
\[
 0\longrightarrow C_\zeta\xrightarrow{1\mapsto h}
 C[C_p]\longrightarrow C[Z]/h(Z)\longrightarrow0.
\]
The kernel is freely generated by the monic polynomial $h$.
There is also an exact order-two relation.
For any dual-number lattice $E$, let $E^-$ be its coefficient twist
by $\epsilon\mapsto-\epsilon$.
The admissible sequence
\[
 0\longrightarrow E^-\longrightarrow
 (C[\epsilon]/(\epsilon^2))\otimes_C\operatorname{Res}_C E
 \longrightarrow E\longrightarrow0
\]
has injection $x\mapsto\epsilon\otimes x-1\otimes\epsilon x$
and evaluation surjection.  Its kernel is a free direct summand in
an underlying coefficient basis.
The underlying $C$-lattice is an extension of two copies of $E_0$.
For $E=E_c$, the matrix $\operatorname{diag}(I_{p-1},-1)$ identifies
$E^-$ with $E$.  It follows that
\begin{equation}\label{eq:odd-order-two}
 2\bigl([E_c]-[(C[\epsilon]/(\epsilon^2))[C_p]]\bigr)=0.
\end{equation}

\emph{Infinite residue fields.}
Suppose $k$ is infinite, and choose one field factor $F$ of $C/sC$.
There are infinitely many primitive elements $\alpha\in F$ such
that $\alpha$ and $a\alpha$ are not $k$-conjugate for any
$a\in\F_p^\times\setminus\{1\}$.
To prove this, work in a normal closure of $F/k$.
Each unwanted equation $\sigma(\alpha)=a\alpha$ defines a proper
$k$-linear subspace of $F$, except for $(\sigma,a)=(1,1)$:
if it held identically, evaluation at $1$ would give $a=1$ and then
$\sigma=1$.
Avoid these finitely many subspaces.  Over an infinite field their
complement has infinitely many points; the equations with $a=1$
also exclude every nonprimitive element.
Choose a countable sequence $\alpha_i$ with pairwise disjoint
scalar-conjugacy orbits, each of which is finite.

Lift $(\alpha_i,0,\ldots,0)\in C/sC$ to $c_i\in C$, putting
$\alpha_i$ in the selected factor.  Let $f_i$ be its minimal
polynomial over $k$.
The $s$-filtration of $C/(s^m)$ has $m$ factors $C/sC$.
In each layer exactly one factor contributes to the $f_i$-primary
part: it is the simple $O[U]$-module $F$, with $U=\alpha_i$.
Every other factor has $U=0$ and disappears.
The primary composition length is therefore the odd integer $m$.

\emph{Finite residue fields.}
Suppose $k=\F_Q$, and choose $r$ such that $\bar G$ splits over
$\F_{Q^r}$.
For strictly increasing positive multiples $n_i$ of $r$, lift a
monic irreducible polynomial of degree $n_i$ over $k$ to $O_1$.
Its quotient $S_i$ is finite free over $O_1$ and has unique maximal
ideal $(s)$ and residue field $k_i=\F_{Q^{n_i}}$; thus it is a
discrete valuation ring.
Set $C_i=S_i[Y]/G(Y)$.
Its reduction is a product of $p-1$ copies of $k_i$.
Choose a generator $\alpha_i$ of $k_i^\times$, lift
$(\alpha_i,0,\ldots,0)$ to $c_i\in C_i$, and let $f_i$ be its
minimal polynomial over $k$.
Again the selected primary length is $m$.
The scalar-conjugacy orbit is free: if
$\alpha_i^{Q^j}=a\alpha_i$ with $0<j<n_i$ and $a\in\F_p^\times$,
then
\[
 Q^{n_i}-1\mid(p-1)(Q^j-1),
\]
although the positive right side is smaller than the left side.
For $j=0$ one obtains $a=1$.
Distinct degrees make the chosen scalar orbits disjoint.

\emph{Induction and splitting.}
In either case write $C_i$ for the chosen coefficient algebra and
$a_i$ for its rank over $O$.
Restrict $E_{c_i}$ to $A$, induce from $C_p$ to $S_p$, and subtract
$a_i[A[S_p]]$; call the resulting class $x_i$.
Restriction preserves freeness and exactness because $C_i$ is finite
free over $O$.
The reduced regular-class equality proves that $x_i$ specializes to
zero over $O$, and~\eqref{eq:odd-order-two} gives $2x_i=0$.
In the Mackey restriction back to $C_p$, the normalizer terms give
each scalar twist $a\in\F_p^\times$ exactly once.
All other intersections are trivial and give Tate-acyclic free
group lattices.  Therefore
\[
 \bigl(\nu_{f_j}\circ\operatorname{Res}_{C_p}^{S_p}\bigr)(x_i)
 =\begin{cases}1,&i=j,\\0,&i\ne j.\end{cases}
\]
This proves nonzero exact order two and independence.
Every finite-length Tate module has finite primary support, so
the direct sum of these invariants retracts the inclusion of the
$x_i$.  The invariants vanish on every Young permutation class.
After inverting $p$, rational nil invariance identifies each class
with its regular reduction, proving the last vanishing assertion.
\end{proof}

\begin{remark}\label{rem:odd-ramification-scope}
The condition in Theorem~\ref{thm:odd-ramification} is sufficient.
When $v_2(e)>v_2(p-1)$, the cyclotomic distance $e/\gcd(e,p-1)$
in this construction is even, and its parity detectors vanish.
Further finite coefficient extensions cannot make that distance odd:
if $e'$ is a multiple of $e$ and $p-1$ divides $e'$, then
$e'/(p-1)$ remains even.
This does not establish a positive result in the complementary range.
\end{remark}

\section{The limit of the primary parity obstruction}
\label{sec:complementary-detector-limit}

The ramification condition in Theorem~\ref{thm:odd-ramification}
is an exact boundary for its primary parity detectors.  Beyond that
condition all these detectors vanish, in every matrix rank.  This
does not decide whether the induction ring is polynomial there.

\begin{theorem}\label{thm:complementary-detector-limit}
Let $O$ be a mixed-characteristic discrete valuation ring with odd
residue characteristic $p$, residue field $k$, and normalized valuation
$v$.  Put $e=v(p)$ and suppose
\[
 v_2(e)>v_2(p-1).
\]
Let $E$ be an $O[\epsilon]/(\epsilon^2)$-free $C_p$-lattice,
let $M=E/\epsilon E$, and let $\delta$ be the coefficient connecting
operator on $\widehat H^\bullet(C_p,M)$.  Set $u=\delta^2$.
For every monic irreducible $f\in k[z]$ other than $z$, the
$f$-primary part of $\widehat H^0(C_p,M)$ has even composition
length over $O[u]$.  Thus every nonzero-primary length-parity
detector of Section~\ref{sec:odd-ramification} is zero.

There is no completeness, finiteness, or perfection assumption on
$k$.  The conclusion is a vanishing theorem for these detectors;
it asserts neither generation nor that the reduction kernel is zero.
\end{theorem}

The proof uses the hypersurface operator and periodicity theorem of
\citet[Proposition~1.3, Corollary~1.4, pp.~40--41;
Theorem~6.1(iii), pp.~52--53]{Eisenbud1980}.  We first prove the
geometric parity calculation to which that theorem will be applied.

\subsection{Symmetric factorizations of every rank}

\begin{lemma}\label{lem:symmetric-factorization-parity}
Assume that $O$ is complete and satisfies the hypotheses of the
theorem.  Put $B=O[[x]]$ and $q(x)=(1+x)^p-1$.
If a square matrix $D$ over $B$ satisfies $D^2=q(x)I$, then
\[
 H(D(0))=\ker D(0)/\operatorname{im}D(0)
\]
has finite $O$-length, and that length is even.
\end{lemma}

\begin{proof}
Consider the double cover
\[
 S=B[y]/(y^2-q(x)).
\]
It is a complete local Cohen--Macaulay domain of dimension two.
Indeed it is finite free over $B$, and $q$ has valuation one at
the prime $(x)$ of $B$, so it is not a square in its fraction field.
A maximal Cohen--Macaulay $S$-module is free over $B$: the parameters
$(\pi,x)$ have depth two on it, and the Auslander--Buchsbaum formula
over the regular local ring $B$ gives projective dimension zero.
Conversely a $B$-free $S$-module is maximal Cohen--Macaulay.
Multiplication by $y$ is therefore a matrix of the displayed kind.

After passing to $K=\operatorname{Frac}(O)$ and completing at $x=0$,
the cover becomes $K[[y]]$, because $q'(0)=p$ is a unit in $K$.
The module is torsion-free there and hence free.  Its reduction
modulo $x$ is free over $K[y]/(y^2)$, so $D(0)$ is exact over $K$.
Its homology over $O$ consequently has finite length.

For a maximal Cohen--Macaulay module $V$ define
\[
 \chi(V)=\ell_O H(y:V/xV\longrightarrow V/xV)\pmod2.
\]
An exact sequence of such modules splits over $B$, so reduction
modulo $x$ remains exact.  The resulting one-periodic homology
sequence is an exact three-cycle of finite-length modules.  Its
length defect is twice an image length.  Thus $\chi$ is additive
modulo two.  It is zero on $S$-free modules: modulo $x$, multiplication
by $y$ on $S$ is a contractible length-two nilpotent block.

This homomorphism extends to $G_0(S)$.  Every finite $S$-module has
a second syzygy that is maximal Cohen--Macaulay, by the depth lemma.
Define its value by the alternating sum in a two-step free
presentation.  Schanuel's lemma makes the value independent of the
presentation, since free summands have value zero; a horseshoe
presentation proves additivity.  This argument does not claim finite
projective dimension for arbitrary $S$-modules.

The special fibre is the domain
\[
 S/(\pi)=k[[x,y]]/(y^2-x^p).
\]
Multiplication by $y$ is injective on it, with quotient
$k[[x]]/(x^p)$.  Hence $p[k]=0$ in $G_0(S)$.  Since $p$ is odd,
$\chi(k)=0$, and every finite-length $S$-module has value zero.
The sequence given by multiplication by $\pi$ also gives
$\chi(S/(\pi))=0$.

Let $P$ be a height-one prime of $S$ not containing $\pi$.
For a nonzero element $a(x)+yb(x)$ of $P$, its norm
$a(x)^2-q(x)b(x)^2$ is a nonzero element of $P\cap B$.
Remove its maximal power of $\pi$.  The resulting primitive series
is a unit times a nonconstant distinguished monic polynomial in
$x$.  One can see this preparation directly: modulo $\pi$, write
the series as $x^d\overline U$.  Given a monic degree-$d$ polynomial
$P_d$ and a unit $U$ correct modulo $\pi^j$, divide the error by
$\overline U$ and split it into a polynomial of degree less than
$d$ and $x^d$ times a series.  These two terms correct $P_d$ and
$U$ modulo $\pi^{j+1}$.  Completeness gives the factorization.
The degree is positive because the series lies in a proper prime.
Consequently $S/P$ is finite over $O$.

Its fraction field $L$ is a finite extension of $K$.  The normalization
$O_L$ is finite over $O$, and $O_L/(S/P)$ has finite length.
Write $\alpha,\beta$ for the images of $x,y$ in $O_L$; both lie
in its maximal ideal and $\beta^2=q(\alpha)$.
Over
\[
 S_L=O_L[[x]][y]/(y^2-q(x)),
\]
the point ideal $I=(x-\alpha,y-\beta)$ has quotient $O_L$.
It is free over $O_L[[x]]$ with the two indicated generators,
on which multiplication by $y$ is
\[
 D_{\alpha,\beta}=
 \begin{pmatrix}
 \beta&c(x)\\ x-\alpha&-\beta
 \end{pmatrix},\qquad
 c(x)=\frac{q(x)-\beta^2}{x-\alpha}.
\]
The division is integral since $q(\alpha)=\beta^2$.
This matrix squares to $q(x)I$.  Since $S_L$ is free over $S$,
the point sequence gives $\chi(O_L)=-\chi(I)$.

We compute the specialized length over $O_L$.  Write
$e_L=v_L(p)$, where $v_L$ is normalized.  It is an integer multiple
of $e$, so $v_2(e_L)>v_2(p-1)$.  Suppose first that
$\alpha\beta\ne0$, and put $l=v_L(\alpha)$ and $s=v_L(\beta)$.
Away from the equality threshold $e_L=(p-1)l$, the binomial expansion
gives
\[
 2s=\min(e_L+l,pl).
\]
In the lower range $2s=pl$, so $l$ is even.  In the upper range
$2s=e_L+l$, and $e_L$ is even, so again $l$ is even.
At the threshold, $l=e_L/(p-1)$ is an even integer.

The specialized matrix has entries
$\beta,\beta^2/\alpha,-\alpha,-\beta$.
It is a nonzero rank-one square-zero matrix.  Its homology length
is its nonzero elementary-divisor valuation,
\[
 m=\min(s,l,2s-l).
\]
In the lower range $m=l$; in the upper range $m=\min(e_L,l)$.
At the threshold cancellation can only increase $2s$ above $pl$,
and $m=l$.  All these integers are even.
If $\alpha=0$, then $\beta=0$ and the nonzero entry is $p$,
giving length $e_L$.  If $\beta=0$ but $\alpha\ne0$, then
$1+\alpha$ is a nontrivial $p$th root of unity, and the length is
$v_L(\alpha)=e_L/(p-1)$, also even.  These cases exhaust the points.

Restriction from $O_L$ to $O$ multiplies a length by the residue
degree, so $\chi(I)=0$.  The finite normalization cokernel then gives
$\chi(S/P)=0$.  Finally, prime filtrations generate $G_0(S)$.
The zero prime contributes $S$, the only height-one prime containing
$\pi$ is $(\pi)$, every other height-one prime was treated above,
and the maximal ideal contributes $k$.  All values are zero.
Thus $\chi$ is identically zero, proving the lemma.
\end{proof}

\subsection{From a primary lattice to a symmetric factorization}

\begin{proof}[Proof of Theorem~\ref{thm:complementary-detector-limit}]
We may first complete $O$.  Tate cohomology commutes with this flat
extension, and its finite primary composition series has unchanged
length and residue factors.  Completion also preserves $v(p)$.

Let $T$ be the action on $E$ and lift it to $O[[\epsilon]]$.
Its stable curvature is
\[
 \widetilde U=(\widetilde T^p-I)/\epsilon^2.
\]
To identify its coefficient action, expand the two-periodic Tate
differential as $d(\epsilon)=d+\epsilon Q+\epsilon^2R+\cdots$.
The identity $d(\epsilon)^2=\epsilon^2\widetilde U$ gives
\[
 \widetilde U(0)=Q^2+dR+Rd.
\]
Since $Q$ induces $\delta$, the curvature induces exactly
$u=\delta^2$ on $\widehat H^\bullet(C_p,M)$.
Integral primary decomposition of the curvature therefore selects
the required cohomology pieces.  On every selected nonzero component
the curvature is a unit, and the cancellation argument of
Lemma~\ref{lem:truncation-primary-acyclic}, valid over any complete
DVR, gives zero Tate cohomology of that component of $E$ itself.

We must keep track of primary composition length, rather than only
ordinary $O$-length.  Write the irreducible residue polynomial as
\[
 f(z)=g(z^{p^a}),
\]
with $g$ separable irreducible.  Pass to a finite separable residue
extension splitting $g$, realized by a finite unramified extension
$O'/O$.  Each resulting irreducible factor is
$z^{p^a}-\alpha_j$, of degree $p^a$: a separable extension is
linearly disjoint from the corresponding purely inseparable field
extension.  Each simple original $O[u]$-factor becomes one
$p^a$-dimensional residue factor in each selected summand.  Hence,
on any one such summand $H_j$,
\[
 \ell_{O'}(H_j)=p^a\ell_{O[u]}(H_f).
\]
The factor $p^a$ is odd.  It is enough to prove even ordinary length
on each selected lattice.  Rename $O'$ as $O$ for the next steps;
its value of $e$ is unchanged.

Let $F$ denote a selected curvature-unit lattice.  Since
$\widehat H^\bullet(C_p,F)=0$, the coefficient sequence for
multiplication by the uniformizer gives zero Tate cohomology of
$F/\pi F$.  A finite $k[C_p]=k[x]/(x^p)$-module is a sum of
$k[x]/(x^j)$, $1\le j\le p$; every summand with $j<p$ has
nonzero Tate cohomology.  Thus $F/\pi F$ is free over $k[C_p]$.
Lift a basis.  Nakayama gives a surjection $O[C_p]^r\to F$, and
equality of $O$-ranks gives injectivity.  Hence $F$ is free over
\[
 R=O[C_p]=O[[x]]/(q(x)),\qquad q(x)=(1+x)^p-1.
\]
Multiplication by $\epsilon$ is an $R$-linear matrix $D$ with
$D^2=0$ and $\ker D=\operatorname{im}D$, because $F$ is free
over $O[\epsilon]/(\epsilon^2)$.

If $D$ has a unit entry, choose $v$ with $w=Dv$ primitive.
The reductions of $v,w$ are independent, since $D^2=0$.
Choose a functional $\lambda$ with $\lambda(w)=1$ and
$\lambda(v)=0$.  The projection
\[
 P(z)=\lambda(Dz)v+\lambda(z)w
\]
commutes with $D$ and is the identity on the span of $v,w$.
It splits off the contractible block $Dv=w$, $Dw=0$.
As an $O[\epsilon]/(\epsilon^2)[C_p]$-module this block is free,
so it contributes zero to the coefficient Tate groups.
After removing all such blocks, the remaining exact one-periodic
free complex is minimal over $R$.

Lift the same matrix $D$ in every degree to $B=O[[x]]$ and write
\[
 \widetilde D^2=q(x)V.
\]
Cancellation of $q$ shows that $V$ commutes with $\widetilde D$.
The matrix $V$ is the degree-two hypersurface operator of the minimal
periodic resolution.  By Eisenbud's Theorem~6.1(iii), it gives the
periodicity and is invertible modulo the maximal ideal.  This also
holds for the chosen lift: the operator is unique up to homotopy,
and on a minimal complex a null-homotopic map is zero modulo the
maximal ideal.  Comparison with a matrix-factorization resolution,
where the operator is the identity, proves invertibility of its
residue.  Thus $V$ is a unit matrix over $B$.

After a finite unramified extension, $V$ has a commuting square root.
Here is the residue-field argument without a perfection assumption.
In the finite commutative algebra $k[\overline V]$, each reduced
field factor is $L=k(v)$, with $v\ne0$.  Write its minimal polynomial
as $h(z^{p^b})$ with $h$ separable.  Then
$L_s=k(v^{p^b})$ is separable over $k$.  Adjoin a square root
$\mu$ of $v^{p^b}$ over $L_s$; this is a separable extension.
A finite separable extension of $k$ can contain the required
conjugate copies of these fields and square roots.  In each resulting
factor,
\[
 w=v^{(p^b+1)/2}/\mu\qquad\text{satisfies}\qquad w^2=v.
\]
The root lifts through the nilpotent radical since $2w$ is a unit.
Thus it is a polynomial square root of $\overline V$.
Realize the finite separable residue extension unramifiedly, lift
the polynomial root to $W_0(V)$, and iterate
\[
 W_{j+1}=(W_j+VW_j^{-1})/2.
\]
The matrices converge in the complete local topology to a unit
$W$ with $W^2=V$.  Every iterate, and therefore the limit, commutes
with $\widetilde D$.  Consequently
\[
 D_*:=\widetilde D W^{-1}\qquad\text{satisfies}\qquad D_*^2=qI.
\]
At $x=0$ its kernel and image are those of $\widetilde D(0)$,
because $W(0)$ is a commuting unit.  Unramified extension preserves
ordinary module length.  Lemma~\ref{lem:symmetric-factorization-parity}
therefore makes $\ell_O H(D(0))$ even.

Finally, let $M_F=F/\epsilon F$.  The coefficient sequence
$0\to M_F\to F\to M_F\to0$ identifies $M_F$ with
$\ker D=\operatorname{im}D$.  On the invariants of the free
$R$-module $F$, multiplication by $\epsilon$ is $D(0)$.
Taking invariants gives
\[
 H^1(C_p,M_F)=\ker D(0)/\operatorname{im}D(0).
\]
Since $F$ has zero Tate cohomology, its coefficient connecting map
identifies the two Tate parities of $M_F$.  Their common ordinary
length is therefore even.  Dividing the earlier length identity by
the odd integer $p^a$ proves the required primary parity.
Completion preserves it, so the conclusion holds over the original
DVR as well.
\end{proof}

The theorem explains why higher-rank searches with the same parity
coordinates cannot settle the complementary ramification range.
A classification there requires information beyond these coordinates.

\section{Truncation exponents prime to the residue characteristic}
\label{sec:prime-to-p-all-dvr}

The divisor-class obstruction extends from dual numbers to every
truncation whose exponent is invertible in the DVR.  We use the
generic-hypersurface equivalence to define the detector on all
coefficient-free lattices.  The order of the detected class is proved
by explicit exact sequences.

\begin{theorem}\label{thm:prime-to-p-all-dvr}
Let \(O\) be a characteristic-zero DVR with residue characteristic
\(p>0\), and let \(n\ge2\) satisfy \(p\nmid n\).  Put
\(A_n=O[t]/(t^n)\).  There is a class of exact order \(n\) in
\[
 \ker\bigl(K_0(\Lat_{A_n}(S_p))
                   \longrightarrow K_0(\Lat_O(S_p))\bigr)
\]
which is zero after extension to \(\operatorname{Frac}(O)\).
An additive detector defined on the whole Grothendieck group
annihilates all Young permutation lattices and detects this class.
Thus \(\Phi_{A_n}\) is not surjective in degree \(p\).
There is no restriction on the ramification or the residue field,
and \(O\) need not be complete.
\end{theorem}

The theorem includes \(p=2\) with \(n\) odd.  It asserts the exact
order of a cyclic subgroup, without asserting that this subgroup
is a direct summand.

\subsection{Finite lattices and the mesh exact sequences}

Fix a generator \(\sigma\) of \(P=C_p\), write \(X=\sigma-1\), and set
\[
 g(X)=(1+X)^p-1=Xh(X),\qquad
 h(X)=p+\binom p2X+\cdots+X^{p-1}.
\]
For \(0\le i\le n\), define \(M_i\) over \(A_n[X]\) by generators
\(a,b\) and relations
\begin{equation}\label{eq:prime-to-p-presentations}
 Xa=t^{n-i}b,\qquad t^ia+h(X)b=0,
 \qquad
 \Phi_i=
 \begin{pmatrix}X&t^i\\-t^{n-i}&h(X)\end{pmatrix}.
\end{equation}
The determinant of \(\Phi_i\) is \(g(X)+t^n\), which is \(g(X)\)
over \(A_n\).  Its adjugate shows that \(g\) annihilates \(M_i\).

\begin{lemma}\label{lem:prime-to-p-mesh}
Each \(M_i\) is a \(P\)-lattice free of rank \(p\) over \(A_n\).
One has \(M_0\simeq M_n\simeq A_n[P]\), and, for \(1\le i<n\),
there is a coefficient-split exact sequence
\begin{equation}\label{eq:prime-to-p-mesh}
 0\longrightarrow M_i
 \xrightarrow{(\beta_{i-1},\alpha_i)}
 M_{i-1}\oplus M_{i+1}
 \xrightarrow{(\alpha_{i-1},-\beta_i)}
 M_i\longrightarrow0.
\end{equation}
Consequently, with \(x_i=[M_i]-[A_n[P]]\),
\begin{equation}\label{eq:prime-to-p-mesh-order}
 x_i=i x_1,\qquad n x_1=0,\qquad x_{n-1}=-x_1.
\end{equation}
For \(1\le i<n\), these classes vanish after coefficient reduction
\(t\mapsto0\) and after extension to \(\operatorname{Frac}(O)\).
\end{lemma}

\begin{proof}
An \(A_n\)-basis is
\begin{equation}\label{eq:prime-to-p-basis}
 a,\qquad X^j b\quad(0\le j\le p-2).
\end{equation}
The first relation specifies \(Xa\), while the second specifies
\(X^{p-1}b\), since \(h\) is monic.  Conversely these prescriptions
define an \(X\)-operator on the free module with basis
\eqref{eq:prime-to-p-basis}; both displayed relations hold there.
The resulting maps between this free module and the presented module
are inverse.  This proves freeness, and \(g(X)=0\) gives the action
of \(P\).  At \(i=0\) eliminate \(a=-hb\); at \(i=n\) eliminate
\(b=Xa\).  In both cases the remaining relation is \(g=0\), proving
the endpoint identifications.

In the bases \eqref{eq:prime-to-p-basis}, put
\[
 \alpha_i:M_i\longrightarrow M_{i+1},\quad
 \alpha_i=\operatorname{diag}(t,I_{p-1}),\qquad
 \beta_i:M_{i+1}\longrightarrow M_i,\quad
 \beta_i=\operatorname{diag}(1,tI_{p-1}).
\]
For example, \(\alpha_i\) sends \(a\) to \(ta\) and \(b\) to \(b\)
in the target, and \(\beta_i\) sends \(a\) to \(a\) and \(b\) to
\(tb\).  Substitution in \eqref{eq:prime-to-p-presentations}
proves equivariance; both composites are multiplication by \(t\).
On the first basis coordinate, the maps in
\eqref{eq:prime-to-p-mesh} are \((1,t)\) and \((t,-1)\);
on each remaining coordinate, they are \((t,1)\) and \((1,-t)\).
These formulas show that their composite is zero and that the kernel
of the second is the image of the first.  For coordinates \((r,v)\)
on the last \(M_i\), a coefficient right inverse is
\[
 (r,v)\longmapsto ((0,v),(-r,0)).
\]
A coefficient left inverse on the first map is
\[
 ((r_L,v_L),(r_R,v_R))\longmapsto(r_L,v_R).
\]
Thus \eqref{eq:prime-to-p-mesh} is admissible.  Its relations are
\(2x_i=x_{i-1}+x_{i+1}\), with \(x_0=x_n=0\), which give
\eqref{eq:prime-to-p-mesh-order}.

For \(0<i<n\), reduction modulo \(t\) is
\[
 O[X]/(X)\oplus O[X]/(h).
\]
The exact sequence
\begin{equation}\label{eq:prime-to-p-reduction}
 0\longrightarrow O[X]/(h)
  \xrightarrow{\ \cdot X\ } O[X]/(Xh)
  \xrightarrow{\ X\mapsto0\ }O\longrightarrow0
\end{equation}
is split over \(O\) by the constant polynomials on the right.
Division in \(O[X]\) verifies injectivity and the stated kernel.
This proves reduction vanishing.

Let \(K=\operatorname{Frac}(O)\).  Since \(h(0)=p\) is a unit in
\(K\), the same reduction is isomorphic to \(K[P]\).
Lift a cyclic generator of that reduction to \(M_i\otimes_OK\).
It defines a homomorphism
\[
 (K[t]/t^n)[P]\longrightarrow M_i\otimes_OK.
\]
Both sides are free of coefficient rank \(p\), and modulo \(t\)
the map is an isomorphism.  Its determinant is a unit in the local
Artinian coefficient ring, so it is an isomorphism.  This proves
generic vanishing.  None of these arguments uses completeness.
\end{proof}

\subsection{Divisor classes of exact prescribed order}

Assume for now that \(O\) is complete, with uniformizer \(\pi\)
and residue field \(k\), and set
\begin{equation}\label{eq:prime-to-p-surface}
 S=O[[X,t]]/(g(X)+t^n),\qquad B=O[[X]],\qquad
 I_a=(X,t^a)\subset S\quad(1\le a<n).
\end{equation}

\begin{lemma}\label{lem:prime-to-p-surface}
The ring \(S\) is a normal two-dimensional complete local domain.
Each \(I_a\) is rank-one reflexive and nonprincipal.  If
\(\mathfrak p=(X,t)\), then
\[
 \operatorname{div}(X)=n\mathfrak p,\qquad
 [I_a]=a[I_1]\quad\text{in }\operatorname{Cl}(S).
\]
In particular \([I_1]\) has exact order \(n\) in
\(\operatorname{Cl}(S)\).  Its class in
\(K_0(D_{\mathrm{sg}}(S)^{\mathrm{ic}})\) has order at least \(n\):
if \(r[I_1]=0\) there, then \(n\mid r\).
\end{lemma}

\begin{proof}
The polynomial \(t^n+g(X)\) is Eisenstein at \((X)\) over
\(B_{(X)}\), because \(g=Xh\) and \(h(0)=p\ne0\).
It is therefore irreducible over \(\operatorname{Frac}(B)\).
Since \(S\) is free of rank \(n\) over \(B\), it is a domain.

Its special fiber is
\[
 k[[X,t]]/(t^n+X^p),
\]
which is also a domain.  Indeed, every root of \(t^n+X^p\) over
an algebraic closure of \(k((X))\) has \(X\)-valuation \(p/n\).
A monic factor of degree \(r\) over \(k((X))\) would have constant
coefficient of valuation \(rp/n\), an integer.  Since \((n,p)=1\),
one has \(n\mid r\), proving irreducibility.  The monic
power-series quotient embeds in this field extension.

The hypersurface \(S\) is Cohen--Macaulay of dimension two.
Its only height-one prime above \(\pi\) is \((\pi)\), which has a
regular localization because its maximal ideal is principal.
At every prime not containing \(\pi\), the derivative
\[
 g'(X)=p(1+X)^{p-1}
\]
is a unit.  The derivation in \(X\) shows that the defining equation
is outside the square of the maximal ideal in the corresponding
regular ambient localization; its quotient is regular.
Thus \(S\) satisfies \(R_1\) and \(S_2\), hence is normal.

In the \(B\)-basis \(1,t,\ldots,t^{n-1}\) of \(S\), one has
\begin{equation}\label{eq:prime-to-p-ideal-basis}
 I_a=BX\oplus BXt\oplus\cdots\oplus BXt^{a-1}
            \oplus Bt^a\oplus\cdots\oplus Bt^{n-1}.
\end{equation}
This is stable under multiplication by \(t\), using \(t^n=-Xh\),
and is exactly the ideal generated by \(X,t^a\).
It is \(B\)-free and torsion free over the normal surface \(S\);
the regular sequence \(\pi,X\) proves depth two.  Thus it is
reflexive.

Its two displayed ideal generators are minimal.  If
\(cX+dt^a\in\mathfrak m_S I_a\), reduction modulo \((\pi,X)\)
in \(k[[t]]/(t^n)\) shows that \(d\in\mathfrak m_S\), because
\(t^a\notin(t^{a+1})\).  Then reduce \(cX\) modulo
\((\pi,t,X^2)\).  This quotient is \(k[X]/(X^2)\), since \(p\ge2\);
the image of \(\mathfrak m_S I_a\) is zero and that of \(X\) is
nonzero.  Hence \(c\in\mathfrak m_S\) as well.  Therefore
\(\dim_k I_a/\mathfrak m_S I_a=2\), and \(I_a\) is not principal.
This argument does not require a coefficient field in \(O\).

The only height-one prime containing \(X\) is
\(\mathfrak p=(X,t)\), with \(S/\mathfrak p=O\).  At that prime
\(h\) is a unit, \(t\) is a uniformizer, and
\(X=-t^n/h\) has valuation \(n\).  Thus
\(\operatorname{div}(X)=n\mathfrak p\).
The ideal \(I_a\) has valuation \(a\) at \(\mathfrak p\) and zero
at all other height-one primes, giving
\([I_a]=a[I_1]\), with the same ideal-class sign convention.
Nonprincipality for every \(1\le a<n\) proves exact order \(n\).
For the normality and reflexivity criteria and the identification of
rank-one reflexive modules with divisor classes, see
\citet[Tags 031S, 0AVB, and 0EBM]{Stacks2026}.

For the last assertion, the first divisor class gives a homomorphism
\[
 c_1:G_0(S)\longrightarrow\operatorname{Cl}(S).
\]
It can be defined by the determinant lattice of the torsion-free
quotient of a module, together with its torsion lengths at height-one
primes.  The determinant convention has the opposite sign to the
valuation of an ideal, while the torsion lengths have positive sign.
Localizing an exact sequence at each height-one DVR proves additivity,
and a change of generic basis changes the divisor by a principal one.
This map takes a rank-one reflexive ideal to its class and takes a
free module to zero.  Since \(S\) is local, it kills every perfect
complex and descends to \(K_0D_{\mathrm{sg}}(S)\).
The inclusion into the idempotent completion induces an injection
on \(K_0\), by \citet[Corollary 2.3, p.~6]{Thomason1997}.
Thus a relation \(r[I_1]=0\) after idempotent completion already
holds before it, and applying \(c_1\) proves \(n\mid r\).
No map from the new idempotent objects to \(\operatorname{Cl}(S)\)
is required.
\end{proof}

\subsection{Specialization of a supported formal disc}

Write
\(D_{\mathrm{sg}}(C)=D^b(\operatorname{coh}C)/\operatorname{Perf}(C)\)
and use a superscript \(\mathrm{ic}\) for idempotent completion.

\begin{lemma}\label{lem:prime-to-p-formal-disc}
Let \(S_0\) be a complete Noetherian local hypersurface, put
\(B_0=S_0[[w]]\), and let \(\mathcal T_w\) be the full subcategory
of \(D_{\mathrm{sg}}(B_0)\) consisting of objects perfect after
inverting \(w\).  There is a triangle functor
\[
 \mathcal E:\mathcal T_w^{\mathrm{ic}}
                 \longrightarrow D_{\mathrm{sg}}(S_0)^{\mathrm{ic}}
\]
which takes an \(S_0\)-module \(N\), regarded as a \(B_0\)-module
killed by \(w\), to the class of that same module \(N\).
\end{lemma}

\begin{proof}
Consider the supported quotient
\[
 \mathcal C_w=
 D^b_{V(w)}(\operatorname{coh}B_0)/
                  \operatorname{Perf}_{V(w)}(B_0).
\]
Its natural functor to \(D_{\mathrm{sg}}(B_0)\) is fully faithful.
Here is the required factorization check.  A supported bounded coherent
complex has a bounded model of coherent \(w\)-power-torsion modules;
one power \(w^N\) kills every term.  Any morphism to it from a bounded
finite-projective complex \(P^\bullet\) is represented by a chain map
and factors through \(P^\bullet/w^NP^\bullet\).  The latter is
represented by tensoring \(P^\bullet\) with the two-term Koszul
complex on \(w^N\), since \(w\) is regular on \(B_0\).  It is perfect
and supported on \(V(w)\).  The Verdier-quotient factorization
criterion now gives full faithfulness.  The supported derived model
and this criterion are \citet[Lemmas 2.1, 2.5, and 2.6,
preprint pp.~3--5]{Orlov2011};
the argument just given is affine and requires no coefficient field.

The image is dense in \(\mathcal T_w\).  Over the local hypersurface
\(B_0\), represent an object of \(\mathcal T_w\), up to shift, by a
finite maximal Cohen--Macaulay module \(L\).  The usual
matrix-factorization description of the singularity category
identifies its endomorphisms with stable module maps
\citep[equation (1.1), p.~2, and Appendix B, pp.~37--40]{BurkeWalker2012}.
Its stable endomorphism module is finite over \(B_0\): it is the
quotient of \(\operatorname{Hom}_{B_0}(L,L)\) by the image of
\(\operatorname{Hom}_{B_0}(L,B_0)\otimes_{B_0}L\).
This description commutes with localization.  After inverting \(w\),
the object is perfect, so its stable identity is zero.  Therefore
some \(w^N\) kills that identity already over \(B_0\).
The cone of \(w^N:L\to L\) has \(w\)-power-torsion cohomology,
and in the singularity category it is isomorphic to \(L\oplus L[1]\).
This proves density, hence
\[
 \mathcal C_w^{\mathrm{ic}}\simeq\mathcal T_w^{\mathrm{ic}}.
\]

Under restriction of scalars, every coherent \(w\)-power-torsion
module is finite over \(S_0\).  This gives a functor from the supported
bounded coherent category to \(D^b(\operatorname{coh}S_0)\).
A supported perfect \(B_0\)-complex has finite flat amplitude over
\(S_0\), because \(B_0\) is \(S_0\)-flat.  Its restricted cohomology
is bounded and finite, so over the Noetherian ring \(S_0\) it is
perfect.  Restriction therefore induces a triangle functor
\(\mathcal C_w\to D_{\mathrm{sg}}(S_0)\).
Extend it to idempotent completions and use the displayed equivalence.
On a module killed by \(w\) this functor is exactly the asserted
restriction of scalars.
\end{proof}

\subsection{A detector defined for every lattice}

Continue with \(O\) complete and use the surface \(S\) in
\eqref{eq:prime-to-p-surface}.  Set
\[
 Q=O[[t,X]],\qquad f=t^n,\qquad
 R=Q/(f,g)=A_n[P].
\]
Here \(f,g\) is a regular sequence.  Indeed \(t^n\) is regular,
and in \(A_n[[X]]\) multiplication by \(g\) is injective: its lowest
nonzero \(X\)-coefficient is the \(O\)-non-zero-divisor \(p\).
Distinguished-polynomial division identifies the quotient with
the finite group ring.

Form the generic hypersurface
\[
 Y=\operatorname{Proj}Q[T_f,T_g]/(fT_f+gT_g).
\]
Let \(\beta:\mathbb P^1_R\hookrightarrow Y\) and
\(\varpi:\mathbb P^1_R\to\operatorname{Spec}R\) be the natural maps.
The mixed-characteristic generic-hypersurface theorem gives mutually
inverse triangle equivalences
\begin{equation}\label{eq:prime-to-p-generic-equivalence}
 F=\beta_*\varpi^*:D_{\mathrm{sg}}(R)\longrightarrow D_{\mathrm{sg}}(Y),
 \qquad
 \Psi=R\varpi_*\beta^\flat,\qquad
 \beta^\flat=L\beta^*(-1)[-1].
\end{equation}
These are \citet[Theorem 2.10, p.~8; the paragraph after Corollary 2.11,
p.~9; and the formula preceding (4.3), p.~19]{BurkeWalker2012}.
Their hypotheses require a regular Noetherian ambient ring of finite
dimension and a regular sequence; they impose no equicharacteristic
assumption.  All those hypotheses hold here.

\begin{lemma}\label{lem:prime-to-p-detector}
There is a homomorphism
\[
 \rho:K_0(\Lat_{A_n}(P))
                  \longrightarrow K_0(D_{\mathrm{sg}}(S)^{\mathrm{ic}})
\]
defined on the full exact category.  It vanishes on \(A_n[P]\)
and on the trivial lattice \(A_n\), and satisfies
\begin{equation}\label{eq:prime-to-p-values}
 \rho([M_i])=-[I_{n-i}]\qquad(0<i<n).
\end{equation}
If \(\alpha_a\) is the automorphism \(\sigma\mapsto\sigma^a\),
\(a\in\mathbb F_p^\times\), then
\begin{equation}\label{eq:prime-to-p-normalizer-values}
 \rho(\alpha_a^*x_{n-1})=
 \begin{cases}
 -[I_1],&a=1,\\
 0,&a\ne1.
 \end{cases}
\end{equation}
\end{lemma}

\begin{proof}
We first verify that a single specialization construction is defined
for every source lattice.  If \(M\) is finite free over \(A_n\), its
twisted cyclic resolution has free \(R\)-terms \(R\otimes_{A_n}M\),
and
\[
 \operatorname{Ext}^q_R(M,M)
  =H^q(P,\operatorname{Hom}_{A_n}(M,M)),\qquad q>0.
\]
These groups are two-periodic, killed by \(p\), and of bounded finite
\(O\)-length.  The even and odd cohomology-operator modules are finite
over \(R[\chi_f,\chi_g]\), by
\citet[Remark 3.19, p.~17]{BurkeWalker2012}.  Discard finitely many
initial degrees so that \(p\) kills both modules; this does not change
their projective support.

Put \(C=R/pR\), an Artinian local ring with residue field \(k\).
The two modules are now finite graded \(C[\chi_f,\chi_g]\)-modules
whose component lengths are bounded.  Reducing modulo the maximal
ideal of \(C\) gives finite graded \(k[\chi_f,\chi_g]\)-modules with
bounded Hilbert function, hence zero-dimensional projective support.
The maximal ideal of \(C\) is nilpotent, so reduction does not change
the underlying support.  The support is therefore a finite set of
closed points of \(\mathbb P^1_k\).  This uses neither a coefficient
field in \(O\) nor any perfection assumption on \(k\).

The stable-Ext sheaf comparison identifies this support with the
nonperfect locus of \(F(M)\)
\citep[Corollaries 3.18 and 3.20, pp.~17--18;
the end of Section 8.1, pp.~33--34]{BurkeWalker2012}.
One can test nonperfectness by the stable self-End sheaf: its stalk
contains the identity precisely when the corresponding local object
is nonzero in the singularity category.  Thus \(F(M)\) has finite
nonperfect support for every lattice \(M\).

On the chart \(T_g=1\), put \(z=T_f/T_g\).  Complete at the rational
closed point \((\pi,t,X,z-1)\), and write \(w=z-1\).  The completed
ring is
\[
 \widehat Y=O[[t,X,w]]/(g(X)+(1+w)t^n).
\]
Because \(n\) is a unit, recursive division by \(n\) gives a unique
\(v(w)\in1+wO[[w]]\) with \(v(w)^n=1+w\).
The change \(T=v(w)t\) identifies this completed ring with
\[
 S[[w]],\qquad S=O[[X,T]]/(g(X)+T^n).
\]
It is the identity in the \(t\)-coordinate on \(w=0\).
The completed \(F(M)\) is perfect outside its closed point, since
its original nonperfect locus was finite and perfection is preserved
under the flat completion map.  In particular it is perfect after
inverting \(w\).  Apply Lemma~\ref{lem:prime-to-p-formal-disc}.
This defines \(\rho\).

The functor \(M\mapsto\beta_*\varpi^*M\) is exact on coherent modules:
\(\varpi\) is flat and \(\beta\) is a closed immersion.  Completion,
the coordinate change, and the functor of
Lemma~\ref{lem:prime-to-p-formal-disc} are triangle functors.
Consequently every coefficient-split exact triple of lattices gives
an additive relation under \(\rho\).  This establishes the
whole-category assertion.

\paragraph{Value on the explicit lattices.}
Over \(S\), set \(N_i=\operatorname{coker}\Phi_i\) for \(0<i<n\).
The map
\[
 a\longmapsto t^{n-i},\qquad b\longmapsto X
\]
identifies \(N_i\) with \(I_{n-i}\).
To verify the full presentation, lift a relation
\(u t^{n-i}+vX=0\) to \(O[[X,t]]\), and subtract a multiple of
\(t^n+Xh\).  The remaining relation on the regular sequence
\(t^{n-i},X\) is a multiple of \((X,-t^{n-i})\).
Thus all relations are generated by the columns
\((X,-t^{n-i})\) and \((t^i,h)\) of \(\Phi_i\).
Since \(S\) is a domain, \(f=t^n\) is regular on \(N_i\), and
the presentation gives \(N_i/fN_i=M_i\).

Let \(i_1:\operatorname{Spec}S\hookrightarrow Y\) be the section
\(T_f=T_g\), or \(z=1\).  It is globally closed and lies entirely
in the chosen chart.  On this chart \(\beta\) is cut out by \(f\).
Thus the inverse in \eqref{eq:prime-to-p-generic-equivalence} gives
\[
 \Psi(i_{1*}N_i)\simeq M_i[-1].
\]
There is no higher Tor because \(f\) is regular on \(N_i\);
\(\mathcal O(-1)\) is trivial on the section, and \(\varpi\) restricts
to an isomorphism on it.  Applying the inverse equivalence yields
\[
 F(M_i)\simeq i_{1*}I_{n-i}[1]
       \quad\text{in }D_{\mathrm{sg}}(Y).
\]
After completion this is \(I_{n-i}[1]\), killed by \(w\).
Lemma~\ref{lem:prime-to-p-formal-disc} gives
\eqref{eq:prime-to-p-values}.  The regular module \(R=A_n[P]\)
has value zero because \(\beta\) is a regular closed immersion.

The trivial \(P\)-module \(A_n=R/(X)\) also has value zero.
Near \(z=1\), \(z\) is a unit, so \(F(A_n)\) is \(Y/(X)\):
the equation modulo \(X\) forces \(zt^n=0\), hence \(f=0\).
The element \(X\) is regular on this chart.  In the regular
ambient ring, reduction of \(g+zt^n\) modulo \(X\) is the nonzero
element \(zt^n\) of the domain \(O[[t]][z]\); cancelling \(X\)
in a lifted relation proves the assertion.  Hence \(Y/(X)\)
has a length-one locally free resolution and is perfect there.

\paragraph{Normalizer separation.}
For the representative \(1\le a\le p-1\), set
\[
 \varphi_a(X)=(1+X)^a-1,\qquad
 U_a(X)=\sum_{r=0}^{a-1}(1+X)^{pr}.
\]
Then \(\varphi_a\) is an invertible formal substitution,
\(U_a\) is a unit with residue \(a\), and
\[
 g(\varphi_a(X))=U_a(X)g(X).
\]
The generic-hypersurface automorphism on the chart is
\[
 (X,z)\longmapsto(\varphi_a(X),\,U_a(X)z);
\]
it multiplies the defining equation by \(U_a\) and respects the
\(\beta,\varpi\) diagram.  Thus \(F\) is natural for this automorphism.
Pullback of support at the residual point \(z=1\) is support at
\(z=a^{-1}\).  For \(a\ne1\) that is a different rational point
over \(k\), so the completed twisted \(F(M_i)\) is perfect at the
selected point.  For \(a=1\) it is the original object.
This proves \eqref{eq:prime-to-p-normalizer-values}.
The inverse exponent results from pullback; it is not replaced
by \(a\) in this argument.
\end{proof}

\subsection{Exact order, induction, and coefficient transfer}

\begin{proof}[Proof of Theorem~\ref{thm:prime-to-p-all-dvr}]
First suppose \(O\) is complete and put \(\eta=x_{n-1}\).
Lemma~\ref{lem:prime-to-p-mesh} proves \(n\eta=0\) and both kernel
assertions.  Lemma~\ref{lem:prime-to-p-detector} gives
\(\rho(\eta)=-[I_1]\).
If \(r\eta=0\), Lemma~\ref{lem:prime-to-p-surface} therefore gives
\(n\mid r\).  Thus \(\eta\) has exact order \(n\).

Set
\[
 \xi=\operatorname{Ind}_{P}^{S_p}\eta,\qquad
 \Theta=\rho\circ\operatorname{Res}_{P}^{S_p}.
\]
In Mackey's formula, double cosets outside the normalizer have
trivial intersection with \(P\).  Their terms factor through
restriction to the trivial group, where \(\eta\) is zero since
both its lattices have coefficient rank \(p\).
The normalizer quotient is \(\mathbb F_p^\times\) and supplies
each twist once.  Hence
\[
 \Theta(\xi)
   =\sum_{a\in\mathbb F_p^\times}\rho(\alpha_a^*\eta)
   =-[I_1].
\]
Induction preserves all the mesh exact sequences, so \(n\xi=0\).
The last equality and Lemma~\ref{lem:prime-to-p-surface} show that
\(\xi\) has exact order \(n\).  Reduction and generic vanishing
commute with induction.

The detector kills every Young permutation lattice in degree \(p\).
For the partition \((p)\), its cyclic restriction is the trivial
module, whose value is zero.  In every other partition of \(p\),
each part is less than \(p\), so the Young subgroup has order
prime to \(p\).  The \(P\)-action on its cosets is free, and its
restricted permutation module is a sum of regular \(A_n[P]\)-modules.
These also have value zero.  Since \(\Theta(\xi)\ne0\), generation
fails in degree \(p\).

For an arbitrary DVR \(O\), use the same finite presentations,
coefficient splittings, and induction to define \(\xi_O\).
Its order divides \(n\), and both kernel properties were established
without completeness in Lemma~\ref{lem:prime-to-p-mesh}.
Completion carries \(\xi_O\) to the class over \(\widehat O\)
just proved to have exact order \(n\).  Thus \(\xi_O\) has
exact order \(n\) as well.  Pullback of the completed detector
provides the asserted detector over \(O\).
\end{proof}

\begin{corollary}\label{cor:prime-to-p-part-all-dvr}
Let \(O\) be any characteristic-zero DVR of residue characteristic
\(p\), and let \(N\ge2\).  Write
\[
 N=p^r d,\qquad (d,p)=1.
\]
If \(d>1\), the degree-\(p\) reduction and generic kernels for
\(A_N=O[t]/(t^N)\) contain a class of exact order \(d\), detected
by a homomorphism annihilating every Young permutation lattice.
In particular generation fails for every exponent having a prime
factor different from \(p\).
\end{corollary}

\begin{proof}
Put \(m=p^r\) and \(A_d=O[u]/(u^d)\).
The map
\[
 A_d\longrightarrow A_N,\qquad u\longmapsto t^m
\]
makes \(A_N\) finite free over \(A_d\) on
\(1,t,\ldots,t^{m-1}\).  Indeed the monomials
\(u^a t^b=t^{ma+b}\), with \(0\le a<d\) and \(0\le b<m\),
are exactly the \(O\)-basis monomials of \(A_N\).
Coefficient extension and coefficient restriction preserve
coefficient-free lattices and their coefficient-split exact sequences.
The composite of restriction after extension is \(m\) copies of the
identity functor, because the displayed coefficient basis has trivial
group action.  It therefore induces multiplication by \(m\) on \(K_0\).

Let \(\xi_d\) be the exact-order-\(d\) class of
Theorem~\ref{thm:prime-to-p-all-dvr}, and extend coefficients to
obtain \(\xi_N\).  It satisfies \(d\xi_N=0\), whereas its coefficient
restriction is \(m\xi_d\), of exact order \(d\) since \((m,d)=1\).
Thus \(\xi_N\) has exact order \(d\).
Coefficient reduction and generic extension commute with this map,
so both kernel properties persist.

Compose coefficient restriction with the detector for \(A_d\).
Restriction sends an \(A_N\)-Young lattice to \(m\) copies of the
corresponding \(A_d\)-Young lattice, and so the resulting detector
kills the Young span.  Its value on \(\xi_N\) is \(m\) times
the value on \(\xi_d\), which has exact order \(d\); this value
is nonzero.  The same proof gives an exact-order-\(\ell^{v_\ell(N)}\)
class for every prime divisor \(\ell\ne p\), by taking that maximal
prime-power divisor in place of \(d\).
\end{proof}

\begin{remark}
These arguments do not convert an order-\(d\) subgroup into a split
summand.  The remaining exponent cases not addressed here are pure
powers of the residue prime.  Combining this result with the
already established obstruction when the truncation exponent does
not divide the ramification index leaves only those pure
\(p\)-power exponents which do divide it; the dual-number case in
residue characteristic two is covered separately.
\end{remark}

\begin{corollary}\label{thm:all-dvr-odd-dual}
For every mixed-characteristic DVR $O$ of odd residue characteristic
$p$, the ring $O[\epsilon]/(\epsilon^2)$ has a nonzero class of exact
order two in both the degree-$p$ reduction and generic kernels.
An additive detector on the whole group annihilates the Young span
and detects this class.  No splitting is asserted.
\end{corollary}
\begin{proof}
Apply Theorem~\ref{thm:prime-to-p-all-dvr} with $n=2$.
\end{proof}
\section{Truncations divisible by the residue characteristic}
\label{sec:ramified-vertical-truncations}

A different integral lattice removes the remaining ramification
restriction for non-generation.  Its invariant lattice has covolume
\(n/p\), independently of the valuation of \(p\).
We use the whole-category primary index from
Corollary~\ref{cor:ramified-all-truncations}; the construction and
its exact-category compatibility are recalled below.

\begin{theorem}\label{thm:ramified-vertical-truncations}
Let \(O\) be a characteristic-zero DVR with residue characteristic
\(p>0\), let \(n\ge2\) be divisible by \(p\), and put
\(A=O[t]/(t^n)\), \(r=n/p\).
There is a class \(\xi\in K_0(\Lat_A(S_p))\) which vanishes both
under coefficient reduction \(t\mapsto0\) and after extension to
\(\operatorname{Frac}(O)\), and a homomorphism
\[
 \Theta:K_0(\Lat_A(S_p))\longrightarrow\Z/n
\]
which annihilates every Young permutation lattice and satisfies
\[
 \Theta(\xi)=r\pmod n.
\]
In particular, \(\Phi_A\) is not surjective in degree \(p\).
There is no restriction on the ramification or the residue field,
and \(O\) need not be complete.
\end{theorem}

Since \(0<r<n\), the displayed value is nonzero and has order \(p\).
Theorem~\ref{thm:vertical-ideal-exact-order} below proves that the
source class itself has exact order \(p\), and that its cyclic
subgroup splits from the whole Grothendieck group when \(p^2\nmid n\).

\begin{corollary}\label{cor:all-dvr-truncations-negative}
For every characteristic-zero DVR \(O\) of positive residue
characteristic and every \(n\ge2\), the ring \(O[t]/(t^n)\)
fails polynomial generation in degree equal to the residue
characteristic.
\end{corollary}

\begin{proof}
If \(p\mid n\), apply
Theorem~\ref{thm:ramified-vertical-truncations}.
If \(p\nmid n\), apply Theorem~\ref{thm:prime-to-p-all-dvr}.
\end{proof}

\subsection{The primary covolume on the whole exact category}

We first work over a complete \(O\), with uniformizer \(\pi\),
residue field \(k\), fraction field \(K\), and normalized valuation
\(v=v_\pi\).  Put \(P=C_p=\langle\sigma\rangle\), \(R=A[P]\), and
let \(\mathcal E=\Lat_A(P)\).
A map is stably zero if it factors through a finite free
\(R\)-module.

\begin{lemma}\label{lem:vertical-primary-index}
There is a homomorphism
\[
 \eta:K_0(\mathcal E)\longrightarrow\Z/n
\]
with the following properties.
It vanishes on every lattice obtained by coefficient extension
from \(O\).  If a lattice \(E\), with generator matrix \(T\), has a
lift \(\widetilde T\) over \(O[[t]]\) for which
\[
 U_E=\frac{\widetilde T^p-I}{t^n}\pmod{t^n}
\]
is the scalar \(-1\), then \(E\) has zero integral Tate cohomology
and
\[
 \eta(E)=v(\det E^P)\pmod n.
\]
The determinant is measured using an \(O\)-basis of \(E^P\) and
an \(A_K\)-basis of \((E_K)^P\), where \(A_K=K[t]/t^n\).
If instead \(U_E=-aI\), with \(a\in\{2,\ldots,p-1\}\), then
\(\eta(E)=0\).
\end{lemma}

\begin{proof}
This is the \(u+1\) coordinate of the construction in
Corollary~\ref{cor:ramified-all-truncations}.
We give its choice and exactness checks for the present setting.

For any \(E\in\mathcal E\), choose a lift \(\widetilde T\) and
define \(U_E\) by the displayed formula.
It is an \(R\)-endomorphism, since
\(\widetilde T^p-I\) commutes with \(\widetilde T\).
For an equivariant map \(F:E\to E'\), write
\[
 \widetilde T'\widetilde F-\widetilde F\widetilde T
       =t^n\widetilde B.
\]
Telescoping gives
\begin{equation}\label{eq:vertical-stable-curvature}
 U_{E'}F-FU_E
 =\sum_{i=0}^{p-1}T'^{\,p-1-i}BT^i
 =\sum_{j=0}^{p-1}T'^jHT^{-j},
 \qquad H=BT^{-1}.
\end{equation}
This group trace factors through \(R\otimes_A E\): use
\[
 v\longmapsto\sum_j\sigma^j\otimes T^{-j}v,\qquad
 \sigma^j\otimes w\longmapsto T'^jH(w).
\]
The same calculation compares lifts and bases.  Also the trace
of an equivariant map is \(p\) times that map.  Thus stable Hom
groups are killed by \(p\), and \(U_E\) is stably natural and
independent of those choices.

Let \(E_-\) be the integral Fitting summand on which \(U_E+I\)
is nilpotent modulo \(\pi\).  Its complement has \(U_E+I\)
invertible.  The Fitting decomposition modulo \(\pi\) lifts
uniquely over the complete DVR by compatible Chinese-remainder
projectors on \(E/\pi^qE\).  These projectors are polynomials in
\(U_E\), commute with the \(A\)- and \(P\)-actions, and give
\(A\)-free summands, since \(A\) is local.  For a sequence with
intertwined endomorphisms, a common annihilating polynomial gives
compatible projectors and an exact sequence of these summands.
This uses only a linear primary polynomial, so it requires no
finiteness or perfection assumption on \(k\).

The resulting summand is independent of the lift up to stable
isomorphism.  Write \(p=u\pi^e\) with \(u\in O^\times\).
On \(E_-\), some power of \(U_E+I\) is divisible by \(\pi\).
Its \(e\)-th power is consequently divisible by \(p\), so is
stably zero.  Division by \(\pi\) or \(p\) preserves equivariance
because the lattices have no such torsion.
On the complementary summand \(U_E+I\) is invertible.
Stable naturality in \eqref{eq:vertical-stable-curvature}
therefore makes every stable map between a selected summand and
an unselected summand zero.  The identity comparing two choices,
and its inverse, are consequently block diagonal in the stable
category and induce stable inverse maps on the selected summands.

The curvature on \(E_-\) is invertible, since its residue is
primary at \(-1\ne0\).  We check integral Tate acyclicity.
Choose an \(A\)-basis adapted to the summand and lift it to
\(O[[t]]\).  Off-diagonal blocks of the conjugated generator
lift are divisible by \(t^n\).  Every contribution leaving a
diagonal block and returning is divisible by \(t^{2n}\).
Thus its diagonal block \(\widetilde T_-\) has curvature
congruent to \(U_E|_{E_-}\) modulo \(t^n\), and this curvature
\(\widetilde U_-\) is a unit over \(O[[t]]\).
Write \(\operatorname{Nm}(Z)=I+Z+\cdots+Z^{p-1}\).
If \((T_--I)x=0\), lift \(x\) and write
\((\widetilde T_--I)x=t^ny\).  Multiplication by
\(\operatorname{Nm}(\widetilde T_-)\) and cancellation of \(t^n\)
give
\[
 \widetilde U_-x=\operatorname{Nm}(\widetilde T_-)y.
\]
Since \(\widetilde U_-\) commutes with the lift, \(x\) is a norm
boundary modulo \(t^n\).
If instead \(\operatorname{Nm}(T_-)x=0\), write
\(\operatorname{Nm}(\widetilde T_-)x=t^ny\) and multiply by
\(\widetilde T_--I\).  This gives
\[
 x=(\widetilde T_--I)\widetilde U_-^{-1}y
 \pmod{t^n}.
\]
Both Tate groups therefore vanish.

For any cohomologically trivial coefficient-free lattice \(V\),
the invariant module \(V^P\) is \(O\)-free.  The projector
\(\operatorname{Nm}(T)/p\) makes \((V_K)^P\) projective, hence
free, over the local ring \(A_K\).  Define its covolume
\[
 \iota(V)=v(\det V^P)\pmod n
\]
using an \(A_K\)-basis and the associated \(K\)-basis
\(1,t,\ldots,t^{n-1}\) in each coordinate.
An \(O\)-basis change has determinant of valuation zero.
For \(B\in\operatorname{GL}_d(A_K)\), the induced map on each
successive \(t\)-adic quotient is its constant matrix \(B_0\),
so
\[
 \det_K(B)=\det(B_0)^n.
\]
Thus the generic basis ambiguity is a multiple of \(n\).

The covolume is additive on admissible sequences of
cohomologically trivial lattices.  Their invariant modules
form a short exact sequence because \(H^1(P,V')=0\) for the
left term.  Compatible generic bases and \(O\)-bases give a
block triangular determinant.  It vanishes on a free \(R\)-module,
whose invariants have group-norm elements as an \(A\)-basis.
It is also invariant under stable isomorphism.
Indeed, for stable inverse maps \(f:V\to W\) and \(g:W\to V\),
write \(gf-I=ba\) through a free module \(Q_0\).
Then \((f,a):V\to W\oplus Q_0\) is split by \((g,-b)\).
The complementary projection is
\[
 \begin{pmatrix}I-fg&fb\\-ag&I+ab\end{pmatrix}.
\]
Every entry is stably zero; hence its image is projective.
The ring \(R\) is local, since its residue algebra is
\(k[P]\simeq k[X]/X^p\).  The complement is therefore free,
so the two lattices become isomorphic after adjoining free
modules, whose covolumes vanish.

Set \(\eta(E)=\iota(E_-)\).
Stable independence of \(E_-\) and the preceding argument prove
choice independence.  For an admissible exact sequence, choose
an \(A\)-splitting and lift the generator in upper triangular
blocks.  The curvature is upper triangular with the endpoint
curvatures on its diagonal.  The compatible Fitting projectors
give an admissible sequence
\[
 0\longrightarrow E'_-\longrightarrow E_-
       \longrightarrow E''_-\longrightarrow0.
\]
All three terms are cohomologically trivial, so their covolumes
add.  Choice independence permits these compatible lifts for
every original sequence.  This proves the homomorphism on the
whole \(K_0(\mathcal E)\).
For a constant \(O\)-lattice, a constant lift has curvature zero,
so the selected summand is zero.
Finally, \(U_E=-I\) selects all of \(E\), whereas
\(U_E=-aI\), \(a\ne1\) in \(\F_p^\times\), selects none.
\end{proof}

\subsection{A finite ideal lattice over every DVR}

The witness construction is defined over \(O\) before completion.
Choose a uniformizer \(\pi\), set \(n=pr\), and form
\[
 C=O[t],\qquad
 F(Y,t)=(1+Y-t^r)^p-1+t^n,\qquad
 S=C[Y]/(F).
\]
Write
\begin{equation}\label{eq:vertical-polynomial-coefficients}
 \begin{split}
 F(Y,t)&=Y^p+\sum_{j=0}^{p-1}q_j(t)Y^j,\\
 q_0(t)&=(1-t^r)^p-1+t^n,\\
 q_j(t)&=\binom pj(1-t^r)^{p-j}\qquad(1\le j<p).
 \end{split}
\end{equation}
Every \(q_j(t)\) is divisible by \(p\), hence by \(\pi\).
For \(q_0\), this follows either by expansion or from
\((1-t^r)^p\equiv1-t^{pr}\pmod p\).
In particular \(F(Y,t)\equiv Y^p\pmod\pi\).
Put
\[
 N=(\pi,Y)\subset S.
\]

\begin{lemma}\label{lem:vertical-ideal-lattice}
The \(C\)-module \(N\) is free with basis
\[
 e_0=\pi,\qquad e_j=Y^j\quad(1\le j\le p-1).
\]
The module \(M=N/t^nN\) is a free \(A\)-module of rank \(p\).
Multiplication by \(1+Y-t^r\) gives a \(P\)-action on \(M\)
and has an integral polynomial lift \(\widetilde T\) satisfying
\begin{equation}\label{eq:vertical-exact-curvature}
 \widetilde T^p=I-t^nI.
\end{equation}
Under \(X=Y-t^r\), there is an equivariant inclusion
\begin{equation}\label{eq:vertical-ideal-inclusion}
 M\hookrightarrow R=A[X]/(g(X))
\end{equation}
with image \((\pi,X+t^r)\), and its quotient is \(k[t]/t^n\).
\end{lemma}

\begin{proof}
Since \(F\) is monic, \(S\) is \(C\)-free on the powers
\(1,Y,\ldots,Y^{p-1}\).
The span of \(\pi,Y,\ldots,Y^{p-1}\) is stable under \(Y\):
the only relation to check is \(Y^p=-\sum_jq_jY^j\), whose
constant term is divisible by \(\pi\).
This span is therefore precisely the ideal \((\pi,Y)\).
Its generators are \(C\)-independent because \(\pi\) is a
nonzerodivisor on \(C\).
This proves both freeness assertions.

The defining equation of \(S\) gives
\((1+Y-t^r)^p=1-t^n\).
It acts on the \(C\)-free module \(N\), proving
\eqref{eq:vertical-exact-curvature}; modulo \(t^n\) this is a
group action.  Explicitly, in the indicated basis,
\[
 \begin{split}
 \widetilde T e_0&=(1-t^r)e_0+\pi e_1,\\
 \widetilde T e_j&=(1-t^r)e_j+e_{j+1}
                       \qquad(1\le j\le p-2),\\
 \widetilde T e_{p-1}
   &=-\frac{q_0}{\pi}e_0
     -\sum_{j=1}^{p-2}q_j e_j
     +(1-t^r-q_{p-1})e_{p-1}.
 \end{split}
\]
The sums indexed from \(1\) to \(p-2\) are empty at \(p=2\).
Every displayed coefficient lies in \(O[t]\).

Since \(S/N=k[t]\), multiplication by \(t^n\) is injective
on \(S/N\).  Tensoring \(0\to N\to S\to S/N\to0\) with
\(C/(t^n)\) therefore preserves injectivity on the left.
The substitution \(X=Y-t^r\) identifies \(S/t^nS\) with
\(A[X]/(g(X))=R\).
It gives \eqref{eq:vertical-ideal-inclusion}, the stated ideal
image, and the quotient \(k[t]/t^n\).
\end{proof}

\subsection{The invariant lattice and the two specialization kernels}

\begin{lemma}\label{lem:vertical-invariant-volume}
For the lattice \(M\) of Lemma~\ref{lem:vertical-ideal-lattice},
\[
 M^P=(\pi,t^r)h(X)\subset A_Kh(X)=(M_K)^P.
\]
Consequently, over complete \(O\),
\begin{equation}\label{eq:vertical-index-value}
 \eta(M)=r\pmod n.
\end{equation}
The cyclic class \(x=[M]-[R]\) vanishes under \(t\mapsto0\)
and over \(K=\operatorname{Frac}(O)\).
\end{lemma}

\begin{proof}
In the regular module \(R=A[X]/(Xh)\), the kernel of \(X\)
is exactly \(Ah(X)\).  For example, divide a representative
by the monic polynomial \(Xh\) to obtain degree less than \(p\).
If its product with \(X\) is zero, that product is a scalar
multiple of \(Xh\); cancellation of \(X\) in \(A[X]\) gives
a scalar multiple of \(h\).  Its coefficient of \(X^{p-1}\)
also proves that \(A\to Ah\) is injective.

Using the inclusion \eqref{eq:vertical-ideal-inclusion},
membership of \(a h(X)\) in \(M\) is tested in
\[
 R/M=k[t]/t^n,\qquad X=-t^r.
\]
Modulo \(\pi\), one has \(h(X)=X^{p-1}\), so its image is
\((-1)^{p-1}t^{r(p-1)}=(-1)^{p-1}t^{n-r}\).
Its annihilator in \(k[t]/t^n\) is \((t^r)\).
Thus \(a h(X)\) belongs to \(M\) exactly when
\(a\in(\pi,t^r)\), proving the formula for the invariants.

After inverting \(\pi\), the ideal \(M\subset R\) becomes all
of \(R_K\); in particular \(h(X)\) is an \(A_K\)-basis for the
invariant module.  Relative to that basis, the quotient of
\(A\) by the invariant lattice is
\[
 A/(\pi,t^r)\simeq k[t]/t^r,
\]
of \(O\)-length \(r\).
Equation~\eqref{eq:vertical-exact-curvature} gives the actual
scalar curvature \(-I\), so
Lemma~\ref{lem:vertical-primary-index} applies and gives
\eqref{eq:vertical-index-value}.

The same generic ideal equality proves \(M_K\simeq R_K\).
For reduction, multiplication by \(t\) is also regular on
\(S/N=k[t]\).  Therefore \(M/tM=N/tN\) injects into
\(O[P]\) as the ideal
\[
 I_\pi=(\pi,X).
\]
Let \(I=XO[P]\) be the augmentation ideal.
The augmentation sequences
\[
 0\longrightarrow I\longrightarrow I_\pi
        \longrightarrow \pi O\longrightarrow0,\qquad
 0\longrightarrow I\longrightarrow O[P]
        \longrightarrow O\longrightarrow0
\]
split over \(O\); the quotients are trivial lattices and
\(\pi O\simeq O\).
Hence \([I_\pi]=[O[P]]\) in \(K_0(\Lat_O(P))\).
This proves the reduction assertion, without asserting that
the two reduced lattices are isomorphic.
\end{proof}

\subsection{Normalizer twists, induction, and descent}

\begin{proof}[Proof of Theorem~\ref{thm:ramified-vertical-truncations}]
First suppose \(O\) is complete.
For \(1\le a\le p-1\), let \(M^{(a)}\) be the twist on which
\(\sigma\) acts by \(T^a\).  Use \(\widetilde T^a\) as a lift.
Equation~\eqref{eq:vertical-exact-curvature} gives
\[
 \frac{\widetilde T^{ap}-I}{t^n}
 =\frac{(1-t^n)^a-1}{t^n}I
 \equiv -aI\pmod{t^n}.
\]
The invariants of \(M^{(a)}\) are the same as those of \(M\),
because \(\sigma^a\) generates \(P\).  Thus
Lemma~\ref{lem:vertical-primary-index} gives
\begin{equation}\label{eq:vertical-normalizer-values}
 \eta(M^{(a)})=
 \begin{cases}
 r\pmod n,&a=1,\\
 0,&a\ne1.
 \end{cases}
\end{equation}
Every twist of \(R\) is a constant lattice, so its value is zero.

Embed \(P\) as a \(p\)-cycle in \(S_p\), and put
\[
 x=[M]-[R],\qquad
 \xi=\operatorname{Ind}_{P}^{S_p}x
     =[\operatorname{Ind}_{P}^{S_p}M]-[A[S_p]],
 \qquad
 \Theta=\eta\circ\operatorname{Res}_{P}^{S_p}.
\]
For a double coset outside the normalizer of \(P\), the
intersection of \(P\) with its conjugate is trivial, since
\(P\) has prime order.
The corresponding Mackey term factors through
\(\operatorname{Res}_{1}^{P}x=0\): both \(M\) and \(R\)
are coefficient-free of rank \(p\).
Inside the normalizer, \(N_{S_p}(P)/P\simeq\F_p^\times\)
gives one copy of every automorphism twist.
Equation~\eqref{eq:vertical-normalizer-values} therefore yields
\[
 \Theta(\xi)=
 \sum_{a\in\F_p^\times}\eta(x^{(a)})
 =r\pmod n.
\]
Changing the convention for conjugate modules replaces the
labels by their inverses, which permutes the same set and
leaves its unique identity contribution unchanged.

Every Young permutation lattice is obtained from an integral
\(O\)-permutation lattice by coefficient extension to \(A\).
The same is true after restriction to \(P\).
Lemma~\ref{lem:vertical-primary-index} therefore gives
\(\Theta=0\) on the Young span.
Restriction and induction preserve coefficient freeness and
coefficient-split exact sequences, so all these maps are maps
on the specified exact Grothendieck groups.
The two kernel statements for \(x\), proved in
Lemma~\ref{lem:vertical-invariant-volume}, commute with
induction and give the assertions for \(\xi\).

For an arbitrary DVR \(O\), construct \(N\), \(M\), and \(\xi\)
by the same finite formulas over \(O[t]\).
The bases, ideal inclusions, augmentation sequences, and
generic equality above do not use completeness.
Extend coefficients to the completion \(\widehat O\).
This is flat, preserves coefficient-split sequences and
coefficient-free objects, and retains the uniformizer,
residue field, and normalized valuations.
Compose the completed \(\Theta\) with this scalar-extension
map.  The invariant quotient remains \(k[t]/t^r\), so the
value on the original \(\xi\) is still \(r\pmod n\).
The Young span is still annihilated.
This proves the theorem without a completeness assumption.
\end{proof}

\begin{example}\label{ex:vertical-pure-power-tests}
The new ideal lattice differs from the standard
\(g+t^n\) mesh of Section~\ref{sec:pure-power-vanishing}.
For \(p=n=3\), \(O=\Z_3[\pi]\), \(\pi^3=3\), its matrix modulo
\(t^3\), in the basis \(\pi,Y,Y^2\), is
\[
 T=
 \begin{pmatrix}
 1-t&0&\pi^2(t-t^2)\\
 \pi&1-t&-3(1-t)^2\\
 0&1&-2(1-t)
 \end{pmatrix}.
\]
The invariant ideal is \((\pi,t)\), and the selected coordinate
is \(1\) modulo \(3\).
For \(p=2,n=4\), \(O=\Z_2[\pi]\), \(\pi^4=2\), a matrix is
\[
 T=
 \begin{pmatrix}
 1-t^2&\pi^3t^2\\
 \pi&-1+t^2
 \end{pmatrix}
 \quad\text{over }O[t]/t^4.
\]
Here the invariant ideal is \((\pi,t^2)\) and the coordinate is
\(2\) modulo \(4\).
For the exact polynomial lift in
\eqref{eq:vertical-exact-curvature}, its upper-right entry is
\(\pi^3(t^2-t^4)\).
These nonzero classes coexist with the vanishing standard mesh
classes proved in Theorem~\ref{thm:pure-power-mesh-vanishing}; they are different integral lattices.
\end{example}

\subsection{Exact order of the vertical-ideal class}

The ideal used in Theorem~\ref{thm:ramified-vertical-truncations}
belongs to a finite chain whose exact sequences determine its order.
This also supplies a split subgroup when the truncation exponent is
not divisible by the square of the residue characteristic.

\begin{theorem}\label{thm:vertical-ideal-exact-order}
Let $O$ be any mixed-characteristic DVR of residue characteristic $p$,
let $\pi$ be a uniformizer, and let $n\geq2$ be divisible by $p$.
Put $s=n/p$, $A=O[t]/t^n$, and $P=C_p$.  Write $X=\sigma-1$ in
$R=A[P]$, and put
\[
 M=(\pi,X+t^s)\subset R,\qquad
 y=[M]-[R]\in K_0(\Lat_A(P)).
\]
Then $y$ and $x=\operatorname{Ind}_P^{S_p}y$ have exact order $p$.
Both lie in the generic and $t$-reduction kernels.  If $p^2\nmid n$,
then each of these classes generates a subgroup $\Z/p$ that is split
in the respective whole lattice Grothendieck group.  For $S_p$ the
retraction can be chosen to annihilate every Young permutation class.
No restriction on the ramification, residue field, or completeness of
$O$ is required.
\end{theorem}

\begin{proof}
Set $Y=X+t^s$ in $R$.  The ring $R$ is free over $A$, with basis
$1,Y,\ldots,Y^{p-1}$, because the change from powers of $X$ is
unitriangular.  Its monic relation in $Y$ has reduction $Y^p$
modulo $\pi$: in residue characteristic $p$ it becomes
\[
 (1+Y-t^s)^p-1=Y^p-t^{ps}=Y^p.
\]
Here $t^{ps}=t^n=0$ in $A$.  Every coefficient below the leading
coefficient of that relation is consequently divisible by $\pi$.

For $0\leq j\leq p$, let
\[
 I_j=(\pi,Y^j)\subset R.
\]
It is free over $A$, with basis
\begin{equation}\label{eq:vertical-ideal-bases}
 \pi Y^q\ (0\leq q<j),\qquad Y^q\ (j\leq q<p).
\end{equation}
To verify this assertion, the indicated span contains $\pi$ and
$Y^j$ and is closed under multiplication by $Y$.  The only last
basis-vector calculation uses the monic relation for $Y^p$; its
coefficients are divisible by $\pi$, as just shown.  Thus the span
is the ideal $I_j$.  Its displayed generators are independent,
since they are diagonal multiples of the basis of $R$ and $\pi$
is a nonzerodivisor on $A$.  The endpoint assertions are
\[
 I_0=R,\qquad I_p=\pi R\cong R.
\]
The second equality again uses $Y^p\in\pi R$.

For $1\leq j<p$ consider the $R$-linear sequence
\begin{equation}\label{eq:vertical-ideal-mesh}
 0\longrightarrow I_j
 \xrightarrow{\ a\mapsto(a,-Ya)\ }
 I_{j-1}\oplus I_{j+1}
 \xrightarrow{\ (b,c)\mapsto Yb+c\ }
 I_j\longrightarrow0.
\end{equation}
Its last map is surjective, since
$YI_{j-1}+I_{j+1}=(\pi,Y^j)=I_j$.
Its first map is injective.  A pair in the kernel of the last map
has the form $(b,-Yb)$, with $b\in I_{j-1}$ and
$Yb\in I_{j+1}$.  Reducing modulo $\pi R$ identifies the ambient
ring with
\[
 (O/\pi)[t,Y]/(t^n,Y^p).
\]
In its free basis $1,Y,\ldots,Y^{p-1}$ over
$(O/\pi)[t]/t^n$, the condition $Y\bar b\in(Y^{j+1})$
forces $\bar b\in(Y^j)$: the only possible extra kernel term
is a multiple of $Y^{p-1}$, which already belongs to $(Y^j)$.
Thus $b\in I_j$, proving exactness in the middle.

The sequence is coefficient-split.  An explicit $A$-linear
section of its last map sends a low basis vector
$\pi Y^q$, $0\leq q<j$, to $(0,\pi Y^q)$, and a high
basis vector $Y^q$, $j\leq q<p$, to $(Y^{q-1},0)$.
These pairs lie in the required ideals and map to their specified
basis vectors.  Therefore \eqref{eq:vertical-ideal-mesh} is an
admissible exact sequence in precisely the lattice category in
question.

Put $z_j=[I_j]-[R]$.  The mesh relations give
\[
 2z_j=z_{j-1}+z_{j+1}\quad(1\leq j<p),\qquad z_0=z_p=0.
\]
Induction on $j$ yields $z_j=jz_1$ and hence $pz_1=0$.
Since $I_1=M$, this says $py=0$.  Exactness of induction gives
$px=0$ as well.

The whole-group detectors in
Theorem~\ref{thm:ramified-vertical-truncations} take $y$ and
$x$ to $s=n/p$ in $\Z/n$.  This value has exact order $p$.
Consequently neither class has smaller order, and both have exact
order $p$.  Their generic and reduction vanishings are the ones
proved in that theorem for this same ideal $M$.

Finally suppose that $p^2\nmid n$, so that $s$ is prime to $p$.
Compose either detector with reduction $\Z/n\to\Z/p$, and
multiply by the inverse of $s$ modulo $p$.  The resulting
homomorphism from the whole Grothendieck group takes the
corresponding class to $1$.  Sending $1\in\Z/p$ to that class
is well-defined because its order is $p$, and is a section of the
homomorphism.  The $S_p$ detector still annihilates the Young span,
so the same is true of this retraction.
\end{proof}

\begin{remark}
When $p^2\mid n$, the homomorphism $\Z/n\to\Z/p$ kills
$n/p$.  The exact-order result still applies, but this composition
does not prove a splitting.  The subgroup generated by this
particular class need not be a direct summand merely because its
order is known.
\end{remark}

\section{Two ramified mesh families with vanishing classes}
\label{sec:pure-power-vanishing}

For the two coefficient rings considered here, the standard mesh
with determinant \(g(X)+t^n\) has the regular-lattice class at every
vertex.  We prove this using explicit coefficient-split sequences.
Polynomial generation nevertheless fails for these rings by
Theorem~\ref{thm:ramified-vertical-truncations}.  The calculation
therefore identifies a specific family of zero classes within those
negative cases.

Throughout the mesh argument, let \(O\) be a complete
characteristic-zero DVR with residue characteristic \(p\), and put
\[
 A=O[t]/(t^n),\qquad P=C_p=\langle\sigma\rangle,
 \qquad R=A[P],\qquad X=\sigma-1,
\]
where \(n\ge2\).  Write
\[
 g(X)=(1+X)^p-1=Xh(X),\qquad
 h(X)=p+\binom p2X+\cdots+X^{p-1}.
\]
For \(0\le i\le n\), let \(M_i\) have generators \(a,b\) and
presentation
\begin{equation}\label{eq:pure-power-mesh-presentation}
 \Phi_i=\begin{pmatrix}X&t^i\\-t^{n-i}&h(X)\end{pmatrix},
 \qquad Xa=t^{n-i}b,\qquad t^ia+h(X)b=0.
\end{equation}

\begin{theorem}\label{thm:pure-power-mesh-vanishing}
In each of the following cases,
\[
 \begin{array}{c|c|c}
 p&n&O\\ \hline
 3&3&\mathbb Z_3[\pi],\quad\pi^3=3\\
 2&4&\mathbb Z_2[\pi],\quad\pi^4=2,
 \end{array}
\]
one has
\[
 [M_i]=[A[C_p]]\quad\text{in }K_0(\Lat_A(C_p))
 \qquad(0\le i\le n).
\]
Their induced differences also vanish in
\(K_0(\Lat_A(S_p))\).
\end{theorem}

The displayed rings are the complete DVRs defined by the indicated
Eisenstein polynomials.  The theorem concerns these particular
presentations and their induction images.

\subsection{The coefficient-split mesh}

The determinant of \eqref{eq:pure-power-mesh-presentation} is
\(g+t^n\), so its cokernel over \(A\) is annihilated by \(g\).
It is \(A\)-free of rank \(p\), with basis
\[
 a,\qquad X^jb\quad(0\le j\le p-2).
\]
The first relation specifies \(Xa\), and the monic second relation
specifies \(X^{p-1}b\).  Conversely, these prescriptions define
the operator on the displayed free module and satisfy both relations,
giving the asserted presentation.  At \(i=0\) eliminate \(a=-hb\);
at \(i=n\) eliminate \(b=Xa\).  The remaining relation is \(g=0\),
so \(M_0\simeq M_n\simeq R\).

In these bases, the maps
\[
 \alpha_i:M_i\longrightarrow M_{i+1},\quad
 \alpha_i=\operatorname{diag}(t,I_{p-1}),\qquad
 \beta_i:M_{i+1}\longrightarrow M_i,\quad
 \beta_i=\operatorname{diag}(1,tI_{p-1})
\]
are equivariant by substitution in the relations.  Their composites
are multiplication by \(t\).  For \(1\le i<n\), they give
\begin{equation}\label{eq:pure-power-mesh}
 0\longrightarrow M_i
 \xrightarrow{(\beta_{i-1},\alpha_i)}M_{i-1}\oplus M_{i+1}
 \xrightarrow{(\alpha_{i-1},-\beta_i)}M_i\longrightarrow0.
\end{equation}
On coefficient vectors \((r,v)\), the maps are
\[
 (r,v)\longmapsto((r,tv),(tr,v)),\qquad
 ((r_L,v_L),(r_R,v_R))\longmapsto(tr_L-r_R,v_L-tv_R).
\]
The kernel equations give exactness.  A coefficient left inverse
of the first map and a coefficient right inverse of the second are
\[
 ((r_L,v_L),(r_R,v_R))\longmapsto(r_L,v_R),\qquad
 (r,v)\longmapsto((0,v),(-r,0)).
\]
Thus \eqref{eq:pure-power-mesh} is admissible.  Setting
\(x_i=[M_i]-[R]\), the mesh and endpoints give
\begin{equation}\label{eq:pure-power-mesh-relations}
 x_i=i x_1,\qquad nx_1=0,\qquad x_{n-1}=-x_1.
\end{equation}

\subsection{An elementary finite-module criterion}

Set
\[
 C=O[[t]],\qquad S=C[X]/(g(X)+t^n),\qquad J=(X,t)\subset S.
\]
The ring \(S\) is \(C\)-free with basis \(1,X,\ldots,X^{p-1}\),
and \(S/t^nS=R\).  Moreover \(J\) is \(C\)-free with basis
\(t,X,\ldots,X^{p-1}\): reducing \(X^p\) produces a constant
multiple of \(t^n\), and these displayed elements are independent
in the \(C\)-basis of \(S\).
The relations on this basis give
\begin{equation}\label{eq:pure-power-ideal-lift}
 J\simeq\operatorname{coker}
 \begin{pmatrix}X&t^{n-1}\\-t&h(X)\end{pmatrix},
 \qquad a\longmapsto t,\quad b\longmapsto X,
 \qquad J/t^nJ\simeq M_{n-1}.
\end{equation}
Write \(\varepsilon:S\to S/J=O\) for evaluation at \(t=X=0\).

\begin{proposition}\label{prop:pure-power-finite-module-criterion}
Suppose a finite \(S\)-module \(N\), free over \(C\), has an
\(S\)-linear endomorphism \(b\) fitting into an exact sequence
\begin{equation}\label{eq:pure-power-normalization-sequence}
 0\longrightarrow N\xrightarrow{b}N
 \xrightarrow{\nu}O\longrightarrow0,
\end{equation}
where \(S\) acts on \(O\) through \(\varepsilon\).
Then \([M_i]=[R]\) in \(K_0(\Lat_A(C_p))\) for every
\(0\le i\le n\).
\end{proposition}

\begin{proof}
Form the finite \(S\)-module
\[
 E=S\times_O N
   =\{(s,u):\varepsilon(s)=\nu(u)\}.
\]
Its two projections give exact sequences
\begin{equation}\label{eq:pure-power-fiber-product-sequences}
 \begin{split}
 0&\longrightarrow J\longrightarrow E\longrightarrow N
       \longrightarrow0,\qquad j\longmapsto(j,0),\\
 0&\longrightarrow N\longrightarrow E\longrightarrow S
       \longrightarrow0,\qquad u\longmapsto(0,b(u)).
 \end{split}
\end{equation}
The second kernel identification uses
\(\ker\nu=\operatorname{im}b\) and the injectivity of \(b\).

Here are coefficient splittings.  Choose a \(C\)-basis
\(u_1,\ldots,u_d\) of \(N\), put \(a_j=\nu(u_j)\in O\),
and define
\[
 \lambda\!\left(\sum_j f_j(t)u_j\right)
       =\sum_j f_j(t)a_j\in C\subset S.
\]
Then \(\varepsilon\lambda=\nu\), so
\(u\mapsto(\lambda(u),u)\) splits \(E\to N\) over \(C\).
Choose \(u_*\in N\) with \(\nu(u_*)=1\).  If
\(s=\sum_{j=0}^{p-1}s_j(t)X^j\), then
\(s\mapsto(s,s_0(t)u_*)\) is a \(C\)-linear section of
\(E\to S\).  Thus both sequences are coefficient-split and
\(E\) is \(C\)-free.

Reduction modulo \(t^n\) preserves these splittings.  Using
\eqref{eq:pure-power-ideal-lift}, write
\(\overline E=E/t^nE\) and \(\overline N=N/t^nN\).
The resulting admissible sequences in \(\Lat_A(C_p)\) are
\[
 \begin{gathered}
 0\longrightarrow M_{n-1}\longrightarrow\overline E
       \longrightarrow\overline N\longrightarrow0,\\
 0\longrightarrow\overline N\longrightarrow\overline E
       \longrightarrow R\longrightarrow0.
 \end{gathered}
\]
Their relations give \([M_{n-1}]=[R]\).  Now
\eqref{eq:pure-power-mesh-relations} gives \(x_1=0\), and then
\(x_i=0\) for every \(i\).
\end{proof}

\subsection{The two finite algebras}

For \(p=n=3\) and \(\pi^3=3\), define
\begin{equation}\label{eq:pure-power-cubic-polynomial}
 P_3(v,t)=v^3+\pi^2(1-t)v^2+\pi(1-t)^2v+t(t-1).
\end{equation}
Direct expansion gives
\begin{equation}\label{eq:pure-power-cubic-identity}
 g(\pi v-t)+t^3=\pi^3P_3(v,t).
\end{equation}
Consequently \(N=C[v]/(P_3)\), which is \(C\)-free of rank three,
is an \(S\)-algebra by \(X\mapsto\pi v-t\).
Multiplication by \(v\) has determinant, up to sign,
\(P_3(0,t)=t(t-1)\ne0\).  Since \(C\) is a domain and \(N\)
is \(C\)-free, this multiplication is injective.  Also
\[
 N/vN=C/(t(t-1))=C/(t)=O,
\]
with \(X=t=0\) on the quotient.  Thus
\begin{equation}\label{eq:pure-power-cubic-normalization}
 0\longrightarrow N\xrightarrow{v}N
       \longrightarrow S/J\longrightarrow0
\end{equation}
satisfies the finite-module criterion.

For \(p=2,n=4\) and \(\pi^4=2\), put
\[
 q(t)=-\pi^2t-(1+\pi^2)t^2
\]
and define
\begin{equation}\label{eq:pure-power-quartic-polynomial}
 P_4(z,t)=z^2+\pi(1+q(t))z
       +t(1+t)\bigl((1+\pi^2)t^2-1\bigr).
\end{equation}
The substitution
\begin{equation}\label{eq:pure-power-quartic-coordinate}
 X=\pi^3z+q(t)
\end{equation}
gives the identity
\begin{equation}\label{eq:pure-power-quartic-identity}
 g(\pi^3z+q(t))+t^4=\pi^6P_4(z,t).
\end{equation}
Indeed its linear coefficient in \(z\) is
\(2\pi^3(1+q)=\pi^7(1+q)\), and its constant coefficient is
\[
 q(t)^2+2q(t)+t^4
       =\pi^6t(1+t)\bigl((1+\pi^2)t^2-1\bigr).
\]
Thus \(N=C[z]/(P_4)\) is a \(C\)-free \(S\)-algebra of rank two.
Its constant term
\begin{equation}\label{eq:pure-power-quartic-constant}
 P_4(0,t)=t(1+t)\bigl((1+\pi^2)t^2-1\bigr)
\end{equation}
is \(t\) times a unit of \(C\).  The multiplication determinant
therefore proves that \(z\) is injective on \(N\), and
\(N/zN\simeq O\), with \(X=t=0\).  We obtain
\begin{equation}\label{eq:pure-power-quartic-normalization}
 0\longrightarrow N\xrightarrow{z}N
       \longrightarrow S/J\longrightarrow0.
\end{equation}
The criterion applies again and proves
Theorem~\ref{thm:pure-power-mesh-vanishing}.  Induction to \(S_p\)
preserves coefficient-free modules and coefficient-split sequences,
and therefore preserves these zero differences.

\subsection[A surface divisor for the unit one plus pi]{A surface divisor for the unit \(1+\pi\)}

The following proposition describes a different integral parameter.
It is a statement about a normal surface and its divisor class group;
the exact-category mesh theorem above has already been proved.

\begin{proposition}\label{prop:pure-power-unit-surface}
Let \(n=p^a\), \(a\ge1\), and let \(O\) be a complete
characteristic-zero DVR with uniformizer \(\pi\) and
\(v_\pi(p)\ge2\).  Set \(c=1+\pi\) and
\[
 S_c=O[[X,t]]/(g(X)+ct^n),\qquad J_c=(X,t).
\]
Then \(S_c\) is a normal two-dimensional local domain, and
\([J_c]\) has exact order \(n\) in \(\operatorname{Cl}(S_c)\).
\end{proposition}

\begin{proof}
As a monic polynomial in \(t\) over \(O[[X]]\), the equation
\(t^n+Xh(X)/c\) is Eisenstein at \((X)\): the element \(h(0)=p\)
is a unit in \(O[[X]]_{(X)}\).  The quotient is finite free over
\(O[[X]]\), so irreducibility over its fraction field proves that
\(S_c\) is a domain.  It is a hypersurface of dimension two and
is Cohen--Macaulay.

Put \(r=p^{a-1}\) and \(Y=X+t^r\).  Modulo \(\pi\), the defining
equation is \(Y^p\).  Thus the unique height-one prime over
\(\pi\) is the image of \((\pi,Y)\).  In the regular ambient
ring localized at this prime, \(t\) is a unit, and
\begin{equation}\label{eq:pure-power-unit-vertical}
 g(X)+(1+\pi)t^n\equiv\pi t^n\pmod{(\pi,Y)^2}.
\end{equation}
All intermediate binomial coefficients lie in \((p)\subseteq
(\pi^2)\).  The pure \(t^n\) terms cancel for odd \(p\); at
\(p=2\), their sum is \(2t^n\in(\pi^2)\).
The right side of \eqref{eq:pure-power-unit-vertical} is outside
the square of the ambient maximal ideal, so this height-one
localization is regular.

At a height-one prime not containing \(\pi\), the derivative
\(g'(X)=p(1+X)^{p-1}\) is a unit.  Here \(1+X\) is already a
unit in the formal ambient ring.  These localizations are therefore
regular as well.  Serre's \(R_1+S_2\) criterion proves normality
\citep[Tag 031S]{Stacks2026}.

For \(1\le b\le n\), put \(I_b=(X,t^b)\).
Its quotient \(S_c/I_b\simeq O[t]/t^b\) has depth one, so the
depth lemma gives depth two for \(I_b\).  At height-one primes
the ideal is torsion-free over a DVR.  Hence it is a rank-one
reflexive module, and represents a divisor class
\citep[Tags 0AVB and 0EBM]{Stacks2026}.

The prime \(J_c\) has quotient \(O\).  At its localization,
\(p\), \(h(X)\), and \(c\) are units, so \(Xh=-ct^n\) gives
\[
 v_{J_c}(t)=1,\qquad v_{J_c}(X)=n.
\]
It is the only height-one prime containing \(X\).  Consequently
\[
 \operatorname{div}(X)=nJ_c,\qquad
 [I_b]=b[J_c]\quad\text{in }\operatorname{Cl}(S_c),
\]
using a consistent ideal-class convention.

For \(1\le b<n\), the element \(t^b\) is nonzero modulo \(X\),
since that quotient is \(O[t]/t^n\).  Also \(X\notin(t^b)\):
setting \(t=0\) leaves the nonzero class of \(X\) in
\(O[[X]]/(Xh)\), because \(h\) is a nonunit.
If \((X,t^b)\) were principal, write \(X=ur\), \(t^b=us\),
and \(u=\alpha X+\beta t^b\) for a generator \(u\ne0\).
Then \(1=\alpha r+\beta s\), so one of \(r,s\) is a unit in
the local ring.  One of \(X,t^b\) would generate the ideal,
contrary to the two preceding observations.  Thus every \(I_b\),
\(1\le b<n\), is nonprincipal.  Together with
\(\operatorname{div}(X)=nJ_c\), this proves exact order \(n\).
\end{proof}

\begin{remark}[Scope of the surface calculation]
\label{rem:pure-power-missing-specialization}
The failure of polynomial generation for the two rings in
Theorem~\ref{thm:pure-power-mesh-vanishing} follows from
Theorem~\ref{thm:ramified-vertical-truncations}.  The separate
question addressed here is whether the fixed surface class of
Proposition~\ref{prop:pure-power-unit-surface} itself yields a
whole-lattice detector through the corresponding parameter chart.
The proposition alone does not construct that specialization functor.

The completed family at this parameter is
\[
 O[[X,t,w]]/(g(X)+(c+w)t^n).
\]
Trivializing it by an integral rescaling of \(t\) would require
\(u(w)^n=1+w/c\) with \(u(w)\in1+wO[[w]]\).  Its linear
coefficient would be \(1/(nc)\), which is not integral.
Derived restriction of a module killed by \(w\) is the
two-term zero-differential complex on that module and has zero
Grothendieck class.

The ramified substitution \(w=c((1+v)^n-1)\) does trivialize the
family after a finite free cover of degree \(n\).  The inverse
image of \(w=0\) has coefficient ring
\[
 O[[v]]/((1+v)^n-1),
\]
free of rank \(n\) over \(O\).  Replacing \(t\) by the unit
multiple \((1+v)t\) preserves the ideal \((X,t^b)\); forgetting
\(v\) therefore gives \(n\) copies of that ideal.
In the Grothendieck group of
\(D_{\mathrm{sg}}(S_c)=D^b(\operatorname{mod}S_c)/
\operatorname{Perf}(S_c)\), the \(t\)-filtration of
\(S_c/(X,t^b)\) gives \([I_b]=b[J_c]\).
Since \(I_n=(X)\simeq S_c\), one has \(n[J_c]=0\).
Thus this cover kills the displayed torsion witness.  These
calculations delimit the two proposed specialization procedures;
they do not change the negative generation conclusion cited above.
\end{remark}

\subsection{The relation between the vertical ideal and the standard mesh}

The vertical-ideal class and the standard mesh respond differently
to ramification.  An exact relation between them explains both a
family of vanishings and a failure of splitting for a particular
nonzero class.

\begin{theorem}\label{thm:vertical-standard-relation}
Let $O$ be a mixed-characteristic DVR of residue characteristic $p$,
with uniformizer $\pi$ and $h=v_\pi(p)$.  Let $n=ps$ with $s\geq1$,
and put $A=O[t]/t^n$, $R=A[C_p]$, $X=\sigma-1$, and $Y=X+t^s$.
Let
\[
 y=[(\pi,Y)]-[R],\qquad z=[M_1]-[R],
\]
where the standard family $M_i$ is defined by
\eqref{eq:pure-power-mesh-presentation}.  Then
\begin{equation}\label{eq:vertical-standard-relation}
 h y=s z\quad\text{in }K_0(\Lat_A(C_p)).
\end{equation}
The same equality holds after induction to $S_p$.

If $p\mid h$, then $s z=0$.  In particular, if $n=p$ and
$p\mid h$, every standard-mesh difference $[M_i]-[R]$ vanishes.
If $p\nmid h$ and $p^2\mid n$, then $y$ and its induction are
nonzero elements of order $p$ which are divisible by $p$ in their
respective whole Grothendieck groups.  The cyclic subgroup generated
by either of these particular classes is consequently not a direct
summand.
\end{theorem}

\begin{proof}
The proof takes place over the original DVR.  No passage to a
normalization or a completion is needed.  For $0\leq a\leq h$ put
\[
 J_a=(\pi^a,Y)\subset R.
\]
The monic relation for $Y$ has every lower coefficient divisible by
$p$: modulo $p$ it is $Y^p-t^n=Y^p$.  It follows that $J_a$ is
$A$-free of rank $p$, with basis
\begin{equation}\label{eq:vertical-power-basis}
 \pi^a,Y,Y^2,\ldots,Y^{p-1}.
\end{equation}
Indeed, reducing modulo $\pi^a$ identifies the ambient algebra
with $(A/\pi^aA)[Y]/Y^p$.  The inverse image of the ideal $(Y)$
consists precisely of elements whose constant coefficient in the
monic $Y$-basis is divisible by $\pi^a$.  This proves spanning in
\eqref{eq:vertical-power-basis}; independence follows because
$\pi$ is a nonzerodivisor on $A$.  We have $J_0=R$ and
$J_1=(\pi,Y)$.

Write $H(X)=g(X)/X$, where $g(X)=(1+X)^p-1$.  There is a polynomial
$V(t)\in O[t]$ with constant term $-1$ such that
\begin{equation}\label{eq:vertical-parameter-unit}
 t^{n-s}-H(-t^s)=pV(t).
\end{equation}
For odd $p$, the leading term of $H(-t^s)$ is $t^{n-s}$,
and every other binomial coefficient in $H$ is divisible by $p$.
For $p=2$, the left side is $2(t^s-1)$, giving $V=t^s-1$.
These computations prove the identity and its asserted constant
term at every prime.  Thus $V$ is a unit in $A$, where $t$ is
nilpotent.  Evaluating the group relation at $X=-t^s$ gives
\[
 g(-t^s)=-t^sH(-t^s)=p t^sV(t),
\]
since $t^n=0$.  Consequently
\begin{equation}\label{eq:vertical-power-quotient}
 R/(Y)\cong A/(p t^s)\cong A/(\pi^h t^s).
\end{equation}

For $1\leq a<h$ the sequence
\begin{equation}\label{eq:vertical-power-mesh}
 0\longrightarrow J_a
 \xrightarrow{u\mapsto(u,\pi u)}J_{a-1}\oplus J_{a+1}
 \xrightarrow{(v,w)\mapsto\pi v-w}J_a
 \longrightarrow0
\end{equation}
is exact.  Injectivity of the first map is immediate, and
surjectivity of the second follows from
$\pi J_{a-1}+J_{a+1}=J_a$.  To check its kernel, let
$v\in J_{a-1}$ satisfy $\pi v\in J_{a+1}$.  Modulo $Y$,
this is a colon computation in the ring
\eqref{eq:vertical-power-quotient}.  In $A$ it reads
\[
 \pi b\in\pi^{a+1}A+\pi^h t^sA.
\]
Cancel the nonzerodivisor $\pi$ in $A$ to obtain
\[
 b\in\pi^aA+\pi^{h-1}t^sA=\pi^aA,
\]
because $a\leq h-1$.  Therefore $v\in J_a$, proving exactness
in the middle.  A coefficient section sends the basis vector
$\pi^a$ to $(\pi^{a-1},0)$ and $Y^q$, $1\leq q<p$, to
$(0,-Y^q)$.  Thus \eqref{eq:vertical-power-mesh} is admissible.

Writing $c_a=[J_a]-[R]$, these relations and $c_0=0$ give
$c_a=a c_1=a y$ for $0\leq a\leq h$.  This also holds when
$h=1$, when there is no interior mesh.  In particular
\begin{equation}\label{eq:vertical-power-endpoint}
 [J_h]-[R]=h y,\qquad J_h=(p,Y).
\end{equation}

We identify that endpoint with a standard-mesh lattice.  Put
\[
 B(X,t)=t^{n-s}-H(X).
\]
Polynomial division by the monic linear polynomial $Y=X+t^s$
and \eqref{eq:vertical-parameter-unit} give
\[
 B=pV+YW\quad\text{for some }W\in A[X].
\]
Since $V$ is a unit, $(B,Y)=(p,Y)=J_h$.
The standard lattice $M_s$ has relations
\[
 Xe=t^{n-s}f,\qquad t^se+H(X)f=0.
\]
The assignment
\begin{equation}\label{eq:vertical-standard-isomorphism}
 e\longmapsto B,\qquad f\longmapsto Y
\end{equation}
respects them, because in $R$
\[
 XB=Xt^{n-s}-g(X)=t^{n-s}Y,\qquad
 t^sB+H(X)Y=t^n+g(X)=0.
\]
Its image is $J_h$.  It is an isomorphism: the source basis is
$e,X^jf$ for $0\leq j\leq p-2$, the target basis is
$p,Y,\ldots,Y^{p-1}$, and the last $p-1$ image vectors are
$Y(Y-t^s)^j$.  They give a triangular change of basis with
unit diagonal on the positive powers of $Y$, while the first
image has constant coordinate $V$ in the target basis.
The determinant of the full coefficient matrix is therefore $V$,
a unit.  This proves \eqref{eq:vertical-standard-isomorphism}
as an isomorphism of coefficient-free $R$-modules.

The standard mesh gives $[M_s]-[R]=s z$ by
\eqref{eq:pure-power-mesh-relations}.  Combining this with
\eqref{eq:vertical-power-endpoint} proves
\eqref{eq:vertical-standard-relation}.  Induction is exact and
preserves this equality.

By Theorem~\ref{thm:vertical-ideal-exact-order}, $py=0$ and
both $y$ and its induced class have exact order $p$.
If $p\mid h$, the left side of
\eqref{eq:vertical-standard-relation} is zero, giving $s z=0$.
For $n=p$, one has $s=1$, so $z=0$ and then every standard
mesh difference is zero.

Finally suppose $p\nmid h$ and $p^2\mid n$.  Choose an integer
$b$ with $bh\equiv1\pmod p$.  Since $py=0$, the relation gives
\[
 y=bh y=bs z
    =p\left(b\frac{s}{p}z\right).
\]
The same equality holds after induction, with a nonzero left side
in both cases.  Every homomorphism to $\Z/p$ kills a
$p$-divisible element.  If the cyclic subgroup generated by
that element were a direct summand, its projection, followed by
identification of its generator with $1\in\Z/p$, would contradict
this fact.  This proves the nonsplitting assertion for the specified
classes.
\end{proof}

\begin{remark}
The nonsplitting assertion concerns the subgroup generated by the
vertical-ideal class.  For example, over $\Z_p$ the whole generic
kernel already contains split copies of $\Z/n$.  If $p^2\mid n$,
the relation is $y=(n/p)z$, so the nonzero subgroup generated by
$y$ lies inside such a cyclic group as a $p$-divisible subgroup.
There is no claim here that the whole Grothendieck group has no
other split subgroup of order $p$.
\end{remark}

\begin{corollary}\label{cor:vertical-primary-split}
In the notation of Theorem~\ref{thm:vertical-standard-relation}, write
$n=p^k m$ with $p\nmid m$, and suppose $p\nmid h$.
The classes $w=mz$ and $w_{S_p}=\operatorname{Ind}_{C_p}^{S_p}w$
have exact order $p^k$ and generate subgroups split in the respective
whole Grothendieck groups.  They lie in both the generic and the
$t$-reduction kernels.  The symmetric-group retraction annihilates
the Young span.  The subgroup generated by $y$, or by its induction,
is the order-$p$ subgroup of the corresponding displayed cyclic
group of order $p^k$.
\end{corollary}

\begin{proof}
Let $\eta$ and $\Theta$ be the whole-group detectors of
Theorem~\ref{thm:ramified-vertical-truncations}, normalized to
take $y$ and its induction to $s=n/p$ in $\Z/n$.
Applying them to $hy=sz$, and to its induced equality, gives
\[
 s(\eta(z)-h)=0,\qquad
 s(\Theta(\operatorname{Ind}z)-h)=0\quad\text{in }\Z/n.
\]
Since $n=ps$, both detector values are congruent to $h$ modulo $p$.
Their reductions modulo $p^k$ are therefore units.  The standard
mesh gives $nz=0$, hence $p^k w=0$, and the same relation after
induction.  Their detector values modulo $p^k$ are units because
$p\nmid m$.  Thus both source classes have exact order $p^k$.
Reducing each detector modulo $p^k$ and multiplying by the inverse
of its value on the corresponding class gives a whole-group
retraction, with section sending $1$ to that class.

For the kernel assertions, all standard lattices become regular
after inverting $p$.  Indeed, $H(X)/p$ is the trivial central
idempotent of $A[1/p][C_p]$.  On that factor $X=0$ and $H=p$,
so the standard relations eliminate $f$; on the complementary
factor $H=0$ and $X$ is invertible, so they eliminate $e$.
Both give a free rank-one module over the corresponding factor.
For $0<i<n$, reduction at $t=0$ gives
$O\oplus O[X]/H(X)$.  Multiplication by $X$ identifies the second
summand with the augmentation ideal of $O[C_p]$.  The
coefficient-split augmentation sequence therefore makes the reduced
Grothendieck class equal to $[O[C_p]]$.  This proves both vanishings
for $z,w$ and their inductions.

Finally choose $b$ with $bh\equiv1\pmod p$.  The exact order of
$y$ and the relation give
\[
 y=bhy=bsz=b p^{k-1}w,
\]
and the same equality after induction.  As $p\nmid b$, these are
the unique order-$p$ subgroups of the displayed cyclic groups.
For $k>1$ their nonsplitting in the whole group was proved in
Theorem~\ref{thm:vertical-standard-relation}.
\end{proof}

\section{All square-zero extensions of the localized integers}
\label{sec:arbitrary-idealizations}

Fix a prime $p$ and put $O=\mathbb Z_{(p)}$.  For an arbitrary
$O$-module $M$, without a finite-generation hypothesis, write
\[
 O\ltimes M=O\oplus M,\qquad
 (a,m)(b,n)=(ab,an+bm).
\]
This is a local characteristic-zero ring with square-zero ideal $M$
and residue field $\mathbb F_p$.

\begin{theorem}\label{thm:arbitrary-idealizations}
For every $O$-module $M$, the polynomial map generated by the trivial
symmetric-group classes is an isomorphism over $O\ltimes M$ if and
only if $M$ is a $\mathbb Q$-vector space.  In every negative case,
surjectivity already fails in degree $p$.

If $M$ has nonzero torsion, then already in symmetric-group degree
$p$ the kernel of augmentation
\[
 K_0\bigl(\operatorname{Lat}_{O\ltimes M}(S_p)\bigr)
 \longrightarrow K_0\bigl(\operatorname{Lat}_{O}(S_p)\bigr)
\]
contains a direct summand isomorphic to a countably infinite direct
sum of copies of $\mathbb Z$.
In particular this assertion includes the divisible torsion module
$M=\mathbb Q/O$.
\end{theorem}

\subsection{An integer detector for the divisible torsion extension}

Put $D=\mathbb Q/O$ and $B=O\ltimes D$.  Fix
$H=C_p=\langle s\rangle$.  If $E$ is a finite $B$-free $B[H]$-module,
put $E_0=E/DE$.  Freeness gives a canonical identification
$DE=D\otimes_OE_0$ and canonical exact sequences
\begin{equation}\label{eq:divisible-two-extensions}
 0\longrightarrow D\otimes_OE_0\longrightarrow E\longrightarrow E_0
 \longrightarrow0,
 \qquad
 0\longrightarrow E_0\longrightarrow\mathbb Q\otimes_OE_0
 \longrightarrow D\otimes_OE_0\longrightarrow0.
\end{equation}
Use the two-periodic cochain complex with even differential $S-I$
and odd differential $I+S+\cdots+S^{p-1}$.  Write its cohomology as
$H^j$, with the indicated identification $H^{j+2}=H^j$.
The two groups $H^0(E_0),H^1(E_0)$ are finite-dimensional
$\mathbb F_p$-spaces: they are finitely generated over $O$ and are
killed by $p$, using the norm on invariants and the identity
\[
 (S-I)\sum_{i=0}^{p-2}(p-1-i)S^i
 =I+S+\cdots+S^{p-1}-pI.
\]

Let $\delta_E^j:H^j(E_0)\to H^{j+1}(D\otimes E_0)$ and
$\beta_E^j:H^j(D\otimes E_0)\to H^{j+1}(E_0)$ be the boundaries
of~\eqref{eq:divisible-two-extensions}.  Define
\begin{equation}\label{eq:divisible-u}
 u_E^j=\beta_E^{j+1}\delta_E^j:
 H^j(E_0)\longrightarrow H^{j+2}(E_0)=H^j(E_0).
\end{equation}
Both boundary maps are canonical.  The second is an isomorphism
because the rational middle complex is acyclic, although only its
naturality is needed below.

We check compatibility with admissible exact sequences explicitly.
Choose a $B$-linear splitting of $0\to E'\to E\to E''\to0$.
The cochain differential has upper-triangular block form; denote its
constant diagonal blocks by $d_0',d_0''$, its constant off-diagonal
block by $U_0$, and its $D$-valued blocks by $d_1',d_1'',U_1$.
The square-zero part of $d^{j+1}d^j=0$ is
\[
 d_1'{}^{j+1}U_0^j+U_0^{j+1}d_1''{}^j
 =-d_0'{}^{j+1}U_1^j-U_1^{j+1}d_0''{}^j.
\]
On an integral cocycle the last term vanishes and the first term
on the right is a boundary in $D\otimes E_0'$.  Thus, writing
$\partial_0,\partial_D$ for the connecting maps of the reduced
and $D$-coefficient sequences, we obtain
\begin{equation}\label{eq:divisible-delta-compatibility}
 \delta_{E'}\partial_0=-\partial_D\delta_{E''}.
\end{equation}
For the second boundary, represent a $D$-coefficient cocycle $z$
by a rational lift $\widetilde z$.  Its differential is an integral
cocycle $w$, which represents $\beta_{E''}[z]$.  The constant block
identity gives
\[
 d_0'{}^{j+1}U_0^j\widetilde z
 =-U_0^{j+1}d_0''{}^j\widetilde z=-U_0^{j+1}w.
\]
Consequently
\begin{equation}\label{eq:divisible-beta-compatibility}
 \beta_{E'}\partial_D=-\partial_0\beta_{E''}.
\end{equation}
The two signs cancel in the composite, so
$u_{E'}\partial_0=\partial_0u_{E''}$.  All other maps commute with
$u$ by functoriality.

For a monic irreducible $f\in\mathbb F_p[t]$, $f\ne t$, define
\begin{equation}\label{eq:divisible-euler-detector}
 \chi_f(E)=
 \ell_{\mathbb F_p[t]}\bigl(H^0(E_0)_{f\text{-primary}}\bigr)
 -\ell_{\mathbb F_p[t]}\bigl(H^1(E_0)_{f\text{-primary}}\bigr),
 \qquad t=u_E.
\end{equation}
Taking primary parts preserves exactness: a common annihilator
$f^a q$, with $(f,q)=1$, gives a common primary projector by Bezout's
identity.  The six-term periodic exact cycle therefore proves that
$\chi_f$ is an integer-valued additive invariant.  It vanishes on
every lattice extended from $O$, because the first sequence
in~\eqref{eq:divisible-two-extensions} then splits equivariantly and
$\delta_E=0$.  Each class has only finitely many nonzero $\chi_f$.

Let $\eta=1/p+O\in D$.  On the rank-one $B[H]$-module with
$s=1+a\eta$, $a\in\mathbb F_p$, the integral reduction is trivial,
with $H^0=\mathbb F_p$ and $H^1=0$.  The first boundary sends $1$ to
$a\eta$, and the second lifts it to $a/p$ and applies the norm $p$.
Thus $u=a$ on the even group.  More generally, for a matrix $C$ over
$\mathbb F_p$, the action $I+\eta C$ has $u=C$ on its even Tate
group and has zero odd integral Tate group.

\subsection{Separating arbitrary nonzero torsion}

Suppose $M$ has nonzero torsion.  Since $O$ is a discrete valuation
ring, some nonzero $m\in M$ satisfies $pm=0$.  Define
$Om\to D$ by $m\mapsto\eta$.  This extends to an $O$-linear map
\begin{equation}\label{eq:torsion-separating-functional}
 \lambda:M\longrightarrow D.
\end{equation}
Here is an elementary proof of the required extension property.
Use Zorn's lemma on partial extensions.  If a partial map is
defined on $N\subset M$ and $x\notin N$, the ideal
$\{a\in O:ax\in N\}$ is principal, say $(d)$.  If $d=0$, choose
the image of $x$ arbitrarily.  If $d\ne0$, divisibility of $D$
provides $y\in D$ with $dy=\lambda(dx)$.  The formula
$\lambda(n+ax)=\lambda(n)+ay$ is then well-defined on $N+Ox$.
Thus a maximal partial extension has domain $M$.

Put $A=O\ltimes M$.  The map~\eqref{eq:torsion-separating-functional}
induces a coefficient homomorphism $A\to B$.  Let $L=A$ with
$s$ acting as $1+m$, and set
\[
 N=\operatorname{Ind}_{H}^{S_p}L,\qquad
 N_0=A\otimes_O\operatorname{Ind}_{H}^{S_p}O.
\]
This is a representation since $(1+m)^p=1$.
Apply base change to $B$, restriction to $H$, and then
$\chi_{t-1}$.  In the Mackey decomposition, a double coset has
intersection $H$ precisely when its representative normalizes
$H$; all other intersections are trivial.  The normalizer quotient
is $\mathbb F_p^\times$, with each automorphism represented once,
and the corresponding rank-one twists have $u=a$ for
$a\in\mathbb F_p^\times$.  The other summands are induced from
the trivial subgroup and have zero reduced Tate cohomology.
Consequently the composite integer invariant takes
\[
 [N]-[N_0]\quad\hbox{to}\quad1.
\]
It vanishes on all Young permutation classes.  This proves failure
of generation in degree $p$, and also proves that the displayed
class spans a split copy of $\mathbb Z$ in the augmentation kernel.

For odd $p$ this class has zero trace at every permutation.  At a
$p$-cycle, the two induced traces differ by
$\sum_{a=1}^{p-1}am=p(p-1)m/2=0$, since $pm=0$.
At the identity the ranks agree; every other nonidentity cycle type
is not conjugate into $H$, so both induced traces vanish.
For the smallest example $M=\mathbb F_p\epsilon$, the cyclic
line has action $1+\epsilon$.  Its induced trace is therefore the
same as the constant induced line when $p$ is odd, despite the
infinite order detected above.


For the countably infinite direct summand, choose strictly
increasing positive integers $r_i$ coprime to $p-1$, and monic
irreducibles $f_i\in\mathbb F_p[t]$ of degree $r_i$ with nonzero
constant term.  Such irreducibles exist in every positive degree.
Let $C_i$ be a companion matrix, and let $L_i=A^{r_i}$ with action
$I+mC_i$.  Put
\[
 x_i=[\operatorname{Ind}_{H}^{S_p}L_i]-r_i[N_0].
\]
Twisting by $s\mapsto s^a$ replaces $C_i$ by $aC_i$ and its
polynomial by $a^{r_i}f_i(t/a)$.  Equality with $f_i$ forces
$a^{r_i}=1$ by comparison of constant terms, hence $a=1$ because
$(r_i,p-1)=1$.  Different chosen degrees have disjoint scaling
orbits.  The induced invariant family therefore satisfies
\[
 \chi_{f_j}\bigl(\operatorname{Res}_{H}^{S_p}(x_i\otimes_A B)\bigr)
 =\delta_{ij}.
\]
It has finite support on each class and so maps to
$\bigoplus_i\mathbb Z$.  Sending its standard basis vectors to
$x_i$ is a section.  All $x_i$ reduce to zero over $O$, proving the
direct-summand assertion in the theorem.

\subsection{The torsion-free alternatives}

\begin{lemma}\label{lem:padic-separation}
Let $O=\Z_{(p)}$, let $M$ be a torsion-free $O$-module, and let
$m\in M\setminus pM$.  There is an $O$-linear map
\[
 \varphi:M\longrightarrow\Z_p,\qquad \varphi(m)=1.
\]
\end{lemma}

\begin{proof}
Choose an $\F_p$-basis $(\overline m_s)_{s\in S}$ of $M/pM$
containing $\overline m$, and choose representatives $m_s\in M$, with
$m_{s_0}=m$.  These representatives are $O$-linearly independent.
Indeed, divide the coefficients in a nonzero finite relation by their
smallest common power of $p$.  Torsion-freeness allows cancellation of
that power in $M$, and reduction modulo $p$ would then give a nonzero
relation among the chosen basis vectors.  Consequently
\[
 F=\sum_{s\in S}O m_s=\bigoplus_{s\in S}O m_s
\]
is a free submodule of $M$, and $F/pF\to M/pM$ is an isomorphism.

For every $n\geq1$ the inclusion induces an isomorphism
\begin{equation}\label{eq:idealization-completion}
 F/p^nF\xrightarrow{\ \sim\ }M/p^nM.
\end{equation}
For surjectivity, successive approximation modulo $p$ writes an arbitrary
$x\in M$ in the form
\[
 x=f_0+pf_1+\cdots+p^{n-1}f_{n-1}+p^n x_n,
 \qquad f_j\in F,\quad x_n\in M.
\]
For injectivity, if $f\in F\cap p^nM$, injectivity modulo $p$ first
gives $f=pf_1$ with $f_1\in F$.  Cancellation of $p$ in $M$ gives
$f_1\in p^{n-1}M$; iteration yields $f\in p^nF$.

The coordinate projection $F\to O$ at $s_0$ induces compatible maps
$F/p^nF\to O/p^nO$.  Transporting them across
\eqref{eq:idealization-completion} and taking their inverse limit gives
\[
 M\longrightarrow\varprojlim_n O/p^nO=\Z_p.
\]
It sends $m$ to $1$.  Only finite linear combinations occur at every
stage, so the proof imposes no restriction on the cardinality of $S$.
\end{proof}


We include the reduction to the flat dual-number obstruction for
completeness.  Suppose $M$ is torsion-free and $M\ne pM$.  Choose $m\notin pM$.  Lemma~\ref{lem:padic-separation} gives
an $O$-linear map $\ell:M\to\mathbb Z_p$ with $\ell(m)=1$.
Hence there are coefficient maps
\[
 O[\epsilon]/(\epsilon^2)\longrightarrow O\ltimes M
 \longrightarrow\mathbb Z_p[\epsilon]/(\epsilon^2),
 \qquad \epsilon\longmapsto m\longmapsto\epsilon.
\]
The explicit rank-$p$ cyclic lattice and its induced degree-$p$
counterexample from the flat dual-number theorem can be extended
along these maps.  Its parity detector is still nonzero over
$\mathbb Z_p$: the two-periodic construction uses finite integral
Tate groups over the residue field, and the displayed boundary
matrices are unchanged by completion.  Equivalently the proof of
Theorem~\ref{thm:dual-numbers-all-primes} applies verbatim to the
DVR $\mathbb Z_p$.  A Young expression over $O\ltimes M$ would
give one after the second base change, a contradiction.

It remains to prove the positive alternative.  A torsion-free
$O$-module with $M=pM$ is exactly a $\mathbb Q$-vector space,
because multiplication by $p$ is then bijective.  Suppose $M$ is
such a vector space and let $G$ be any finite group.  In an
$A$-basis, every representation can be written
\[
 \rho(g)=(I+c_g)\rho_0(g),\qquad c_g\in\operatorname{Mat}_r(M),
\]
where $\rho_0$ is its reduction over $O$.  Conjugation by $\rho_0$
gives the action on the matrix space, and the representation
identity is the cocycle equation $c_{gh}=c_g+g\cdot c_h$.
Define
\[
 X=\frac1{|G|}\sum_{g\in G}c_g.
\]
Then $g\cdot X=X-c_g$.  Since the ideal is square-zero,
conjugation by $I-X$ gives
\[
 (I-X)\rho(g)(I+X)=\rho_0(g).
\]
Thus every lattice is isomorphic to a scalar extension from $O$.
The polynomial theorem over $O$, together with the augmentation
retraction, proves the polynomial isomorphism over $A$.
The torsion and torsion-free cases exhaust all modules, completing
the proof of Theorem~\ref{thm:arbitrary-idealizations}.

\begin{corollary}\label{cor:all-square-zero-extensions}
Let $I^2=0$ in a commutative ring $A$ and suppose $A/I\simeq
\mathbb Z_{(p)}$.  Then the polynomial map over $A$ is an
isomorphism if and only if $I$ is a $\mathbb Q$-vector space under
its natural $\mathbb Z_{(p)}$-module structure.
If $I$ is finitely generated, this holds if and only if $I=0$.
In particular, for every $n\ge1$,
\[
 \mathbb Z_{(p)}[\epsilon]/(p^n\epsilon,\epsilon^2)
\]
has a split countably generated free abelian summand in the
degree-$p$ augmentation kernel.
\end{corollary}

\begin{proof}
Every integer prime to $p$ is a unit modulo $I$ and therefore is
a unit in $A$.  The localization property supplies a unique
section $\mathbb Z_{(p)}\to A$ of the quotient map.  Thus
$A\simeq O\ltimes I$, and the theorem applies.  If $I$ is finitely
generated and $I=pI$, Nakayama's lemma gives $I=0$.
For the displayed cyclic torsion ideal, one may take
$m=p^{n-1}\epsilon$ in the construction; it has order $p$.
\end{proof}

No classification of nonsplit higher nilpotent extensions or of
arbitrary singular local rings is asserted here.

\section{Quadratic orders and infinitely many torsion classes}
\label{sec:singular-orders}

Failure of generation can already occur over a one-dimensional local
domain whose generic fibre is a field.  The obstruction in this section
is a family of additive invariants of integral involutions.  It records
information that disappears from every ordinary character.

\begin{theorem}\label{thm:singular-order}
Let $O=\Z_{(2)}$ or $\Z_2$, let $r,s\in O$, and suppose
$t^2+2st+4r$ is irreducible over $\operatorname{Frac}(O)$.  Put
\[
 A=O[t]/(t^2+2st+4r).
\]
Then $A$ is a one-dimensional noetherian local characteristic-zero
$\Z$-flat domain, and $\Phi_A$ is not surjective.  More precisely, the
kernel of
\[
 K_0(\Lat_A(C_2))\longrightarrow
 K_0(\operatorname{Rep}_{\operatorname{Frac}(A)}(C_2))
\]
contains a direct summand isomorphic to a countable direct sum of
copies of $\Z/2$.
\end{theorem}

For $r=1$ and $s=0$, one nonzero class of exact order two is
\[
 [L]-[A[C_2]],\qquad L=A^2,\qquad
 T_L=\begin{pmatrix}1+t&2\\2-t&-1-t\end{pmatrix}.
\]


We first construct the invariants, including their exact-sequence
compatibility.  All reductions in the proof are modulo $2$ in $O$;
write $k=\F_2$.

\begin{lemma}\label{lem:order-detectors}
Let $\mathcal F_s$ be the monic irreducible polynomials in $k[z]$
other than $z$, also excluding $z+1$ when $s$ is odd.  For every
$f\in\mathcal F_s$, there is an
additive homomorphism
\[
 \nu_f:K_0(\Lat_A(C_2))\longrightarrow\Z/2
\]
that vanishes on every lattice obtained by extending scalars from $O$.
\end{lemma}

\begin{proof}
Put $n=\operatorname{rank}_A E$ and choose an $A$-basis.  Its involution has a unique
expression $T=P+tQ$ with $P,Q$ matrices over $O$.  Since $1,t$ is an
$O$-basis of $A$, the equation $T^2=I$ gives
\[
 P^2-4rQ^2=I,\qquad PQ+QP=2sQ^2.
\]
Set $D=P-I$.  Then
\begin{equation}\label{eq:order-curvature}
 D^2+2D=4rQ^2,\qquad DQ+QD=2sQ^2-2Q.
\end{equation}
The reduction $d$ of $D$ has square zero on $V=k^n$.
Put $H(E)=\ker d/\operatorname{im}d$.  The reduction of $Q$ induces an
operator $B$ on $H(E)$.  For a class represented by $v\in O^n$ with
$Dv\in2O^n$, define
\[
 C[v]=[Dv/2\bmod2].
\]
This is well-defined.  The vector $Dv/2$ is a cycle because
$D(Dv/2)=2rQ^2v-Dv$ is even.  Changing $v$ by $2w$ changes its image by
the boundary $Dw$; changing its homology class by the boundary $Dw$
changes the image by $D^2w/2=2rQ^2w-Dw$, again a boundary modulo two.
Equation~\eqref{eq:order-curvature} now gives
\begin{equation}\label{eq:order-operators}
 C^2+C=\overline r B^2,\qquad CB+BC=B+\overline s B^2.
\end{equation}
In particular $Z=B^2$ commutes with both operators.

These operators are natural.  If an equivariant map is written
$F=F_0+tF_1$, comparison of coefficients gives
\[
 P'F_0-F_0P=4r(Q'F_1-F_1Q),
 \]
 \[
 Q'F_0-F_0Q=F_1P-P'F_1+2s(Q'F_1-F_1Q).
\]
Thus $\overline F_0$ induces a map on $H$ commuting with $B$; the second
identity makes its commutator a boundary on cycles.  The first identity,
after division by two on cycles, proves commutation with $C$.
In particular the construction is independent of the chosen basis.

Fix $f\in\mathcal F_s$ of degree $a$, and let $H_f(E)$ be the
$f$-primary subspace for $Z$.  Its dimension is divisible by $2a$.
Indeed, filtration by powers of $f(Z)$ reduces the assertion to
modules over $F=k[z]/(f)$ on which $Z$ is $\alpha=z\bmod f$.
Set $U=B+\overline s\alpha$.  Equations~\eqref{eq:order-operators}
give
\[
 CU+UC=U,\qquad U^2=\eta:=\alpha+\overline s\alpha^2,
 \qquad C^2+C=\overline r\alpha.
\]
The exclusions in $\mathcal F_s$ ensure $\eta\ne0$.
This operator algebra is $\operatorname{Mat}_2(F)$: with
$\beta^2=\eta$ in the finite field $F$, send
\[
 U/\beta\longmapsto\begin{pmatrix}1&0\\1&1\end{pmatrix},\qquad
 C\longmapsto\begin{pmatrix}0&1\\\overline r\alpha&1\end{pmatrix}.
\]
The defining relations hold, and the four spanning elements
$1,U,C,UC$ map to independent matrices.
Every finite module therefore
has $F$-dimension divisible by two.  Define
\begin{equation}\label{eq:order-invariant}
 \nu_f(E)=\frac{\dim_k H_f(E)}{2\deg f}\pmod2.
\end{equation}

It remains to prove additivity; dimension of homology alone would not
suffice.  In an $A$-split exact sequence
$0\to E'\to E\to E''\to0$, choose compatible bases and write
\[
 P=\begin{pmatrix}P'&U\\0&P''\end{pmatrix},\qquad
 Q=\begin{pmatrix}Q'&W\\0&Q''\end{pmatrix}.
\]
The square-zero reductions give the exact three-cycle
\begin{equation}\label{eq:order-three-cycle}
 H(E')\longrightarrow H(E)\longrightarrow H(E'')
       \xrightarrow{\partial}H(E').
\end{equation}
The first two maps commute with $B,C$.  The connecting map satisfies
\begin{equation}\label{eq:order-connecting}
 B'\partial=\partial B'',\qquad
 C'\partial=\partial(C''+I).
\end{equation}
For a cycle $v$ in the last term, $\partial[v]=[Uv]$.  The off-diagonal
parts of~\eqref{eq:order-curvature} are
\[
 D'U+UD''+2U=4r(Q'W+WQ''),
 \]
 \[
 Q'U+UQ''=-D'W-WD''-2W+2s(Q'W+WQ'').
\]
Division of the first identity by two on a cycle proves the formula
for $C$ in~\eqref{eq:order-connecting}; the second proves the formula
for $B$ modulo boundaries.

All three maps preserve the primary subspaces for $Z$, and all their
images are invariant under $B,C$.  Their image dimensions are therefore
multiples of $2a$, by the matrix-algebra calculation.  Primary
decomposition preserves exactness in~\eqref{eq:order-three-cycle}.
If its three image dimensions are $u,v,w$, then the sum of its three
term dimensions is $2(u+v+w)$, a multiple of $4a$.
Dividing by $2a$ modulo two gives
$\nu_f(E)=\nu_f(E')+\nu_f(E'')$.
Finally, for a lattice defined over $O$ one has $Q=0$, hence $B=0$ and
$H_f=0$ for every $f\ne z$.
\end{proof}

\begin{proof}[Proof of Theorem~\ref{thm:singular-order}]
For each $f\in\mathcal F_s$, let $a=\deg f$ and choose the
matrix $D_f$ with entries $0,1$ lifting the companion matrix of $f$.  Define
\[
 B_f=\begin{pmatrix}D_f&0\\sD_f^2-D_f&-D_f\end{pmatrix},\qquad
 C_f=\begin{pmatrix}0&I\\rD_f^2&-I\end{pmatrix},\qquad
 T_f=I+2C_f+tB_f.
\]
Direct multiplication gives
$C_f^2+C_f=rB_f^2$ and $B_fC_f+C_fB_f=sB_f^2-B_f$, so that $T_f^2=I$ over
$A$.  Let $L_f=A^{2a}$ with this action.  Its square-zero reduction is
zero, and on $H(L_f)=k^{2a}$ the central operator is
\[
 Z=\operatorname{diag}(\overline D_f^{\,2},\overline D_f^{\,2}).
\]
Squaring is an automorphism of the field $k[z]/(f)$, so
$\overline D_f^{\,2}$ still has irreducible polynomial $f$.  Consequently
\[
 \nu_g(L_f)=\begin{cases}1,&g=f,\\0,&g\ne f.\end{cases}
\]

For every lattice $E$ there is an exact sequence
\[
 0\longrightarrow E\otimes\mathrm{sgn}
 \longrightarrow A[C_2]\otimes_A E
 \longrightarrow E\longrightarrow0,
\]
with diagonal action in the middle.  This middle lattice is a free
$A[C_2]$-module of rank $\operatorname{rank}_A E$, by untwisting the
action on the regular factor.  The matrix
\[
 J_f=\begin{pmatrix}I&0\\-I-tD_f&-I\end{pmatrix}
\]
is invertible and satisfies $J_fT_f=-T_fJ_f$.  Hence
$L_f\simeq L_f\otimes\mathrm{sgn}$ and
\[
 \delta_f=[L_f]-a[A[C_2]],\qquad 2\delta_f=0.
\]
The invariants $\nu_g$ show that these classes are independent and
nonzero.  For each lattice only finitely many $\nu_f$ are nonzero, so
$(\nu_f)_f$ takes values in the direct sum and retracts the inclusion
of the classes $\delta_f$.  They therefore form a direct summand, also
inside the ordinary-character kernel.  There are infinitely many monic
irreducibles in $\mathcal F_s$.
Over the fraction field each $T_f$ has $a$ eigenvalues $1$ and $a$
eigenvalues $-1$: it squares to one and has trace zero.  Thus all
$\delta_f$ have zero ordinary character.  When $r=1$, $s=0$, and $f=z+1$, this is the displayed rank-two example.

In degree two the subring generated by the trivial classes is spanned
by $\rho_2$ and $\rho_1^2=[A[C_2]]$.  Both have action matrices over
$O$, and all $\nu_f$ vanish on them, whereas $\nu_f(L_f)=1$ for every $f\in\mathcal F_s$.
This proves the failure of surjectivity.

Finally, irreducibility over $\operatorname{Frac}(O)$ and the monic
presentation make $A$ a domain, finite free of rank two over $O$.
Its quotient by $2$ is $k[t]/(t^2)$, so it is local with maximal ideal
$\mathfrak m=(2,t)$ and dimension one.  The defining equation has no
linear term in the cotangent space: its embedding dimension is two,
so the ring is singular.  Freeness over $O$ gives $\Z$-flatness.
\end{proof}

\begin{corollary}\label{cor:quadratic-order-classification}
Let $K/\Q_2$ be a quadratic field extension, and let $A$ be a
$\Z_2$-order in $K$.  Then $\Phi_A$ is an isomorphism if and only if
$A$ is the maximal order.  For every nonmaximal order, the
ordinary-character kernel in degree two has a direct summand
isomorphic to a countable direct sum of copies of $\Z/2$.
\end{corollary}

\begin{proof}
Write the maximal order as $O_K=\Z_2\oplus\Z_2\omega$.
Such a basis exists because $1$ is primitive in the free rank-two
$\Z_2$-module $O_K$.  The finite-index submodule
$A/\Z_2\subseteq O_K/\Z_2$ is $2^n\Z_2\omega$ for a unique $n\ge0$.
Consequently $A=\Z_2+2^nO_K$.  If $n=0$, it is a discrete valuation
ring and the positive theorem applies.  If $n\ge1$, the generator
$t=2^n\omega$ has polynomial
\[
 t^2-2^n\operatorname{Tr}_{K/\Q_2}(\omega)t
       +2^{2n}\operatorname{N}_{K/\Q_2}(\omega).
\]
The linear coefficient is divisible by two and the constant
coefficient by four.  The polynomial is irreducible over $\Q_2$,
so Theorem~\ref{thm:singular-order} applies.
\end{proof}

\section{Every nonmaximal quadratic order at an odd prime}
\label{sec:odd-quadratic-orders}

A second coefficient boundary strengthens the parity invariant when the
coefficient ring is a quadratic order.  The useful extra operator is the
canonical homotopy for multiplication by the order of a cyclic group.
It forces certain primary lengths to be even, so that their halves
provide a new invariant modulo two.

\begin{theorem}\label{thm:odd-quadratic-orders}
Let $p$ be an odd prime, let $O=\Z_{(p)}$ or $\Z_p$, and let $d\in p^2O$.
Put $A=O[t]/(t^2-d)$.  Then the degree-$p$ kernel of generic
specialization contains a direct summand isomorphic to
\[
 \bigoplus_{i\ge1}\mathbb Z/2.
\]
Every displayed generator is represented by the difference between an
induced explicit lattice and a multiple of the regular lattice.
In particular $\Phi_A$ is not surjective.
If $O=\Z_p$ and $d$ is nonzero and nonsquare in $\Q_p$, the same direct
summand lies in the kernel of extension to the normalization of $A$.
\end{theorem}

\begin{corollary}\label{cor:odd-quadratic-orders}
For every odd prime $p$, every nonmaximal $\mathbb Z_p$-order in a
quadratic extension of $\mathbb Q_p$ fails polynomial generation.
Its degree-$p$ normalization kernel contains a split countable direct
sum of copies of $\mathbb Z/2$.
\end{corollary}

We first construct the invariant, keeping all coefficient and
exact-sequence compatibilities explicit.  Write $b=d/p^2$.

\subsection{Three odd operators on Tate cohomology}

Let $L$ be an $A$-free $C_p$-lattice and set $V=L/(p,t)L$.
The assumption $d\in p^2O$ gives two natural exact sequences
\begin{align}
 0\longrightarrow V\xrightarrow{p}L/(t,p^2)L
       \longrightarrow V\longrightarrow0,\label{eq:order-beta-seq}\\
 0\longrightarrow V\xrightarrow{t}L/pL
       \longrightarrow V\longrightarrow0.\label{eq:order-delta-seq}
\end{align}
Indeed, the middle coefficient rings are respectively
$O/p^2O$ and $\mathbb F_p[t]/t^2$, and $L$ is free.
Let $\beta$ and $\delta$ be their degree-one Tate connecting maps.
All Tate groups in this subsection are finite-dimensional over
$\mathbb F_p$.

Choose an $A$-basis and write the differentials of the complete
period-two Tate complex as
\[
 d_i=X_i+tY_i,
\]
where $X_i,Y_i$ have entries in $O$.  The differentials alternate
between $T-1$ and $N(T)=1+T+\cdots+T^{p-1}$.  The identity
$d_{i+1}d_i=0$ gives
\begin{equation}\label{eq:order-XY-relations}
 X_{i+1}X_i=-p^2bY_{i+1}Y_i,
 \qquad X_{i+1}Y_i+Y_{i+1}X_i=0.
\end{equation}
For a cycle represented by an integral column $z$ with $X_i z=pw$,
the two connecting maps are
\[
 \beta[z]=[w],\qquad \delta[z]=[Y_i z].
\]
They are the connecting maps of
\eqref{eq:order-beta-seq}--\eqref{eq:order-delta-seq}, so the formulas
are independent of lifts and bases.  For example, replacing $z$ by
$z+ph$ changes $w$ by the boundary $X_i h$; changes by coboundaries
are handled by \eqref{eq:order-XY-relations}.
Those same identities give
\begin{equation}\label{eq:order-beta-delta}
 \beta^2=-\overline b\,\delta^2,
 \qquad \beta\delta+\delta\beta=0.
\end{equation}
In fact $X_{i+1}w=-pbY_{i+1}Y_i z$, while
$X_{i+1}Y_i z=-pY_{i+1}w$, and division by $p$ gives the two formulas.

Both operators anticommute with the Tate connecting map $\partial$
of an admissible sequence of $A$-lattices.  Here is a direct check of
the signs.  Split the sequence over $A$ and write the upper block of
$d_i$ as $Z_i+tW_i$.  The upper blocks of $d_{i+1}d_i=0$ are
\begin{align*}
 X'_{i+1}Z_i+Z_{i+1}X''_i
   &=-p^2b(Y'_{i+1}W_i+W_{i+1}Y''_i),\\
 X'_{i+1}W_i+Y'_{i+1}Z_i+Z_{i+1}Y''_i+W_{i+1}X''_i&=0.
\end{align*}
On $X''_i z=pw$, the first formula, divided by $p$ and reduced
modulo $p$, gives $\beta'\partial=-\partial\beta''$.
The second says that
$\delta'\partial+\partial\delta''$ is null-homotopic, with homotopy
given by the upper blocks $W_i$.

For the third operator put
\[
 Q(z)=\frac{N(z)-p}{z-1}\in\mathbb Z[z].
\]
Define a degree-minus-one homotopy on the complete complex by
\[
 h_i=I\quad(i\text{ even}),\qquad
 h_i=-Q(T)\quad(i\text{ odd}).
\]
The polynomial identity $N(T)-(T-1)Q(T)=pI$ gives
\begin{equation}\label{eq:order-p-homotopy}
 d_{i-1}h_i+h_{i+1}d_i=pI.
\end{equation}
After reduction to $V$, this is an odd cohomology operator, using
periodicity to identify degrees minus one and plus one.  It is natural
in $V$, since its matrices are polynomials in the group action.
The group generator acts as the identity on Tate cohomology, whence
\begin{equation}\label{eq:order-h-square}
 h^2=-Q(1)=-p(p-1)/2=0.
\end{equation}
This uses the assumption that $p$ is odd.

Its relations with the two coefficient boundaries are
\begin{equation}\label{eq:order-h-boundaries}
 \beta h+h\beta=I,\qquad \delta h+h\delta=0.
\end{equation}
For the first, use \eqref{eq:order-p-homotopy} modulo $p^2$ in the
coefficient complex defining $\beta$.  If $dz=pw$, then
$d(hz)=p(z-hw)$, giving the stated formula after division by $p$.
The constant terms of the full homotopy identity differ from this
calculation only by multiples of $d=p^2b$, which vanish after that
division and reduction.
For the second, expand the full homotopy as $h_i=H_i+tK_i$.
The coefficient of $t$ in \eqref{eq:order-p-homotopy} is
\[
 X_{i-1}K_i+Y_{i-1}H_i+H_{i+1}Y_i+K_{i+1}X_i=0.
\]
Thus $\delta h+h\delta$ is null-homotopic, with the indicated $K_i$.

Finally $h$ anticommutes with $\partial$ for an admissible sequence.
For completeness, reduce the sequence modulo $(p,t)$ and write
its differential and homotopy in upper triangular blocks.  If $R_i$
is the upper block of $h_i$, the upper block of $dh+hd=0$ is
\[
 X'_{i-1}R_i+Z_{i-1}h''_i+h'_{i+1}Z_i+R_{i+1}X''_i=0.
\]
It gives $h'\partial+\partial h''=0$ on cohomology.

\subsection{An additive half-length invariant}

Put $u=\delta^2$ and take the invertible Fitting part of each Tate
group.  Fitting extraction is exact on finite-dimensional modules.
On this part $\delta$ is invertible.  Define even operators
\[
 C=\delta^{-1}\beta,\qquad J=\delta h.
\]
They commute with $u$ and with every map in the Tate long exact
sequence: each of their two odd factors anticommutes with connecting
maps.  The preceding identities give
\begin{equation}\label{eq:order-CJ}
 C^2=\overline b,\qquad J^2=0,\qquad CJ+JC=-I.
\end{equation}
In particular $\ker J=\operatorname{im}J$.  If $Jv=0$, then
$v=-JCv$; the reverse inclusion follows from $J^2=0$.
This holds on every $u$-primary submodule and every connecting-image
submodule, because all these submodules are stable under $C$ and $J$.

For a monic irreducible $f\in\mathbb F_p[u]$, $f\ne u$, let
$H^0_f$ denote the $f$-primary part.  The exact sequence
\[
 0\longrightarrow\operatorname{im}J\longrightarrow H^0_f
   \xrightarrow{J}\operatorname{im}J\longrightarrow0
\]
shows that its $\mathbb F_p[u]$ composition length is even.  Hence
\begin{equation}\label{eq:order-half-length}
 \mu_f(L)=\frac12\operatorname{length}_{\mathbb F_p[u]}H^0_f
                    \pmod2
\end{equation}
is defined.  It is additive.  Indeed, exact primary extraction gives a
six-term exact cycle.  Its two connecting-image modules have equal
ordinary composition length, because $\delta$ identifies them, and
each length is even by \eqref{eq:order-CJ}.  The even-parity length
defect is the negative of the sum of these lengths.  Dividing by two
therefore gives an even integer.

Every lattice extended from $O$ has $\delta=0$, so all these invariants
vanish on such lattices and, in particular, on Young permutation
lattices.  Only finitely many primary factors occur for any one lattice.

\subsection{The lattice and its exact order}

Choose a primitive element $\eta$ of the residue field
$k_s=\F_{p^s}$, and lift its monic irreducible polynomial over
$\F_p$ to a monic polynomial $F\in O[X]$.  Put $O_s=O[X]/(F)$.
This is a finite free local DVR over $O$, of degree $s$ and residue
field $k_s$.  Indeed, its reduction modulo $p$ is a field and
integrality forces every maximal ideal to contract to $pO$, so its
unique maximal ideal is $pO_s$.  The monic polynomial $F$ is
irreducible over $\operatorname{Frac}(O)$: a monic factorization
would lie in $O[X]$, since $O$ is integrally closed, and would
contradict irreducibility after reduction.  Thus $O_s$ is a domain.
Every nonzero element is a power of $p$ times a unit, by taking
the minimum valuation of its coefficients in the monic basis.
This proves the DVR assertion without assuming completeness.

We choose units $\alpha,\kappa\in O_s$ with
$\overline\kappa^{-1}=\eta$ and
\begin{equation}\label{eq:order-alpha}
 \alpha^p=1-d/\kappa.
\end{equation}
If $d=0$, take $\alpha=1$ and any lift of $\eta^{-1}$ for $\kappa$.
Otherwise write $d=p^m b_0$ with $m\ge2$ and $b_0\in O^\times$.
Choose a unit $y\in O_s$ lifting $-\eta\overline b_0$, and put
\[
 \alpha=1+p^{m-1}y,\qquad
 \kappa=\frac{d}{1-\alpha^p}.
\]
The quotient belongs to $O_s^\times$.  In fact
\[
 1-\alpha^p=-p^m yU,\qquad
 U=1+\sum_{j=2}^{p}\binom pj p^{j(m-1)-m}y^{j-1}
       \in1+pO_s.
\]
For $2\le j<p$ the summand coefficient has valuation
$(j-1)(m-1)\ge1$; for $j=p$ its valuation is
$p(m-1)-m\ge1$, since $p$ is odd.  Therefore
$\kappa=-b_0/(yU)$ is a unit,
$\overline\kappa^{-1}=-\overline y/\overline b_0=\eta$,
and \eqref{eq:order-alpha} holds exactly.

Let $\Lambda_s=O_s[z]/(z^p-1)$ and put
\begin{equation}\label{eq:order-D-matrix}
 q(z)=\frac{z^p-\alpha^p}{z-\alpha},\qquad
 D=\begin{pmatrix}0&z-\alpha\\\kappa q(z)&0\end{pmatrix}.
\end{equation}
On $E=\Lambda_s^2$, let the generator of $C_p$ act by $z$ and let
$t$ act by $D$.  Then
\[
 D^2=\kappa(z^p-\alpha^p)I=dI,
\]
so this is a module over $A_s=O_s[t]/(t^2-d)$.
It is free of rank $p$ over $A_s$.  Its quotient modulo $(p,t)$ is
\[
 k_s[z]/(z-1)\ \oplus\ k_s[z]/(z-1)^{p-1},
\]
which has dimension $p$.  Nakayama gives a surjection $A_s^p\to E$.
Both modules are free of rank $2p$ over $O_s$, so its kernel is zero.

This also gives an explicit $p$ by $p$ matrix over $A_s$.
Let $C_\alpha$ be the companion matrix of $q$, let $v$ be its first
coordinate column, and let $w$ be the last coordinate row.  In the
basis consisting of the first $\Lambda_s^2$ generator and
$1,z,\ldots,z^{p-2}$ times the second, the group generator is
\begin{equation}\label{eq:order-T-matrix}
 T=\begin{pmatrix}
       \alpha&\kappa^{-1}t w\\
       t v&C_\alpha
    \end{pmatrix}.
\end{equation}
Its characteristic polynomial is
$(z-\alpha)q(z)-\kappa^{-1}t^2=z^p-1$, which verifies $T^p=I$ directly.

For the coefficient boundary of $E$, both Tate parities of
$V=E/(p,t)E$ have dimension two over $k_s$.  In even degree take
$v_1,(z-1)^{p-2}v_2$, and in odd degree take $v_1,v_2$.
Modulo $p$ the off-diagonal entries of $D$ are
$z-1$ and $\overline\kappa(z-1)^{p-1}$.  The difference and norm
calculations therefore give, in both directions,
\[
 \delta=\begin{pmatrix}0&\overline\kappa^{-1}\\1&0\end{pmatrix},
 \qquad u=\overline\kappa^{-1}I.
\]

Let $\sigma$ send $t$ to $-t$.  The quadratic coefficient sequence is
\begin{equation}\label{eq:order-sign-sequence}
 0\longrightarrow E^\sigma\longrightarrow
 A_s\otimes_{O_s}\operatorname{Res}_{O_s}E
 \longrightarrow E\longrightarrow0.
\end{equation}
The final map is evaluation, and the kernel map is
$v\mapsto t\otimes v-1\otimes Dv$.
In the two coefficient coordinates its image is $(-Dv,v)$.
Multiplication by $t$ on the kernel corresponds to $-D$ on the source,
because $D^2=dI$; the sequence is split over $A_s$, since its quotient
is free.  The matrix $\operatorname{diag}(1,-1)$ anticommutes with
$D$ and commutes with $z$, so $E^\sigma\simeq E$.
Its underlying $O_s[C_p]$-module is free of rank two, whence
\begin{equation}\label{eq:order-exact-two}
 2\bigl([E]-[A_s[C_p]]\bigr)=0.
\end{equation}

\subsection{Induction and the direct summand}

Choose distinct positive integers $s_i$ tending to infinity.  In
$k_i=\mathbb F_{p^{s_i}}$ choose a primitive multiplicative element
$\eta_i$, and choose $\alpha_i,\kappa_i$ by the preceding construction.
Let $f_i$ be the
minimal polynomial of $\eta_i$ over $\mathbb F_p$.
For the corresponding lattice, the even Tate group is $k_i^2$ with
$u=\eta_i$, so its $f_i$-primary composition length is two and
$\mu_{f_i}=1$.

These scalar orbits are free.  If
$j\eta_i=\eta_i^{p^l}$ with $j\in\mathbb F_p^\times$ and
$0<l<s_i$, then
\[
 p^{s_i}-1\mid(p-1)(p^l-1),
\]
although the positive integer on the right is smaller.  If $l=0$,
the equality forces $j=1$.  Distinct degrees give disjoint orbits.

Restrict the constructed lattice from $A_{s_i}$ to $A$, call it $L_i$,
and put
\[
 x_i=[\operatorname{Ind}_{C_p}^{S_p}L_i]-s_i[A[S_p]].
\]
Equation~\eqref{eq:order-exact-two} gives $2x_i=0$.
The generator-change comparison for cyclic Tate complexes scales
$u=\delta^2$ by the corresponding element of $\mathbb F_p^\times$.
In the Mackey restriction, the normalizer supplies each such twist
once; every other intersection is trivial and supplies a free
$A[C_p]$-lattice.  Consequently
\[
 (\mu_{f_j}\circ\operatorname{Res}_{C_p}^{S_p})(x_i)=\delta_{ij}.
\]
The finite-support observation following
\eqref{eq:order-half-length} makes the direct sum of these invariants
a retraction of the map sending the standard basis of
$\bigoplus_i\mathbb Z/2$ to the $x_i$.  This proves both exact order
two and the direct-summand assertion.  Their nonzero invariants,
which vanish on every Young lattice, prove non-surjectivity.

For generic vanishing, write $K=\operatorname{Frac}(O)$.
Then $A[1/p]$ is a field, a product of two fields,
or the dual-number algebra over $K$, according as $d$ is
nonsquare, nonzero square, or zero.  The characteristic-zero
coefficient result gives torsion-free lattice Grothendieck groups
in the local cases.  In the product case representations decompose
componentwise and the classes of coefficient-free objects form a
subgroup of the product of the two characteristic-zero representation
groups, so there is again no torsion.  Thus all $x_i$ vanish generically.
When $O=\Z_p$ and $A$ is a domain, its normalization is a DVR and its
symmetric-group lattice Grothendieck group is torsion-free by
Theorem~\ref{thm:regular-local}; normalization extension therefore also
kills the $x_i$.  This completes the proof of
Theorem~\ref{thm:odd-quadratic-orders}.

\begin{proof}[Proof of Corollary~\ref{cor:odd-quadratic-orders}]
Let $O_K$ be the ring of integers of a quadratic extension of
$\mathbb Q_p$.  Since $2$ is a unit, the trace projection
$\frac12\operatorname{Tr}_{K/\mathbb Q_p}$ splits
$O_K=O\oplus O\theta$, where $\theta$ has trace zero.
Thus $\theta^2=c\in O$.  The valuation of $c$ is zero or one:
if $p^2\mid c$, the element $\theta/p$ would be integral and hence
belong to $O_K$, contradicting this integral basis.
If $c$ is a unit, its residue is nonsquare, since otherwise Hensel's
lemma would make $c$ a square in $\mathbb Q_p$.

For any order $B\subset O_K$, the same trace projection shows that
its trace-zero part is a submodule $p^rO\theta$, for some $r\ge0$.
Hence $B=O\oplus p^rO\theta$.  The nonmaximal orders have $r\ge1$
and are isomorphic to
\[
 O[t]/(t^2-p^{2r}c),
\]
whose constant term is divisible by $p^2$.  Apply the theorem.
\end{proof}

\begin{corollary}\label{cor:localized-odd-quadratic-orders}
Let $p$ be odd.  Every nonmaximal finite $\Z_{(p)}$-order in a
quadratic number field is local and fails polynomial generation.
Its degree-$p$ generic-specialization kernel contains a split
countably infinite direct sum of copies of $\Z/2$.
\end{corollary}

\begin{proof}
Put $O=\Z_{(p)}$, and let $D$ be its integral closure in the quadratic
field.  This is finite free of rank two over the DVR $O$.
The trace projection splits $D=O\oplus O\theta$ with $\theta$ of
trace zero, so $\theta^2=c\in O$.  If $p^2$ divided $c$, then
$\theta/p$ would be integral and belong to $D$, contradicting this
basis.  Thus $v_p(c)$ is zero or one.

For any finite order $B\subset D$ containing $O$, half the trace
sends $B$ to $O\subset B$.  Its trace-zero part is therefore
$p^rO\theta$ for some $r\ge0$, and
$B=O\oplus p^rO\theta$.  Nonmaximality means $r\ge1$.
Consequently
\[
 B\simeq O[t]/(t^2-p^{2r}c),
\]
and the quadratic constant lies in $p^2O$.  This finite algebra is
local, since its quotient modulo $p$ is $\F_p[t]/(t^2)$.
Apply Theorem~\ref{thm:odd-quadratic-orders}.

The maximal order $D$ itself can be semilocal when $p$ splits in
the number field.  No nonsquare-residue claim about a unit $c$ is
needed here.
\end{proof}

\begin{example}\label{ex:cubic-quadratic-order}
Let $O=\Z_{(3)}$ or $\Z_3$ and take the nonmaximal ramified
quadratic order $A=O[t]/(t^2-27)=O[3\sqrt3]$.
The explicit choices $\alpha=-8$ and $\kappa=1/19$ satisfy
$\alpha^3=1-27/\kappa$.  They give
\[
 T=\begin{pmatrix}
 -8&0&19t\\
 t&0&-64\\
 0&1&8
 \end{pmatrix},\qquad T^3=I.
\]
The cyclic difference $[A^3,T]-[A[C_3]]$ has exact order two.
Its induction to $S_3$ is detected by $\mu_{u-1}$ after restriction
to $C_3$, since the other normalizer twist has eigenvalue $-1$.
\end{example}

\section{Congruent coefficient algebras, after Maranda}
\label{sec:congruent-coefficients}

Sufficiently close coefficient algebras have the same lattice
Grothendieck group for a fixed finite group.  This transfers the
truncated-polynomial obstruction to equations with small nonzero
coefficients, including totally ramified field orders.
The mechanism is Maranda's averaging and successive approximation:
\citet[Theorem~2, pp.~347--348]{Maranda1953} proves rigidity of
integral representations from sufficiently close reductions, and
his Theorem~3, pp.~348--351, lifts modular subconstituents with the
precision improvement used below.  We give the coefficient-algebra
comparison and its exact-sequence compatibility explicitly, since
these are the assertions needed for the present Grothendieck groups.

\begin{theorem}\label{thm:congruent-coefficients}
Let $A$ and $B$ be nonzero commutative algebras finite free over
$\Z_p$, let $G$ be a finite group, and put $v=v_p(|G|)$.
For an integer $r>2v$, fix a unital algebra isomorphism
\[
 \alpha:A/p^rA\xrightarrow{\sim}B/p^rB.
\]
It determines a canonical rank-preserving bijection between the
isomorphism classes in $\Lat_A(G)$ and $\Lat_B(G)$.  This bijection
preserves direct sums and the triples of classes represented by
coefficient-split short exact sequences.  Hence it induces an
isomorphism
\[
 K_0(\Lat_A(G))\xrightarrow{\sim}K_0(\Lat_B(G)).
\]
The isomorphism respects tensor and exterior powers and representations
defined over $\Z_p$.  It commutes with restriction and induction for
subgroups of $G$.  The assertion concerns isomorphism classes and
exact-sequence relations; no equivalence of the categories of all
morphisms is asserted.
\end{theorem}

\subsection{Correction and rigidity}

We first establish the two estimates that make the comparison
independent of choices.  Let $R$ be a finite free $\Z_p$-algebra,
write $m=|G|$, and suppose that matrices $\rho_g\in\mathrm{GL}_d(R)$
satisfy $\rho_1=I$ and
\[
 D(g,h)=\rho_g\rho_h-\rho_{gh}\in p^s\mathrm{M}_d(R),
 \qquad s>2v.
\]
The matrices need not yet define an action.  Set
\[
 a_g=\frac1m\sum_{k\in G}D(g,k)\rho_k^{-1},\qquad
 \rho'_g=\rho_g-a_g.
\]
Division by $m$ is legitimate here: every numerator entry is
divisible by $p^s$, and $m/p^v$ is a unit.  Thus
$a_g\in p^{s-v}\mathrm{M}_d(R)$.  The new matrices are invertible
and $\rho'_1=I$.

Associativity gives the exact identity
\[
 \rho_gD(h,k)-D(gh,k)+D(g,hk)-D(g,h)\rho_k=0.
\]
Multiply on the right by $\rho_k^{-1}$ and average.  If
\[
 E(g,h)=a_g\rho_h-\frac1m\sum_kD(g,hk)\rho_k^{-1},
\]
reindexing the first sum shows that
\begin{align*}
 E(g,h)
 &=\frac1m\sum_kD(g,hk)
       (\rho_{hk}^{-1}\rho_h-\rho_k^{-1})\\
 &=\frac1m\sum_kD(g,hk)\rho_{hk}^{-1}D(h,k)\rho_k^{-1}.
\end{align*}
Consequently $E(g,h)\in p^{2s-v}\mathrm{M}_d(R)$, and the new
defect satisfies
\begin{equation}\label{eq:congruence-improvement}
 \rho'_g\rho'_h-\rho'_{gh}
   =-E(g,h)+a_ga_h\in p^{2s-2v}\mathrm{M}_d(R).
\end{equation}
This calculation uses no cohomology action associated with an
approximate representation.

Iteration replaces $s$ by
$s_j=2v+2^j(s-2v)$, which tends to infinity.  Completeness and
separatedness of $R$ give an exact action $\sigma$ satisfying
\begin{equation}\label{eq:congruence-corrected}
 \sigma_g\equiv\rho_g\pmod {p^{s-v}}.
\end{equation}
The construction preserves block diagonal and block upper triangular
matrix algebras.  In either case its diagonal blocks are exactly
the separate constructions on the original diagonal blocks.

For rigidity, let $\sigma,\tau$ be exact actions, and suppose an
invertible matrix $C$ satisfies
\[
 \tau_gC-C\sigma_g\in p^s\mathrm{M}_d(R),\qquad s>v.
\]
The averaged intertwiner
\begin{equation}\label{eq:congruence-intertwiner}
 F=\frac1m\sum_{g\in G}\tau_gC\sigma_g^{-1}
\end{equation}
is integral and congruent to $C$ modulo $p^{s-v}$.  It is therefore
invertible.  Reindexing the sum proves $\tau_hF=F\sigma_h$ for
every $h\in G$.  In particular exact actions congruent modulo
$p^s$, for $s>v$, are isomorphic.

\subsection{Transport of objects and every exact relation}

Lift $\alpha$ to a unital $\Z_p$-linear isomorphism
$\theta:A\to B$.  Such a lift exists: $1$ is primitive in each
finite free $\Z_p$-module, since $1\in pA$ would imply $A=pA$,
contrary to Nakayama's lemma.  Extend $1$ to bases and lift the
matrix of $\alpha$, fixing its first column.  Its determinant is a
unit.  The inverse map lifts $\alpha^{-1}$, and
\begin{equation}\label{eq:congruence-multiplication}
 \theta(ab)-\theta(a)\theta(b)\in p^rB.
\end{equation}

Choose an $A$-basis for a lattice and apply $\theta$ entrywise to
its action matrices $\rho_g$.  The resulting matrices are
invertible, since their products with the entrywise images of
$\rho_g^{-1}$ are congruent to $I$ modulo $p^r$.
Their multiplication defect lies in $p^r$.  The correction above
produces an exact $B$-action with
\[
 \sigma_g\equiv\theta(\rho_g)\pmod {p^{r-v}}.
\]
Since $r-v>v$, rigidity makes its isomorphism class unique among
actions with this congruence.  Two lifts of $\alpha$ differ by
$p^r$, so their corrected actions are isomorphic.  For a change
of basis $C$, the invertible matrix $\theta(C)$ is an intertwiner
modulo $p^{r-v}$ by \eqref{eq:congruence-multiplication};
\eqref{eq:congruence-intertwiner} makes it an exact intertwiner.
Thus the transported class depends only on $\alpha$ and the
original isomorphism class.

Transporting back with $\theta^{-1}$ gives an exact action
congruent to $\rho$ modulo $p^{r-v}$: the first correction and the
second correction both have this divisibility.  Rigidity identifies
the round trip with the original action.  The two transports are
therefore inverse bijections.

Now choose an $A$-splitting of an admissible short exact sequence.
In its adapted basis the action matrices are block upper triangular.
Entrywise transport and every correction preserve that form, and
their diagonal blocks correct the actions on the subobject and
quotient separately.  The corrected matrices thus give an admissible
short exact sequence with the three transported classes.
The inverse transport has the same property.  Block diagonal
matrices prove direct-sum compatibility.  Passing from the free
abelian groups on isomorphism classes to their quotients by all
exact-sequence relations proves the claimed $K_0$ isomorphism.

Tensor and exterior power matrices are integral polynomials in
the entries.  Applying $\theta$ before or after those polynomials
changes the result by $p^r$.  The corrected exact actions are
therefore congruent modulo $p^{r-v}$ and are isomorphic by rigidity.
Matrices over $\Z_p$ are fixed throughout.  For a subgroup
$H\le G$, put $v_H=v_p(|H|)\le v$.  Restriction followed by
correction, or correction followed by restriction, gives exact
$H$-actions congruent modulo $p^{r-v}$, which exceeds $v_H$.
For induction, fixed coset representatives give matrices with blocks
equal to action matrices or zero.  The two induced $G$-actions are
congruent modulo $p^{r-v}$, which exceeds $v$.
Rigidity proves both compatibilities and completes the theorem.


\section{Completion and the generic specialization kernel}
\label{sec:completion-generic-kernel}

Completion need not algebraize every coefficient-free representation.
It is injective on the entire lattice Grothendieck group and preserves
the entire generic specialization kernel.
The distinction is useful here: after adding one relative syzygy,
every class in that kernel can be represented by lattices whose
generic group-algebra modules are free, and those lattices admit a
direct algebraization by approximation of their coefficient bases.

\begin{theorem}\label{thm:completion-generic-kernel}\label{thm:completion-dvr-kernel}
Let $O$ be any characteristic-zero DVR, with uniformizer $\pi$,
fraction field $K$, completion $\widehat O$, and completed fraction
field $\widehat K$.  Let $A$ be a nonzero commutative algebra finite
free over $O$, and let $G$ be a finite group.  Put
\[
 \widehat A=A\otimes_O\widehat O,\qquad
 D=A[1/\pi],\qquad \widehat D=\widehat A[1/\pi].
\]
Completion induces an injection
\[
 K_0(\Lat_A(G))\hookrightarrow K_0(\Lat_{\widehat A}(G))
\]
and an isomorphism
\[
 \ker\!\left(K_0(\Lat_A(G))\longrightarrow K_0(\Lat_D(G))\right)
 \xrightarrow{\ \sim\ }
 \ker\!\left(K_0(\Lat_{\widehat A}(G))
                   \longrightarrow K_0(\Lat_{\widehat D}(G))\right).
\]
Moreover, any subgroup of the right-hand kernel that is split in
$K_0(\Lat_{\widehat A}(G))$ lifts to a subgroup of the left-hand
kernel that is split in $K_0(\Lat_A(G))$.
\end{theorem}

Here $\Lat_A(G)$ consists of finite free $A$-modules with
$A$-linear $G$-action, and its admissible sequences are the
coefficient-split short exact sequences.  A free module has one common
finite rank over the whole coefficient ring, including when that ring
is nonlocal or has several components.  Neither locality nor
reducedness of $A$ is assumed, and no Henselian, excellence,
ramification, or residue-field condition is imposed on $O$.  The theorem makes no assertion that all
completed lattices descend, nor that completion is an equivalence
on morphism groups.

\subsection{Generic projectives and algebraization}

We use first an elementary observation about the generic categories.
If $C$ is a finite-dimensional commutative algebra over a
characteristic-zero field, every $C$-free $G$-module is projective
over $C[G]$: a coefficient-linear section of an epimorphism onto it
becomes equivariant after averaging over $G$.  Thus every admissible
sequence in $\Lat_C(G)$ splits equivariantly.  The algebra $C[G]$
is Artinian, and its finitely generated projective modules have
unique multiplicity vectors in the finite set of indecomposable
projective classes.  The coefficient-free projectives form a
submonoid of this free commutative monoid.  Its group completion
embeds in the corresponding free abelian group.  Consequently
\begin{equation}\label{eq:completion-generic-cancellation}
 K_0(\Lat_C(G))\text{ is torsion-free},\qquad
 [U]=[V]\text{ in }K_0(\Lat_C(G))\Longrightarrow U\cong V.
\end{equation}
This includes nonreduced coefficient algebras $C$.

\begin{lemma}\label{lem:completion-generic-free-algebraization}
Let $M\in\Lat_{\widehat A}(G)$.  If
$M[1/\pi]\cong\widehat D[G]^s$, then there exists
$L\in\Lat_A(G)$ with
\[
 \widehat A\otimes_A L\cong M,
 \qquad L[1/\pi]\cong D[G]^s.
\]
\end{lemma}

\begin{proof}
Choose the displayed generic isomorphism and view $M$ inside the
standard representation $\widehat D[G]^s$.  With $k=s|G|$, a
coefficient basis of $M$ gives
\[
 M=P\widehat A^k,\qquad P\in\mathrm{GL}_k(\widehat D).
\]
The algebra $D$ is dense in $\widehat D$: in an $O$-basis of
$A$ this is the density of $K$ in $\widehat K$.  Choose
$P_0\in\mathrm M_k(D)$ sufficiently close to $P$ that
\[
 P^{-1}P_0\in I+\pi\mathrm M_k(\widehat A).
\]
Then $P_0\widehat A^k=M$.  The matrix $P_0$ is invertible over
$D$, since it is invertible after the faithfully flat scalar
extension $K\to\widehat K$.

Set $L=P_0A^k\subset D[G]^s$.  Let $R_g$ denote the standard
regular-action matrix of $g\in G$.  Stability of $M$ gives
\[
 P_0^{-1}R_gP_0\in\mathrm M_k(D)\cap
                         \mathrm M_k(\widehat A)
                       =\mathrm M_k(A).
\]
The last equality follows coordinatewise from
$K\cap\widehat O=O$ in any $O$-basis of $A$: membership in
the completed DVR is nonnegative valuation, and the valuation on $K$
is the original DVR valuation.
Applying the same statement to $g^{-1}$ shows that these matrices
are invertible over $A$.  Thus $L$ is a coefficient-free
$G$-lattice, its completion is $M$, and its generic module is the
chosen standard $D[G]^s$.
\end{proof}

\subsection{Reflection of isomorphisms and exact triples}

The completion facts used below are finite-module exactness,
flatness, faithful flatness when the completed ideal lies in the
Jacobson radical, and equality of finite quotients
\citep[Lemmas~10.97.1--10.97.4, Tags 00MA, 00MB, 00MC, 031C]{Stacks2026}.

\begin{lemma}\label{lem:completion-reflects-triples}
For coefficient-free $A$-lattices, completion reflects
isomorphisms and the existence of coefficient-split short exact
sequences with three prescribed isomorphism classes.
\end{lemma}

\begin{proof}
For $L,N\in\Lat_A(G)$, equivariant homomorphisms are the kernel
of the finitely many linear equivariance equations on
$\operatorname{Hom}_A(L,N)$.  Flatness of $\widehat O$ over $O$
therefore gives
\begin{equation}\label{eq:completion-hom-base-change}
 \operatorname{Hom}_{A[G]}(L,N)\otimes_O\widehat O
 \cong
 \operatorname{Hom}_{\widehat A[G]}(\widehat L,\widehat N).
\end{equation}
The integral Hom module is finite and torsion-free over the DVR
$O$, hence dense in its completion.  If $\widehat L\cong\widehat N$,
approximate an equivariant completed isomorphism modulo $\pi$ by an
integral equivariant map $U$.  Its reduction modulo $\pi$ is
invertible.  Since $A$ is finite over $O$, the ideal $\pi A$ lies in
its Jacobson radical, and $A/\pi A=\widehat A/\pi\widehat A$.
Therefore $\det U$ is a unit in $A$, so $U$ is an isomorphism.

For the exact-sequence assertion, coefficient-split extensions of
$N$ by $L$ are classified by
\[
 H^1\!\left(G,\operatorname{Hom}_A(N,L)\right),
\]
where $G$ acts by conjugation.  Indeed, a coefficient-linear
splitting writes the action in upper block-triangular form, whose
off-diagonal entries give a one-cocycle; changing the splitting
adds a coboundary.

This cohomology group is finite over $O$ and is killed by $|G|$.
For an explicit check of the latter statement, if $c_g$ is a
one-cocycle and $S=\sum_hc_h$, summing
$c_{gh}=c_g+gc_h$ over $h$ yields
\[
 |G|c_g=S-gS.
\]
Since $O$ has characteristic zero, $|G|$ is nonzero in $O$; it is
a unit times a power of $\pi$.  Thus this cohomology group has finite
length over $O$ (and is zero if $|G|$ is a unit), so it is unchanged
by completion.  Flat base change of the finite group cochain complex
also gives
\[
 H^1\!\left(G,\operatorname{Hom}_A(N,L)\right)\otimes_O\widehat O
 \cong
 H^1\!\left(G,
       \operatorname{Hom}_{\widehat A}(\widehat N,\widehat L)
       \right).
\]
Every completed extension class consequently descends to a
coefficient-split extension over $A$.  If its completed middle
term is isomorphic to the completion of a prescribed lattice,
the isomorphism-reflection argument just proved identifies the
integral middle term with that lattice.  This proves reflection
of exact triples.  Preservation follows from flat base change.
\end{proof}

\subsection{Cofinality and exact-sequence relations}

We use the following published cofinality lemma
\citep[Chapter~II, Lemma~7.2, p.~62 of the author's chapter
version]{Weibel2013}.

\begin{lemma}[Cofinality]\label{lem:completion-elementary-cofinality}
Let $\mathcal F$ be a full, isomorphism-closed, extension-closed
exact subcategory of an exact category $\mathcal E$, with the
induced exact structure.  Suppose every $X\in\mathcal E$ has a
supplement $X'\in\mathcal E$ such that
$X\oplus X'\in\mathcal F$.  Then
$K_0(\mathcal F)\to K_0(\mathcal E)$ is injective.
\end{lemma}

The extension-closure hypothesis controls exact-sequence relations
as well as stable isomorphisms.

\subsection{Identification of the kernels}

\begin{proof}[Proof of Theorem~\ref{thm:completion-generic-kernel}]
Write $\mathcal E_A=\Lat_A(G)$ and let $\mathcal F_A$ be its
full subcategory of lattices $L$ for which $L[1/\pi]$ is free over
$D[G]$.  Define $\mathcal E_{\widehat A}$ and
$\mathcal F_{\widehat A}$ in the same way.  Both subcategories
$\mathcal F$ are extension-closed: a generic extension of free
group-algebra modules splits and is free.

Lemmas~\ref{lem:completion-generic-free-algebraization}
and~\ref{lem:completion-reflects-triples} show that completion
gives a bijection between the isomorphism classes in
$\mathcal F_A$ and $\mathcal F_{\widehat A}$ and preserves and
reflects their exact triples.  Comparing the generators and
relations defining their Grothendieck groups gives
\begin{equation}\label{eq:completion-free-subcategory-kzero}
 K_0(\mathcal F_A)\xrightarrow{\sim}
 K_0(\mathcal F_{\widehat A}).
\end{equation}

These are cofinal subcategories in the sense of
Lemma~\ref{lem:completion-elementary-cofinality}.  Indeed, for
$M\in\mathcal E_A$ of coefficient rank $a$, the action map
\[
 A[G]\otimes_A M\longrightarrow M,
 \qquad g\otimes m\longmapsto gm,
\]
is a coefficient-split epimorphism.  Its kernel $\Omega M$ is
coefficient-free: an explicit coefficient isomorphism from
$\bigoplus_{g\ne1}M$ onto that kernel is
\[
 (m_g)_{g\ne1}\longmapsto
 \sum_{g\ne1}\bigl(g\otimes m_g-1\otimes gm_g\bigr).
\]
The middle term is $A[G]^a$.  Generically $M[1/\pi]$ is projective
by averaging, so
\[
 M[1/\pi]\oplus(\Omega M)[1/\pi]\cong D[G]^a.
\]
Consequently $M\oplus\Omega M\in\mathcal F_A$.  The same
argument applies after completion.  The lemma now gives
injections
\begin{equation}\label{eq:completion-cofinal-injections}
 K_0(\mathcal F_A)\hookrightarrow K_0(\mathcal E_A),\qquad
 K_0(\mathcal F_{\widehat A})\hookrightarrow
                         K_0(\mathcal E_{\widehat A}).
\end{equation}

We can now prove injectivity on the full Grothendieck groups.
Let $\mathcal I_A$ be the full, isomorphism-closed subcategory of
$\mathcal E_{\widehat A}$ consisting of objects isomorphic to
completions of objects of $\mathcal E_A$.  The extension-descent
argument in Lemma~\ref{lem:completion-reflects-triples} shows that
$\mathcal I_A$ is extension-closed.  Reflection of isomorphisms
and preservation and reflection of exact triples identify the
generators and relations, giving
\[
 K_0(\mathcal E_A)\xrightarrow{\sim}K_0(\mathcal I_A).
\]
For every $M\in\mathcal E_{\widehat A}$, the explicit syzygy above
gives $M\oplus\Omega M\in\mathcal F_{\widehat A}$.
Lemma~\ref{lem:completion-generic-free-algebraization} gives
$\mathcal F_{\widehat A}\subseteq\mathcal I_A$.  Thus
$\mathcal I_A$ is cofinal in $\mathcal E_{\widehat A}$, and
Lemma~\ref{lem:completion-elementary-cofinality} yields
\[
 K_0(\mathcal E_A)\xrightarrow{\sim}K_0(\mathcal I_A)
 \hookrightarrow K_0(\mathcal E_{\widehat A}).
\]
This composite is completion.  The argument uses algebraization
only for regular generic type and makes no descent assertion for
an arbitrary completed lattice.

The generic free-group-algebra rank defines homomorphisms
\[
 \nu_A:K_0(\mathcal F_A)\to\Z,
 \qquad
 \nu_{\widehat A}:K_0(\mathcal F_{\widehat A})\to\Z,
\]
and the isomorphism
\eqref{eq:completion-free-subcategory-kzero} preserves them.
Under \eqref{eq:completion-cofinal-injections},
$\ker\nu_A$ is exactly the generic specialization kernel in
$K_0(\mathcal E_A)$.  To see surjectivity onto that kernel,
write any class as $[M]-[N]$.  If its generic image is zero,
\eqref{eq:completion-generic-cancellation} gives
$M[1/\pi]\cong N[1/\pi]$.  Adding $\Omega M$ to both terms
represents the same class by
\[
 [M\oplus\Omega M]-[N\oplus\Omega M],
\]
and both lattices are in $\mathcal F_A$ with the same generic
free-group-algebra rank.  Conversely, a difference of rank zero
in $K_0(\mathcal F_A)$ has zero generic class.  The same reasoning
identifies the completed kernel with $\ker\nu_{\widehat A}$.
Equation~\eqref{eq:completion-free-subcategory-kzero} therefore
proves the asserted isomorphism of kernels.

Finally, let $s:H\to K_0(\mathcal E_{\widehat A})$ and
$q:K_0(\mathcal E_{\widehat A})\to H$ satisfy $qs=1$, with
$s(H)$ contained in the completed generic kernel.  Lift $s$
through the kernel isomorphism, obtaining
$s_A:H\to K_0(\mathcal E_A)$, and put
\[
 q_A=q\circ\text{completion}.
\]
Then $q_As_A=1$.  The lifted subgroup is consequently split in
the entire integral lattice Grothendieck group, as claimed.
\end{proof}

\subsection{Deep monic deformations over the localized integers}

\begin{corollary}\label{cor:deep-monogenic-orders}
Let $O=\Z_p$ or $\Z_{(p)}$, and let $F\in O[t]$ be monic of
degree $n\ge2$ with
\[
 F\equiv t^n\pmod {p^3}.
\]
For $B=O[t]/F$, the generic specialization kernel in
$K_0(\Lat_B(S_p))$ contains a subgroup isomorphic to
$\bigoplus_{i\ge1}\Z/n$ that is split in the whole Grothendieck
group.  A countable such summand has a coordinate retraction
annihilating the Young permutation span.  In particular,
$\Phi_B$ is not surjective in degree $p$.
\end{corollary}

\begin{proof}
First let $O=\Z_p$, and put $B_0=\Z_p[t]/t^n$.
The ring maps $\Z_p\to B_0\to\Z_p$, with $t\mapsto0$ in the
second map, split the reduction homomorphism on lattice
Grothendieck groups.  Its kernel is consequently a direct summand
of the whole group.  The split countable $\Z/n$ subgroup of
Theorem~\ref{thm:all-truncations} is therefore split in
$K_0(\Lat_{B_0}(S_p))$.

The coefficient algebra $B$ is finite free of rank $n$, and the
hypothesis on $F$ gives $B/p^3B\cong B_0/p^3B_0$.
Since $v_p(|S_p|)=1$, Theorem~\ref{thm:congruent-coefficients}
applies with $r=3>2$.  Transport the preceding subgroup,
together with its section and retraction.  Its exact orders
and its splitting are preserved.  The generic target is torsion-free
by \eqref{eq:completion-generic-cancellation}, applied to
$B[1/p]$, so this torsion subgroup lies in the generic kernel.

For completeness, the Young-span assertion does not require a
particular normalization of the original retraction.  There are
only finitely many degree-$p$ Young permutation lattices.  Each
has finitely supported image in $\bigoplus_i\Z/n$.  Delete the
finite union of those supports and restrict the section to the
remaining coordinates.  Composing the retraction with their
coordinate projection gives a retraction of a countable direct
summand that annihilates every Young permutation lattice.
If the retraction from Theorem~\ref{thm:all-truncations} is already
so normalized, no coordinates need be deleted.

Now let $O=\Z_{(p)}$.  The completion of $B$ is
$\widehat B=\Z_p[t]/F$, so the complete case supplies a split
countable subgroup in its generic kernel.
Theorem~\ref{thm:completion-generic-kernel} lifts that subgroup
to the generic kernel over $B$ and preserves its splitting in
the full lattice Grothendieck group.  Compose the completed
coordinate retraction with completion.  Each integral Young
permutation lattice completes to its corresponding completed
Young lattice, so this retraction still annihilates the Young
span.  It takes the lifted subgroup onto the nonzero coordinate
vectors.  Since the degree-$p$ image of $\Phi_B$ is the Young
permutation span, $\Phi_B$ cannot be surjective.
\end{proof}

\section{Distinguished monic coefficient rings over unramified DVRs}
\label{sec:unramified-distinguished-monic}\label{sec:singular-distinguished-monic}

The order of the obstruction is determined by the first coefficient
of valuation one.  We prove the result for every characteristic-zero
DVR with uniformizer \(p\), including arbitrary residue fields and
noncomplete bases.  The source lattices and their exact relations are
finite polynomial constructions over the original DVR.  Completion
is used only for the additive detectors.  We include the supported
length construction and every categorical calculation used below.
Throughout, lattice exact sequences split over the coefficient ring;
they need not split equivariantly.

\begin{theorem}\label{thm:unramified-distinguished-monic}\label{thm:singular-distinguished-monic}
Let \(O\) be a characteristic-zero DVR with residue field \(k\) of
characteristic \(p>0\), and suppose that its normalized valuation
satisfies \(v_O(p)=1\).  Let
\[
 F(t)=t^n+\sum_{i=0}^{n-1}a_it^i,
 \qquad n\ge2,\quad a_i\in pO,\quad a_0\in p^2O,
 \qquad A=O[t]/(F).
\]
Set
\[
 r=\min\{1\le i<n:a_i\notin p^2O\},\qquad N=\min(n,pr),
\]
with \(r=\infty\), \(N=n\) if this set is empty.
For \(G=C_p\) and \(G=S_p\), there are homomorphisms
\[
 s_G:\bigoplus_{j\ge1}\mathbb Z/N\longrightarrow K_0(\Lat_A(G)),
 \qquad
 \eta_G:K_0(\Lat_A(G))\longrightarrow\bigoplus_{j\ge1}\mathbb Z/N
\]
with \(\eta_Gs_G=1\), and with the image of \(s_G\) in the kernel
of extension to \(A[1/p]\).  For \(G=S_p\), the map \(\eta_G\)
annihilates every Young permutation lattice.  Thus the distinguished
section classes have exact order \(N\), and \(\Phi_A\) is not
surjective in degree \(p\).
The residue field may be finite, infinite, or imperfect, and \(O\)
need not be complete.  No reducedness of \(A[1/p]\) is assumed.
The section classes admit polynomial presentation matrices over finite
\'etale local extensions of \(O\); when \(k\) is infinite these
extensions may all be taken to be \(O\) itself.
\end{theorem}

\subsection{Length on the supported singularity category}

For a Noetherian ring \(S\), write
\(D_{\mathrm{sg}}(S)=D^b(\operatorname{coh}S)/\operatorname{Perf}(S)\).
All categories below carry their usual differential graded, or equivalently
stable, enhancements.  A superscript \(\mathrm{ic}\) denotes idempotent
completion.

\begin{lemma}\label{lem:singular-monic-supported-length}
Let \(T\) be a complete regular local ring of dimension two, let
\(S=T[t]/(t^N)\), and let \(Z\) be its closed point.  Put
\[
 \mathcal Q_Z=
 \left(D^b_Z(\operatorname{coh}S)/\operatorname{Perf}_Z(S)\right)^{\mathrm{ic}}.
\]
There is a homomorphism
\[
 \lambda_N:K_0(\mathcal Q_Z)\longrightarrow\mathbb Z/N
\]
which sends a finite-length \(S\)-module \(L\) to
\(\operatorname{length}_T(L)\bmod N\).
Moreover \(\mathcal Q_Z\) is equivalent to the idempotent completion
of the full subcategory of \(D_{\mathrm{sg}}(S)\) consisting of objects
perfect outside \(Z\).
\end{lemma}

\begin{proof}
Write
\[
 \mathcal P_S=\operatorname{Perf}_Z(S),\quad
 \mathcal C_S=D^b_Z(\operatorname{coh}S),\quad
 \mathcal P_T=\operatorname{Perf}_Z(T),
\]
and let \(j:\mathcal P_S\to\mathcal C_S\) be the inclusion.
Restriction of scalars gives an exact functor
\(R:\mathcal C_S\to\mathcal P_T\): the restricted complex is bounded
coherent with finite-length cohomology and is perfect because \(T\)
is regular.  There is also an exact functor
\[
 G:\mathcal P_S\longrightarrow\mathcal P_T,
       \qquad G(P)=P\otimes_S^{\mathbf L}T.
\]
This derived reduction is used only on perfect \(S\)-complexes.

For \(0\le a\le N\), define exact functors
\[
 L_a(P)=\operatorname{Res}_T(P\otimes_S^{\mathbf L}t^aS).
\]
They take values in \(\mathcal P_T\): in a bounded finite-projective
model for \(P\), every term becomes finite projective over \(T\),
since \(t^aS\) is \(T\)-free.  Localization verifies the support
condition.  The exact sequences
\(0\to t^{a+1}S\to t^aS\to T\to0\) give natural cofiber sequences
\(L_{a+1}\to L_a\to G\).  Their endpoints are \(L_0=Rj\) and
\(L_N=0\).  Additivity therefore supplies a homotopy of maps of
nonconnective \(K\)-theory spectra
\begin{equation}\label{eq:singular-monic-spectrum-filtration}
 \mathbf K(R)\mathbf K(j)\simeq N\mathbf K(G).
\end{equation}
This assertion is at the level of spectra, rather than only of Euler
characteristics.

Nonconnective localization and Morita invariance give a cofiber sequence
\[
 \mathbf K(\mathcal P_S)\longrightarrow\mathbf K(\mathcal C_S)
      \longrightarrow\mathbf K(\mathcal Q_Z).
\]
Use the homotopy in \eqref{eq:singular-monic-spectrum-filtration} in
the square whose horizontal maps are \(\mathbf K(j)\) and
\(N:\mathbf K(\mathcal P_T)\to\mathbf K(\mathcal P_T)\), and whose
vertical maps are \(\mathbf K(G)\) and \(\mathbf K(R)\).
Its map on cofibers is
\begin{equation}\label{eq:singular-monic-cofiber-map}
 \mathbf K(\mathcal Q_Z)\longrightarrow
 \operatorname{cofib}\bigl(N:\mathbf K(\mathcal P_T)
                                  \to\mathbf K(\mathcal P_T)\bigr).
\end{equation}
The localization and additivity inputs, including the use of the
idempotent-completed quotient, are
\citet[Theorem 9.8, p.~808; the localization consequence on p.~814;
Remark 9.33 and Theorem 9.34, pp.~815--816]{BlumbergGepnerTabuada2013}.

Regularity identifies \(\mathcal P_T\) with
\(D^b_Z(\operatorname{coh}T)\).  Its ordinary cohomological
\(t\)-structure is bounded, with heart the finite-length \(T\)-modules.
Truncation triangles and composition series identify
\(K_0(\mathcal P_T)\) with \(\mathbb Z\) by Euler length, and
\citet[Theorem 1.1, author p.~2]{AntieauGepnerHeller2019} gives
\(K_{-1}(\mathcal P_T)=0\).  The long exact sequence of the target
cofiber in \eqref{eq:singular-monic-cofiber-map} therefore gives
\(\pi_0=\mathbb Z/N\).  The resulting map \(\lambda_N\) has the
stated values because restriction sends a finite-length module to
its underlying \(T\)-module.  This construction is already defined
on \(\mathbf K(\mathcal Q_Z)\); no vanishing of
\(K_{-1}(\operatorname{Perf}_Z(S))\), and no extension from an
uncompleted Grothendieck group, has been assumed.

We verify the last assertion to identify the domain that will occur
in the proof.  Choose regular parameters \(u,v\) of \(T\); they form
an \(S\)-regular sequence because \(S\) is \(T\)-free.  A bounded
coherent complex supported at \(Z\) has a bounded model of
finite-length modules.  Some common powers \(u^a,v^a\) kill all its
terms.  A map to it from a bounded projective complex \(P\) factors
through \(P/(u^a,v^a)P\).  The latter is represented by tensoring
\(P\) with the Koszul complex on \(u^a,v^a\), so it is perfect and
supported at \(Z\).  The supported-derived model and the
Verdier-quotient factorization criterion are
\citet[Lemmas 2.1, 2.5, and 2.6, author pp.~3--5]{Orlov2011}.
They show that the supported quotient is fully faithful in
\(D_{\mathrm{sg}}(S)\).

For density, represent an object perfect outside \(Z\), up to shift,
by a finite maximal Cohen--Macaulay module \(M\).  The hypersurface
description identifies singularity-category maps with stable module
maps \citep[equation (1.1), p.~2, and Appendix B,
pp.~37--40]{BurkeWalker2012}.  The stable endomorphism module of
\(M\) is the quotient of \(\operatorname{Hom}_S(M,M)\) by the image
of \(\operatorname{Hom}_S(M,S)\otimes_SM\).  It is finite, and this
description commutes with localization.  Its support is contained
in \(Z\), so powers \(u^a,v^a\) kill the stable identity.
The cone of multiplication by \(u^a\) on \(M\) is therefore
\(M\oplus M[1]\) in the singularity category.  Multiplication by
\(v^a\) is zero on this cone there, since scalar multiplication is
natural.  Taking its cone produces an object containing \(M\) as
a direct summand and represented by the Koszul complex of \(M\)
on \(u^a,v^a\).  That complex has finite-length cohomology.
Thus the fully faithful functor has dense image, and becomes an
equivalence after idempotent completion.
\end{proof}

\subsection{Parameter choices over the residue field}

Put \(g(X)=(1+X)^p-1\).  We choose finite \'etale local extensions
\(D_j/O\), units \(c_j\in D_j\), and \(q_j\in pD_j\) such that
\begin{equation}\label{eq:singular-monic-constant-root}
 g(q_j)+c_ja_0=0.
\end{equation}
The associated closed points of \(\mathbb A^1_k\) will have
pairwise disjoint \(\mathbb F_p^\times\)-orbits, each orbit being
free.  If \(r<\infty\), put \(\beta=a_r/p\in O^\times\).
In the case \(n=pr\) we also require
\begin{equation}\label{eq:unramified-monic-bad-parameters}
 \bar c_j+(-\bar c_j\bar\beta)^p\ne0.
\end{equation}

Here are the exact constant formulas.  If \(a_0=0\), take
\(q_j=0\) and any chosen unit \(c_j\).  Otherwise write
\(a_0=p^\nu u\), with \(\nu\ge2\), \(u\in O^\times\), and take
\begin{equation}\label{eq:unramified-monic-parameter-formulas}
 q_j=p^{\nu-1}\theta_j,\qquad c_j=-g(q_j)/a_0.
\end{equation}
This quotient is integral.  Except when \(p=2,\nu=2\), its residue
is \(-\bar\theta_j/\bar u\).  In that exceptional case its residue
is \(-\bar\theta_j(\bar\theta_j+1)/\bar u\).

Suppose first that \(k\) is finite.  Choose distinct prime integers
\(m_j>p\), and monic polynomials \(Q_j\in O[Z]\) lifting
irreducible polynomials of degree \(m_j\) over \(k\).  Set
\(D_j=O[\theta_j]/(Q_j)\).  It is finite free, with residue field
\(k_j\) of degree \(m_j\) over \(k\), and is local with maximal
ideal \(pD_j\): every maximal ideal of a finite integral extension
contracts to the maximal ideal of \(O\), and its special fiber is
a field.  Gauss's lemma shows that \(D_j\) is a domain.
It is consequently a DVR.  The residue polynomial is separable,
so \(Q_j'(\theta_j)\) is a unit and \(D_j/O\) is finite \'etale.
These assertions do not require \(O\) to be Henselian.

For \(a_0=0\), choose \(c_j=\theta_j\); otherwise use
\eqref{eq:unramified-monic-parameter-formulas}.
The residue \(\bar c_j\) has degree \(m_j\) over \(k\).
In the exceptional quadratic expression, the degree of
\(k(\bar\theta_j)\) over \(k(\bar\theta_j^2+\bar\theta_j)\)
is at most two and divides the odd prime \(m_j\), hence is one.
In all other cases the full degree is immediate.  Therefore
\eqref{eq:unramified-monic-bad-parameters} holds: its left side is
the value of a nonzero degree-\(p\) polynomial over \(k\), whereas
\(\bar c_j\) has degree greater than \(p\).

These closed points have distinct degrees over \(k\).
If \(a\bar c_j\) is a \(k\)-Galois conjugate of \(\bar c_j\),
with \(a\in\mathbb F_p^\times\), taking the norm from \(k_j\)
to \(k\) gives \(a^{m_j}=1\).  Since
\(\gcd(m_j,p-1)=1\), one has \(a=1\).
This proves the required orbit separation.  The degrees and norms
in this argument are relative to \(k\), not to \(\mathbb F_p\).

Now suppose that \(k\) is infinite.  Take \(D_j=O\).
If \(a_0=0\), the possible residues \(\bar c_j\) are all of
\(k^\times\).  The same is true using
\eqref{eq:unramified-monic-parameter-formulas} outside the case
\(p=2,\nu=2\), because its residual expression is a nonzero linear
function of \(\bar\theta_j\).  In that exceptional case the possible
residues are a nonzero scalar multiple of the Artin--Schreier image
\[
 \{x^2+x:x\in k\}.
\]
This image is infinite: each fiber contains at most two elements,
whereas \(k\) is infinite.  Remove zero and, when necessary, the
finitely many roots excluded by
\eqref{eq:unramified-monic-bad-parameters}.  Inductively choose
countably many remaining elements in distinct
\(\mathbb F_p^\times\)-orbits.  At each step only finitely many
additional elements are excluded.  Every nonzero rational point
has a free scalar orbit.  Lift the selected residues, or their
chosen preimages \(\bar\theta_j\), to \(O\), and use the formulas
above.  This proves all parameter requirements without any
perfection assumption on \(k\).

\subsection{Finite polynomial lattices over the original DVR}\label{subsec:singular-monic-polynomial-sections}

Fix one of the resulting \(D=D_j,c=c_j,q=q_j\), and put
\(v=(1+q)^{p-1}\in D^\times\).  The following finite recursion
is over the original ring \(D\):
\begin{equation}\label{eq:singular-monic-finite-jet}
 \gamma_0=q,\qquad
 \gamma_d=-\frac{[t^d]\bigl(g(\sum_{h<d}\gamma_ht^h)+cF\bigr)}{pv}
 \quad(1\le d<N),\qquad
 \alpha_c(t)=\sum_{d<N}\gamma_dt^d.
\end{equation}
Before degree \(r\), all earlier coefficients lie in \(pD\),
and \(g(pC)\) is divisible by \(p^2\).  Since \(a_d\in p^2O\),
the numerator is in \(p^2D\), giving \(\gamma_d\in pD\).
At degree \(r<\infty\), it gives
\(\bar\gamma_r=-\bar c\bar\beta\ne0\).
Thereafter the residual partial polynomial has \(t\)-order \(r\).
Its \(p\)-th power has no term in degree \(d<N\le pr\), and
\(g\) reduces to that \(p\)-th power.  The degree-\(d\) term
of \(cF\) is also divisible by \(p\), since \(d<n\).
Thus every division by \(p\) in
\eqref{eq:singular-monic-finite-jet} is integral.  If \(r=\infty\),
the first argument applies throughout.

The new coefficient changes degree \(d\) by precisely
\(g'(q)\gamma_d=pv\gamma_d\).  The constant equation is
\eqref{eq:singular-monic-constant-root}.  We obtain an exact
polynomial identity
\begin{equation}\label{eq:singular-monic-polynomial-root}
 g(\alpha_c(t))+cF(t)=t^NU_c(t),\qquad
 \alpha_c,U_c\in D[t],\quad\deg\alpha_c<N.
\end{equation}
Its constant residue is
\begin{equation}\label{eq:singular-monic-unit-residue}
 \bar U_c(0)=
 \begin{cases}
 \bar c,&r=\infty\text{ or }n<pr,\\
 (-\bar c\bar\beta)^p,&pr<n,\\
 \bar c+(-\bar c\bar\beta)^p,&n=pr.
 \end{cases}
\end{equation}
It is nonzero by the parameter choices.  The algebra
\(A_D=D[t]/(F)\) is local with maximal ideal \((p,t)\), because
its special fiber is \((D/p)[t]/t^n\).  Hence \(U_c(t)\) is a
unit in \(A_D\), including when \(D\) is not complete.

Set
\[
 h_c(X,t)=\frac{g(X)-g(\alpha_c(t))}{X-\alpha_c(t)}.
\]
For \(0\le i\le N\), define \(M_i(c)\) by generators \(e,f\)
and the two relations given by the columns of
\begin{equation}\label{eq:singular-monic-factorizations}
 \begin{pmatrix}
 X-\alpha_c(t)&t^i\\-U_c(t)t^{N-i}&h_c(X,t)
 \end{pmatrix}.
\end{equation}
The determinant is \(g+cF\).  Equivalently, the relations are
\begin{equation}\label{eq:singular-monic-relations}
 (X-\alpha_c)e=U_ct^{N-i}f,\qquad t^ie+h_c(X,t)f=0.
\end{equation}
They give an \(A_D\)-basis
\begin{equation}\label{eq:singular-monic-free-basis}
 e,\qquad X^jf\quad(0\le j\le p-2).
\end{equation}
Indeed the first relation specifies \(Xe\), while monicity of
\(h_c\) specifies \(X^{p-1}f\).  These prescriptions define an
operator on the free module with the displayed basis, satisfy both
relations, and give inverse maps to the presented module.  The same
argument first works over \(D[t]\).  The adjugate of the matrix
annihilates its cokernel by \(g+cF\), so after imposing \(F=0\)
the operator \(\sigma=1+X\) defines a \(C_p\)-action.

At \(i=0\) eliminate \(e=-h_cf\); the remaining relation is
\((g+cF)f=0\).  At \(i=N\), use the unit \(U_c\) in \(A_D\)
to eliminate \(f=U_c^{-1}(X-\alpha_c)e\), leaving
\((g+cF)e=0\).  Consequently
\begin{equation}\label{eq:singular-monic-endpoints}
 M_0(c)\simeq M_N(c)\simeq A_D[C_p].
\end{equation}

In these bases define
\[
 \alpha_i:M_i\longrightarrow M_{i+1},\quad
 \alpha_i=\operatorname{diag}(t,I_{p-1}),\qquad
 \beta_i:M_{i+1}\longrightarrow M_i,\quad
 \beta_i=\operatorname{diag}(1,tI_{p-1}).
\]
Substitution in \eqref{eq:singular-monic-relations} proves
equivariance; both composites are multiplication by \(t\).
For \(1\le i<N\), the sequence
\begin{equation}\label{eq:singular-monic-mesh}
 0\longrightarrow M_i
 \xrightarrow{(\beta_{i-1},\alpha_i)}M_{i-1}\oplus M_{i+1}
 \xrightarrow{(\alpha_{i-1},-\beta_i)}M_i\longrightarrow0
\end{equation}
is coefficient-split exact.  Write a coefficient vector as \((a,v)\),
where \(a\) is its first coordinate.  The two maps are
\[
 (a,v)\longmapsto((a,tv),(ta,v)),\qquad
 ((a_L,v_L),(a_R,v_R))\longmapsto
                  (ta_L-a_R,v_L-tv_R).
\]
The kernel equations give exactness.  A coefficient left inverse
of the first map and a coefficient right inverse of the second are
\[
 ((a_L,v_L),(a_R,v_R))\longmapsto(a_L,v_R),\qquad
 (a,v)\longmapsto((0,v),(-a,0)).
\]
These are actual splittings even when \(t\) is a zero-divisor.
With \(x_i(c)=[M_i(c)]-[A_D[C_p]]\), the mesh and endpoints give
\begin{equation}\label{eq:singular-monic-annihilation}
 x_i(c)=i x_1(c),\qquad Nx_1(c)=0.
\end{equation}
Thus the annihilation relation holds over the original coefficient ring.

The characteristic polynomial of the \(X\)-operator is
\[
 (Z-\alpha_c)h_c(Z,t)+U_ct^N=g(Z)+cF(t).
\]
This follows from the companion block for \(h_c\); the coupling
to the first coordinate contributes \(U_ct^N\).  The characteristic
polynomial of \(\sigma\) on \(M_i\) is therefore \(Z^p-1\).
Put \(B=A_D[1/p]\), a commutative Artinian algebra, possibly
nonreduced.  Averaging a \(B\)-linear splitting makes
\(M_i[1/p]\) projective over \(B[C_p]\).  This group algebra is
finite \'etale over \(B\).  Over every residue field of \(B\),
the separable characteristic polynomial \(Z^p-1\), with every
irreducible factor occurring once, gives the regular representation.
On each local Artinian factor of \(B[C_p]\) the projective module
therefore has rank one and is free.  Hence
\begin{equation}\label{eq:singular-monic-generic-zero}
 M_i(c)[1/p]\simeq B[C_p],\qquad x_i(c)[1/p]=0.
\end{equation}

\subsection{The detector after completion}

We verify the categorical hypotheses over the complete coefficient
ring without imposing a restriction on its residue field.
Let \(\widehat O\) be the completion, put
\(\widehat A=\widehat O[t]/F\), and set
\[
 Q=\widehat O[[t,X]],\qquad R=Q/(F,g)=\widehat A[C_p].
\]
The ring \(Q\) is regular Noetherian of dimension three.
Distinguished division gives the displayed quotient, and \(F,g\)
is a regular sequence: \(F\) is nonzero, and the lowest nonzero
\(X\)-coefficient of \(g\) on \(\widehat A[[X]]\) is the
non-zero-divisor \(p\).  Form
\[
 Y=\operatorname{Proj}Q[T_F,T_g]/(FT_F+gT_g),
\]
with the natural maps
\(\beta:\mathbb P^1_R\hookrightarrow Y\) and
\(\pi:\mathbb P^1_R\to\operatorname{Spec}R\).
The generic-hypersurface equivalence and its inverse are
\begin{equation}\label{eq:singular-monic-generic-equivalence}
 \mathcal F=\beta_*\pi^*:D_{\mathrm{sg}}(R)
        \simeq D_{\mathrm{sg}}(Y),\qquad
 \mathcal F^{-1}=R\pi_*L\beta^*(-1)[-1].
\end{equation}
The precise inputs are \citet[Theorem 2.10, p.~8; the paragraph
after Corollary 2.11, p.~9; and the formula preceding (4.3),
p.~19]{BurkeWalker2012}.  The ambient ring is regular Noetherian
of finite dimension and the sequence is regular, as just checked.
No coefficient-field, equicharacteristic, or residue-perfectness
hypothesis is needed.

For every \(\widehat A\)-free cyclic lattice \(E\), its positive
self-Ext groups are cyclic-group cohomology with coefficients
\(\operatorname{Hom}_{\widehat A}(E,E)\).  They are two-periodic,
killed by \(p\), and of bounded finite dimension over \(k\).
The even and odd operator modules are finite over
\(R[\chi_F,\chi_g]\)
\citep[Remark 3.19, p.~17]{BurkeWalker2012}.
Their high-degree tails are therefore finite graded modules over
\[
 (R/pR)[\chi_F,\chi_g],\qquad
 R/pR=k[t,X]/(t^n,X^p),
\]
with bounded Hilbert functions after reducing the nilpotent maximal
ideal of \(R/pR\).  Their projective support is a finite set of
closed points of \(\mathbb P^1_k\).  The stable-Ext comparison
identifies it with the nonperfect locus of the generic-hypersurface
image \citep[Corollaries 3.18 and 3.20, pp.~17--18, and the end of
Section 8.1, pp.~33--34]{BurkeWalker2012}.  This argument works over
an arbitrary field \(k\), including an imperfect one.

Complete the chart \(z=T_F/T_g\) after extension to
\(\widehat D_j=D_j\otimes_O\widehat O\), at the selected residue
point \(z=\bar c_j\).  Write \(z=c_j+w\).
The formal constant root is obtained from
\[
 G(B)=g(pB)/p^2,\qquad
 G(B)+z(a_0/p^2)=0,
 \qquad G'(B)=(1+pB)^{p-1}.
\]
The chosen constant \(q_j/p\) is an exact solution at \(w=0\),
and the derivative is a unit.  Formal implicit-function iteration,
followed by multiplication by \(p\), gives the chosen branch
\(\gamma_0(z)\in p\widehat D_j[[w]]\)
for the root of \(g(\gamma_0)+za_0=0\).
Repeat the finite recursion
\eqref{eq:singular-monic-finite-jet} over \(\widehat D_j[[w]]\).
It gives
\begin{equation}\label{eq:singular-monic-root-factorization}
\begin{gathered}
 g(\alpha(z,t))+zF=t^NU(z,t),\qquad
 U\in\widehat D_j[[w,t]]^\times,\\
 \alpha(c_j,t)=\alpha_{c_j}(t).
\end{gathered}
\end{equation}
All coefficient divisibility checks remain valid modulo \(p\),
regardless of perfection of the residue field.

The root equation defines the closed immersion
\begin{equation}\label{eq:singular-monic-root-section}
 \iota_j:\operatorname{Spec}S_j\longrightarrow
 \operatorname{Spec}\widehat D_j[[w,t,X]]/(g+(c_j+w)F),
 \qquad S_j=\widehat D_j[[w,t]]/(t^N).
\end{equation}
It has Tor dimension one.  In the regular ambient power-series ring,
the quotient by \(X-\alpha\) is the domain \(\widehat D_j[[w,t]]\)
and the potential reduces to the nonzero element \(t^NU\).
To check regularity on the hypersurface directly, an equality
\((X-\alpha)b=(g+zF)d\) forces \(d\) to be divisible by
\(X-\alpha\) after reduction modulo that element.  Cancellation
in the ambient domain then gives \(b\in(g+zF)\).
Derived root pullback is consequently bounded and preserves perfect
objects.
The image of every source lattice is perfect off the target
closed point, by the finite-support argument above.
Lemma~\ref{lem:singular-monic-supported-length} applies with
\(T=\widehat D_j[[w]]\), a complete regular local ring of
dimension two.  That lemma has no residue-field hypothesis and
already includes the idempotent-completion and supported-category
arguments.  Composing its length map with these functors and with
completion defines
\begin{equation}\label{eq:singular-monic-detector}
 \rho_j:K_0(\Lat_A(C_p))\longrightarrow\mathbb Z/N.
\end{equation}
Every step sends arbitrary coefficient-split exact sequences to
triangles.  This is therefore a whole-category additive detector.

There are only finitely many nonzero selected coordinates on any
lattice: its original nonperfect support is finite, and the chosen
closed points are distinct.  Thus \((\rho_j)_j\) takes values in
the direct sum, even when \(k\) is uncountable.
It kills the regular lattice and the trivial lattice.  The former
is perfect; the latter has generic-hypersurface image \(Y/(X)\)
near a selected unit parameter, where \(X\) is regular since the
potential modulo \(X\) is \(zF\ne0\).

\subsection{Evaluation and split sections}

Work over \(\widehat D_j\), abbreviating \(c_j\) by \(c\).
Before imposing \(F\), the source matrix defines a module \(L_i\)
over \(\widehat D_j[[t,X]]\), free over \(\widehat D_j[[t]]\)
in the basis \eqref{eq:singular-monic-free-basis}.  The polynomial
operator extends to the formal-power-series action: modulo \((p,t)\)
it is nilpotent, including at the endpoints by the eliminations
already displayed.  Thus the asserted formal module and basis are
well-defined.  Multiplication by \(F\) on \(L_i\) is injective,
and \(M_i=L_i/FL_i\).

Let \(j_c\) be the closed constant section \(z=c\) of \(Y_{\widehat D_j}\),
with source \(\operatorname{Spec}\widehat D_j[[t,X]]/(g+cF)\).
On this chart the immersion \(\beta\) is defined by \(F\), since
\(g=-zF\).  Its derived pullback of \(j_{c*}L_i\) has no higher
Tor because \(F\) is regular on \(L_i\), and gives \(M_i\)
on the constant section of \(\mathbb P^1_R\).  The twist
\((-1)\) is trivial on that section; its projection to
\(\operatorname{Spec}R\) is the identity and has no higher direct
images.  The inverse in \eqref{eq:singular-monic-generic-equivalence}
therefore gives \(M_i[-1]\), so
\begin{equation}\label{eq:singular-monic-image-shift}
 \mathcal F(M_i)\simeq j_{c*}L_i[1].
\end{equation}

Multiplication by \(X-\alpha_c(t)\) on \(L_i\) has determinant,
up to sign, \(t^NU_c(t)\), by its characteristic polynomial.
It is injective over the domain \(\widehat D_j[[t]]\).
Since \(\alpha(z,t)\) specializes to the exact finite polynomial
\(\alpha_c(t)\), the root pullback of \(j_{c*}L_i\) has no
higher Tor.  It is killed by \(w\) and over
\(S_0=\widehat D_j[[t]]/(t^N)\) has presentation
\begin{equation}\label{eq:singular-monic-root-presentation}
 -U_ct^{N-i}f=0,\qquad t^ie+p v(t)f=0,
 \quad v(t)=(1+\alpha_c(t))^{p-1}\in S_0^\times.
\end{equation}
It is the ideal \(J_i=(p,t^i)\subset S_0\), through
\[
 e\longmapsto p v(t),\qquad f\longmapsto-t^i.
\]
Here is the complete syzygy check.  If \(a p v-bt^i=0\) in
\(S_0\), reduction modulo \(p\) gives
\(b=t^{N-i}u+p b_1\).  The element \(p\) is regular on
\(S_0\), so the original equality gives \(av=b_1t^i\).
Consequently
\[
 (a,b)=(b_1/v)(t^i,pv)+u(0,t^{N-i}).
\]
These are precisely the displayed relation columns, since \(U_c\)
is a unit.  This calculation includes \(i=0,N\).

As an \(S_j\)-module killed by \(w\), the ideal fits in
\[
 0\longrightarrow J_i\longrightarrow S_j/(w)
      \longrightarrow S_j/(p,w,t^i)\longrightarrow0.
\]
The middle term is perfect, because \(w\) is regular on \(S_j\).
The last term has length \(i\) over \(\widehat D_j[[w]]\).
This uses that \(p\) is a uniformizer, and introduces no
ramification multiplier.  The shift \([1]\) in
\eqref{eq:singular-monic-image-shift} reverses the negative sign
of the ideal class.  Hence the detector starting over
\(\widehat D_j\) satisfies
\begin{equation}\label{eq:singular-monic-witness-values}
 \rho_c([M_i])=i\pmod N.
\end{equation}

If \(k\) is infinite, \(D_j=O\), so there is no coefficient-degree
factor.  If \(k\) is finite, the completed extension
\(\widehat D_j/\widehat O\) is Galois with residue extension
\(k_j/k\).  For clarity, this follows directly from Hensel lifting:
the reduction of \(Q_j\) splits into \(m_j\) distinct roots over
\(k_j\); each lifts uniquely in \(\widehat D_j\), giving the
\(m_j\) automorphisms.  Distinct roots have unit differences, and
the Chinese remainder theorem gives
\[
 \widehat D_j\otimes_{\widehat O}\widehat D_j
          \simeq\prod_{\tau\in\operatorname{Gal}(k_j/k)}\widehat D_j.
\]
The polynomial recursion commutes with these automorphisms.
Consequently restriction to \(O\), followed by completion and
extension back to \(\widehat D_j\), splits the witness into
conjugates supported at \(\tau(\bar c_j)\).
Exactly one survives at the selected point.  Thus, with
\(m_j=[D_j:O]\) and \(m_j=1\) in the infinite case, define
\[
 y_j=[\operatorname{Res}_{D_j/O}M_1(c_j)]-m_j[A[C_p]].
\]
The evaluation is
\begin{equation}\label{eq:singular-monic-coordinate-values}
 \rho_i(y_j)=\delta_{ij}\pmod N.
\end{equation}
No residue degree is divided out, so no condition on
\(\gcd(m_j,N)\) is required.  By
\eqref{eq:singular-monic-annihilation}, \(Ny_j=0\), and the
generic regularity argument shows that \(y_j\) is generically zero.
These give the cyclic split section.
For completeness, coefficient restriction does not enlarge the original
support point: the finite \'etale coefficient map has exact pushforward
preserving perfect complexes, flat base change commutes with \(\pi^*\),
and the closed immersion \(\beta_*\) commutes with that base change.
Thus the generic-hypersurface image of a restricted witness is the
finite pushforward of its image before restriction.  Its nonperfect
support maps to its selected closed point over \(k\).  Distinct
degrees in the finite-field case, or distinct rational points in the
infinite-field case, therefore give the stated zero off-diagonal values.

We fix the normalizer convention explicitly.  For
\(a\in\{1,\ldots,p-1\}\), put
\[
 \phi_a(X)=(1+X)^a-1,\qquad
 V_a(X)=\sum_{h=0}^{a-1}(1+X)^{ph}.
\]
Then \(g(\phi_a(X))=V_a(X)g(X)\), with \(V_a(0)=a\).
The formal change
\[
 (X,z)\longmapsto(\phi_a(X),V_a(X)z)
\]
is an automorphism, multiplies the potential by a unit, and commutes
with the generic-hypersurface diagram.  Under its geometric pullback,
support at \(\bar c\) moves to \(a^{-1}\bar c\).

If its induced ring automorphism is \(h_a(\sigma)=\sigma^a\),
geometric pullback of a module is \(R\otimes_{h_a,R}M\).
The identification \(r\otimes m\mapsto h_a^{-1}(r)m\) makes
the generator act by the inverse group twist.  Thus the concrete
twist \(M^{(a)}\) whose generator acts as \(\mathcal T^a\), when
\(\mathcal T\) acts on \(M\), has its support multiplied by \(a\).
These conventions give inverse bijections of
\(\mathbb F_p^\times\), and both preserve the identity label.

The selected nonzero rational points have disjoint free scalar orbits
when \(k\) is infinite.  In the finite case the relative norm
argument above excludes every nonidentity scalar-conjugate collision,
and the distinct degrees exclude collisions between different indices.
Therefore for the concrete twists
\begin{equation}\label{eq:singular-monic-normalizer-values}
 \rho_i(y_j^{(a)})=
 \begin{cases}1,&i=j\text{ and }a=1,\\0,&\text{otherwise}.
 \end{cases}
\end{equation}
In particular the inverse labeling convention in Mackey's formula
does not change the unique surviving term.

Set \(z_j=\operatorname{Ind}_{C_p}^{S_p}y_j\), and
\(\eta_{S_p}=(\rho_i)_i\operatorname{Res}_{C_p}^{S_p}\).
For a double coset outside the normalizer, the intersection of
\(C_p\) with its conjugate is trivial, since it has prime order.
Its Mackey contribution is induced from the underlying coefficient
class of \(y_j\), which is zero by its virtual coefficient rank.
The quotient \(N_{S_p}(C_p)/C_p=\mathbb F_p^\times\) supplies
each scalar twist once.  Therefore
\(\eta_{S_p}(z_j)=e_j\).  Induction preserves the actual
annihilation relation and generic vanishing.  This gives the
symmetric-group section and proves exact order \(N\).

Finally, a degree-\(p\) Young lattice restricts to a trivial cyclic
lattice when its partition is \((p)\); every other Young subgroup
has order prime to \(p\), so its restriction is a sum of regular
cyclic lattices.  Both types are killed.  All sources, meshes and
generic isomorphisms were constructed over the original \(O\).
Only the detectors used completion.  This proves
Theorem~\ref{thm:unramified-distinguished-monic} in full.

\begin{corollary}\label{cor:distinguished-monic-regularity}
Let \(O\) be any characteristic-zero DVR with \(v_O(p)=1\), let \(n\ge2\), and
let \(F=t^n+\sum_{i<n}a_it^i\) with every \(a_i\in pO\).
For \(A=O[t]/(F)\), the following are equivalent:
\[
 \Phi_A\text{ is an isomorphism};\qquad\Phi_A\text{ is surjective};\qquad A\text{ is regular};\qquad
 a_0\in pO\setminus p^2O.
\]
In the regular case \(A\) is a DVR.  In every other case
Theorem~\ref{thm:singular-distinguished-monic} gives the stated
split infinite torsion obstruction already in degree \(p\).
\end{corollary}

\begin{proof}
The algebra \(A\) is finite free over \(O\), has dimension one,
and is local with maximal ideal \((p,t)\), since its special
fiber is \(k[t]/(t^n)\).  If \(a_0=pu\) with \(u\)
a unit, write
\(F=t^n+p(u+\sum_{i\ge1}(a_i/p)t^i)\).  The expression in
parentheses is a unit in \(A\), so \(p\in(t^n)\) and the maximal
ideal is generated by \(t\).  Thus \(A\) is a one-dimensional
regular local ring, hence a DVR; alternatively \(F\) is Eisenstein.
Its polynomial-generation property follows from
Theorem~\ref{thm:regular-local}.
If \(a_0\in p^2O\), then \(F\in(p,t)^2\) in the regular local
ring \(O[t]_{(p,t)}\).  Thus the cotangent space of \(A\) has
dimension two while \(\dim A=1\), so \(A\) is singular.  Here
\(N\ge2\), and the preceding theorem proves nonsurjectivity.
\end{proof}


\begin{corollary}\label{cor:unramified-monic-completed-field}\label{cor:singular-monic-generic-field-order}
Under the hypotheses of Theorem~\ref{thm:unramified-distinguished-monic},
suppose that \(\widehat K[t]/(F)\) is a field, where
\(\widehat K=\operatorname{Frac}(\widehat O)\).
Then \(N=n\), so the split section has countably many generators
of exact order \(n\).  This includes arbitrary residue fields.
\end{corollary}

\begin{proof}
The constant term is nonzero, since \(F\) is irreducible of
degree \(n\ge2\).  Write \(\nu=v_O(a_0)\ge2\).
A complete discretely valued field has a unique extension of its
absolute value to a finite extension; hence that valuation is
Galois invariant \citep[Theorem 4.7 and Corollary 4.8,
pp.~54--55]{WilliamsWuthrich2025}.  The cited statements require
completeness, not finiteness or perfection of the residue field.
Thus every conjugate root of \(F\) has valuation \(\nu/n\).
The elementary symmetric expressions give
\[
 v_O(a_i)\ge\frac{(n-i)\nu}{n}.
\]
If \(r<\infty\), its valuation-one coefficient forces
\(r\ge n(1-1/\nu)\ge n/2\), so \(pr\ge n\).
If \(r=\infty\), \(N=n\) by definition.  This proves the claim.
For a noncomplete base the hypothesis deliberately concerns the
completed generic algebra.
\end{proof}

\begin{example}\label{ex:singular-monic-p2-lower-coefficient}
Let \(p=2\), \(F=t^3+2t+4\), and \(A=O[t]/(F)\), with \(O=\mathbb Z_2\) or \(O=\mathbb Z_{(2)}\).  Then \(r=1\) and \(N=2<n=3\).
Take
\[
 D=O[\theta]/(\theta^3+\theta+1),\quad
 q=2\theta,\quad c=-\theta(\theta+1),\quad s=1+q,
\]
and put
\[
 a(t)=q-\frac{c}{s}t,\qquad
 U(t)=\frac{c^2}{s^2}+ct.
\]
The polynomial defining \(D\) has no root modulo two, hence is
irreducible and separable there.  Thus \(D\) is finite \'etale local
of degree three over \(O\); over \(\mathbb Z_2\) it is the
unramified cubic extension.  Both \(c\) and \(s\) are units.
Direct expansion, using \(q^2+2q+4c=0\), gives
\[
 g(a(t))+cF(t)=t^2U(t),\qquad g(X)=X^2+2X.
\]
Consequently
\[
 \mathcal T=
 \begin{pmatrix}1+a(t)&-t\\ U(t)t&-1-a(t)\end{pmatrix}
 \quad\text{satisfies}\quad
 \mathcal T^2=I-cF(t)I=I
 \quad\text{over }A_D.
\]
Let \(L\) be the corresponding rank-two \(A_D\)-lattice.
The residue \(\bar c=\bar\theta^2+\bar\theta\) has degree three,
so the one-conjugate evaluation above gives
\[
 \left[\operatorname{Res}_{D/O}L\right]-3[A[C_2]]
\]
exact order two and generic value zero.  The same finite formulas
work over \(\mathbb Z_{(2)}\), because \(s\) remains a unit.
Since \(S_2=C_2\), this symmetric-group class also lies outside the
Young span.  No assertion that the generic algebra is a field is needed.
\end{example}

\begin{remark}\label{rem:distinguished-completed-field-hypothesis}
The completed-field condition cannot be replaced by irreducibility
only over \(\Q\) for the localized base.  The polynomial
\(t^3+2t+4\) has none of its possible integer roots
\(\pm1,\pm2,\pm4\), so it is irreducible over \(\Q\).
But \(F(2u)=4(2u^3+u+1)\), and the latter cubic has a simple
root modulo two.  Hensel's lemma gives a root in \(\Q_2\).
Here \(r=1\) and \(N=2<3\).
\end{remark}

\section{Coefficient rings with a one-dimensional tangent space in the special fibre}
\label{sec:curvilinear-special-fibre}

The residue field need not be the prime field for the preceding
criterion to apply.  A Cohen coefficient ring separates the residue
field from the nilpotent direction in the special fibre.  We make
this reduction explicit, including its consequence before completion.
For an Artinian local ring $(C,\mathfrak n,k)$, write
\[
 \operatorname{edim}(C)=\dim_k\mathfrak n/\mathfrak n^2.
\]
Thus $\operatorname{edim}(C)\leq1$ means that its maximal ideal
has at most one generator.  This condition concerns the special
fibre, rather than the embedding dimension of the coefficient ring
itself.

\begin{theorem}\label{thm:curvilinear-special-fibre}
Let $(A,\mathfrak m,k)$ be a complete noetherian local ring of mixed
characteristic $(0,p)$.  Suppose that multiplication by $p$ is
injective on $A$, that $A/pA$ is Artinian, and that
\[
 \operatorname{edim}(A/pA)\leq1.
\]
Then the following conditions are equivalent:
\[
 \Phi_A\text{ is an isomorphism},\qquad
 A\text{ is regular},\qquad A\text{ is a DVR}.
\]
More precisely, there is a complete DVR $O\subset A$, with
uniformizer $p$ and residue field $k$, such that, for
$n=\operatorname{length}_{A/pA}(A/pA)$,
\begin{equation}\label{eq:curvilinear-presentation}
 A\cong O[t]/F,\qquad
 F=t^n+\sum_{i=0}^{n-1}a_it^i,\qquad a_i\in pO.
\end{equation}
If $n=1$, then $A=O$.  If $A$ is singular, then $n\geq2$ and
$a_0\in p^2O$.  Let $r$ be the least $1\leq i<n$ for which
$a_i\notin p^2O$, with $r=\infty$ if there is none, and put
\[
 N=\min(n,pr),\qquad N=n\text{ when }r=\infty.
\]
The degree-$p$ generic specialization kernel contains a subgroup
isomorphic to $\bigoplus_{\mathbb N}\Z/N$ which is split in
$K_0(\Lat_A(S_p))$.  Its retraction is defined on that whole
Grothendieck group and annihilates every Young permutation class.
The same assertion holds for $C_p$.  In the singular case, if $A[1/p]$ is a field,
then $N=n$.
\end{theorem}

The integer $n$ in the presentation is the length of the special
fibre over its own residue field.  When that field is finite of
degree $f$ over $\F_p$, the rank of $A$ over $\Z_p$ is $fn$.
It is the relative degree $n$, rather than $fn$, that appears in
the field-case obstruction.

\subsection{The coefficient presentation}

\begin{lemma}\label{lem:curvilinear-cohen-presentation}
Under the hypotheses of Theorem~\ref{thm:curvilinear-special-fibre},
the presentation \eqref{eq:curvilinear-presentation} exists.
\end{lemma}

\begin{proof}
The Cohen structure theorem supplies a coefficient ring of $A$
\citep[Theorem~10.160.8, Tag~032A]{Stacks2026}.  Since $A$ has
characteristic zero, this coefficient ring is a complete DVR $O$
with uniformizer $p$, and its residue field maps isomorphically
to $k$.  Equivalently, the map from the Cohen DVR into $A$ is
injective: its kernel is an ideal of that DVR, and a nonzero
kernel would contain a power of $p$, contrary to injectivity of
multiplication by $p$ on the nonzero ring $A$.

Put $C=A/pA$.  The map $O/pO\to C$ is injective.  Indeed, a
nonzero residue class from $O/pO$ maps to an element of $A$
whose image in $k$ is nonzero; that element is a unit, and
cannot belong to $pA$.  This identifies $k$ with a coefficient
field of $C$.

The Artinian hypothesis gives an integer $c\geq1$ with
$\mathfrak m^c\subset pA$.  Consequently
\[
 p^jA\subset\mathfrak m^j,\qquad
 \mathfrak m^{cj}\subset p^jA\quad(j\geq1).
\]
The $p$-adic and maximal-ideal-adic topologies on $A$ are
therefore equivalent.  In particular $A$ is complete and
separated for the $p$-adic topology.

Choose a $k$-basis $\bar b_1,\ldots,\bar b_d$ of $C$ and
lifts $b_1,\ldots,b_d\in A$.  These lifts form an $O$-basis
of $A$.  For surjectivity, express any element of $A$ modulo
$pA$ as a $k$-linear combination of the $\bar b_i$, lift its
coefficients to $O$, subtract, and divide the difference by
$p$.  Repeating this operation gives coefficients in $O$ as
convergent $p$-adic series, and represents the original element
as their linear combination of the $b_i$.  For independence,
reduce an equality $\sum_i c_i b_i=0$ modulo $pA$.  The basis
property forces every $c_i\in pO$.  Since $p$ is a
nonzerodivisor on $A$, divide the equality by $p$ and repeat.
Each $c_i$ then belongs to $\bigcap_jp^jO=0$.  Thus $A$ is
finite free over $O$.

Write $\mathfrak n$ for the maximal ideal of $C$.  If
$\operatorname{edim}(C)=0$, Nakayama's lemma gives
$\mathfrak n=0$, so $C=k$ and the preceding basis argument,
with $b_1=1$, gives $A=O$.

Otherwise choose $\bar t\in\mathfrak n$ whose image spans
$\mathfrak n/\mathfrak n^2$.  Nakayama's lemma gives
$\mathfrak n=(\bar t)$.  Let $n$ be the least positive
integer with $\bar t^n=0$.  Repeatedly subtracting a residue
in the coefficient field and dividing by $\bar t$ shows that
\[
 1,\bar t,\ldots,\bar t^{n-1}
\]
span $C$ over $k$.  They are independent: in a nonzero
relation, the lowest nonzero coefficient would express
$\bar t^j$ times a unit as zero for some $j<n$.  This
contradicts the definition of $n$.  Hence
$C\cong k[t]/t^n$, and both its $k$-dimension and its
composition length are $n$.

Choose any lift $t\in\mathfrak m$ of $\bar t$.  The same
basis argument now proves that $1,t,\ldots,t^{n-1}$ is an
$O$-basis of $A$.  There is a unique relation
\[
 t^n+\sum_{i=0}^{n-1}a_it^i=0,\qquad a_i\in O.
\]
Reducing modulo $p$ proves $a_i\in pO$ for every $i$.
The homomorphism $O[T]/(T^n+\sum_i a_iT^i)\to A$ carries
one displayed $O$-basis onto the other and is therefore an
isomorphism.  This proves the lemma.
\end{proof}

\subsection{Regularity and the obstruction}

\begin{proof}[Proof of Theorem~\ref{thm:curvilinear-special-fibre}]
The lemma gives the asserted presentation and proves that $A$
is finite flat over the DVR $O$.  Therefore $\dim A=1$.
When $n=1$ we have $A=O$ and the positive regular-ring theorem
applies.

Suppose $n\geq2$.  The local ring $O[t]/F$ has maximal ideal
$(p,t)$.  The localization of $O[t]$ at $(p,t)$ is a regular
local ring of dimension two.  In its cotangent space, the
image of $F$ is the image of $a_0$: the term $t^n$ and every
$a_it^i$ with $i\geq1$ belong to $(p,t)^2$.  Thus the
quotient has embedding dimension one exactly when
$a_0\notin p^2O$, and embedding dimension two otherwise.
Since it has dimension one, this is exactly the stated
regularity criterion.  A one-dimensional noetherian regular
local ring is a DVR.

In the regular case the positive theorem for regular local
coefficient rings proves that $\Phi_A$ is an isomorphism.
In the singular case,
Theorem~\ref{thm:unramified-distinguished-monic} applies to
the coefficient DVR $O$ and the polynomial $F$.  Its section
and whole-group Young-killing retraction give the claimed
split subgroup, with $N=\min(n,pr)$.  Here $N\geq2$,
because $n\geq2$, $p\geq2$, and every finite $r$ is
positive.  The retraction is nonzero, so $\Phi_A$ is not
surjective.  This proves both implications in the criterion.

It remains to verify the field-case exponent.  Suppose that
$L=A[1/p]$ is a field.  Then $F$ is irreducible over
$K=\operatorname{Frac}(O)$.  Put $s=v_O(a_0)$; the
singular case gives $s\geq2$.  Since $O$ is complete, its
valuation has a unique extension to each finite algebraic
extension of $K$.  One way to see the needed uniqueness is
to take the integral closure in a finite Galois splitting
field.  It is finite by the trace-pairing argument, is local
because $O$ is henselian, and is a DVR.  Its unique valuation
is consequently Galois invariant
\citep[Lemmas~10.153.9 and~10.161.8, Tags~04GM and~032L;
Remark~15.113.6, Tag~09E8]{Stacks2026}.
All $n$ roots of $F$ therefore have the same normalized
valuation $s/n$.  Each product of $n-i$ roots has valuation
$(n-i)s/n$, so the elementary-symmetric expression for $a_i$
gives
\[
 v_O(a_i)\geq \frac{(n-i)s}{n}.
\]
If $v_O(a_i)=1$, then
$i\geq n(1-1/s)\geq n/2$.  Hence $pi\geq n$ for every
such index.  This proves $pr\geq n$, or $r=\infty$, and
therefore $N=n$.  No assumption on the cardinality or
perfection of $k$ was used.
\end{proof}

\subsection{Finite algebras over an arbitrary DVR}

The complete criterion descends along completion for a finite free
algebra over any DVR.  Its hypothesis uses the special fibre modulo
the rational prime, which matters when the given coefficient DVR is
ramified.

\begin{corollary}\label{cor:arithmetic-curvilinear}\label{cor:arbitrary-dvr-curvilinear}
Let $O$ be any mixed-characteristic DVR of residue characteristic
$p$, and let $A$ be a nonzero commutative local algebra finite free
over $O$.  Suppose that
\[
 \operatorname{edim}(A/pA)\leq1.
\]
Then $\Phi_A$ is an isomorphism if and only if $A$ is regular.
If $A$ is singular, its degree-$p$ generic kernel contains a
countable sum $\bigoplus_{\mathbb N}\Z/N$ that is split in the
whole lattice Grothendieck group, with a retraction annihilating
the Young permutation span.  Here $N$ is computed from a Cohen
coefficient presentation of $\widehat A$ as in
Theorem~\ref{thm:curvilinear-special-fibre}.
If $\widehat A[1/p]$ is a field, then
\[
 N=\operatorname{length}_{A/pA}(A/pA).
\]
There is no completeness or residue-field hypothesis on $O$.

In particular, when $v_O(p)=1$, this criterion applies to every
finite free local monogenic $O$-algebra, including arbitrary and
imperfect residue fields.  For a ramified coefficient DVR the
hypothesis remains the displayed condition on $A/pA$.
\end{corollary}

\begin{proof}
Let $\pi$ be a uniformizer of $O$, write $p=u\pi^h$ with $u$
a unit, and let $\mathfrak m$ be the maximal ideal of $A$.
The ring $A$ is noetherian of dimension one.  The quotient
$A/\pi A$ is an Artinian local ring, so
$\mathfrak m^c\subset\pi A$ for some $c$.  It follows that
the $\mathfrak m$-adic and $\pi$-adic topologies on $A$ agree;
they also agree with the $p$-adic topology.  Its completion is
\[
 \widehat A=A\otimes_O\widehat O.
\]
This ring is local and complete, and
$\widehat A/p\widehat A=A/pA$.  In particular the special fibre,
its length, and its embedding dimension are unchanged.
Freeness over $O$ makes $p$ a nonzerodivisor on $A$, and faithful
flatness gives the same conclusion for $\widehat A$.
Both rings have dimension one; their maximal-ideal cotangent
spaces agree.  Thus $A$ is regular if and only if $\widehat A$
is regular, also the general completion criterion
\citep[Lemma~15.44.4, Tag~07NY]{Stacks2026}.

If $A$ is regular, the positive theorem for regular local
coefficient rings gives the assertion.  If it is singular,
Theorem~\ref{thm:curvilinear-special-fibre} gives the stated
split subgroup $H$ of the generic kernel over $\widehat A$.
Let $s:H\to K_0(\Lat_{\widehat A}(S_p))$ and $q$ be its
section and retraction.  Inverting $p$ is the same as inverting
$\pi$, so Theorem~\ref{thm:completion-dvr-kernel} identifies
the completed generic kernel with the one over $A$.
Lift $s$ through that isomorphism to $s_A$, and compose $q$
with completion on the whole group to obtain $q_A$.
Then $q_As_A=1_H$.  Completion preserves every Young lattice,
so $q_A$ still annihilates the Young span.  It is nonzero because
$N\geq2$ in the singular case, proving failure of generation.
The completed-field formula follows from the complete theorem
and the unchanged special-fibre length.

Finally suppose $v_O(p)=1$ and write $k_0=O/pO$.  If $A$ is
monogenic over $O$, then $A/pA$ is a finite-dimensional local
quotient of $k_0[T]$.  Its defining ideal is $(f^e)$ for a
monic irreducible polynomial $f$ and an integer $e\geq1$.
Its maximal ideal is generated by the image of $f$, so its
embedding dimension is at most one.  This argument does not
require $f$ to be separable.  The length over the residue field
$k_0[T]/f$ is $e$, which is the field-case exponent; the absolute
$O$-rank $e\deg f$ need not equal it.  The Cohen coefficient
ring used after completion need not extend the original
coefficient DVR, and no such compatibility is required.
\end{proof}


\begin{remark}
For a ramified coefficient DVR with $v_O(p)\geq2$, the monogenic
algebra $A=O[t]/t^2$ illustrates the distinction between the two
special fibres.  In the regular local ring $O[t]_{(\pi,t)}$, both
$p$ and $t^2$ belong to $(\pi,t)^2$.  Thus $A/pA$ has two
independent cotangent directions, represented by $\pi$ and $t$,
whereas $A/\pi A$ has a principal maximal ideal.  The hypothesis
of Corollary~\ref{cor:arbitrary-dvr-curvilinear} is the former one.
\end{remark}

\begin{example}\label{ex:curvilinear-residue-field}
Let $D=\Z_2[\omega]$ with $\omega^2+\omega+1=0$, and set
\[
 A=D[t]/(t^2+4).
\]
The coefficient DVR $D$ has residue field $\F_4$.  The
ring $A$ has rank four over $\Z_2$, while its special fibre
is $\F_4[t]/t^2$ and has length two over its residue field.
Its maximal ideal is $(2,t)$ and its embedding dimension
is two, so it is singular.  Here $r=\infty$ and $N=2$.
The theorem gives a split countable sum of exact-order-two
classes in the degree-two generic kernel, with a
whole-group Young-killing retraction.  The nontrivial
residue field creates no additional multiplicity in this
cyclic order.
\end{example}

\section{Consequences and explicit coefficient families}
\label{sec:monic-depth-two-split}

\begin{corollary}\label{thm:monic-depth-two-split}
Let $p$ be a prime, $n\ge2$, $O=\Z_p$ or $\Z_{(p)}$, and
\[
 F(t)=t^n+\sum_{i=0}^{n-1}a_it^i,\qquad a_i\in p^2O,
 \qquad A=O[t]/F(t).
\]
There are homomorphisms
\[
 s:\bigoplus_{i\ge1}\Z/n\longrightarrow K_0(\Lat_A(S_p)),
 \qquad
 \eta:K_0(\Lat_A(S_p))\longrightarrow\bigoplus_{i\ge1}\Z/n
\]
with $\eta s=1$, with the image of $s$ contained in the generic
specialization kernel, and with $\eta$ annihilating the degree-$p$
Young permutation span.  Thus this kernel contains a countable
direct sum of exact order-$n$ classes that splits from the whole
Grothendieck group.  The analogous assertion holds for $C_p$.
In particular $\Phi_A$ is not surjective in degree $p$.
\end{corollary}


\begin{proof}
Every nonconstant coefficient has valuation at least two, so
\(r=\infty\) and \(N=n\) in
Theorem~\ref{thm:singular-distinguished-monic}.
\end{proof}

The next construction gives finite companion matrices and a second
description of the coordinates for binomial equations.  The general
classification above includes all these coefficient rings.

\subsection{The binomial construction}
\label{sec:binomial-depth-two-all}

The coefficient Bockstein gives an orientation on a stable module over
$\F_p[t]/t^n$.  Its length modulo $n$ supplies the additive invariant
for every truncation exponent, including when the generic coefficient
algebra is a totally ramified field.  The construction uses all finite
coefficient test modules, so its additivity is established on the full
lattice category.

\begin{theorem}\label{thm:binomial-depth-two-all-exponents}
Let $p$ be a prime, $n\ge2$, $O=\Z_p$ or $\Z_{(p)}$, and
$d\in p^2O$.  Put $A=O[t]/(t^n-d)$.  The kernel
\[
 \ker\!\left(K_0(\Lat_A(S_p))
       \longrightarrow K_0(\Lat_{A[1/p]}(S_p))\right)
\]
contains a subgroup isomorphic to
$\bigoplus_{i\ge1}\Z/n$ that is split in the whole lattice
Grothendieck group.  It has a coordinate retraction annihilating
the degree-$p$ Young permutation span.  Consequently $\Phi_A$
is not surjective in degree $p$.
The same split torsion assertion holds for the cyclic precursors.
\end{theorem}

The case $d=0$ is Theorem~\ref{thm:all-truncations}.
At $p=2$, the cases $d\in8O$ follow from
Corollary~\ref{cor:deep-monogenic-orders}.
We prove the remaining assertions over $O=\Z_p$, and then pass
to $\Z_{(p)}$ by completion.

\subsection{Curvature components and exact residual Tate complexes}

Write $F=t^n-d$, $k=\F_p$, and $R=k[t]/t^n$.
For an $A$-free $C_p$-lattice $E$ with generator matrix $T$,
choose a polynomial lift $\widetilde T$ and put
\[
 \widetilde U=\frac{\widetilde T^p-I}{F},
 \qquad U_E=\widetilde U\bmod F.
\]
The division is possible in $O[t]$.  Cancellation of $F$ shows
that $U_E$ is an $A[C_p]$-endomorphism.
A map is called stably zero if it factors through a finite free
$A[C_p]$-lattice.

\begin{lemma}\label{lem:binomial-all-primary}
The stable class of $U_E$ is independent of the lift and basis
and is natural in $E$.  Stable Hom groups are killed by $p$.
For each irreducible $f\in k[z]$, the integral $f$-primary
summand $E_f$ is $A$-free and is well-defined up to stable
isomorphism.  If $f\ne z$, then $E_f$ and $E_f/pE_f$ have zero
Tate cohomology.
\end{lemma}

\begin{proof}
For an equivariant map $B:E\to E'$, lift it and write
$\widetilde T'\widetilde B-\widetilde B\widetilde T=F\widetilde C$.
Telescoping gives
\[
 U_{E'}B-BU_E
 =\sum_{i=0}^{p-1}T'^{\,p-1-i}CT^i.
\]
This is a group trace and factors through a regular lattice.
The same calculation compares two lifts.  The group trace of
an equivariant map is $p$ times that map.

The usual integral primary projectors for an endomorphism of a
finite free $O$-module are obtained by Chinese remainders modulo
$p^a$ and passage to the inverse limit.  They are $A[C_p]$-linear.
Their summands are $A$-free because $A$ is local, and the
projectors are exact on intertwined short exact sequences.
On the $f$-part some power of $f(U_E)$ is divisible by $p$,
and is therefore stably zero.  Bezout's identity over $k[z]$
makes every stable intertwining map between distinct primary
parts zero.  Comparing two choices using the identity map and
its inverse consequently gives stable isomorphisms of their
corresponding primary parts.

For $f\ne z$, $U_E$ is invertible on $E_f$.  Lift a coefficient
basis adapted to this summand to the local ring
$O[t]_{(p,t)}$.  The off-diagonal action blocks are divisible
by $F$; every term leaving and returning to a diagonal block
of $\widetilde T^p$ is divisible by $F^2$.
The curvature of the diagonal lift is therefore the restricted
curvature, and is a unit in the local ring.

If $(T-I)x=0$ modulo $F$, lift and write
$(\widetilde T-I)x=Fy$.  Multiplication by $N(\widetilde T)$
and cancellation of $F$ give
$\widetilde Ux=N(\widetilde T)y$.
Thus $x$ is a norm boundary.  Reversing the two factors proves
the other Tate exactness assertion.  Multiplication by $p$
on $E_f$ is injective; its coefficient long exact sequence
then proves the same vanishing for $E_f/pE_f$.
\end{proof}

For a compatible action lift on an admissible sequence, choose
upper triangular blocks.  Its curvature is upper triangular
with the endpoint curvatures on the diagonal.  The primary
projectors therefore give admissible sequences of the selected
components.  This observation will supply exactness before any
coordinate is evaluated.

\subsection{A stable module representing all coefficient tests}

Let $E$ be a selected cohomologically trivial component, and put
$\overline E=E/pE$.  Its exact, two-periodic Tate complex has
free $R$-modules as terms and differentials
\[
 D^0=T-I,\qquad D^1=N(T).
\]
Write $(-)^*=\operatorname{Hom}_R(-,R)$ and define
\[
 M_E=\operatorname{coker}
 \bigl((D^0)^*:\overline E^*\longrightarrow\overline E^*\bigr).
\]

The coefficient-of-$t^{n-1}$ pairing identifies $R$ with its
$k$-linear dual, so $R$ is a symmetric Frobenius algebra.
Its finite-module stable category
$\mathcal S_n=\underline{\operatorname{mod}}(R)$ is triangulated,
with triangles induced by short exact sequences and suspension
given by an injective embedding
\citep[Section~13.4, Theorem~13.7, p.~50 of the author's
version]{Buhler2010}.

\begin{lemma}\label{lem:binomial-all-representability}
For every finite $R$-module $X$ there is a natural isomorphism
\begin{equation}\label{eq:binomial-all-representable}
 \widehat H^0(C_p,\overline E\otimes_R X)
 \cong \underline{\operatorname{Hom}}_R(M_E,X).
\end{equation}
The assignment $E\mapsto M_E$ is contravariant.
An admissible sequence of selected components gives an actual
short exact sequence
\begin{equation}\label{eq:binomial-all-transpose-sequence}
 0\longrightarrow M_{E''}\longrightarrow M_E
   \longrightarrow M_{E'}\longrightarrow0,
\end{equation}
and the resulting stable Hom long exact sequence agrees with
the Tate long exact sequence under \eqref{eq:binomial-all-representable}.
\end{lemma}

\begin{proof}
Dualizing the exact Tate complex preserves exactness because
$R$ is self-injective.
Ordinary $\operatorname{Hom}_R(M_E,X)$ identifies with the
kernel of $D^0$ on $\overline E\otimes_R X$.
The dual complex embeds $M_E$ as the image of $(D^1)^*$ in
a free module.  Every map from $M_E$ to a projective extends
across this embedding, since projectives are injective.
Maps factoring through projectives are consequently exactly
the image of $D^1$ on $\overline E\otimes_R X$.  This proves
\eqref{eq:binomial-all-representable}.

An equivariant map $E\to E'$ dualizes to a map of complexes
in the opposite direction and hence gives $M_{E'}\to M_E$.
For an admissible sequence, the dual complexes form a
degreewise split short exact sequence of exact complexes.
The induced map on cycles is surjective: express a quotient
cycle as a differential, lift its preimage, and apply the
differential.  The cokernel sequence is therefore
\eqref{eq:binomial-all-transpose-sequence}.
Applying Hom to these complete resolutions identifies its
connecting map with the Tate connecting map, as asserted.
\end{proof}

\subsection{Lifting every test-module map and its suspension}

Put
\[
 \widetilde R=(O/p^2)[t]/t^n=A/p^2A.
\]
Every finite $R$-module is a direct sum of
$X_i=R/(t^i)$, $1\le i\le n$.  Use the standard lifts
$\widetilde X_i=(O/p^2)[t]/t^i$ and their direct sums.
Every map $X_i\to X_j$ lifts: its value at $1$ is divisible
by $t^{\max(j-i,0)}$, and the remaining coefficients lift
arbitrarily.  This also proves the assertion for matrices
of maps between direct sums.  An isomorphism lifts to an
isomorphism because it is invertible modulo $p$ between
finite free $O/p^2$-modules.

For a test module $X$ these lifts give a coefficient sequence
\[
 0\longrightarrow\overline E\otimes_R X
 \xrightarrow{p}(E/p^2E)\otimes_{\widetilde R}\widetilde X
 \longrightarrow\overline E\otimes_R X\longrightarrow0.
\]
It is exact since $E/p^2E$ is $\widetilde R$-free and
$\widetilde X$ is $O/p^2$-free.  Let $\beta_X$ be its Tate
Bockstein.  Lifting each test-module map proves naturality
on the entire category of finite $R$-modules.  A functorial
choice of lifts of all maps is unnecessary: any lift gives
the required naturality square with the prescribed reduction.
In particular, changes of cyclic decomposition of $X$
do not change the resulting operator.

Set
\[
 Q_p(z)=\frac{N(z)-p}{z-1},\qquad
 h_{\rm even}=I,\qquad h_{\rm odd}=-Q_p(T).
\]
The identity $dh+hd=pI$ makes $h$ an odd Tate operator modulo
$p$.  If a lifted cycle satisfies $dz=pw$, then
$d(hz)=p(z-hw)$.  Division by $p$ proves
\begin{equation}\label{eq:binomial-all-beta-h}
 \beta h+h\beta=I.
\end{equation}
All these operators are natural in both $E$ and $X$.
Define the even operator $P_X=\beta_Xh_X$.

\begin{lemma}\label{lem:binomial-all-boundaries}
The operator $P$ commutes with every Tate connecting map of
an admissible lattice sequence.  It also commutes with the
coefficient connecting maps implementing the stable suspension
in the variable $X$.
\end{lemma}

\begin{proof}
For a lattice sequence, reduce modulo $p^2$, tensor with
$\widetilde X$, and choose an $O/p^2$-linear splitting.
The differentials have upper triangular blocks.  Their
off-diagonal block $Z$ is an odd cochain map modulo $p^2$.
On a cycle with $dz=pw$, division by $p$ therefore gives
$\beta'\partial+\partial\beta''=0$.
The upper block of $dh+hd=pI$ gives
$h'\partial+\partial h''=0$ modulo $p$.
Their product commutes with $\partial$.

For suspension use
\[
 0\longrightarrow X_i\xrightarrow{t^{n-i}}R
 \longrightarrow X_{n-i}\longrightarrow0.
\]
These sequences lift to exact sequences of the standard
$\widetilde R$-modules and are split over $O/p^2$.
The identical two-boundary calculation shows that both
$\beta$ and $h$ anticommute with their coefficient connecting
maps.  Hence $P$ commutes with them.  The middle test module
$R$ has zero Tate cohomology, so these connecting maps
identify the two parities with stable suspension.
Direct sums cover every finite test module.
\end{proof}

By Yoneda and Lemma~\ref{lem:binomial-all-representability},
$P^0_X$ is induced by a unique stable endomorphism
$P_E$ of $M_E$.  The suspension assertion identifies $P^1$
with the shifted endomorphism.  Consequently the first
assertion of Lemma~\ref{lem:binomial-all-boundaries} says that
these endomorphisms commute with all three arrows of the
triangle \eqref{eq:binomial-all-transpose-sequence},
including its connecting arrow.

\subsection{Stable summands and additive length}

The category $\mathcal S_n$ is idempotent complete.
Indeed, every finite $R$-module is a sum of the cyclic
$X_i$.  Remove the free summands $X_n$.
On $X_i$, $i<n$, a map factoring through a free module is
divisible by $t^{n-i}$ and belongs to the radical of its
endomorphism ring.  The projective-factor ideal on a finite
sum is therefore contained in the radical of its finite-dimensional
endomorphism algebra.  Idempotents lift through this nilpotent
ideal and split as module summands.

The function
\[
 \lambda(M)=\dim_k M\pmod n
\]
is well-defined on stable isomorphism classes.
For completeness, stable inverse maps give a split inclusion
$M\to N\oplus Q$ with $Q$ projective; its complementary
idempotent factors through projectives, so its complement
is projective.  Thus a stable isomorphism becomes an actual
isomorphism after adding free $R$-modules.  Such modules
have dimensions divisible by $n$.
The function $\lambda$ is additive on distinguished triangles
because they are induced by short exact sequences.

We also use that a direct summand of a distinguished triangle
is distinguished
\citep[Lemma~1.6 and proof, pp.~821--822]{BalmerSchlichting2001}.
Thus a compatible idempotent on a triangle can be split
without losing the relation needed for $\lambda$.

\subsection{Odd primes}

Suppose $p$ is odd.  The group generator acts trivially on
Tate cohomology, and
\[
 h^2=-Q_p(1)=-p(p-1)/2=0.
\]
Equation~\eqref{eq:binomial-all-beta-h} gives
\[
 P^2=\beta(1-\beta h)h=P-\beta^2h^2=P.
\]
No assertion that $\beta^2=0$ is used.
Let $M_E^+$ denote the image of the stable idempotent $P_E$.
The preceding triangle compatibilities and direct-summand
lemma imply
\[
 \lambda(M_E^+)=\lambda(M_{E'}^+)+\lambda(M_{E''}^+).
\]

For each irreducible $f\in k[z]$, $f\ne z$, define
\begin{equation}\label{eq:binomial-all-odd-coordinate}
 \eta_f(L)=-\lambda(M_{L_f}^+)\in\Z/n.
\end{equation}
These are independent of the curvature choices.
In fact, a map factoring through $A[C_p]^a$ acts as zero on
every test functor: after reduction and tensoring with $X$
the regular module is free over $k[C_p]$.
Stable equivalences of the primary components therefore
give isomorphisms of their representable stable Hom functors,
intertwining the natural Bockstein and $h$.
Yoneda gives compatible stable isomorphisms of their
oriented summands.

For an admissible sequence choose the compatible primary
projectors from Lemma~\ref{lem:binomial-all-primary}.
The full triangle argument now proves additivity of
\eqref{eq:binomial-all-odd-coordinate}.
Thus these are homomorphisms on the full lattice $K_0$.
Each class has finite primary support.  Constant $O$-lattices
have a constant action lift with curvature zero, so every
$\eta_f$ annihilates them.

\subsection{The prime two at valuation two}

Suppose $p=2$ and $d=4b$ with $b\in O^\times$.
Let $E$ be the entire invertible Fitting part of the curvature.
Here $h^2=-I$ and $P$ is not assumed idempotent.
For any standard test module $X_i$, lift to $O[t]/t^i$.
The products of the lifted Tate differentials are
$-4b\widetilde U(t)$, since $i\le n$.
If $D_i z=2w$, then
$D_{i+1}w=-2b\widetilde U(t)z$.
Dividing twice by two gives, on every test module,
\[
 \beta^2=-\overline b\,U(t)=U(t),\qquad
 P^2+P=U(t).
\]
This uses the full curvature polynomial on $X_i$, not merely
its value at $t=0$.

The stable endomorphism $P_E$ of $M_E$ still commutes with
all arrows of the associated triangle.
For an irreducible $q\in\F_2[z]$, $q\ne z,z+1$, let
$M_{E,q}$ be its primary stable summand.
For a triangle, choose a common annihilating polynomial
of the three finite-dimensional stable endomorphisms.
Chinese remainders give one polynomial primary projector
on all three vertices; it is an idempotent endomorphism
of that triangle.  Define
\begin{equation}\label{eq:binomial-all-two-coordinate}
 \eta_q(L)=\lambda(M_{E,q})\in\Z/n.
\end{equation}
The direct-summand lemma proves additivity.
The same stable-equivalence argument used at odd primes
proves choice independence.  These coordinates have finite
support and vanish on constant lattices.
The relation $P(P+I)=U(t)$ places the selected $q$-part
over invertible curvature.

\subsection{Binomial witnesses and their exact annihilation}

Let $O_s/O$ be a finite unramified extension, and choose
units $\alpha,c\in O_s$ with
\begin{equation}\label{eq:binomial-all-parameter}
 \alpha^p+cd=1.
\end{equation}
Put $A_s=O_s\otimes_O A$, let $C_\alpha$ be the companion
of $q_\alpha(z)=(z^p-\alpha^p)/(z-\alpha)$,
and let $a$ be its last coordinate row and $b_0$ its first
coordinate column.  Define
\[
 T_j=
 \begin{pmatrix}
 \alpha&ct^j a\\ t^{n-j}b_0&C_\alpha
 \end{pmatrix},
 \qquad 0\le j\le n.
\]
The companion identity gives
\[
 \det(zI-T_j)=z^p-\alpha^p-ct^n=z^p-1,\qquad
 \widetilde T_j^p-I=c(t^n-d)I.
\]
Thus $L_j=(A_s^p,T_j)$ is a lattice with curvature $c$.

Set $D_-=\operatorname{diag}(1,tI_{p-1})$ and
$D_+=\operatorname{diag}(t,I_{p-1})$.
For $1\le j<n$ they give coefficient-split exact sequences
\[
 0\longrightarrow L_j
 \xrightarrow{(D_-,D_+)}L_{j-1}\oplus L_{j+1}
 \xrightarrow{(D_+,-D_-)}L_j\longrightarrow0.
\]
On each coordinate these are the split sequences
$x\mapsto(x,tx)$, $(v,w)\mapsto tv-w$, or their transposes.
The endpoint lattices are regular: at $j=0$ the unit upper
attachment gives a cyclic companion vector, and at $j=n$
the unit lower attachment does so.
Writing $r_j=[L_j]$, the relations
$2r_j=r_{j-1}+r_{j+1}$ and $r_n=r_0$ yield
\begin{equation}\label{eq:binomial-all-annihilator}
 n\bigl([L_1]-[A_s[C_p]]\bigr)=0.
\end{equation}
This is an actual exact-category relation, independent
of the coordinate calculations.

\subsection{The transpose cokernel and its selected summand}

We give an explicit presentation.
All residue calculations below are over
$R_s=k_s[t]/t^n$, where $k_s=O_s/pO_s$.
Since $\alpha\equiv1\pmod p$, put $m=p-1$ and write
\[
 \overline q_\alpha(z)
 =z^m+q_{m-1}z^{m-1}+\cdots+q_0=(z-1)^m.
\]
Present $M_{L_j}=\operatorname{coker}(T_j-I)^*$ by
generators $x,y_1,\ldots,y_m$.
The top action row gives $\overline c\,t^j y_m=0$.
For $2\le r\le m$, the companion rows give
\[
 y_{r-1}-y_r-q_{r-1}y_m=0.
\]
Eliminating $y_1,\ldots,y_{m-1}$ with these unit-coefficient
relations leaves the first companion-row relation
\[
 t^{n-j}x-\bigl(1+q_{m-1}+\cdots+q_0\bigr)y_m
 =t^{n-j}x-\overline q_\alpha(1)y_m=t^{n-j}x=0.
\]
For $m=1$ the elimination list is empty and the same final
relation holds directly.  Since $\overline c$ is a unit,
\begin{equation}\label{eq:binomial-all-cokernel}
 M_{L_j}\cong R_s/(t^{n-j})\,x\ \oplus\
                    R_s/(t^j)\,y_m.
\end{equation}
The first generator corresponds under
\eqref{eq:binomial-all-representable} at $X=k_s$ to the
even trivial Tate class; the second corresponds to the
even cyclotomic Tate class.

At odd $p$, choose $d=p^v b$, $v\ge2$, and a unit $y\in O_s$.
Set
\[
 \alpha=1+p^{v-1}y,\qquad c=(1-\alpha^p)/d.
\]
The binomial expansion shows that $c$ is a unit and
$\overline c=-\overline y/\overline b$, so every nonzero
residue parameter is available.
On the trivial class $h_{\rm even}=I$ and
$\beta_{\rm odd}=I$, because
\[
 \frac{N(\alpha)}p\equiv1\pmod p.
\]
For the latter identity expand
$N(1+s)=\sum_{\ell=1}^p\binom p\ell s^{\ell-1}$ with
$s\in pO_s$; after division by $p$ every term after the
constant term is divisible by $p$.
On the even cyclotomic class, represented by
$(C_\alpha-I)^{p-2}b_0$ modulo $p$, the map $h_{\rm even}=I$
has odd image in $\operatorname{im}(C_\alpha-I)$, and hence
is zero in cohomology.
Thus $P$ selects the trivial even class and kills the
cyclotomic even class.

Under \eqref{eq:binomial-all-cokernel}, the selected stable
summand is consequently $R_s/(t^{n-j})$.
Indeed, the induced idempotent has diagonal residues $1,0$
in the semisimple quotient of the endomorphism algebra
of these cyclic modules.  Krull--Schmidt and idempotent
lifting determine its image up to stable isomorphism.
For $j=1$ this also covers $n=2$, when the two cyclic
summands have equal length and the selected multiplicity
is one.  If $\overline c$ has degree $s$ and minimal
polynomial $f$, coefficient restriction to $O$ therefore gives
\begin{equation}\label{eq:binomial-all-odd-value}
 \eta_f(L_1)=s\pmod n.
\end{equation}

At $p=2$, $d=4b$, choose an odd prime degree $s>n$ and
a primitive element $\overline y\in k_s^\times$, with
an integral lift $y$.  Put
\[
 \alpha=1+2y,\qquad c=-y(1+y)/b.
\]
Both $y$ and $1+y$ are units and
\eqref{eq:binomial-all-parameter} holds.
On even residual Tate cohomology,
\[
 P_{\rm even}=\operatorname{diag}(1+\overline y,\overline y),
\]
since $h_{\rm even}=I$ and
$\beta_{\rm odd}=(T_j(0)+I)/2$.
The elements $\overline y$ and $1+\overline y$ have
different absolute traces because $s$ is odd, so their
minimal polynomials are distinct.
If $q$ is the minimal polynomial of $\overline y$,
its primary projector in \eqref{eq:binomial-all-cokernel}
selects $R_s/(t^j)y_m$.
The same endomorphism-algebra argument just used makes
this a statement about the actual stable summand.
Thus
\begin{equation}\label{eq:binomial-all-two-value}
 \eta_q(L_1)=s\pmod n.
\end{equation}

\subsection{Induction, retractions, and the two coefficient bases}

\begin{proof}[Proof of Theorem~\ref{thm:binomial-depth-two-all-exponents}]
For odd $p$ and $d\ne0$, choose distinct prime degrees
$s_i>n(p-1)$ and nonzero residue parameters
$\overline c_i$ of those degrees.
Let $f_i$ be their minimal polynomials and restrict the
corresponding $L_1$ from $A_{s_i}$ to $A$, obtaining $L_i$.
The scalar orbits of the parameters under $\F_p^\times$
are free.  If $a\overline c_i$ is Frobenius-conjugate to
$\overline c_i$, taking the field norm gives $a^{s_i}=1$;
coprimality with $p-1$ forces $a=1$.
Different degrees give disjoint orbits.

Changing the cyclic generator to its $a$th power multiplies
the stable curvature by $a$, as follows by expanding
$(I+F\widetilde U)^a$ and dividing by $F$.
Mackey restriction of induction from $C_p$ to $S_p$ supplies
each normalizer twist once.  Every other intersection is
trivial and gives a regular cyclic lattice.
All nonidentity normalizer twists have zero $f_i$-primary
part before applying $P$, so no invariance of $P$ under
generator changes is required.
For
\[
 x_i=[\operatorname{Ind}_{C_p}^{S_p}L_i]
                            -s_i[A[S_p]]
\]
we obtain
\[
 \bigl(\eta_{f_j}\operatorname{Res}_{C_p}^{S_p}\bigr)(x_i)
 =\begin{cases}s_i\pmod n,&i=j,\\0,&i\ne j.\end{cases}
\]
Normalize each coordinate by $s_i^{-1}\in\Z/n$.
Equation~\eqref{eq:binomial-all-annihilator}, preserved by
coefficient restriction and induction, gives $nx_i=0$.
The normalized coordinates, with finite support on each
class, are a retraction of the homomorphism sending the
standard generators of $\bigoplus_i\Z/n$ to the $x_i$.
This proves exact order $n$ and splitting in the whole group.
The cyclic assertion follows before induction.

For $p=2$ and $v_2(d)=2$, choose distinct odd prime degrees
$s_i>n$ and the parameters in
\eqref{eq:binomial-all-two-value}.
The two primary parameters for each witness have degree
$s_i$, so different degrees separate all witnesses.
The values $s_i$ are units modulo $n$.
The same annihilation relation and normalized coordinates
give the split direct sum.  Here $C_2=S_2$.

The remaining nonzero case is $p=2$ and $d=2^v b$ with
$v\ge3$ and $b$ a unit.  Choose the unramified residue parameters
of Theorem~\ref{thm:all-truncations}, and set
$\alpha=1+2^{v-1}y$, $c=(1-\alpha^2)/d$.
Then $c$ is a unit with residue $-\overline y/\overline b$,
and $\alpha\equiv1\pmod4$.  The matrices $T_j$ above reduce
modulo four to the truncated-polynomial witnesses with the same
parameter $c$.  The coefficient comparison of
Theorem~\ref{thm:congruent-coefficients}, at precision three for
$C_2$, transports those witnesses to exact actions with the same
reduction modulo four.  Its rigidity argument identifies the
transported actions with $T_j$, since four has valuation strictly
larger than $v_2(|C_2|)=1$.  Thus the transported coordinates take
the same values on these explicit companions.  The mesh relation
\eqref{eq:binomial-all-annihilator} gives their exact order $n$.
This also supplies explicit generators in the deeper case.

For the constructed witnesses, $t$ is invertible after
inverting $p$, since $d\ne0$.  The diagonal maps relating
successive $L_j$ become isomorphisms, and $L_0$ is regular.
Hence the classes vanish generically.
Every Young permutation lattice, and its restriction to
$C_p$, is extended from $O$.  Its curvature is zero, so
all coordinates annihilate the Young span.
This proves non-surjectivity in degree $p$.

The cases $d=0$ and $p=2,\ d\in8O$ were covered at the
start by the truncation and deep-monogenic results.
Finally let $O=\Z_{(p)}$.  The completion of $A$ is the
corresponding binomial algebra over $\Z_p$.
Theorem~\ref{thm:completion-generic-kernel} lifts the
completed split subgroup through the generic-kernel
isomorphism.  Compose its completed retraction with
completion to obtain a retraction on the whole integral
Grothendieck group.  Young permutation lattices complete
to the corresponding Young lattices, so this retraction
still annihilates their span.

The companion witnesses constructed above can also be chosen
over $\Z_{(p)}$ itself.  For each residue parameter choose a
monic lift of its separable minimal polynomial and let $O_s$
be the resulting finite free local etale $\Z_{(p)}$-algebra.
Its completion is the unramified extension used above.
For odd $p$, and also for $p=2$ with $v_2(d)\ge3$, the formulas
$\alpha=1+p^{v-1}y$ and $c=(1-\alpha^p)/d$ lie in $O_s$:
the numerator is divisible by $p^v$ and $d/p^v$ is a unit.
For $p=2$, $v_2(d)=2$, use
$\alpha=1+2y$ and $c=-y(1+y)/(d/4)$.
The companion matrices, coefficient restriction, induction,
and mesh sequences are then defined over $\Z_{(p)}$.
Their completed detector values are exactly those computed
above, so they give the same sections over the localized integers.
This proves all assertions.
\end{proof}

\begin{example}\label{ex:quintic-depth-two-twentyfive}
Let $p=n=5$, $A=\Z_5[t]/(t^5-25)$, $\alpha=-4$ and $c=41$.
Then $\alpha^5+25c=1$, and
\[
 q_\alpha(z)=z^4-4z^3+16z^2-64z+256.
\]
With $C_\alpha$ its companion, last coordinate row $a$,
and first coordinate column $b_0$, the matrix
\[
 T=\begin{pmatrix}-4&41t\,a\\t^4b_0&C_\alpha\end{pmatrix}
\]
satisfies $T^5=I$.
Its curvature has residue $1$, whose scalar orbit under
$\F_5^\times$ is free.  The degree-one coordinate in
\eqref{eq:binomial-all-odd-value} shows that
$[\operatorname{Ind}_{C_5}^{S_5}L]-[A[S_5]]$ has exact
order five.
If $\pi^5=5$, then $t=\pi^2$ and
$A=\Z_5[\pi^2]$ has index $25$ in the totally ramified
maximal order $\Z_5[\pi]$.
Thus the conclusion includes the exponent equal to
the residue characteristic and a totally ramified generic field.
\end{example}

\subsection{Arithmetic domain examples}

The following small matrices exhibit the obstruction in nonmaximal
orders of totally ramified fields.

\begin{example}\label{ex:cubic-depth-two-nine}
Take $p=3$, $d=9$, $\alpha=4$ and $c=-7$.  Then
$\alpha^3+9c=1$, and over $A=\Z_3[t]/(t^3-9)$ the matrix
\[
 T=\begin{pmatrix}
 4&0&-7t\\ t^2&0&-16\\0&1&-4
 \end{pmatrix}
\]
satisfies $T^3=I$.  The residue parameter is $2$, whose degree-one
minimal polynomial has a free scalar orbit under
$\F_3^\times$.  The corresponding coordinate shows that
$[\operatorname{Ind}_{C_3}^{S_3}L]-[A[S_3]]$ has exact order
three.  Finite unramified parameter extensions give the
countable family for this same coefficient ring.

If $\pi^3=3$, then $t=\pi^2$ has valuation $2/3$, so
$t^3-9$ is irreducible and its field is the totally ramified
cubic field $\Q_3(\pi)$.  Its order is
\[
 \Z_3[\pi^2]=\Z_3\cdot1\oplus\Z_3\cdot\pi^2
                         \oplus\Z_3\cdot3\pi,
\]
of index three in the maximal order $\Z_3[\pi]$.
\end{example}

\begin{example}\label{ex:deep-cubic-binomial}
Take $p=3$, $n=3$, and $d=81$.  With $\alpha=28$ and $c=-271$,
\[
 T=\begin{pmatrix}
 28&0&-271t\\ t^2&0&-784\\0&1&-28
 \end{pmatrix},\qquad T^3=I.
\]
Indeed $28^3+81(-271)=1$.  The induced class minus the regular
class has exact order three: modulo $3$ its parameter is $2$,
and the corresponding truncated-polynomial coordinate is one.
The coefficient ring is a totally ramified cubic local domain.
If $\pi^3=3$, then $t=3\pi$ satisfies $t^3=81$ and has valuation
$4/3$; its degree over $\Q_3$ is therefore three.  The infinite
family in the corollary is obtained by the indicated finite
coefficient extensions.
\end{example}

\begin{example}\label{ex:quartic-binomial-order}
Let $O=\Z_{(3)}$ or $\Z_3$, and put
\[
 A=O[t]/(t^4-27).
\]
The choices $\alpha=-8$ and $\kappa=1/19$ give
\[
 T=\begin{pmatrix}
 -8&0&19t\\
 t^3&0&-64\\
 0&1&8
 \end{pmatrix},\qquad T^3=I.
\]
Indeed $\alpha^3=1-27\cdot19$, and its characteristic polynomial
is $z^3+512-19t^4=z^3-1$.  Restriction to
$O[s]/(s^2-27)$, $s=t^2$, has the nonzero coordinate $\mu_{u-1}$,
since $19\equiv1\pmod3$.  The induced degree-three class is nonzero,
has exact order four and vanishes generically.  Indeed, the
stable-module coordinate in Theorem~\ref{thm:binomial-depth-two-all-exponents}
is one modulo four: its residue parameter is $1$, of degree one,
and its scalar orbit under $\mathbb F_3^\times$ is free.

This is a quartic local-domain example.  A root of $t^4-27$ has
$3$-adic valuation $3/4$, so its field degree is divisible by four;
the polynomial is therefore irreducible over $\Q_3$, and hence also
over $\Q$.  Its normalization is $O[\pi]$, where
$\pi=3/t$, $\pi^4=3$, and $t=\pi^3$.  The ring $O[\pi]$ is a DVR
with uniformizer $\pi$.
\end{example}


\section{Shallow next-to-leading coefficients}
\label{sec:shallow-top-coefficient}

\begin{corollary}\label{cor:shallow-top-split}
Let \(p\) be any prime, \(n\ge2\), and \(O=\Z_p\) or
\(\Z_{(p)}\).  If
\[
 F=t^n+pb\,t^{n-1}+\sum_{i=0}^{n-2}a_it^i,
 \qquad b\in O,\quad a_i\in p^2O,
\]
the generic kernel for \(O[t]/F\) contains a countable direct sum
of exact order-\(n\) classes which splits from the whole
\(K_0(\Lat_{O[t]/F}(S_p))\).  The retraction annihilates the
Young span.  The analogous cyclic-group statement also holds.
\end{corollary}
\begin{proof}
If \(b\in pO\), then \(r=\infty\); otherwise \(r=n-1\).
In either case \(N=n\), because \(p(n-1)\ge n\).
Apply Theorem~\ref{thm:singular-distinguished-monic}.
\end{proof}

For odd primes the following direct calculation describes another
coordinate system and an explicit cubic companion lattice.
Its particular preimages are not identified with the mesh section
in the preceding corollary.

\begin{theorem}\label{thm:shallow-top-coefficient}
Let $p$ be odd, $n\ge2$, $O=\Z_p$ or $\Z_{(p)}$, and
\[
 F(t)=t^n+pb\,t^{n-1}+\sum_{i=0}^{n-2}a_it^i,
 \qquad b\in O,\quad a_i\in p^2O.
\]
Put $A=O[t]/F(t)$.  There is a homomorphism
\[
 K_0(\Lat_A(S_p))\longrightarrow\bigoplus_{j\ge1}\Z/n
\]
which annihilates the degree-$p$ Young permutation span and whose
restriction to the kernel of generic specialization is surjective.
The analogous assertion holds for $C_p$.
In particular, $\Phi_A$ is not surjective in degree $p$.
No annihilation relation or splitting of the preimages is asserted.
\end{theorem}

The proof changes the coefficient test modules, so that every test
map and its suspension still lift.  It does not apply the binomial
Bockstein to unavailable truncated-polynomial lifts.

\begin{proof}
If $n=2$, replacing $t$ by $u-pb/2$ identifies $A$ with
$O[u]/(u^2-d)$, where $d=p^2b^2/4-a_0\in p^2O$.
Theorem~\ref{thm:binomial-depth-two-all-exponents} proves the result.
Hence assume $n\ge3$, initially over $O=\Z_p$.
Write $k=\F_p$, $W=O/p^2O$, $R=k[t]/t^n$, and
$\widetilde R=A/p^2A$.

\paragraph{A palindromic chain of coefficient lifts.}
Choose monic linear polynomials $\ell_1,\ldots,\ell_n\in W[t]$.
For $n=2m+1$, let all of them be $t$, except
$\ell_{m+1}=t+pb$.
For $n=2m$, let all of them be $t$, except
$\ell_m=\ell_{m+1}=t+pb/2$.
Their sequence is palindromic, and
\[
 \prod_{r=1}^n\ell_r=t^n+pb\,t^{n-1}=F\pmod{p^2}.
\]
For $0\le i\le n$ put $f_i=\prod_{r=1}^i\ell_r$, with
$f_0=1$.  Then
\begin{equation}\label{eq:shallow-palindrome}
 f_i\mid f_j\quad(i\le j),\qquad
 f_i f_{n-i}=f_n,
 \qquad \widetilde R=W[t]/f_n.
\end{equation}
Let $X_i=R/(t^i)$ and
$\widetilde X_i=W[t]/f_i$, for $1\le i\le n$.
These are $W$-free lifts of the $X_i$.

Every $R$-map $X_i\to X_j$ lifts to a
$\widetilde R$-map.  If $i\le j$, its value at $1$ has the
form $t^{j-i}h(t)$; lift it to
$(f_j/f_i)\widetilde h(t)\bmod f_j$, lifting each coefficient
of $h$ arbitrarily to $W$.
If $i\ge j$, lift the arbitrary image $h(t)$ of $1$ to
$\widetilde h(t)\bmod f_j$; the divisibility $f_j\mid f_i$
makes this well-defined.
Matrices of these maps cover arbitrary finite direct sums.
An isomorphism lifts to an isomorphism, since its reduction
is invertible between $W$-free modules of the same rank.

The coefficient syzygy also lifts exactly:
\begin{equation}\label{eq:shallow-lifted-syzygy}
 0\longrightarrow\widetilde X_i
 \xrightarrow{\ f_{n-i}\ }\widetilde R
 \longrightarrow\widetilde X_{n-i}\longrightarrow0.
\end{equation}
For injectivity, if $f_{n-i}a\in(f_i f_{n-i})$ in $W[t]$,
cancel the monic polynomial $f_{n-i}$ to get $a\in(f_i)$.
The cokernel is visibly $W[t]/f_{n-i}$.
The sequence is split over $W$ because the quotient is
$W$-free.  Its reduction is
$0\to X_i\xrightarrow{t^{n-i}}R\to X_{n-i}\to0$.

\paragraph{The full-category Bockstein and its suspension.}
For an $A$-free cyclic lattice $E$, use curvature
$U_E=(\widetilde T^p-I)/F\bmod F$.
The stable naturality, integral primary extraction, and
cohomological triviality of every component $E_f$, $f\ne z$,
in Lemma~\ref{lem:binomial-all-primary} remain valid.
Indeed, their proofs use trace telescoping and cancellation
of the monic $F$, with $F\bmod p=t^n$; no congruence modulo
$p^2$ enters those arguments.
For a selected component $E$, its exact residual Tate
complex has free $R$-terms.  Thus
Lemma~\ref{lem:binomial-all-representability} supplies
\[
 M_E=\operatorname{coker}(T-I)^*,\qquad
 \underline{\operatorname{Hom}}_R(M_E,X)
 \cong\widehat H^0(C_p,(E/pE)\otimes_R X),
\]
and the contravariant short exact sequence of transpose
cokernels associated to every admissible sequence.

For a direct sum of the chosen test lifts, use
\[
 0\longrightarrow(E/pE)\otimes_R X
 \xrightarrow{p}(E/p^2E)\otimes_{\widetilde R}\widetilde X
 \longrightarrow(E/pE)\otimes_R X\longrightarrow0.
\]
It is exact because $E/p^2E$ is $\widetilde R$-free and
$\widetilde X$ is $W$-free.  Its Tate boundary is $\beta_X$.
The preceding lift of every map gives naturality on the
entire finite $R$-module category.  No functorial choice of
map lifts is needed: every such lift gives the naturality
square with the same prescribed maps on its two reductions.
A change of cyclic decomposition lifts to an isomorphism,
so the resulting operator is independent of that presentation.

Use the cyclic homotopy
$h_{\rm even}=I$, $h_{\rm odd}=-Q_p(T)$, with
$Q_p(z)=(N(z)-p)/(z-1)$.
The identity $dh+hd=pI$ gives $\beta h+h\beta=I$.
Since $p$ is odd, $h^2=-Q_p(1)=0$ on Tate cohomology, and
$P=\beta h$ is idempotent.

We check the boundary signs with the new lifts explicitly.
For an admissible lattice sequence, or for
\eqref{eq:shallow-lifted-syzygy} after tensoring with $E/p^2E$,
choose a degreewise $W$-linear splitting of the complexes.
Write the off-diagonal differential block as $Z$.
Then $d'Z+Zd''=0$ over $W$.
For a lifted cycle with $d''z=pw$, we have
$d'(Zz)=-pZw$; dividing by $p$ shows
$\beta'\partial+\partial\beta''=0$.
Writing the off-diagonal block of $h$ as $Y$, the off-diagonal
part of $dh+hd=pI$ is
\[
 d'Y+Yd''+h'Z+Zh''=0.
\]
It gives $h'\partial+\partial h''=0$ on cohomology.
Consequently $P'\partial=\partial P''$ in both kinds of
sequence.  The middle coefficient test $R$ has zero Tate
cohomology, so the reduced syzygy boundaries identify the
two parities with stable suspension.

Yoneda therefore gives a stable idempotent $P_E$ of $M_E$
which commutes with all three arrows of the transpose
triangle, including its suspension arrow.
Let $M_E^+$ be its image in the idempotent-complete stable
category of $R$.  The stable-length and direct-summand
arguments of Section~\ref{sec:binomial-depth-two-all} now
prove that
\[
 \eta_f(L)=-\dim_k M_{L_f}^+\pmod n\qquad(f\ne z)
\]
is additive on every admissible sequence.
These arguments use the Frobenius stable-category theorem
and the direct-summand theorem for triangles in exactly the
scopes already established there
\citep[Section~13.4, Theorem~13.7, author p.~50]{Buhler2010}
\citep[Lemma~1.6 and proof, pp.~821--822]{BalmerSchlichting2001}.
Choice independence holds on the full category: maps
factoring through $A[C_p]$ act as zero on every test functor,
and the natural $\beta,h$ intertwine the stable isomorphisms
between choices of curvature-primary components.
Each class has finite primary support.  Constant $O$-lattices
have zero curvature and hence all coordinates zero.

\paragraph{Witnesses and their actual selected stable module.}
Write $F=a_0+tH(t)$, and work over a finite unramified
extension $O_s/O$.  If $a_0=p^v u\ne0$, take
\[
 \alpha=1+p^{v-1}y,\qquad c=(\alpha^p-1)/a_0,
\]
where $y$ is a unit.  Then $c$ is a unit and
$\overline c=\overline y/\overline u$ can be prescribed
arbitrarily in $k_s^\times$.
If $a_0=0$, take $\alpha=1$ and any unit $c$.
In both cases $\alpha^p-ca_0=1$.
Let $C_\alpha$ be the companion of
$q_\alpha(z)=(z^p-\alpha^p)/(z-\alpha)$,
with last coordinate row $a$ and first coordinate column $b_0$.
The matrix
\[
 T=\begin{pmatrix}\alpha&ct\,a\\H(t)b_0&C_\alpha\end{pmatrix}
\]
has characteristic polynomial $z^p-1-cF(t)$ before imposing
$F=0$, so $\widetilde T^p-I=cF(t)I$.
It defines a lattice $L$ of scalar curvature $c$.

Here $H\bmod p=t^{n-1}$.
Present the residual transpose cokernel by
$x,y_1,\ldots,y_{p-1}$.
If $m=p-1$ and
$\overline q_\alpha=z^m+q_{m-1}z^{m-1}+\cdots+q_0=(z-1)^m$,
the top row gives $\overline c\,t y_m=0$; the subsequent
companion rows give
$y_{r-1}-y_r-q_{r-1}y_m=0$ for $2\le r\le m$.
Eliminate $y_1,\ldots,y_{m-1}$.
The remaining relation is
\[
 t^{n-1}x-(1+q_{m-1}+\cdots+q_0)y_m
 =t^{n-1}x-\overline q_\alpha(1)y_m=t^{n-1}x=0.
\]
Thus, over $R_s=k_s[t]/t^n$,
\begin{equation}\label{eq:shallow-witness-transpose}
 M_L\cong R_s/(t^{n-1})\,x\ \oplus\ R_s/(t)\,y_m.
\end{equation}

Since $n\ge3$, the palindromic chain has $f_1=t$.
Consequently the chosen lift of the residue test $k_s$ has
$t=0$ even modulo $p^2$.
The upper off-diagonal block vanishes there, as does the
lower one modulo $p^2$, because $H(0)=a_1\in p^2O$.
On the trivial even Tate class, $h=I$ and
$\beta_{\rm odd}=N(\alpha)/p\equiv1\pmod p$.
The last congruence follows by expanding $N(1+w)$ with
$w\in pO_s$ and dividing by $p$.
The even cyclotomic class is represented by
$(C_\alpha-I)^{p-2}b_0$ modulo $p$; since $p>2$, its image
under $h_{\rm even}=I$ belongs to
$\operatorname{im}(C_\alpha-I)$ and is zero in odd cohomology.
Thus $P$ has values $1,0$ on the two residue-test classes.
The idempotent-lifting and Krull--Schmidt argument applied
to \eqref{eq:shallow-witness-transpose} identifies its image
as $R_s/(t^{n-1})$ up to stable isomorphism.
Therefore the negative-length coordinate of $L$ is
$-s(n-1)=s\pmod n$.

The witness is generically regular without a reducedness
assumption on the generic algebra.  Put $B=A_s[1/p]$.
The $B$-free representation is projective over $B[C_p]$
by averaging.  Over every residue field of $B$, its
characteristic polynomial $z^p-1$ is separable and gives
the regular representation.
The commutative Artinian algebra $B[C_p]$ is finite etale
over $B$.  The projective module has rank one on each
of its local factors, so is free of rank one over $B[C_p]$.
Thus $[L]-[A_s[C_p]]$ vanishes generically.

Choose distinct prime degrees $s_j>n(p-1)$, and residue
parameters of those degrees.  The scalar orbits of these
parameters under $\F_p^\times$ are free: equality between
$a\overline c$ and a Frobenius conjugate of $\overline c$
would give $a^{s_j}=1$ after taking the norm, hence $a=1$.
Different degrees give disjoint orbits.
Restrict coefficients to $A$ and induce to $S_p$.
Changing the cyclic generator to its $a$th power multiplies
the stable curvature by $a$; Mackey restriction supplies
each normalizer twist once, and all remaining terms are
regular cyclic lattices.
Accordingly, the $f_j$-coordinates of the generic-zero
classes $[\operatorname{Ind}L_j]-s_j[A[S_p]]$ are diagonal,
with value $s_j\pmod n$ on the diagonal.
Normalize by $s_j^{-1}$.  Finite primary support gives the
claimed map to the direct sum, and these witnesses prove
surjectivity from its generic kernel.
Every Young lattice and its cyclic restriction is constant,
so the Young span is annihilated.

For $O=\Z_{(p)}$, compose the completed coordinates with
completion.  Lift the separable residue parameter polynomial
to obtain a finite free local etale $O$-algebra $O_s$.
The finite formulas for $\alpha,c$ are integral over $O_s$
and define the same companions there.
The preceding Artinian generic argument proves their generic
vanishing before completion, and their completed coordinate
values are unchanged.  They therefore give the desired
localized preimages.  No mesh relation is used.
\end{proof}

\begin{example}\label{ex:shallow-wild-cubic}
At $p=3$, put $A=\Z_3[t]/(t^3+3t^2+9)$.
Its Newton polygon has a single segment of slope $-2/3$.
The coprime-denominator argument proves irreducibility and
total ramification of degree three
\citep[Corollary~3.18, p.~49]{WilliamsWuthrich2025}.
Here is also an explicit normalization calculation.
Set $\pi=t^2/3$.  The defining equation gives
\[
 \pi(t+3)=-3,\qquad
 \pi^3=3(\pi+1)^2,
 \qquad t=-\pi^2+3\pi+3.
\]
The polynomial $\pi^3-3\pi^2-6\pi-3$ is Eisenstein;
thus the normalization is $S=\Z_3[\pi]$
\citep[Theorem~5.9, pp.~67--68]{WilliamsWuthrich2025}.
Since $t^2=3\pi$, the two bases satisfy
\[
 A=\Z_3\oplus3\Z_3\pi\oplus\Z_3\pi^2\subset S.
\]
The index is three, so this is a proper local order in a
totally ramified cubic field.

Choose $\alpha=4$, $c=7$, so $\alpha^3-9c=1$.
One explicit witness is
\[
 T=\begin{pmatrix}
 4&0&7t\\t^2+3t&0&-16\\0&1&-4
 \end{pmatrix},\qquad T^3=I.
\]
Its curvature residue is $1$, with free scalar orbit under
$\F_3^\times$.
The induced class $[\operatorname{Ind}_{C_3}^{S_3}L]-[A[S_3]]$
vanishes generically and has coordinate $1\pmod3$.
This proves it is nonzero; its exact order is not asserted.
The countable family is obtained with the unramified
parameters in the proof.
\end{example}

\section{Monogenic orders with nontrivial generic residue degree}
\label{sec:monogenic-residue-degree}

The curvature construction also applies to a distinguished defining
polynomial whose generic quotient is a field.  The resulting modulus
is the residue degree of that field: this is exactly the ambiguity in
the valuation of a determinant under a change of generic basis.

\begin{theorem}\label{thm:monogenic-residue-degree}
Let $O=\Z_p$ or $\Z_{(p)}$, and let $F\in O[t]$ be monic of degree
$N\ge2$, with
$F\bmod p=t^N$.  Suppose that
\[
 E=\Q_p[t]/(F)
\]
is a field of residue degree $f>1$ over $\Q_p$.  Put
$A=O[t]/(F)$.  There is a homomorphism
\[
 \Theta:K_0(\Lat_A(S_p))\longrightarrow
             \bigoplus_{i\ge1}\Z/f
\]
which vanishes on the Young span and whose restriction to the kernel
of specialization to $\operatorname{Frac}(O)\otimes_O A$ is
surjective.  In particular $\Phi_A$
fails to be surjective in degree $p$.

The proof supplies explicit generic-zero preimages of the coordinate
vectors.  It asserts neither a torsion order for those preimages nor
a splitting of this quotient.
\end{theorem}

\begin{corollary}\label{cor:scaled-unramified-orders}
Let $E/\Q_p$ be unramified of degree $N\ge2$, choose a unit
$\theta$ with $O_E=\Z_p[\theta]$, and let $a\ge1$.  Every proper
local order $A=\Z_p[p^a\theta]$ satisfies the theorem with modulus
$f=N$.  This includes every degree $N\ge2$ and every prime.
\end{corollary}

\subsection{Curvature for a distinguished defining polynomial}

We first work over $O=\Z_p$; descent of the explicit witnesses to
$\Z_{(p)}$ is proved at the end.  The ring $A$ is finite free and
complete over $O$, and is local with
maximal ideal $(p,t)$.  Formal division by the distinguished polynomial
identifies $A$ with $R/(F)$ for $R=O[[t]]$.  Since its generic quotient
is a field, it is a domain and $F(0)\ne0$.

Let $M$ be an $A$-free $C_p$-lattice, with generator matrix $T$.
Choose a lift $\widetilde T$ over $R$ and put
\[
 \widetilde U=\frac{\widetilde T^p-I}{F},\qquad
 U=\widetilde U\bmod F.
\]
The numerator is divisible by $F$, and $U$ commutes with $T$.
The stable-curvature identity of
Lemma~\ref{lem:truncation-stable-curvature} remains valid with $F$
in place of $t^n$.  Indeed, if an equivariant map $H:M\to M'$ is
lifted and
$\widetilde T'\widetilde H-\widetilde H\widetilde T=F\widetilde B$,
the same telescoping calculation gives
\begin{equation}\label{eq:monogenic-stable-curvature}
 U'H-HU=\sum_{j=0}^{p-1}(T')^{p-1-j}BT^j.
\end{equation}
This is the group trace appearing in that lemma, with the same
factorization through an induced lattice.  Consequently the stable
endomorphism $U$ is natural and independent of lifts and bases.

Use the integral primary decomposition of the finite free $O$-module
$M$ for $U$, indexed by monic irreducibles $q\in\F_p[u]$.
The projectors commute with $A$ and $T$, so their images $M_q$ are
$A$-free direct summands.  The complete-local primary construction and
its exactness from Section~\ref{sec:all-truncations} apply without
change: a common annihilating polynomial gives compatible projectors
throughout a short exact sequence.

We verify the altered acyclicity step explicitly.  For $q\ne u$,
$U$ is a unit on $M_q$.  Choose an $A$-basis adapted to this summand
and lift it to an invertible matrix over $R$.  The off-diagonal blocks
of the lifted action are divisible by $F$.  Every term leaving a
block and returning in its $p$th power contains two such blocks.
Thus the diagonal action lift $\widetilde T_q$ satisfies
\[
 \frac{\widetilde T_q^p-I}{F}\bmod F=U|_{M_q}.
\]
Its lifted curvature $\widetilde U_q$ is a unit over $R$, and
\begin{equation}\label{eq:monogenic-factorization}
 (\widetilde T_q-I)N(\widetilde T_q)=F\widetilde U_q.
\end{equation}
If $(\widetilde T_q-I)x=Fy$, multiplication by the norm, cancellation
of the nonzerodivisor $F$, and inversion of $\widetilde U_q$ give
$x=N(\widetilde T_q)\widetilde U_q^{-1}y$ modulo $F$.
Starting with $N(\widetilde T_q)x=Fy$ gives the other Tate parity.
Hence
\[
 \widehat H^\bullet(C_p,M_q)=0\qquad(q\ne u).
\]
This proves the general-$F$ replacement for
Lemma~\ref{lem:truncation-primary-acyclic}.

\subsection{The determinant ambiguity and the additive coordinates}

For a cohomologically trivial $A$-free cyclic lattice $M$, its invariant
module $M^{C_p}$ is finite free over $O$, and its generic fibre is
an $E$-vector space, say $E^r$.  Choose a generic $E$-basis.  Using
$1,t,\ldots,t^{N-1}$ in each copy of $E$, compare an $O$-basis of
$M^{C_p}$ with the resulting $\Q_p$-basis and define
\[
 \iota_f(M)=v_p(\det M^{C_p})\pmod f.
\]
Integral basis changes have unit determinant.  If $B\in\mathrm{GL}_r(E)$
changes the generic basis, its determinant over $\Q_p$ is
$\operatorname{Norm}_{E/\Q_p}(\det_E B)$, and
\begin{equation}\label{eq:monogenic-norm-valuation}
 v_p\bigl(\operatorname{Norm}_{E/\Q_p}(x)\bigr)=f\,v_E(x).
\end{equation}
Here $v_E$ is normalized by an $E$-uniformizer.  To see the formula,
units preserve the valuation ring, while multiplication by a
uniformizer has cokernel the residue field, of $O$-length $f$.
Determinants and multiplicativity give the formula for every nonzero
$x$.  Thus $\iota_f$ is independent of the generic basis modulo $f$.

This index is additive on admissible sequences of cohomologically
trivial lattices.  Taking invariants is exact because the first term
has vanishing $H^1$; compatible integral and generic bases make the
determinant block triangular.  It is zero on induced projectives:
their invariants are finite projective, hence free, over $A$.
The stabilization argument of
Lemma~\ref{lem:truncation-stable-index} therefore proves invariance
under relative stable isomorphism.  That argument uses only the
locality of $A$, the actual stabilization isomorphism, and the
vanishing of the index on projectives, all of which hold here.

For $q\ne u$ set
\begin{equation}\label{eq:monogenic-primary-index}
 \eta_q(M)=\iota_f(M_q)\in\Z/f.
\end{equation}
We recall the descent to the Grothendieck group to specify exactly
what is reused.  Two choices of curvature agree stably by
\eqref{eq:monogenic-stable-curvature}.  Over the complete local ring
$A$, a common product of their characteristic polynomials gives one
Hensel--Bezout primary projector for both endomorphisms.  The two
projectors agree stably, so their images are stably isomorphic.
The preceding stable invariance proves that \eqref{eq:monogenic-primary-index}
does not depend on choices.  For an admissible sequence, choose an
$A$-splitting and upper triangular action lifts.  Its curvature is
upper triangular, and exact primary extraction gives an admissible
sequence of the $q$-components.  Their Tate acyclicity and index
additivity prove that $\eta_q$ is a homomorphism.

Constants have constant action lifts with zero curvature, so every
$\eta_q$ vanishes on constant $O$-lattices.  Each class has finite
primary support.  These observations include the stable-primary,
exactness, and finite-support assertions needed below.

\subsection{A parameter and its companion lattice}

Let $\tau$ denote the image of $t$ in $E$.  It is integral, and its
reduction is zero because $F\bmod p=t^N$.  Thus $v_E(\tau)\ge1$.
Since its norm is $F(0)$ up to sign,
\[
 v_p(F(0))=f\,v_E(\tau)\ge f\ge2.
\]
The valuation needed below is therefore automatic except for the
additional parameter calculation when $p=2$ and $v_p(F(0))=2$.
In that case $f=2$.

Write $F(0)=p^v b_0$, with $b_0$ a unit.  First suppose that $p$ is
odd, or that $p=2$ and $v\ge3$.  Let $O_s/O$ be the
unramified extension of degree $s$, and prescribe a nonzero residue
$\eta\in\F_{p^s}$.  Choose a unit $y\in O_s$ reducing to
$\overline b_0\eta$, and put
\begin{equation}\label{eq:monogenic-parameter}
 \alpha=1+p^{v-1}y,\qquad c=\frac{\alpha^p-1}{F(0)}.
\end{equation}
The binomial expansion gives
\[
 \alpha^p-1=p^v y\,U_0,\qquad U_0\in1+pO_s.
\]
For odd $p$ this uses $v\ge2$; for $p=2$ the factor is exactly
$U_0=1+2^{v-2}y$, and $v\ge3$ suffices.  Hence $c$ is a unit
with residue $\eta$, and
\[
 \alpha^p=1+cF(0),\qquad v_p(\alpha-1)=v-1.
\]

For the remaining case $p=2$, $F(0)=4b_0$, choose an odd prime
$s\ge3$ and a nonzero trace-zero element $\eta\in\F_{2^s}$.
Its degree is $s$, because the only proper subfield is $\F_2$,
whose only trace-zero element is zero for odd $s$.
The map $y\mapsto y^2+y$ has kernel $\F_2$ and image the
trace-zero hyperplane.  Choose $\overline y$ with
$\overline y^2+\overline y=\eta$ and lift it to $O_s$.
Both $y$ and $1+y$ are units.  Set
\[
 \alpha=1+2y,\qquad c=\frac{y(1+y)}{b_0}.
\]
Then $\alpha^2=1+cF(0)$, $\overline c=\eta$, and both
$\alpha-1=2y$ and $q_\alpha(1)=\alpha+1=2(1+y)$ have
valuation one.  Here $q_\alpha$ is the polynomial defined next.
Thus the same companion and ideal calculations apply.

Put
\[
 G(t)=\frac{F(t)-F(0)}t,\qquad
 q_\alpha(z)=\frac{z^p-\alpha^p}{z-\alpha}.
\]
Let $C_\alpha$ be the companion of $q_\alpha$, let $a$ be its last
coordinate row, and let $b$ be its first coordinate column.  Over
$A_s=A\otimes_O O_s$, take the rank-$p$ lattice with generator
\begin{equation}\label{eq:monogenic-witness}
 T=\begin{pmatrix}
       \alpha&ct\,a\\
       G(t)b&C_\alpha
    \end{pmatrix}.
\end{equation}
The companion cofactor identity gives, over $O_s[[t]]$,
\[
 \det(zI-\widetilde T)
   =(z-\alpha)q_\alpha(z)-ctG(t)
   =z^p-1-cF(t).
\]
Cayley--Hamilton therefore gives the exact identity
\begin{equation}\label{eq:monogenic-exact-curvature}
 \widetilde T^p-I=cF(t)I_p.
\end{equation}
Thus \eqref{eq:monogenic-witness} is a cyclic group action, and its
curvature is scalar multiplication by $c$.  Restriction of coefficients
gives an $A$-free lattice of rank $ps$.

\subsection{The ideal calculation}

We compute its coinvariants; the norm identifies coinvariants with
invariants because the selected lattice is cohomologically trivial.
This is an $A$-linear isomorphism, so the covolume modulo $f$ agrees
with the invariant index above.

Let $a_0=\alpha-1$ and
\[
 \lambda=\frac{q_\alpha(1)}c=\frac{F(0)}{a_0}.
\]
Its valuation is one, so $\lambda$ is $p$ times a unit.  The companion
columns identify all cyclotomic coinvariant coordinates with one
generator $w$.  With $v_0$ the remaining coordinate, the presentation
is
\begin{equation}\label{eq:monogenic-coinvariant-presentation}
 tv_0=\lambda w,\qquad a_0v_0+G(t)w=0.
\end{equation}
The map $v_0\mapsto\lambda$, $w\mapsto t$ identifies this module
with $(p,t)A_s$.  To check completeness of the relations, suppose in
$O_s[[t]]$ that $\lambda a+tb=Fc_0$.  Since
$F=a_0\lambda+tG$, this gives
\[
 \lambda(a-a_0c_0)=t(Gc_0-b).
\]
Setting $t=0$ shows $a-a_0c_0=tq_0$, and then
$b=Gc_0-\lambda q_0$.  Every relation is therefore generated by
$(t,-\lambda)$ and $(a_0,G)$, as in
\eqref{eq:monogenic-coinvariant-presentation}.

Relative to an integral basis of $O_s/O$, the module $A_s$ is free
of rank $s$ over $A$, and
\[
 A_s/(p,t)A_s=\F_{p^s}
\]
has $O$-length $s$.  If $\eta$ has degree $s$ over $\F_p$ and minimal
polynomial $q$, the whole witness lies in its $q$-primary component.
Consequently
\begin{equation}\label{eq:monogenic-witness-index}
 \eta_q(L)=s\pmod f,
\end{equation}
and its other primary coordinates vanish.  This computation does not
require $E\otimes_{\Q_p}\operatorname{Frac}(O_s)$ to be a field.

\subsection{The countable quotient and the generic kernel}

Outside the valuation-two case at $p=2$, choose distinct positive
degrees $s_i$ coprime to $f(p-1)$, and monic
irreducibles $q_i\in\F_p[u]$ of those degrees with nonzero constant
terms.  Such irreducibles exist in every degree.  Choose a root
$\eta_i$ of each.  In the valuation-two case, use trace-zero
parameters of distinct odd prime degrees from the preceding
construction; these degrees are prime to $f=2$.
Form the corresponding coefficient-restricted lattices $L_i$ and put
\[
 x_i=[\operatorname{Ind}_{C_p}^{S_p}L_i]-s_i[A[S_p]].
\]
Under $g\mapsto g^j$, the lift $\widetilde T^j$ has residual
curvature $j\eta_i$.  If this scalar twist preserves $q_i$, comparison
of constant terms gives $j^{s_i}=1$, so $j=1$ by the degree condition.
Different degrees have disjoint scalar orbits.

Mackey restriction contributes each of the $p-1$ normalizer twists
once.  Every other intersection is trivial, and gives a relatively
projective cyclic lattice.  Therefore
\[
 (\eta_{q_j}\circ\operatorname{Res}_{C_p}^{S_p})(x_i)
       =s_i\delta_{ij}\pmod f.
\]
Normalize the $i$th coordinate by $s_i^{-1}$ modulo $f$.  Finite
primary support gives the homomorphism $\Theta$ to the direct sum,
and these equations show surjectivity.  Every Young lattice is
constant over $O$, so $\Theta$ kills the Young span.

To verify the asserted generic vanishing, work over
$E_s=E\otimes_{\Q_p}\operatorname{Frac}(O_s)$.  The elements $t$
and $c$ are units.  Starting with the first cyclotomic coordinate in
\eqref{eq:monogenic-witness}, its successive images first span all
the cyclotomic coordinates, and the next image has remaining
coordinate $ct$, a unit.  They form a cyclic basis.  Since the
characteristic polynomial is $z^p-1$, the generic lattice is
$E_s[C_p]$.  Restriction of coefficients and induction show that
$x_i$ is generically zero.  Thus $\Theta$ is surjective already on
the generic kernel in the complete case.

\subsection{Descent to the localized integers}

Now let $O=\Z_{(p)}$.  The hypothesis that $\Q_p[t]/F$ is a field
implies that $F$ is irreducible over $\Q$.  The ring $A$ is finite
free and local over $O$, with maximal ideal $(p,t)$, and its
completion is $\widehat A=\Z_p[t]/F$.  Compose the complete
coordinate homomorphisms with coefficient extension to $\widehat A$.
Admissible sequences remain exact because they split over $A$.
This gives finite-support homomorphisms on the original lattice
group, vanishing on its Young span.

Every required witness is already defined over $O$.  Lift a monic
irreducible polynomial for $\F_{p^s}/\F_p$ to $O[z]$ and take its
quotient $O_s$.  This is a finite etale local $O$-algebra: its
residue algebra is the separable field $\F_{p^s}$, and its maximal
ideal is $pO_s$.  It is therefore a DVR, free of rank $s$ over $O$,
and its completion is the unramified extension used above.
Lift each chosen residue $y$ to this ring.  Formulas
\eqref{eq:monogenic-parameter}, or the valuation-two formula
$c=y(1+y)/b_0$, are finite identities in $O_s$; the valuation
calculations prove that their quotients are units there.
No Hensel root in $\Z_{(p)}$ is needed.

The resulting lattices complete to witnesses with the computed
coordinate vectors, so the pulled-back map is surjective.
Their generic cancellation holds over $\Q$ itself.  Indeed
$A\otimes_O\Q$ is a number field, and after adjoining
$\operatorname{Frac}(O_s)$ it is a product of characteristic-zero
fields.  In every factor $t$ and $c$ are units, and the explicit
cyclic basis above identifies the witness with the regular
representation.  Restriction and induction give generic-zero
preimages over the original ring.  This proves the theorem over
$\Z_{(p)}$ without assuming that arbitrary completed lattices descend.

For Corollary~\ref{cor:scaled-unramified-orders}, let $H$ be the
minimal polynomial of $\theta$.  It has degree $N$ and unit constant
term.  The minimal polynomial of $p^a\theta$ is
\[
 F(t)=p^{aN}H(t/p^a).
\]
It is monic integral, reduces to $t^N$, and has constant-term
valuation $aN$.  The field has residue degree $N$, so all hypotheses
hold.  The order is proper because its residue
field is $\F_p$, whereas that of the maximal order is $\F_{p^N}$.

\begin{example}\label{ex:unramified-cubic-order}
The polynomial $H(z)=z^3-z-1$ has no root over $\F_3$, so it defines
an unramified cubic extension of $\Q_3$.  If $\theta$ is a root, the
order
\[
 A=\Z_3[3\theta]
   \simeq\Z_3[t]/(t^3-9t-27)
\]
has a degree-three generic kernel mapping onto
$\bigoplus_{i\ge1}\Z/3$.  One explicit cyclic witness is obtained
from $\alpha=10$ and $c=-37$:
\[
 T=\begin{pmatrix}
       10&0&-37t\\
       t^2-9&0&-100\\
       0&1&-10
    \end{pmatrix},\qquad T^3=I.
\]
Its induced class minus the regular class has coordinate one modulo
three for the primary label $u-2$.  The theorem supplies independent
coordinates through the unramified parameter extensions; it does not
assign an exact order to this displayed class.
\end{example}


\section{The split Witt--Green ring}

This section answers the separate question on scan 6 using the split
ring exactly as defined there.  Let $k=\overline\F_3$ and $G=C_3$.
Write $J_i$ for the unipotent indecomposable $k[G]$-module of
dimension $i$, $1\leq i\leq3$.  The split Green group is free
abelian on $1=[J_1]$, $y=[J_2]$, $x=[J_3]$, by the Jordan
classification.  Direct tensor and exterior-power computations give
\begin{equation}\label{eq:jordan}
 x^2=3x,\quad xy=2x,\quad y^2=x+1,\qquad
 \lambda^2x=x,\quad\lambda^3x=1,\quad
 \lambda^2y=1,\quad\lambda^3y=0.
\end{equation}
Here $J_3=k[G]$, and the equivariant isomorphism
$k[G]\otimes_k M\simeq k[G]^{\dim M}$ follows by sending
$g\otimes m$ to $g\otimes g^{-1}m$.  For $J_2\otimes J_2$,
the nilpotent operator $g-1$ has ranks two and one in its first
two powers and has cube zero, giving blocks of sizes three and
one.  On $\bigwedge^2J_3$, a cyclic exterior basis is permuted
transitively, and both top exterior powers have determinant one.
This proves every identity in~\eqref{eq:jordan}.

Let $X$ be $J_3$ with the nondegenerate invariant dot product on
its permutation basis, and let $Z=H(J_2)$ be the hyperbolic
symmetric module.  Forgetting forms sends $X$ to $x$ and $Z$ to
$2y$, and commutes with exterior powers.  Newton's third operation
on a pre-$\lambda$-ring is
\[
 \psi^3(a)=a^3-3a\lambda^2(a)+3\lambda^3(a).
\]
It is additive, by the logarithmic-derivative identity for
$\lambda_t$.  A special $\lambda$-ring makes it multiplicative.
The displayed relations give
\[
 \psi^3(x)=3,\qquad \psi^3(y)=2x-2y.
\]
Consequently, after forgetting forms,
\[
 \psi^3(XZ)\longmapsto12,\qquad
 \psi^3(X)\psi^3(Z)\longmapsto12x-12y.
\]
Their difference is nonzero in the free abelian group on $1,y,x$.
These integer coefficients are not reduced modulo three.
The split Witt--Green ring is therefore not special, although two
is invertible in $k$.

The iterate identity also fails.  For $a=x+2y$, realized by
$X\perp H(J_2)$, multiplication of exterior series gives
\[
 \lambda^2a=6x+3,\quad
 \lambda^3a=\lambda^4a=10x+2y+1,\quad
 \lambda^5a=6x+3,\quad\lambda^6a=x+2y.
\]
The universal symmetric polynomial for $\lambda^3(\lambda^2a)$ is
\[
 Q_{3,2}=e_1^2e_4-2e_2e_4-e_1e_5+e_6+e_3^2,
 \qquad e_i=\lambda^i(a).
\]
To verify it, expand $e_3[e_2]$ in six variables by choosing three
distinct edges of a complete graph.  The exponent partitions
$(3,1,1,1)$, $(2,2,1,1)$, $(2,2,2)$, $(2,1,1,1,1)$ and
$(1,1,1,1,1,1)$ have respective coefficients $1,2,1,6,15$;
the displayed polynomial has exactly these coefficients.
Specialization to six variables is injective in symmetric degree
six.  In the split Green ring the two evaluations are
\[
 \lambda^3(6x+3)=441x+7,\qquad Q_{3,2}(a)=445x-4y+3,
\]
whose difference is $-4(x-y-1)\ne0$.

The exact representation group sends $x$ to $3$ and $y$ to $2$,
so both defects vanish there.  Zibrowius's theorem applies after
the metabolic relations defining his Grothendieck--Witt ring
\citep[Definition~1.2, p.~1060]{Zibrowius2015}.  The distinction
between split and exact relations is precisely the distinction
needed to reconcile the two conclusions.

\section{The tensor-power formula and its virtual coefficients}

The formula on scan 28 is an identity of exact equivariant
Grothendieck groups.  An integral Cauchy filtration proves it on any
commutatively ringed topos, without a semisimplicity assumption.
Let $(X,\mathcal A)$ be such a topos, and let $K_0(X)$ be formed
from finite locally free $\mathcal A$-modules of bounded rank and
all short exact sequences between them.  Let $K_n(X)$ be the
analogous exact group with an $S_n$-action.  Induction makes
$\bigoplus_nK_n(X)$ a graded ring, on which $K_0(X)$ acts by
tensoring with modules with trivial action.

For a partition $\nu=(\nu_1,\ldots,\nu_l)$ of $n$, put
$\rho_\nu=\rho_{\nu_1}\cdots\rho_{\nu_l}$, and let $m_\nu$
be the monomial symmetric function, expressed as its integral
polynomial in elementary symmetric functions.  Evaluation at a
class of $K_0(X)$ uses its exterior powers.

\begin{theorem}\label{thm:tensor-formula}
For every $\xi\in K_0(X)$ one has
\[
 \tau^n(\xi)=\sum_{\nu\vdash n}m_\nu(\xi)\rho_\nu
 \quad\text{in }K_n(X).
\]
Here $\tau^n([E])=[E^{\otimes n}]$ with place permutations on
an effective class; the total series $\tau$ extends to virtual
classes by sending addition to the induction product.
\end{theorem}

\begin{proof}
We use the integral Schur functors $S_\lambda$ and the natural
Cauchy filtration
\[
 \operatorname{gr}\operatorname{Sym}^n(F\otimes G)
 \simeq\bigoplus_{\lambda\vdash n}S_\lambda(F)\otimes S_\lambda(G).
\]
Its existence over arbitrary commutative rings, and the freeness
and base-change compatibility of its factors, are the standard
basis and Cauchy theorems of
\citet[Theorem~II.2.16, p.~232; Theorem~III.1.4,
pp.~245--246]{ABW1982}.  Their notation has $L_{(a)}=\bigwedge^a$;
our convention is $S_\lambda=L_{\lambda'}$, where $\lambda'$
is the conjugate partition.  Its weight multiplicities are the
usual Kostka numbers $K_{\lambda\nu}$.
The construction consists of natural exterior diagonalizations
and symmetric multiplications.  The filtration and its free
quotients commute with scalar extension and changes of basis,
so they apply locally to free sheaves and glue on the topos.

Set $W=\Z^n$, with its diagonal torus and its permutation-matrix
$S_n$-action, and let
\[
 C_\lambda=(S_\lambda(W))_{(1,\ldots,1)}.
\]
This is a finite free integral $S_n$-lattice.  Weight extraction
is exact over $\Z$: a comodule for a split diagonal torus is a
direct sum of its weight modules, and all comodule maps preserve
that decomposition.  Thus extracting a weight from our filtration
preserves exactness and local freeness.

The $W$-weight $(1,\ldots,1)$ of
$\operatorname{Sym}^n(E\otimes W)$ is $E^{\otimes n}$, by the map
\[
 e_1\otimes\cdots\otimes e_n\longmapsto
 (e_1\otimes w_1)\cdots(e_n\otimes w_n).
\]
Permuting the $w_i$ corresponds to place permutations.  The first
weight extraction therefore gives the exact identity
\begin{equation}\label{eq:schur-tensor}
 [E^{\otimes n}]=\sum_{\lambda\vdash n}[S_\lambda(E)][C_\lambda].
\end{equation}
Apply this construction to $E=\Z^n$ with a second diagonal torus,
and take its weight $\nu$, padded by zeroes.  The tensor words of
that weight form the transitive Young coset set; the corresponding
weight space of $S_\lambda(E)$ has rank $K_{\lambda\nu}$.
Consequently, already over the integers,
\begin{equation}\label{eq:integral-Young}
 \rho_\nu=\sum_{\lambda\vdash n}K_{\lambda\nu}[C_\lambda].
\end{equation}
This equality follows from an exact filtration; equality of
rational characters is not being substituted for it.

For a third extraction, take the $W$-weight $\nu$ directly from
the Cauchy filtration.  Its left side is
$\bigotimes_j\operatorname{Sym}^{\nu_j}(E)$, whence
\[
 h_\nu(E)=\sum_{\lambda\vdash n}K_{\lambda\nu}[S_\lambda(E)].
\]
The identity $[\operatorname{Sym}^a(E)]=h_a(E)$ follows from
the exterior--symmetric Koszul complex, whose positive-degree
strands are exact locally on a free basis.  The Kostka matrix is
triangular with diagonal entries one in dominance order.  Comparing
the last display with the integral symmetric-function identity
$h_\nu=\sum_\lambda K_{\lambda\nu}s_\lambda$ and inverting this
matrix gives $[S_\lambda(E)]=s_\lambda(E)$ in $K_0(X)$.
Now $s_\lambda=\sum_\nu K_{\lambda\nu}m_\nu$, so
\eqref{eq:schur-tensor} and \eqref{eq:integral-Young} give
\[
 [E^{\otimes n}]
 =\sum_\nu m_\nu(E)\sum_\lambda K_{\lambda\nu}[C_\lambda]
 =\sum_\nu m_\nu(E)\rho_\nu.
\]

It remains to justify virtual classes.  For an exact sequence
$0\to E'\to E\to E''\to0$, the filtration of $E^{\otimes n}$
by the number of factors in $E'$ is $S_n$-stable and has factors
\[
 \operatorname{Ind}_{S_a\times S_{n-a}}^{S_n}
 \bigl((E')^{\otimes a}\otimes(E'')^{\otimes(n-a)}\bigr).
\]
Local splittings show that its submodules and quotients are
finite locally free.  Therefore $\tau(E)=\tau(E')\tau(E'')$
in the degree-completed induction ring.  Its constant coefficient
is one, so inversion of series extends $\tau$ to $K_0(X)$.
The symmetric-function identity
\[
 m_\nu(u+v)=\sum_{\alpha\cup\beta=\nu}m_\alpha(u)m_\beta(v)
\]
and $\rho_\alpha\rho_\beta=\rho_{\alpha\cup\beta}$ show that
the proposed right-hand total series has the same additivity
law.  These integral identities evaluate on $K_0(X)$, whose
special $\lambda$-ring structure is
\citeauthor{SGA6}'s theorem for commutatively ringed topoi
\citep[Expos\'e~VI, Theorem~3.3, pp.~392--393]{SGA6}.
The two series agree on effective classes and hence on all of
$K_0(X)$.
\end{proof}

For $n=2$, the formula is
\[
 [E^{\otimes2}]=\psi^2(E)\rho_2+
 [\bigwedge^2E]\rho_1^2,
 \qquad\psi^2(E)=[E]^2-2[\bigwedge^2E].
\]
The coefficients need not be effective.  For the standard
representation of $\mathrm{GL}_2$ in characteristic zero,
$\psi^2(E)=[\operatorname{Sym}^2E]-[\bigwedge^2E]$.
These are distinct irreducible representations, so the negative
coefficient precludes effectivity.  Thus the formula does not
give the direct-summand interpretation suggested in the margin.


\section{What remains open}

The exact answer for all commutative local rings is not obtained
here.  Positive characteristic is excluded, polynomial
independence is automatic in characteristic zero, and all
coherent regular local rings and local $\Q$-algebras are covered,
as are their filtered local colimits and extensions by uniquely
divisible nil ideals.  These include noncoherent rings.
The unresolved case is generation over singular
mixed-characteristic local rings.  The examples show that
reducedness, absence of integer torsion and the rank-one
determinant condition do not settle it.  Generic eigenbundles can
already obstruct generation over a domain.
The exact answer for every square-zero extension of $\Z_{(p)}$ is
now available: unique divisibility of the ideal is necessary and
sufficient.  Every truncated polynomial ring $O[e]/(e^n)$, $n\ge2$, also fails,
with a split countable $\Z/n$ subgroup in its degree-$p$ reduction
kernel.  More general finite nilpotent coefficient extensions remain
to be classified; the principal truncated family is now settled.  For arbitrary mixed-characteristic DVR coefficients,
Corollary~\ref{cor:all-dvr-truncations-negative} settles every
principal truncation: generation always fails for $n\ge2$.
For exponents divisible by the residue characteristic, the
uniformizer ideal gives a nonzero detector value of order $p$;
the order of the source class and a general splitting remain open.
Theorem~\ref{thm:pure-power-mesh-vanishing} proves that a different,
standard mesh is zero in two pure-power cases, so those particular
classes cannot describe the whole kernel.  The normal surface in
Proposition~\ref{prop:pure-power-unit-surface} also has a nonzero
divisor class; constructing its full-category specialization is a
separate unresolved question, no longer a generation obstruction.
The square-zero classification includes arbitrary
torsion and arbitrary cardinality;
every nonzero torsion ideal produces infinite-order classes.  Flat dual numbers have nonzero
obstruction in every residue characteristic, with an infinite split
two-torsion subgroup in degree $p$.  The same conclusion holds for
all ramified residue-characteristic-two DVRs.  For arbitrary DVRs
of residue characteristic $p>2$,
Theorem~\ref{thm:odd-ramification} gives infinite split two-torsion
when $v_2(e)\le v_2(p-1)$, and
Corollary~\ref{thm:all-dvr-odd-dual} gives a nonzero exact-order-two
class for every ramification index.  Whether the latter class
splits off in the complementary range is a different question.
A square-zero ideal
killed by $p$ instead gives a countably generated free direct summand.
Orders over $\Z_p$ in quadratic field extensions of $\Q_p$, at every prime, are
completely classified: precisely the
maximal ones are positive.  Every nonmaximal one has infinite two-torsion in its
ordinary-character kernel, and extension to the DVR normalization
kills them.  Binomial coefficient rings at every prime with constant
term in $p^2O$ are negative at every exponent $n\ge2$, with split
countable exact order-$n$ generic torsion and explicit domain examples.  The distinguished monic classification covers all equations reducing
to $t^n$: precisely their regular members are positive.  In the
singular completed-field case it gives exact order $n$, including
all shallow totally ramified equations.  Cohen coefficient rings
extend the criterion to complete $p$-torsion-free rings whose special
fibre is Artinian with a principal maximal ideal, over arbitrary
residue fields.  Completion extends the criterion to every finite free local algebra
over any mixed-characteristic DVR with a principal special fibre
modulo $p$.  It therefore classifies every finite free local
monogenic algebra over an unramified DVR, even with an inseparable
residue-field extension.
The remaining finite-flat order questions include special fibres
with at least two independent tangent directions.
None of these explicit direct summands
is claimed to exhaust the respective kernel.

For a chosen characteristic-zero field specialization $A\to F$,
the split injection identifies the remaining obstruction with
the kernel of $R(A)\to R(F)$.  This is a useful separation of
the problem, not an intrinsic classification of coefficient
rings.  A complete answer must determine that kernel from the
ring and explain when its lattice classes vanish by admissible
exact sequences.  The regular-base argument fails precisely
when modular multiplicity modules no longer admit the required
finite free resolutions.

\section*{Note on method}

The mathematics and text were produced by GPT-6 Astra (ChatGPT Pro
mode).  The main assertions were subjected to separate checks
by the same model; the distinction between proof checks and
independent re-proofs is recorded at \url{\recordurl}.
Earlier recorded checks of the degree-two and split quadratic
examples were made by Claude Opus 5.5 at effort high.
No human has checked this paper.  Our text is dedicated to the
public domain under CC0 1.0.  The repository address for the
mathematical records and the verification program
\emph{Symmetric-lattice checks} is \url{\workbenchurl}.

\begingroup\small\interlinepenalty=10000\setlength{\bibsep}{4pt}
\begin{thebibliography}{99}
\bibitem[Akin et al.(1982)]{ABW1982}
Akin, K., Buchsbaum, D. A., \& Weyman, J. (1982).
Schur functors and Schur complexes.
\emph{Advances in Mathematics, 44}(3), 207--278.
\url{https://doi.org/10.1016/0001-8708(82)90039-1}.
\bibitem[Antieau et al.(2019)]{AntieauGepnerHeller2019}
Antieau, B., Gepner, D., \& Heller, J. (2019).
K-theoretic obstructions to bounded t-structures.
\emph{Inventiones Mathematicae, 216}(1), 241--300.
\url{https://doi.org/10.1007/s00222-018-00847-0}.

\bibitem[Balmer and Schlichting(2001)]{BalmerSchlichting2001}
Balmer, P., \& Schlichting, M. (2001). Idempotent completion of
triangulated categories. \emph{Journal of Algebra, 236}(2), 819--834.
Lemma~1.6, pp.~821--822.
\url{https://doi.org/10.1006/jabr.2000.8529}.
\bibitem[Barthel et al.(2025)]{BarthelEtAl2025}
Barthel, T., Benson, D., Iyengar, S. B., Krause, H., \& Pevtsova, J.
(2025). Lattices over finite group schemes and stratification.
\emph{Compositio Mathematica, 161}, 2911--2946.
\url{https://doi.org/10.1112/S0010437X25102844}.
\bibitem[Berthelot et al.(1971)]{SGA6}
Berthelot, P., Grothendieck, A., \& Illusie, L. (Eds.). (1971).
\emph{Th\'eorie des intersections et th\'eor\`eme de Riemann--Roch
(SGA 6)}. Lecture Notes in Mathematics, 225. Springer.
Expos\'e VI, \S3. \url{https://doi.org/10.1007/BFb0066283}.
\bibitem[Blumberg et al.(2013)]{BlumbergGepnerTabuada2013}
Blumberg, A. J., Gepner, D., \& Tabuada, G. (2013).
A universal characterization of higher algebraic K-theory.
\emph{Geometry \& Topology, 17}(2), 733--838.
\url{https://doi.org/10.2140/gt.2013.17.733}.

\bibitem[B\"uhler(2010)]{Buhler2010}
B\"uhler, T. (2010). Exact categories.
\emph{Expositiones Mathematicae, 28}(1), 1--69.
Section~13.4 and Theorem~13.7, p.~50 of the author version.
\url{https://arxiv.org/abs/0811.1480}.

\bibitem[Burke and Walker(2015)]{BurkeWalker2012}
Burke, J., \& Walker, M. E. (2015). Matrix factorizations in higher codimension.
\emph{Transactions of the American Mathematical Society, 367}(5), 3323--3370.
\url{https://doi.org/10.1090/S0002-9947-2015-06323-5}.
Page locators refer to the 2012 author version, arXiv:1205.2552v1:
\url{https://arxiv.org/abs/1205.2552v1}.
\bibitem[Burroughs(1974)]{Burroughs1974}
Burroughs, J. (1974). Operations in Grothendieck rings and the
symmetric group. \emph{Canadian Journal of Mathematics, 26}(3),
543--550. \url{https://doi.org/10.4153/CJM-1974-050-3}.
\bibitem[Clauwens(2010)]{Clauwens2010}
Clauwens, F. J.-B. J. (2010). \emph{The nilpotence degree of torsion
elements in lambda-rings}. arXiv:1004.0829v1. Proposition~2, p.~2.
\url{https://arxiv.org/abs/1004.0829}.
\bibitem[Eisenbud(1980)]{Eisenbud1980}
Eisenbud, D. (1980). Homological algebra on a complete intersection,
with an application to group representations.
\emph{Transactions of the American Mathematical Society, 260}(1), 35--64.
\url{https://doi.org/10.1090/S0002-9947-1980-0570778-7}.
\bibitem[Gersten(1973)]{Gersten1973}
Gersten, S. M. (1973). Some exact sequences in the higher
$K$-theory of rings. In H. Bass (Ed.),
\emph{Algebraic K-theory I: Higher K-theories}
(Lecture Notes in Mathematics, Vol. 341, pp.~211--243).
Springer. The Addendum on p.~242 records Swan's comparison theorem.

\bibitem[GPT-6(2026)]{GPT6Grinberg2026}
GPT-6. (2026). \emph{The anti-arithmetic product of symmetric functions}
(D. Grinberg, Ed.). October~3, 2026. Proof originating with GPT-5.5.
\url{https://darijgrinberg.gitlab.io/algebra/ariprod-gpt.pdf}.

\bibitem[Grothendieck(n.d.)]{Grothendieck133}
Grothendieck, A. (n.d.). \emph{Manuscript on $\lambda$-rings}.
Montpellier archive, cote 133, scans 6, 28--29
(manuscript pages 5, 27--28).
\bibitem[James(1978)]{James1978}
James, G. D. (1978). \emph{The representation theory of the
symmetric groups}. Lecture Notes in Mathematics, 682. Springer.
\url{https://doi.org/10.1007/BFb0067708}.
\bibitem[Maranda(1953)]{Maranda1953}
Maranda, J.-M. (1953).
On $\mathfrak P$-adic integral representations of finite groups.
\emph{Canadian Journal of Mathematics, 5}, 344--355.
\url{https://doi.org/10.4153/CJM-1953-040-2}.

\bibitem[Orlov(2011)]{Orlov2011}
Orlov, D. (2011). Formal completions and idempotent completions of triangulated
categories of singularities. \emph{Advances in Mathematics, 226}(1), 206--217.
\url{https://doi.org/10.1016/j.aim.2010.06.016}.
Lemmas 2.1, 2.5, and 2.6 are cited from author pp.~3--5:
\url{https://arxiv.org/abs/0901.1859v1}.
\bibitem[Pappas(2015)]{Pappas2015}
Pappas, G. (2015). Adams operations and Galois structure.
\emph{Algebra \& Number Theory, 9}(6), 1477--1515.
\url{https://doi.org/10.2140/ant.2015.9.1477}.
\bibitem[Quillen(1973)]{Quillen1973}
Quillen, D. (1973). Higher algebraic $K$-theory: I.
In H. Bass (Ed.), \emph{Algebraic $K$-theory I: Higher
$K$-theories} (pp.~85--147). Lecture Notes in Mathematics, 341.
Springer. \S4, Theorem~3 and Corollary~1; \S5, Theorem~5.
\url{https://doi.org/10.1007/BFb0067053}.
\bibitem[Roggenkamp(1971)]{Roggenkamp1971}
Roggenkamp, K. W. (1971). Grothendieck groups of Bass orders.
\emph{Canadian Journal of Mathematics, 23}(1), 103--115.
\url{https://doi.org/10.4153/CJM-1971-011-9}.
\bibitem[Swan(1963)]{Swan1963}
Swan, R. G. (1963). The Grothendieck ring of a finite group.
\emph{Topology, 2}, 85--110.
\url{https://doi.org/10.1016/0040-9383(63)90025-9}.
\bibitem[The Stacks Project Authors(2026)]{Stacks2026}
{\raggedright
The Stacks Project Authors. (2026). \emph{The Stacks Project}.
Lemmas 10.51.4, 10.97.1--10.97.4, 10.119.7, 10.153.3,
10.153.7, 10.153.9, 10.157.4, 15.24.18, and 31.30.2.
Tags 00IP, 00MA, 00MB, 00MC, 031C, 00PD, 04GG, 04GK, 04GM, 031S, 0AVB, 0EBM.
\url{https://stacks.math.columbia.edu/tag/00IP};
\url{https://stacks.math.columbia.edu/tag/00PD};
\url{https://stacks.math.columbia.edu/tag/031S};
\url{https://stacks.math.columbia.edu/tag/04GG};
\url{https://stacks.math.columbia.edu/tag/04GK};
\url{https://stacks.math.columbia.edu/tag/04GM};
\url{https://stacks.math.columbia.edu/tag/0BNH};
\url{https://stacks.math.columbia.edu/tag/0AVB};
\url{https://stacks.math.columbia.edu/tag/0EBM}.\par}

\bibitem[Thomason(1997)]{Thomason1997}
Thomason, R. W. (1997). The classification of triangulated subcategories.
\emph{Compositio Mathematica, 105}(1), 1--27. Corollary 2.3, p.~6.
\url{https://doi.org/10.1023/A:1017932514274}.
\bibitem[Weibel(2013)]{Weibel2013}
Weibel, C. A. (2013). \emph{The K-book: An introduction to
algebraic $K$-theory}. Graduate Studies in Mathematics, 145.
American Mathematical Society. Chapter~II, Lemma~7.2,
p.~62 in the separately paginated author chapter version.
\url{https://sites.math.rutgers.edu/~weibel/Kbook/Kbook.II.pdf}.
\bibitem[Williams and Wuthrich(2025)]{WilliamsWuthrich2025}
Williams, C., \& Wuthrich, C. (2025). \emph{Local fields}.
MAGIC037 lecture notes, Spring 2025. Corollary 3.18, p.~49;
Theorem 4.7 and Corollary 4.8, pp.~54--55; Theorem 5.9, pp.~67--68.
\url{https://www.maths.nottingham.ac.uk/plp/pmzcw/download/localfields.pdf}.

\bibitem[Zibrowius(2015)]{Zibrowius2015}
Zibrowius, M. (2015). Symmetric representation rings are
$\lambda$-rings. \emph{New York Journal of Mathematics, 21},
1055--1092. \url{https://nyjm.albany.edu/j/2015/21-49.html}.
\end{thebibliography}
\endgroup
\end{document}
